\documentclass{article}%
\usepackage[T1]{fontenc}%
\usepackage[utf8]{inputenc}%
\usepackage{lmodern}%
\usepackage{textcomp}%
\usepackage{lastpage}%
%
\usepackage{amsmath}%
\usepackage{amssymb}%
\usepackage{amsfonts}%
\usepackage{graphicx}%
\usepackage{hyperref}%
\usepackage{tikz}%
\usepackage{pgfplots}%
\usetikzlibrary{math}%
\title{Unicamp 2027 - 1ª fase (Conhecimentos Gerais, prova Q/Y)\thanks{This exam was solved and automatically generated with gpt-resolve: \url{https://github.com/lgabs/gpt-resolve}}}%
\author{claude{-}opus{-}5{-}5}%
\date{October 30, 2026}%
%
\begin{document}%
\normalsize%
\maketitle%
\section*{Questão 1}

\section*{Solução}

\textbf{1. O que o trecho afirma.} O texto descreve duas fases na relação entre a Igreja e o Estado romano:
\begin{itemize}
  \item \emph{Primeiro}, a Igreja ``se beneficiou do apoio imperial''. O édito de tolerância de Galério (311) havia encerrado a Grande Perseguição. O acordo de Milão (313), entre Constantino e Licínio, garantiu liberdade de culto e a devolução dos bens confiscados aos cristãos. Vieram depois o patrocínio imperial a basílicas, a isenção de encargos para o clero, a convocação de concílios pelo imperador (Niceia, 325) e o Édito de Tessalônica (380), com Teodósio, que fez do cristianismo niceno, chamado ali de ``católico'', a religião oficial do Império.
  \item \emph{Depois}, a Igreja ``ajudou a preservar o legado de Roma'' durante as transformações do século V: a deposição de Rômulo Augústulo (476) e a formação dos reinos germânicos. Bispos e mosteiros conservaram o latim, o direito, a cultura escrita e a organização em dioceses, herdada das circunscrições imperiais.
\end{itemize}
A relação descrita é, portanto, de \emph{reforço mútuo}. O Estado deu à Igreja proteção, recursos e status oficial. A Igreja deu ao poder imperial legitimação religiosa e, mais tarde, garantiu a continuidade da herança romana.

\textbf{2. Análise das alternativas.}

\textbf{a) Falsa.} O texto fala de apoio e continuidade, não de ``ruptura'' nem de ``separação''. A separação entre Igreja e Estado é uma noção moderna e seria anacrônica aqui. No Oriente, o vínculo ficou até \emph{mais estreito}: o imperador Marciano convocou o Concílio de Calcedônia (451), base do chamado cesaropapismo bizantino. Por fim, nem o texto nem a historiografia atribuem uma suposta separação a ``disputas entre classes sociais''.

\textbf{b) Verdadeira.} ``Relação simbiótica, na qual ambas as instituições saíram fortalecidas'' resume as duas fases do trecho: a Igreja cresceu com o apoio imperial, e o poder e o legado romanos se apoiaram na Igreja. O ``por vezes'' é uma ressalva correta, porque houve atritos \emph{dentro} dessa aliança. Exemplos: Constâncio II, ao apoiar o arianismo, exilou Atanásio; Juliano tentou restaurar o paganismo (361--363); Ambrósio, bispo de Milão, impôs penitência a Teodósio após o massacre de Tessalônica (390). Os ``determinados grupos sociais'' favorecidos foram o alto clero e as aristocracias. Muitos bispos vinham da elite senatorial (Ambrósio, Sidônio Apolinário) e ganharam patrimônio, isenções e jurisdição civil (a \emph{episcopalis audientia}). O nome ``Igreja Católica Apostólica Romana'' é anacrônico, pois a distinção entre católicos romanos e ortodoxos só se fixa com o Cisma de 1054. Trata-se, porém, de uma imprecisão de nomenclatura que não altera o conteúdo da alternativa, que se refere à Igreja cristã do Império (o próprio Édito de 380 já usa o termo ``católica''). Não há motivo para descartá-la por isso.

\textbf{c) Falsa.} Contradiz diretamente o trecho (``inicialmente se beneficiou do apoio imperial''). A expansão decisiva ocorreu \emph{com} o apoio das autoridades imperiais, ainda no século IV, e não depois do colapso nem de forma autônoma.

\textbf{d) Falsa.} É o distrator mais forte. O texto diz que a Igreja \emph{preservou} o legado romano, não que \emph{substituiu} as instituições romanas. De fato, no Ocidente pós-476 alguns bispos assumiram funções urbanas. Mas isso foi consequência, e não causa, de uma consolidação que já tinha ocorrido sob o Império, graças ao apoio imperial. Além disso, o Império do Oriente continuou existindo, sem ser substituído pela Igreja.

\textbf{3. Conclusão.} Só a alternativa (b) traduz a ideia central do trecho: uma aliança de benefício mútuo, com atritos ocasionais, entre a Igreja e o Estado romano. Essa aliança fortaleceu as duas instituições e certas elites da Antiguidade Tardia.

\medskip
ANSWER: (B) Relação por vezes simbiótica entre Igreja e Estado romano, que fortaleceu as duas instituições e favoreceu determinados grupos sociais.%
\newpage%
\section*{Questão 2}

\section*{Solução}

\textbf{Ideia central do texto.} O trecho (adaptado de M. Norton) traz três afirmações:
\begin{enumerate}
\item muitos povos indígenas \emph{não separavam rigidamente} humanos e não humanos;
\item o critério principal era a oposição entre \emph{seres selvagens} e \emph{seres familiarizados}. Os familiarizados são parentes humanos e animais \emph{integrados ao convívio cotidiano}, e transformar inimigos, presas ou animais selvagens em seres próximos ao grupo era uma prática valorizada;
\item mesmo com os animais trazidos pelos europeus, esses povos os incorporaram segundo \emph{suas próprias referências culturais}, \emph{diferentes} das concepções europeias de pecuária e de propriedade animal.
\end{enumerate}

\textbf{Análise das alternativas.}

\textbf{a) Incorreta.} Ela inverte o conceito. Para o texto, os seres familiarizados são justamente os \emph{integrados} ao convívio cotidiano, e não os que ``não faziam parte'' dele. A segunda metade (papéis sociais distintos dos europeus) é compatível com o texto, mas não salva a definição errada.

\textbf{b) Incorreta.} O texto diz que as comunidades, ``em vez de separar rigidamente humanos e não humanos'', distinguiam ``sobretudo'' selvagens e familiarizados. Portanto, a oposição humano/não humano \emph{não} era o eixo dessa relação. Além disso, a incorporação dos animais ao grupo era uma prática concreta de convívio, e não apenas ``simbólica''.

\textbf{c) Correta.} A alternativa diz que, para alguns povos, o mundo não se organizava pela oposição humano/não humano, o que corresponde ao item 1. Ela também afirma que a distinção entre familiarizados e selvagens gerava formas de organização social diferentes das europeias, o que corresponde aos itens 2 e 3. O texto aproxima parentes humanos e animais do convívio, valoriza a familiarização de seres de fora do grupo e opõe essa lógica à da pecuária e da propriedade animal europeias. A expressão ``alguns povos'' também respeita a cautela do texto, que fala em ``muitos povos'' e ``diversas comunidades'', não em todos.

\textbf{d) Incorreta.} A primeira parte (a convivência próxima favorecia a integração ao grupo) está de acordo com o texto. A conclusão, porém, o contradiz: os ameríndios incorporavam os animais segundo referências ``diferentes das concepções europeias de pecuária e propriedade animal''. Essas práticas \emph{não} se aproximavam da criação, da domesticação e da propriedade animal europeias.

\textbf{Conclusão.} Só a alternativa (c) reproduz fielmente as duas ideias do texto: a oposição selvagem/familiarizado no lugar da oposição humano/não humano e a diferença em relação ao modelo europeu.

ANSWER: (C)%
\newpage%
\section*{Questão 3}

\section*{Solução}

\textbf{1. O que o trecho estabelece.} O art. 94 diz quem pode ser \emph{Eleitor}, isto é, quem vota na eleição de deputados, senadores e membros dos Conselhos de Província. A Constituição de 1824 previa eleições \emph{indiretas}, em dois graus:
\begin{itemize}
\item nas Assembleias Paroquiais, os \emph{votantes} escolhiam os eleitores. Pelo art. 92, para isso era preciso ter renda líquida anual de pelo menos 100 mil réis;
\item os \emph{eleitores} escolhiam os deputados e senadores. Pelo art. 94, citado no enunciado, ficavam excluídos desse grau, entre outros, os que não tivessem renda líquida anual de 200 mil réis ``por bens de raiz, indústria, comércio, ou emprego'', os libertos e os criminosos pronunciados.
\end{itemize}
O critério central é, portanto, a \textbf{renda}: quanto mais alto o grau eleitoral, maior a renda mínima exigida. Um voto que depende de renda ou de patrimônio mínimo chama-se \textbf{voto censitário}. A palavra vem do \emph{census} romano, o registro dos bens dos cidadãos.

\textbf{2. Análise das alternativas.}
\begin{itemize}
\item[a)] \textbf{Falsa.} O voto não era universal. O próprio trecho exclui quem não tinha a renda mínima, os libertos e os criminosos pronunciados. Além disso, mulheres e escravizados não tinham direitos políticos.
\item[b)] \textbf{Falsa.} A Constituição de 1824 não exigia grau de instrução. Os analfabetos podiam votar desde que tivessem a renda exigida. A exclusão dos analfabetos só veio com a Lei Saraiva (1881). O filtro era econômico, e não intelectual.
\item[c)] \textbf{Verdadeira.} O inciso I condiciona a participação a uma renda líquida mínima de 200 mil réis, vinda de bens de raiz (propriedades), indústria, comércio ou emprego. Os incisos II e III acrescentam outras restrições legais (libertos e criminosos pronunciados). Isso corresponde exatamente a ``comprovação de renda mínima ou propriedades, entre outros temas legais''.
\item[d)] \textbf{Falsa.} O texto não usa a cor ou a origem étnica como critério. A exclusão dos libertos se baseia na condição jurídica de quem nasceu escravizado e foi alforriado. Um homem livre de nascimento, qualquer que fosse sua cor, podia ser eleitor se cumprisse o requisito de renda. A sociedade imperial era profundamente desigual do ponto de vista racial, mas o critério \emph{legal} explícito era a renda (e a condição de liberto), e não a raça.
\end{itemize}

\textbf{3. Conclusão.} Na Constituição de 1824, o voto era censitário (e indireto). Para votar e, mais ainda, para ser eleitor, era preciso comprovar renda mínima, o que restringia a participação política às camadas proprietárias.

ANSWER: (C) censitário, pois condicionava o exercício do voto à comprovação de renda mínima ou propriedades, entre outros temas legais.%
\newpage%
\section*{Questão 4}

\section*{Solução}

\textbf{Ideia central do texto.} Sollai (2024) descreve a ``Batalha do Trigo'' (1925) como um programa de modernização agrária. A proposta era trocar as variedades tradicionais de trigo por cultivares ``de elite'', obtidas por melhoramento genético, para que a Itália deixasse de depender do trigo importado. O texto destaca três pontos:
\begin{enumerate}
  \item a política deu certo no norte, onde as condições ambientais e econômicas eram mais favoráveis;
  \item fracassou no sul semiárido e pobre, porque os agricultores \emph{resistiram}: guardaram as próprias sementes e mantiveram as formas tradicionais de cultivo;
  \item essa resistência levou o Estado a \emph{financiar} estratégias alternativas de ``adaptação profunda'', que aproveitavam a agrobiodiversidade local em vez de substituí-la.
\end{enumerate}
Por isso o autor conclui que a ``Revolução Verde'' fascista foi \emph{desigual e contraditória}. Ela não formou um bloco homogêneo de soluções, ao contrário do que costuma acontecer na Revolução Verde do pós-guerra.

\textbf{Contexto histórico.} A Revolução Verde do pós-Segunda Guerra (décadas de 1940 a 1970) difundiu sementes de alto rendimento obtidas por pesquisa científica e melhoramento genético. Essas sementes vinham num ``pacote'' tecnológico relativamente padronizado, com fertilizantes, agrotóxicos, irrigação e mecanização, aplicado em contextos tão diferentes quanto o México, a Índia e o Brasil.

\textbf{Análise das alternativas.}
\begin{itemize}
  \item[a)] \textbf{Incorreta.} A padronização é típica da Revolução Verde do pós-guerra. O texto afirma justamente que o caso fascista \emph{não} foi homogêneo e que se desdobrou de forma desigual entre o norte e o sul.
  \item[b)] \textbf{Incorreta.} A Revolução Verde do pós-guerra se apoiou exatamente em pesquisa científica e em melhoramento genético, como as variedades de trigo e de arroz de alto rendimento. Portanto, esse não é um traço que a diferencie da Batalha do Trigo; pelo contrário, é algo que as duas têm em comum.
  \item[c)] \textbf{Correta.} A alternativa reproduz o argumento do texto. O projeto enfrentou resistência no sul semiárido e pobre, e essa resistência levou o Estado a financiar as estratégias de ``adaptação profunda'', ou seja, alternativas ajustadas às condições e à agrobiodiversidade locais.
  \item[d)] \textbf{Incorreta.} A alternativa inverte a relação de causa e efeito. O Estado não abandonou o sul, e foi isso que provocou a resistência. A ordem dos fatos foi outra: o Estado tentou impor as novas variedades, os agricultores resistiram e, a partir daí, o Estado passou a desenvolver estratégias alternativas para a região.
\end{itemize}

\textbf{Conclusão.} Só a alternativa (c) está de acordo tanto com o texto (resistência regional seguida de estratégias estatais adaptadas ao local) quanto com o contexto histórico. As demais contradizem o texto, como (a) e (d), ou o conhecimento histórico sobre a Revolução Verde, como (b).

ANSWER: (C) O projeto de modernização agrícola fascista encontrou resistência em determinadas regiões da Itália, levando o Estado a desenvolver estratégias alternativas adaptadas às condições locais.%
\newpage%
\section*{Questão 5}

\section*{Solução}

\textbf{Tema:} religiosidade e cultura material no mundo atlântico escravista (séculos XVII--XVIII), com base no excerto de Laura de Mello e Souza, \emph{O Diabo e a terra de Santa Cruz}.

\textbf{Passo 1 -- O que o texto afirma.}
\begin{itemize}
\item \emph{Função:} eram amuletos usados ao pescoço para ``fechar'' o corpo contra agressões físicas e espirituais. Portanto, eram recursos espirituais de proteção.
\item \emph{Composição:} reuniam elementos ``fruto do contato entre diferentes tradições culturais'': itens da flora e da fauna americanas; partes do corpo humano, ligadas a práticas mágicas de matriz africana; objetos católicos (pedra d'ara, hóstias) e orações e cruzes escritas, de origem europeia.
\item \emph{Difusão:} circulavam por redes que ligavam Minas Gerais, Recife, Rio de Janeiro, Lisboa, Costa da Mina e Angola, isto é, por todo o Atlântico escravista. O uso cresceu no século XVIII.
\end{itemize}
Conclusão: as bolsas eram uma prática \emph{nova}, nascida da mistura (sincretismo) de tradições africanas, europeias e americanas. Não eram a simples continuação de uma delas.

\textbf{Passo 2 -- Análise das alternativas.}

\textbf{(a) Incorreta.} Diz que as bolsas ficaram ``imunes'' a elementos europeus e americanos. O texto afirma o contrário: hóstias, pedra d'ara, orações e cruzes são europeus e católicos, e plantas e animais são americanos.

\textbf{(b) Correta.} Junta os três pontos do texto: proteção espiritual (``fechar o corpo''), uso amplo no Atlântico escravista (redes entre América, Europa e África) e articulação de diferentes tradições em uma prática religiosa nova.

\textbf{(c) Incorreta.} Os objetos católicos foram \emph{incorporados} a uma prática mágica, que não era cristã, e não \emph{substituíram} as práticas africanas. Isso é ressignificação, não cristianização. Tanto que os mandingueiros eram presos pela Inquisição, o que mostra que a Igreja via essas práticas como heterodoxas, e não como sinal de conversão.

\textbf{(d) Incorreta.} O texto não fala de matriz islâmica nem de leitura do Alcorão. Os escritos citados são orações e símbolos cristãos. Além disso, a alternativa reduz as bolsas a uma única tradição, o que contraria a ideia central do excerto, que é a mistura de culturas. O nome ``mandinga'' não autoriza essa conclusão.

\textbf{Conclusão:} só a alternativa (b) está de acordo com o excerto: as bolsas de mandinga eram amuletos de proteção espiritual difundidos no Atlântico escravista, nascidos da combinação de elementos africanos, europeus e americanos em uma nova prática religiosa.

ANSWER: (B)%
\newpage%
\section*{Questão 6}

\section*{Solução}

\textbf{Ideia central do texto.} O trecho do relatório da Comissão Nacional da Verdade (2014) aponta três pontos sobre a política indigenista da ditadura:
\begin{enumerate}
\item \emph{Falta de controle social}: houve arbitrariedades de colonos e fazendeiros, conluio de chefes locais com a polícia e os militares, e políticas decididas em Brasília. A censura e a propaganda ufanista dificultavam as denúncias.
\item \emph{Base ideológica}: a doutrina de Segurança Nacional, ``etnocêntrica, racista e desenvolvimentista'', servia para ``legitimar'' violações de direitos humanos.
\item \emph{Conclusão explícita}: ``a sobrevivência dos povos indígenas e a existência de suas culturas foram subordinadas aos grandes projetos de colonização do país''. O exemplo é a rodovia BR-80. Ela cortou o Parque Nacional do Xingu, dividiu os Tukarramãe e levou-lhes doenças como o sarampo. O presidente da Funai, o próprio órgão indigenista, respondeu: ``o Parque Nacional do Xingu não pode impedir o progresso do país''.
\end{enumerate}
Esse contexto histórico corresponde ao Plano de Integração Nacional (1970), à Transamazônica e à ocupação da Amazônia sob o lema ``integrar para não entregar''. Nele, o indígena era visto como alguém a ser ``integrado'' à sociedade nacional, e não como sujeito de direitos territoriais.

\textbf{Análise das alternativas.}

\textbf{a) Incorreta.} Segundo o texto, a integração territorial não foi compatibilizada com a autonomia indígena: a sobrevivência dos povos foi \emph{subordinada} aos projetos. Prova disso é que nem uma área protegida como o Xingu foi poupada pela BR-80.

\textbf{b) Correta.} A alternativa parafraseia a tese do texto. Os projetos de ocupação seguiam uma concepção desenvolvimentista ligada à Segurança Nacional, na qual a permanência dos indígenas em suas terras ficava abaixo dos interesses do Estado. A frase do general Bandeira de Mello (``não pode impedir o progresso'') mostra essa hierarquia de prioridades.

\textbf{c) Incorreta.} O texto fala em regime autoritário, censura e decisões tomadas em Brasília, sem controle social. Não houve participação indígena nas decisões sobre os projetos; pelo contrário, a própria Funai justificou a obra que dividiu os Tukarramãe.

\textbf{d) Incorreta.} As rodovias e a colonização não protegeram as terras indígenas. Elas as fragmentaram, como a BR-80 fez com o Xingu, e agravaram os conflitos com colonos e fazendeiros, além de espalhar doenças.

\textbf{Conclusão.} Somente a alternativa (b) está de acordo com o trecho, que afirma a subordinação da sobrevivência e dos territórios indígenas aos grandes projetos desenvolvimentistas do Estado.

ANSWER: (B)%
\newpage%
\section*{Questão 7}

\section*{Solução}

\textbf{Ideia central.} A questão pede que se cruzem duas fontes: a \emph{imagem}, que mostra uma vitrine de museu com dezenas de consoles, controles, pistolas de jogo e portáteis de várias gerações, e o \emph{excerto}, que trata do papel do console doméstico. A alternativa correta precisa ser sustentada pelas duas e não pode trazer afirmações que nenhuma delas autoriza.

\medskip
\textbf{Passo 1 -- O que diz o excerto.}
\begin{itemize}
\item O console (Magnavox Odyssey em 1972, Atari 2600, Nintendo e Sega nos anos 1980) ``deslocou o jogo eletrônico dos laboratórios, arcades e espaços públicos para o interior das casas''. Mudam o espaço e as formas de convívio em que se joga, ou seja, a \emph{sociabilidade}, e o jogo vira um bem de \emph{consumo} doméstico vendido em massa.
\item Ligado à televisão, ``a tela mais importante da cultura audiovisual de massa do pós-guerra'', o console transformou ``uma superfície antes predominantemente expositiva em um espaço interativo e responsivo''. Assim, mudou a \emph{relação com a tela da TV}: o espectador, antes passivo, passa a agir sobre ela.
\end{itemize}

\textbf{Passo 2 -- O que mostra a imagem.} São objetos físicos (aparelhos, controles, acessórios) de diferentes épocas, organizados como acervo e expostos num museu (ACMI, 2012). Isso revela uma \emph{cultura material} própria dos \emph{games}, que já se tornou patrimônio e objeto de \emph{memória}. O título do artigo citado, ``videogames como fontes históricas'', aponta na mesma direção.

\textbf{Passo 3 -- Análise das alternativas.}
\begin{itemize}
\item[a)] \textbf{Correta.} Cada parte tem apoio nas fontes. A ``cultura material específica'' aparece na imagem. A ``relação histórica com as telas da televisão'' corresponde à passagem da tela expositiva para a tela interativa, descrita no excerto. A sociabilidade e o consumo mudam com a ida do jogo para dentro de casa e com a comercialização em massa dos consoles. A memória aparece na musealização dos objetos e no uso dos jogos como fontes históricas.
\item[b)] \textbf{Incorreta.} O texto fala de uma nova relação entre ver e agir \emph{do lado do jogador}. Ele não afirma que o console mudou o modo de \emph{filmar e exibir} do cinema e da TV, isto é, a produção audiovisual dessas indústrias.
\item[c)] \textbf{Incorreta.} Nenhuma das fontes menciona uma adaptação do console às técnicas de Hollywood. O excerto também não fala em aprofundar as diferenças entre a linguagem do cinema e a da TV. Pelo contrário, destaca que os \emph{games} \emph{articulam} visualidade, som, regra e performance.
\item[d)] \textbf{Incorreta.} O console precisava do televisor para funcionar, o que deu um novo uso à tela doméstica e reforçou sua centralidade. O excerto não fala em ``recusa'' da linguagem dos \emph{games} pela indústria televisiva nem em perda de espaço da TV dentro de casa.
\end{itemize}

\textbf{Conclusão.} Só a alternativa (a) junta a dimensão material mostrada na imagem com a mudança na relação com a tela descrita no excerto, e suas consequências sobre sociabilidade, consumo e memória. As outras atribuem às fontes afirmações que elas não fazem.

\medskip
ANSWER: (A)%
\newpage%
\section*{Questão 8}

\section*{Solução}

\textbf{1) Antes da liberação do gás.} A água do frasco é praticamente pura: os únicos íons vêm da autoionização ($\mathrm{H_2O \rightleftharpoons H^+ + OH^-}$, com $[\mathrm{H^+}]=[\mathrm{OH^-}]=10^{-7}\ \mathrm{mol/L}$). Por isso a condutividade é muito baixa, cerca de $0{,}1\ \mu\mathrm{S/cm}$. A pressão é a do ar no frasco fechado, $1{,}0\ \mathrm{atm}$. Os quatro gráficos concordam nesse trecho.

\textbf{2) O que acontece com a pressão logo após a abertura.} O balão está inflado, ou seja, o gás dentro dele está a uma pressão maior que a do frasco. Ao abrir a torneira, esse gás entra no frasco fechado. Pela equação dos gases ($pV=nRT$, com $V$ e $T$ praticamente constantes), o aumento de $n$ faz a pressão subir de forma brusca. Isso vale para \emph{qualquer} gás. Portanto os gráficos com pressão constante em $1{,}0\ \mathrm{atm}$ (alternativas b e d) são fisicamente incoerentes.

\textbf{3) O que diferencia CO$_2$ e He.}
\begin{itemize}
\item \textbf{He}: é um gás nobre, inerte e muito pouco solúvel em água. Ele não reage com a água e não gera íons. Após o salto, a pressão ficaria praticamente constante no valor mais alto, e a condutividade continuaria em $\approx 0{,}1\ \mu\mathrm{S/cm}$. Nenhum dos gráficos mostra esse comportamento, então as alternativas c e d estão descartadas.
\item \textbf{CO$_2$}: é bastante solúvel em água e forma um ácido fraco:
\[
\mathrm{CO_2(g) \rightleftharpoons CO_2(aq)}, \qquad
\mathrm{CO_2(aq) + H_2O(l) \rightleftharpoons H^+(aq) + HCO_3^-(aq)}.
\]
\end{itemize}
Com o CO$_2$, a observação fica assim:
\begin{itemize}
\item \emph{Pressão:} primeiro há o salto (item 2). Em seguida, o CO$_2$ passa aos poucos da fase gasosa para a solução. Assim, $n_{\text{gás}}$ diminui e a pressão cai gradualmente, de $\approx 1{,}08$ para $\approx 1{,}01\ \mathrm{atm}$.
\item \emph{Condutividade:} os íons $\mathrm{H^+}$ e $\mathrm{HCO_3^-}$ produzidos aumentam a condutividade. A dissolução é lenta, pois depende da transferência de massa pela interface gás--líquido, sem agitação. Por isso a condutividade cresce de forma \emph{contínua} e tende a um patamar quando o sistema se aproxima do equilíbrio. Ela não pode subir de repente e depois cair, como no gráfico b: nada removeria os íons formados.
\end{itemize}

\textbf{Conferência da ordem de grandeza.} As constantes a seguir são valores tabelados usuais, não fornecidos na prova. Considere água saturada com CO$_2$ a $p_{\mathrm{CO_2}}\lesssim 1\ \mathrm{atm}$, com constante de Henry $k_H\approx 3{,}4\times10^{-2}\ \mathrm{mol\,L^{-1}\,atm^{-1}}$ e $K_{a1}\approx4{,}5\times10^{-7}$. Então:
\[
[\mathrm{H^+}]\approx\sqrt{K_{a1}\,k_H\,p}\approx\sqrt{4{,}5\times10^{-7}\cdot3{,}4\times10^{-2}}\approx1{,}2\times10^{-4}\ \mathrm{mol/L}\quad(\mathrm{pH}\approx3{,}9).
\]
As condutividades molares iônicas são $\lambda_{\mathrm{H^+}}\approx350$ e $\lambda_{\mathrm{HCO_3^-}}\approx45\ \mathrm{S\,cm^2\,mol^{-1}}$. Daí:
\[
\kappa\approx(350+45)\cdot1{,}2\times10^{-7}\ \mathrm{S/cm}\approx5\times10^{-5}\ \mathrm{S/cm}\approx 50\ \mu\mathrm{S/cm}.
\]
O valor é de algumas dezenas de $\mu\mathrm{S/cm}$, coerente com os $\approx 60\ \mu\mathrm{S/cm}$ do gráfico a. É um aumento de mais de duas ordens de grandeza, o que justifica o uso da escala logarítmica.

\textbf{Conclusão.} O gás é o CO$_2$. O gráfico correto mostra a pressão com um salto seguido de queda gradual, por causa da absorção do CO$_2$ pela água. Mostra também a condutividade crescendo continuamente, por causa da formação de $\mathrm{H^+}$ e $\mathrm{HCO_3^-}$. Esse gráfico é o da alternativa a.

ANSWER: (A) CO$_2$; pressão sobe e depois cai gradualmente, e a condutividade cresce continuamente%
\newpage%
\section*{Questão 9}

\section*{Solução}

\textbf{1) O que é a substituição galvânica.} O metal sólido (Ag) se oxida e passa para a solução. Ao mesmo tempo, o íon metálico da solução se reduz e se deposita sobre a partícula. O processo só ocorre espontaneamente se o íon tiver potencial de redução maior que o do par Ag$^{+}$/Ag ($+0{,}80$ V). Nesse caso, a fem da pilha formada é positiva e $\Delta G^\circ=-nFE^\circ<0$.

\textbf{2) Teste dos dois íons das alternativas.}

\emph{Au$^{3+}$:}
\[
\begin{aligned}
\text{oxidação:}&\quad 3\,\text{Ag}(s)\rightarrow 3\,\text{Ag}^{+}(aq)+3e^{-}\\
\text{redução:}&\quad \text{Au}^{3+}(aq)+3e^{-}\rightarrow \text{Au}(s)\\
\text{global:}&\quad 3\,\text{Ag}(s)+\text{Au}^{3+}(aq)\rightarrow 3\,\text{Ag}^{+}(aq)+\text{Au}(s)
\end{aligned}
\]
\[
E^\circ=E^\circ_{\text{red}}(\text{Au}^{3+}/\text{Au})-E^\circ_{\text{red}}(\text{Ag}^{+}/\text{Ag})=0{,}93-0{,}80=+0{,}13\ \text{V}>0,
\]
logo a reação é espontânea. Há ainda $\Delta G^\circ=-3\cdot 96500\cdot 0{,}13\approx -38$ kJ por mol de Au.

\emph{Cu$^{2+}$:} na reação $2\,\text{Ag}+\text{Cu}^{2+}\rightarrow 2\,\text{Ag}^{+}+\text{Cu}$,
\[
E^\circ=0{,}34-0{,}80=-0{,}46\ \text{V}<0 .
\]
Essa reação não é espontânea. O Cu$^{2+}$ não consegue oxidar a prata (a reação espontânea seria a inversa: Cu reduzindo Ag$^{+}$). Por isso as alternativas (c) e (d) estão descartadas. Pelo mesmo critério, Pt$^{2+}$ também serviria e Zn$^{2+}$ não, mas esses íons não aparecem nas alternativas.

\textbf{3) Mudança de morfologia.} A estequiometria mostra que \emph{3 átomos de Ag saem} da partícula para cada \emph{1 átomo de Au que entra}. Ag e Au têm a mesma estrutura cristalina (cúbica de faces centradas) e volumes atômicos praticamente iguais (parâmetro de rede $\approx 4{,}08$ \AA\ nos dois). Assim, o volume de metal que permanece é cerca de $1/3$ do volume de prata que foi oxidada:
\[
\frac{V_{\text{Au depositado}}}{V_{\text{Ag removida}}}\approx\frac{1}{3}.
\]
O ouro se deposita na superfície do cubo e forma com a prata uma casca de liga Au--Ag que mantém o contorno cúbico. Já a prata é consumida pelo interior da partícula. O resultado é uma estrutura oca, com paredes porosas, cheias de furos: as chamadas \emph{nanogaiolas}. Isso também concorda com o enunciado, que fala em alterações ``na sua composição \emph{e morfologia}''. Por isso as figuras com o cubo maciço inalterado, (a) e (c), estão descartadas.

\textbf{4) Conclusão.} O íon precisa ter $E^\circ_{\text{red}}>0{,}80$ V, e nas alternativas só o Au$^{3+}$ cumpre isso. A troca de 3 Ag por 1 Au deixa o cubo oco e poroso. A única alternativa que reúne as duas condições é a (b).

\medskip
ANSWER: (B) Au$^{3+}$; o nanocubo maciço vira uma nanogaiola oca e porosa de liga Au--Ag.%
\newpage%
\section*{Questão 10}

\section*{Solução}

\textbf{Reação (lida na figura):}
\[
\tfrac{1}{2}\,\mathrm{N_2(g)} + \tfrac{3}{2}\,\mathrm{H_2(g)} \rightleftharpoons \mathrm{NH_3(g)}
\]
O gráfico mostra $K$ em função de $T$, com a pressão fixa em 300 atm.

\medskip
\textbf{1) Efeito da temperatura.}
No gráfico, a curva só desce. Em $T \approx 250$ K temos $K \approx 10^{5}$; em $T \approx 880$ K temos $K \approx 10^{-4}$. Portanto $K$ cai cerca de 9 ordens de grandeza.

Um valor menor de $K$ indica menor proporção de produto no equilíbrio. Logo, \emph{aumentar a temperatura desfavorece a formação de $\mathrm{NH_3}$}.

Isso é coerente com a equação de van 't Hoff:
\[
\frac{d\ln K}{dT}=\frac{\Delta H^\circ}{RT^2}.
\]
Como $K$ diminui com $T$, temos $\Delta H^\circ<0$, isto é, a síntese é exotérmica. Pelo princípio de Le Chatelier, aquecer desloca o equilíbrio no sentido endotérmico, que é a decomposição da amônia.

\medskip
\textbf{2) Efeito da pressão.}
A uma dada temperatura, $K$ não depende da pressão. O que muda com a pressão é a \emph{composição} de equilíbrio. Contando os mols de gás:
\[
\Delta n_{\text{gás}} = 1-\left(\tfrac{1}{2}+\tfrac{3}{2}\right) = -1 < 0 .
\]
A formação do produto reduz o número de mols de gás.

Para gases ideais, escrevemos $p_i = x_i P$, em que $x_i$ é a fração molar e $P$ a pressão total. Então:
\[
K=\frac{x_{\mathrm{NH_3}}\,P}{\left(x_{\mathrm{N_2}}P\right)^{1/2}\left(x_{\mathrm{H_2}}P\right)^{3/2}}
=\frac{x_{\mathrm{NH_3}}}{x_{\mathrm{N_2}}^{1/2}\,x_{\mathrm{H_2}}^{3/2}}\cdot P^{-1}
\quad\Longrightarrow\quad
\frac{x_{\mathrm{NH_3}}}{x_{\mathrm{N_2}}^{1/2}\,x_{\mathrm{H_2}}^{3/2}} = K\,P .
\]
Com $T$ fixa, $K$ é constante. Assim, aumentar $P$ aumenta a razão entre as frações molares, ou seja, aumenta a fração de $\mathrm{NH_3}$ no equilíbrio.

É o mesmo resultado dado por Le Chatelier: a compressão desloca o equilíbrio para o lado com menos mols de gás. Logo, \emph{aumentar a pressão favorece a formação de $\mathrm{NH_3}$}. Por isso o processo Haber--Bosch opera a pressões altas, como os 300 atm do gráfico.

\medskip
\textbf{Observação:} na indústria se usa temperatura alta mesmo ela sendo desfavorável ao equilíbrio. O motivo é cinético: a reação fica mais rápida, junto com o efeito do catalisador. Isso não contradiz a conclusão termodinâmica acima.

\medskip
\textbf{Análise das alternativas.}
\begin{itemize}
\item a) Errada: o aumento de $T$ desfavorece o produto.
\item b) Errada nos dois efeitos.
\item c) Errada: o aumento de $P$ favorece o produto.
\item d) Correta: o aumento de $T$ desfavorece o produto ($K$ decrescente) e o aumento de $P$ o favorece ($\Delta n_{\text{gás}}<0$).
\end{itemize}

ANSWER: (D) desfavorecida com o aumento da temperatura, mas favorecida com o aumento da pressão.%
\newpage%
\section*{Questão 11}

\section*{Solução}

\textbf{1. O que a aplicação exige.} Para captar água do ar, os poros da MOF precisam atrair fortemente as moléculas de $\mathrm{H_2O}$. A água é muito polar e forma ligações de hidrogênio: é doadora (H ligado a O) e aceptora (pares de elétrons livres do O). Logo, o melhor ligante é o mais \emph{hidrofílico}, ou seja, o que tem mais grupos capazes de formar ligações de hidrogênio e interações dipolo--dipolo com a água. Grupos apolares (anel aromático, $-\mathrm{CH_3}$) interagem com a água apenas por forças de dispersão (London), que são fracas, e tornam o material hidrofóbico.

\textbf{2. Análise dos ligantes.}

\begin{itemize}
\item \textbf{A (ácido 2-sulfotereftálico):} além dos dois grupos $-\mathrm{COOH}$, tem o grupo sulfônico $-\mathrm{SO_3H}$, que é fortemente polar e ácido. Ele tem três oxigênios aceptores e um $\mathrm{O-H}$ doador, que é muito ácido e pode até ionizar-se a $-\mathrm{SO_3^-}$, um grupo que se hidrata intensamente (íon--dipolo). Na MOF, os carboxilatos se ligam aos centros metálicos, mas o $-\mathrm{SO_3H}$ fica \emph{livre}, voltado para dentro dos poros. Por isso A é o ligante de maior afinidade pela água.

\item \textbf{B (ácido tereftálico):} tem dois grupos $-\mathrm{COOH}$, com quatro oxigênios, que formam ligações de hidrogênio como doadores e como aceptores. Mesmo coordenados ao metal, os oxigênios dos carboxilatos ainda podem aceitar ligações de hidrogênio. Por outro lado, não tem nenhum grupo polar além desses, e o anel benzênico é apolar. A afinidade pela água é intermediária.

\item \textbf{C (2-metilimidazol):} tem apenas um $\mathrm{N-H}$ doador e um N aceptor. O restante da molécula ($-\mathrm{CH_3}$ e a ligação $\mathrm{C=C}$ do anel) é apolar. Na MOF (ZIF-8), o $\mathrm{N-H}$ é desprotonado e os \emph{dois} nitrogênios se coordenam ao $\mathrm{Zn^{2+}}$. Assim, não sobra nenhum sítio livre para ligação de hidrogênio, e o poro fica revestido por grupos metila e C--H. Isso explica por que o ZIF-8 é conhecido por ser hidrofóbico. É o pior dos três para captar água.
\end{itemize}

\textbf{3. Ordem de preferência.} Considerando o número e a força das interações com a água (íon--dipolo e ligação de hidrogênio $>$ dipolo--dipolo $>$ dispersão de London), temos
\[
\text{afinidade por } \mathrm{H_2O}:\quad \mathbf{A}\ (-\mathrm{SO_3H} + 2\,-\mathrm{COOH}) \;>\; \mathbf{B}\ (2\,-\mathrm{COOH}) \;>\; \mathbf{C}\ (\text{N coordenados}, -\mathrm{CH_3}\ \text{apolar}).
\]
Portanto a primeira opção é \textbf{A} e a segunda é \textbf{B}.

\textbf{4. Conferência com as alternativas.} A única alternativa que coloca A em primeiro lugar é a (c). As alternativas (a) e (d) põem B em primeiro, o que deixa de lado o grupo sulfônico, muito mais hidrofílico. A alternativa (b) põe C em primeiro, mas C é o ligante mais hidrofóbico dos três.

ANSWER: (C) A e B%
\newpage%
\section*{Questão 12}

\section*{Solução}

\textbf{1. O que acontece no organismo.} O metanol, em si, é pouco tóxico. O problema está em seu metabolismo no fígado, que ocorre em duas oxidações sucessivas:
\[
\mathrm{CH_3OH}\ \xrightarrow{\text{álcool desidrogenase}}\ \mathrm{HCHO}\ \xrightarrow{\text{aldeído desidrogenase}}\ \mathrm{HCOOH}.
\]
O ácido fórmico ($K_a \approx 1{,}8\times10^{-4}$, $\mathrm{p}K_a\approx 3{,}75$) está praticamente todo ionizado no pH do sangue, que é cerca de $7{,}4$:
\[
\mathrm{HCOOH} \longrightarrow \mathrm{HCOO^-} + \mathrm{H^+}.
\]
Portanto a intoxicação gera ácido. O tratamento confirma isso. O etanol e o fomepizol competem com o metanol pela álcool desidrogenase (o fomepizol a inibe), o que bloqueia a formação de $\mathrm{HCOOH}$. A hemodiálise retira do sangue o metanol e o formiato.

\textbf{2. Efeito sobre o tampão do sangue.} O principal tampão do plasma é o par $\mathrm{H_2CO_3/HCO_3^-}$. O $\mathrm{H^+}$ liberado pelo ácido fórmico é neutralizado pelo bicarbonato:
\[
\mathrm{HCO_3^-} + \mathrm{H^+} \longrightarrow \mathrm{H_2CO_3} \longrightarrow \mathrm{CO_2} + \mathrm{H_2O}.
\]
Assim, $[\mathrm{HCO_3^-}]$ diminui. Pela equação de Henderson--Hasselbalch,
\[
\mathrm{pH} = 6{,}1 + \log\frac{[\mathrm{HCO_3^-}]}{0{,}03\,p_{\mathrm{CO_2}}},
\]
uma queda de $[\mathrm{HCO_3^-}]$ faz o pH cair. Nos valores normais ($[\mathrm{HCO_3^-}]\approx 24$ mmol/L e $p_{\mathrm{CO_2}}\approx 40$ mmHg), obtemos $\mathrm{pH}\approx 6{,}1+\log 20\approx 7{,}40$.

Como a perturbação vem de um ácido produzido pelo metabolismo, e não de um problema respiratório, o quadro é de \textbf{acidose metabólica}, que o enunciado chama de ``acidez metabólica''. Ela se caracteriza por:
\begin{itemize}
\item bicarbonato baixo, abaixo do normal de cerca de 22 a 26 mmol/L; o critério é $[\mathrm{HCO_3^-}] < 15$ mmol/L;
\item pH arterial baixo, abaixo da faixa normal de 7,35 a 7,45; o critério é $\mathrm{pH} < 7{,}30$.
\end{itemize}

\textbf{3. Concentração de metanol.} O tratamento é indicado quando o tóxico está \emph{acima} de um limite, ou seja, $[\mathrm{CH_3OH}] > 50$ mg/dL. Conferindo a ordem de grandeza: $50$ mg/dL $= 0{,}50$ g/L, e $\dfrac{0{,}50\ \text{g/L}}{32{,}04\ \text{g/mol}} \approx 1{,}6\times10^{-2}$ mol/L $\approx 16$ mmol/L. É uma quantidade de metanol capaz de gerar, ao ser metabolizada, ácido suficiente para consumir boa parte do bicarbonato do sangue. O valor é coerente.

\textbf{4. Análise das alternativas.}
\begin{itemize}
\item[a)] Errada. A produção de ácido causa acidose, e não alcalose. Além disso, a alternativa se contradiz: chama de alcalose um quadro com $\mathrm{pH}<7{,}30$.
\item[b)] Errada. $[\mathrm{CH_3OH}]<50$ mg/dL é uma concentração menor que o limite. Além disso, $[\mathrm{HCO_3^-}]>15$ mmol/L e $\mathrm{pH}>7{,}30$ não indicam acidose grave.
\item[c)] Correta. Metanol acima do limite, bicarbonato consumido ($<15$ mmol/L) e pH baixo ($<7{,}30$): é a acidose metabólica causada pelo ácido fórmico.
\item[d)] Errada. A concentração de metanol está abaixo do limite, o quadro indicado é de alcalose e o critério $\mathrm{pH}>7{,}30$ não corresponde ao quadro tóxico.
\end{itemize}

ANSWER: (C) $[\mathrm{CH_3OH}] > 50$ mg/dL, acidose metabólica ($\mathrm{HCO_3^-} < 15$ mmol/L e pH arterial $< 7{,}30$)%
\newpage%
\section*{Questão 13}

\section*{Solução}

\textbf{1) Gases produzidos na eletrólise da solução aquosa de NaCl.}
Na água da piscina há os íons $\mathrm{Na^+}$ e $\mathrm{Cl^-}$ do sal dissolvido, além da própria água. Por isso é preciso ver qual espécie reage em cada eletrodo.

\emph{Cátodo (redução).} As espécies que podem ser reduzidas são $\mathrm{Na^+}$ e $\mathrm{H_2O}$. Como $E^\circ(\mathrm{Na^+/Na}) = -2{,}71\ \mathrm{V}$ é muito menor que o potencial de redução da água ($-0{,}83\ \mathrm{V}$ em meio básico), quem se reduz é a água:
\[
\mathrm{2\,H_2O(l) + 2\,e^- \longrightarrow H_2(g) + 2\,OH^-(aq)}
\]

\emph{Ânodo (oxidação).} As espécies que podem ser oxidadas são $\mathrm{Cl^-}$ e $\mathrm{H_2O}$. Pelos potenciais padrão, o $\mathrm{O_2}$ seria um pouco favorecido ($1{,}23\ \mathrm{V}$ contra $1{,}36\ \mathrm{V}$). Porém a formação de $\mathrm{O_2}$ tem uma sobretensão muito grande, e na prática o íon cloreto é que se oxida, como no processo cloro-álcali industrial:
\[
\mathrm{2\,Cl^-(aq) \longrightarrow Cl_2(g) + 2\,e^-}
\]
O próprio enunciado confirma isso: um dos produtos da eletrólise ``reage gerando $\mathrm{HClO}$''. Esse produto precisa conter cloro, e o único que contém é o $\mathrm{Cl_2}$.

Somando as semirreações, a reação global é
\[
\mathrm{2\,NaCl(aq) + 2\,H_2O(l) \longrightarrow 2\,NaOH(aq) + H_2(g) + Cl_2(g)}
\]
Logo, os gases produzidos são $\mathrm{H_2}$ e $\mathrm{Cl_2}$. Com isso ficam eliminadas as alternativas (b) e (c), que indicam $\mathrm{O_2}$.

\textbf{2) Equação de formação do HClO.}
O $\mathrm{Cl_2}$ formado reage com a água em um desproporcionamento: o cloro, com Nox $0$, passa ao mesmo tempo a $+1$ (no $\mathrm{HClO}$) e a $-1$ (no $\mathrm{Cl^-}$):
\[
\mathrm{Cl_2(g) + H_2O(l) \rightleftharpoons HClO(aq) + H^+(aq) + Cl^-(aq)}
\]
Conferindo o balanceamento:
\begin{center}
\begin{tabular}{l|c|c}
 & reagentes & produtos \\ \hline
Cl & 2 & $1 + 1 = 2$ \\
H & 2 & $1 + 1 = 2$ \\
O & 1 & 1 \\
carga & 0 & $0 + 1 - 1 = 0$ \\
\end{tabular}
\end{center}
Os átomos e as cargas estão conservados. Os elétrons também: um Cl ganha 1 elétron (de $0$ para $-1$) e o outro perde 1 elétron (de $0$ para $+1$).

\textbf{3) Análise das alternativas.}
\begin{itemize}
\item (a) Os gases estão corretos, mas a equação $\mathrm{Cl_2 + H_2O \rightleftharpoons HClO + Cl^-}$ não está balanceada: há 2 H à esquerda e 1 H à direita, e a carga passa de $0$ para $-1$. Falta o $\mathrm{H^+}$.
\item (b) $\mathrm{O_2}$ não é produto da eletrólise. Além disso, a equação não fecha: são 2 H à esquerda e 3 H à direita, e a carga passa de $-1$ para $0$.
\item (c) $\mathrm{O_2}$ não é produto da eletrólise. A equação também não fecha: são 3 O à esquerda e 1 O à direita, e a carga passa de $-1$ para $+1$.
\item (d) Os gases ($\mathrm{H_2}$ e $\mathrm{Cl_2}$) estão corretos e a equação está balanceada em átomos e em carga. \textbf{Esta é a alternativa correta.}
\end{itemize}

Observação: o $\mathrm{HClO}$, ao agir como bactericida, volta a ser reduzido a $\mathrm{Cl^-}$. Assim o sal é regenerado e não é consumido, o que explica por que esse tratamento é econômico.

ANSWER: (D) $\mathrm{H_2}$ e $\mathrm{Cl_2}$; $\mathrm{Cl_2(g) + H_2O(l) \rightleftharpoons HClO(aq) + H^+(aq) + Cl^-(aq)}$%
\newpage%
\section*{Questão 14}

\section*{Solução}

\textbf{1) Absorção $\times$ metabolização (curva contínua).}
A concentração plasmática $C$ varia de acordo com o balanço
\[
\frac{dC}{dt} = v_{\text{absorção}} - v_{\text{metabolização}} .
\]
Logo após cada dose (a cada 24 h), $C$ sobe de forma quase vertical. Na 1ª dose, por exemplo, vai de $0$ a cerca de $305\ \text{ng/mL}$ em poucas horas, ou seja, a uma taxa da ordem de dezenas de $\text{ng mL}^{-1}\,\text{h}^{-1}$. Durante essa subida, $dC/dt>0$, portanto a entrada do fármaco no plasma é mais rápida que a sua retirada:
\[
v_{\text{absorção}} > v_{\text{metabolização}} .
\]
Já a queda que vem depois (só metabolização/eliminação) é lenta: de cerca de $305$ para $185\ \text{ng/mL}$ em quase 20 h, isto é, apenas alguns $\text{ng mL}^{-1}\,\text{h}^{-1}$. Outra evidência é o acúmulo: o fármaco não é eliminado por completo antes da dose seguinte, e por isso os vales ($185 \to 265 \to 305 \to \dots \to 325\ \text{ng/mL}$) e os picos ($305 \to 460 \to \dots \to 575\ \text{ng/mL}$) crescem até se estabilizarem. Com isso, as alternativas (c) e (d) ficam descartadas.

\textbf{2) Dose diária (curva tracejada).}
A dose acumulada sobe em degraus iguais de cerca de $35\ \text{mg}$ a cada 24 h ($35, 70, 105, \dots, 350\ \text{mg}$ em 10 dias). Portanto:
\[
d = 35\ \text{mg/dia} \quad\Rightarrow\quad \frac{35\ \text{mg}}{50\ \text{kg}} = 0{,}7\ \text{mg kg}^{-1}\,\text{dia}^{-1},
\]
valor que está dentro da faixa da bula ($0{,}5$ a $1{,}0\ \text{mg kg}^{-1}$), o que confirma a leitura do gráfico.

\textbf{3) Duração mínima do tratamento.}
A meta é uma dose acumulada entre $120$ e $150\ \text{mg kg}^{-1}$. Para $m=50\ \text{kg}$:
\[
D_{\min} = 120 \cdot 50 = 6000\ \text{mg}, \qquad D_{\max} = 150 \cdot 50 = 7500\ \text{mg}.
\]
O tempo \emph{mínimo} corresponde a atingir o limite inferior:
\[
t_{\min} = \frac{6000\ \text{mg}}{35\ \text{mg/dia}} \approx 171\ \text{dias} \approx 170\ \text{dias}.
\]
Os 215 dias correspondem ao limite \emph{superior} da faixa: $7500/35 \approx 214$ dias. Esse seria o tempo para atingir $150\ \text{mg kg}^{-1}$, e não o mínimo. Por isso as alternativas (a) e (c) estão erradas na duração.

\textbf{Conclusão.} A velocidade de absorção é maior que a de metabolização, e a paciente precisa tomar o remédio por no mínimo cerca de 170 dias. Só a alternativa (b) reúne as duas afirmações corretas.

ANSWER: (B)%
\newpage%
\section*{Questão 15}

\section*{Solução}

\textbf{Modelagem.} Sejam $c$ o número de pizzas de calabresa e $b$ o número de pizzas de brigadeiro vendidas ($c,b$ inteiros não negativos). Pelos dados do enunciado:
\begin{itemize}
\item total de pizzas: $c+b=205$;
\item total arrecadado (em reais): $45c+40b=8750$.
\end{itemize}

\textbf{Resolução do sistema.} Da primeira equação, $c=205-b$. Substituindo na segunda:
\begin{align*}
45(205-b)+40b &= 8750\\
9225-45b+40b &= 8750\\
-5b &= 8750-9225=-475\\
b &= 95.
\end{align*}
Logo $c=205-95=110$.

\textbf{Interpretação alternativa (mesmo resultado).} Se as 205 pizzas fossem todas de calabresa, a arrecadação seria $45\cdot 205=9225$ reais, isto é, $9225-8750=475$ reais a mais que o real. Cada pizza de calabresa trocada por uma de brigadeiro diminui o total em $45-40=5$ reais; portanto houve $475/5=95$ pizzas de brigadeiro.

\textbf{Verificação.}
\[
110+95=205 \quad\text{e}\quad 45\cdot 110+40\cdot 95=4950+3800=8750 \text{ reais}.
\]
As duas condições são satisfeitas e os valores são inteiros positivos, como devem ser.

\textbf{Análise das alternativas.} Com o total fixo em 205 pizzas, cada pizza de brigadeiro a mais reduz a arrecadação em 5 reais, então só um valor de $b$ serve. Para as outras alternativas: $b=100$ daria $8725$ reais, $b=105$ daria $8700$ reais e $b=110$ daria $8675$ reais, todos diferentes de $8750$ reais. A alternativa d) (110) é o número de pizzas de \emph{calabresa}, uma armadilha comum. Portanto foram vendidas 95 pizzas de brigadeiro, alternativa a).

ANSWER: (a) 95%
\newpage%
\section*{Questão 16}

\section*{Solução}

\textbf{Ideia.} A distância do centro de uma pedra em $(x,y)$ ao centro do alvo $(0,0)$ é $r=\sqrt{x^2+y^2}$. Como $r\geq 0$, comparar distâncias é o mesmo que comparar seus quadrados $r^2=x^2+y^2$, o que evita raízes. Os raios dos círculos do alvo não interferem: a regra usa apenas as distâncias dos centros das pedras.

\textbf{Equipe azul:}
\begin{align*}
(1,2)&: & r^2&=1+4=5,\\
\left(-\tfrac{3}{2},-1\right)&: & r^2&=\tfrac{9}{4}+1=\tfrac{13}{4}=3{,}25,\\
(-2,-3)&: & r^2&=4+9=13,\\
(4,-1)&: & r^2&=16+1=17.
\end{align*}

\textbf{Equipe vermelha:}
\begin{align*}
\left(\tfrac{1}{2},\tfrac{1}{5}\right)&: & r^2&=\tfrac{1}{4}+\tfrac{1}{25}=\tfrac{29}{100}=0{,}29,\\
(2,0)&: & r^2&=4,\\
(-1,1)&: & r^2&=1+1=2,\\
(2,-2)&: & r^2&=4+4=8.
\end{align*}

\textbf{Vencedora da rodada.} A menor distância de todas é a da pedra vermelha em $\left(\tfrac{1}{2},\tfrac{1}{5}\right)$: $r^2=0{,}29$, isto é, $r\approx 0{,}54$. Ela está menos de $1$ unidade do centro, dentro do botão. A pedra azul mais próxima tem $r^2=3{,}25$. Logo, a equipe \textbf{vermelha} vence a rodada.

\textbf{Pontuação.} A pedra adversária (azul) mais próxima do centro é a de $\left(-\tfrac{3}{2},-1\right)$, então
\[
d=\sqrt{\tfrac{13}{4}}=\frac{\sqrt{13}}{2}\approx 1{,}80,\qquad d^2=3{,}25 .
\]
Contam ponto as pedras vermelhas com $r<d$, ou seja, com $r^2<3{,}25$:
\begin{itemize}
\item $\left(\tfrac{1}{2},\tfrac{1}{5}\right)$: $0{,}29<3{,}25$ \quad (conta);
\item $(-1,1)$: $2<3{,}25$ \quad (conta, pois $r=\sqrt{2}\approx 1{,}41$);
\item $(2,0)$: $4>3{,}25$ \quad (não conta, pois $r=2$);
\item $(2,-2)$: $8>3{,}25$ \quad (não conta, pois $r=2\sqrt{2}\approx 2{,}83$).
\end{itemize}
Nenhuma pedra está exatamente à distância $d$, então não há caso de empate a tratar. A equipe vermelha faz, portanto, $2$ pontos.

\textbf{Conferência.} Pondo todas as pedras em ordem crescente de $r^2$, temos: vermelha ($0{,}29$), vermelha ($2$), azul ($3{,}25$), vermelha ($4$), $\ldots$ As duas primeiras são vermelhas e a terceira é azul. Isso confirma que o vermelho vence com exatamente $2$ pontos.

\textbf{Alternativas.} As letras (a) e (b) estão erradas, porque quem vence é o vermelho. A letra (c) também está errada, pois a pedra em $(-1,1)$ fica mais perto do centro que qualquer pedra azul. Resta a letra (d).

ANSWER: (D) Vermelho venceu com 2 pontos.%
\newpage%
\section*{Questão 17}

\section*{Solução}

\textbf{1) Cálculo do determinante.} Fazemos a expansão de Laplace pela 1ª coluna, que tem um zero:
\begin{align*}
\det A(x) &= (2-x)\begin{vmatrix} -3-x & 9\\ 3 & 1-x\end{vmatrix}
-(-7)\begin{vmatrix} 1 & -3\\ 3 & 1-x\end{vmatrix} + 0\\
&= (2-x)\big[(-3-x)(1-x)-27\big] + 7\big[(1-x)+9\big]\\
&= (2-x)(x^2+2x-30) + 7(10-x)\\
&= (-x^3+34x-60) + (70-7x).
\end{align*}
Logo
\[
y=\det A(x) = -x^3+27x+10 .
\]
\emph{Conferência:} para $x=0$ obtemos $\det A(0)=2(-3-27)-1(-7)+(-3)(-21)=-60+7+63=10$, que coincide com o termo independente.

\textbf{2) Grau e coeficiente líder.} $y$ é um polinômio de grau $3$, e não $4$: só o produto da diagonal principal, $(2-x)(-3-x)(1-x)$, contém $x^3$, com coeficiente $(-1)^3=-1$. Isso elimina as alternativas (c) e (d), que mostram quárticas. Como o coeficiente líder é negativo, $y\to+\infty$ quando $x\to-\infty$ e $y\to-\infty$ quando $x\to+\infty$: a curva vem de cima à esquerda e desce à direita. Isso elimina a alternativa (a).

\textbf{3) Extremos e intercepto.}
\[
y'=-3x^2+27=-3(x-3)(x+3)=0 \iff x=\pm 3 .
\]
Temos $y'<0$ para $x<-3$, $y'>0$ para $-3<x<3$ e $y'<0$ para $x>3$. Portanto:
\begin{itemize}
\item mínimo local em $x=-3$ (à esquerda do eixo $y$): $y(-3)=27-81+10=-44<0$, abaixo do eixo $x$;
\item máximo local em $x=3$ (à direita do eixo $y$): $y(3)=-27+81+10=64>0$, acima do eixo $x$;
\item a curva corta o eixo $y$ em $y(0)=10>0$, ou seja, um pouco acima da origem (é pequeno em comparação com os valores dos extremos, $-44$ e $64$).
\end{itemize}

\textbf{4) Raízes.} Como $y(-5)=125-135+10=0$, $x=-5$ é raiz. Dividindo por $(x+5)$:
\[
y=-(x+5)(x^2-5x-2) \quad\Longrightarrow\quad x=-5,\qquad x=\frac{5\pm\sqrt{33}}{2}\approx -0{,}37 \ \text{e}\ 5{,}37 .
\]
São três raízes reais. Uma fica à esquerda do mínimo, outra fica entre o mínimo e a origem (perto da origem, no lado negativo) e a terceira fica à direita do máximo. Isso concorda com a mudança de sinal: $y(-3)<0<y(0)$.

\textbf{Conclusão.} O gráfico é uma cúbica com coeficiente líder negativo. Ela desce até um mínimo abaixo do eixo $x$ à esquerda do eixo $y$, sobe e corta o eixo $y$ ligeiramente acima da origem, atinge um máximo acima do eixo $x$ à direita do eixo $y$ e depois desce, com três raízes reais. Essa é exatamente a descrição da alternativa (b).

ANSWER: (B) $y=\det A(x)=-x^3+27x+10$%
\newpage%
\section*{Questão 18}

\section*{Solução}

\textbf{Hipótese.} Como $p$ é um polinômio mônico de grau 3, ele tem três raízes $r_1, r_2, r_3$ em $\mathbb{C}$, contadas com multiplicidade. O enunciado diz que todas elas são inteiras. Essa é a leitura pretendida: se fossem admitidas raízes não reais, $b$ nem seria limitado. Por exemplo, com $r_1=-n$ e as outras duas raízes conjugadas de soma $0$ e produto $8/n$, teríamos $b=n$ para todo $n$ inteiro positivo.

\textbf{Relações de Girard.} Para $p(x)=x^3+bx^2+cx+8=(x-r_1)(x-r_2)(x-r_3)$:
\[
r_1+r_2+r_3=-b, \qquad r_1r_2+r_1r_3+r_2r_3=c, \qquad r_1r_2r_3=-8 .
\]
Assim, maximizar $b=-(r_1+r_2+r_3)$ equivale a \emph{minimizar a soma} de três inteiros cujo produto é $-8$. Qualquer terna desse tipo gera um $c$ inteiro, logo real, e portanto é admissível.

\textbf{Análise dos casos.} Como o produto é negativo, o número de raízes negativas é ímpar.

\begin{itemize}
\item \emph{Três raízes negativas:} $\{|r_1|,|r_2|,|r_3|\}$ é uma terna de inteiros positivos com produto $8$, ou seja, $\{1,1,8\}$, $\{1,2,4\}$ ou $\{2,2,2\}$. As somas são
\[
-1-1-8=-10,\qquad -1-2-4=-7,\qquad -2-2-2=-6 .
\]
\item \emph{Uma raiz negativa $-m$ e duas positivas $u, v$, com $muv=8$:} a soma é $u+v-m\ge 1+1-8=-6$, pois $u,v\ge 1$ e $m\le 8$. A igualdade ocorre na terna $(-8,1,1)$.
\end{itemize}

A menor soma possível é, portanto, $-10$, obtida com as raízes $-1,-1,-8$. Logo
\[
b_{\max}=-(-10)=10 .
\]

\textbf{Verificação.}
\[
(x+1)^2(x+8)=(x^2+2x+1)(x+8)=x^3+10x^2+17x+8,
\]
de modo que $b=10$, $c=17$ (real) e o termo independente é $8$, como exigido.

\textbf{Alternativas.}
\begin{itemize}
\item O valor $b=13$ exigiria três inteiros de produto $-8$ e soma $-13$, o que não existe, já que a menor soma é $-10$.
\item Os valores $6$ e $7$ são possíveis, com raízes $(-8,1,1)$ e $(-1,-2,-4)$ respectivamente, mas não são os maiores.
\end{itemize}
Portanto, a alternativa correta é a (c).

ANSWER: (C) 10%
\newpage%
\section*{Questão 19}

\section*{Solução}

\textbf{Ideia.} A altura é $A(t)=\dfrac{7}{10}\,\mathrm{sen}\,\theta(t)+\dfrac{13}{10}$, com $\theta(t)=\dfrac{\pi}{12}(2t+3)$. Como $\dfrac{7}{10}>0$, $A$ é uma função crescente de $\mathrm{sen}\,\theta$. Logo $A$ é mínima exatamente quando $\mathrm{sen}\,\theta=-1$, e então
\[
A_{\min}=-\frac{7}{10}+\frac{13}{10}=\frac{6}{10}=0{,}6\ \text{m}.
\]

\textbf{Intervalo de $\theta$.} Para $0\le t\le 24$:
\[
\theta(0)=\frac{3\pi}{12}=\frac{\pi}{4},\qquad \theta(24)=\frac{51\pi}{12}=\frac{17\pi}{4},
\]
e $\theta$ é crescente em $t$. Então $\theta\in\left[\dfrac{\pi}{4},\dfrac{17\pi}{4}\right]$.

\textbf{Mínimos.} Temos $\mathrm{sen}\,\theta=-1 \iff \theta=\dfrac{3\pi}{2}+2k\pi$, com $k\in\mathbb{Z}$. No intervalo acima, só servem
\[
\theta=\frac{3\pi}{2}\quad\text{e}\quad\theta=\frac{7\pi}{2},
\]
porque $\dfrac{11\pi}{2}>\dfrac{17\pi}{4}$ e $-\dfrac{\pi}{2}<\dfrac{\pi}{4}$. Resolvendo para $t$:
\[
\frac{\pi}{12}(2t+3)=\frac{3\pi}{2}\;\Rightarrow\;2t+3=18\;\Rightarrow\;t=7{,}5\ \text{h},
\]
\[
\frac{\pi}{12}(2t+3)=\frac{7\pi}{2}\;\Rightarrow\;2t+3=42\;\Rightarrow\;t=19{,}5\ \text{h}.
\]
Como $0{,}5$ h $=30$ min, os horários são \textbf{07h30} e \textbf{19h30}.

\textbf{Verificações.}
\begin{itemize}
\item Período: $\omega=\dfrac{2\pi}{12}=\dfrac{\pi}{6}$ rad/h, logo $T=\dfrac{2\pi}{\pi/6}=12$ h. É coerente: os dois mínimos estão separados por $19{,}5-7{,}5=12$ h, o que corresponde a duas marés baixas por dia (maré semidiurna).
\item Valor em $t=7{,}5$: $A(7{,}5)=0{,}7\,\mathrm{sen}\!\left(\dfrac{18\pi}{12}\right)+1{,}3=0{,}7\cdot(-1)+1{,}3=0{,}6$ m. As outras alternativas não dão o mínimo. Por exemplo, em $t=7$ temos $\theta=\dfrac{17\pi}{12}$, logo $\mathrm{sen}\,\theta\approx-0{,}97\neq-1$ e $A(7)\approx 0{,}62$ m $>0{,}6$ m. Do mesmo modo, $t=6$ e $t=6{,}5$ dão $\mathrm{sen}\,\theta>-1$.
\item As alturas variam entre $0{,}6$ m e $2{,}0$ m, valores plausíveis para uma maré.
\end{itemize}

\textbf{Conclusão.} A maré está no ponto mais baixo às 07h30 e às 19h30, que é a alternativa (d).

ANSWER: (D) 07h30 e 19h30%
\newpage%
\section*{Questão 20}

\section*{Solução}

\textbf{1. Área de cada retângulo.} O retângulo de índice $n$ ($n=0,1,2,\dots$) fica sobre o intervalo $[n,\,n+1]$, tem base $1$ e altura
\[
f(n)=\frac{1}{3^{n}} .
\]
Portanto, sua área é
\[
R_n = 1\cdot f(n)=\left(\frac{1}{3}\right)^{n}.
\]

\textbf{2. Área de $B$ como soma.} Dois retângulos vizinhos só compartilham um lado vertical, que tem área nula. Por isso a área da união é a soma das áreas:
\[
\text{Área}(B)=\sum_{n=0}^{\infty}\left(\frac{1}{3}\right)^{n}
=1+\frac{1}{3}+\frac{1}{9}+\frac{1}{27}+\cdots
\]

\textbf{3. Série geométrica.} Essa soma é uma série geométrica de primeiro termo $a_0=1$ e razão $q=\dfrac{1}{3}$. Como $|q|<1$, a série converge e sua soma é
\[
S=\frac{a_0}{1-q}=\frac{1}{1-\frac{1}{3}}=\frac{1}{\frac{2}{3}}=\frac{3}{2}.
\]
Logo, $\text{Área}(B)=\dfrac{3}{2}=1{,}5$ (unidades de área).

\textbf{4. Verificações.}
\begin{itemize}
\item \emph{Contenção $A\subset B$:} $f$ é decrescente. Assim, para $x\in[n,\,n+1]$ vale $f(x)\le f(n)$, e a curva fica abaixo do topo de cada retângulo. Então $B$ de fato dá uma estimativa por excesso.
\item \emph{Ordem de grandeza:} a área exata de $A$ é
\[
\int_{0}^{\infty}3^{-x}\,dx=\left[-\frac{3^{-x}}{\ln 3}\right]_{0}^{\infty}=\frac{1}{\ln 3}\approx 0{,}91 .
\]
Ela é de fato menor que $1{,}5$, como deve ser para uma estimativa por excesso.
\item \emph{Somas parciais:} $1+\frac{1}{3}+\frac{1}{9}=\frac{13}{9}\approx 1{,}44$ e $1+\frac{1}{3}+\frac{1}{9}+\frac{1}{27}=\frac{40}{27}\approx 1{,}48$. Os valores crescem e se aproximam de $1{,}5$, o que confirma o limite.
\end{itemize}

\textbf{5. Análise das alternativas.}
\begin{itemize}
\item (a) $\frac{9}{10}=0{,}9$ é menor que a área de $A$ ($\approx 0{,}91$). Não pode ser uma estimativa por excesso.
\item (b) $\frac{3}{9}=\frac{1}{3}$ é apenas a área do segundo retângulo.
\item (d) $\frac{40}{27}$ é a soma só dos quatro primeiros retângulos ($n=0$ a $3$). Faltam os infinitos retângulos restantes.
\item (c) $\frac{3}{2}$ é a soma de todos os retângulos, ou seja, a área de $B$.
\end{itemize}

\textbf{Conclusão:} a área da região $B$ é $\dfrac{3}{2}$, alternativa (c).

ANSWER: (C) $\dfrac{3}{2}$%
\newpage%
\section*{Questão 21}

\section*{Solução}

\textbf{Parametrização da PA.} Como $a,b,c,d$ formam, nessa ordem, uma progressão aritmética de razão $r$ (com $r<0$), escrevemos
\[
b=a+r,\qquad c=a+2r,\qquad d=a+3r .
\]

\textbf{Traço.} O traço é a soma dos elementos da diagonal principal:
\[
\operatorname{tr}M=a+d=a+(a+3r)=2a+3r=-4. \tag{1}
\]

\textbf{Determinante.} Para uma matriz $2\times 2$, $\det M=ad-bc$. Logo
\[
\det M=a(a+3r)-(a+r)(a+2r)=(a^2+3ar)-(a^2+3ar+2r^2)=-2r^2 .
\]
Repare que o determinante só depende da razão, e não de $a$. Com $\det M=-8$:
\[
-2r^2=-8\;\Longrightarrow\; r^2=4\;\Longrightarrow\; r=\pm 2 .
\]
Como a razão é negativa por hipótese, $r=-2$. (O valor $r=2$ também satisfaz a condição do determinante. Por isso o sinal da razão é necessário para que a resposta seja única: com $r=2$ teríamos $a=-5$, que corresponde à alternativa d, uma armadilha.)

\textbf{Cálculo de $a$.} Substituindo $r=-2$ em (1):
\[
2a+3(-2)=-4\;\Longrightarrow\; 2a=2\;\Longrightarrow\; a=1 .
\]

\textbf{Verificação.} A PA é $1,\,-1,\,-3,\,-5$, de razão $-2<0$, e
\[
M=\begin{pmatrix} 1 & -1\\ -3 & -5\end{pmatrix},\qquad
\det M=1\cdot(-5)-(-1)(-3)=-5-3=-8,\qquad
\operatorname{tr}M=1+(-5)=-4 .
\]
As duas condições do enunciado são satisfeitas. Portanto $a=1$, que é a alternativa a.

ANSWER: (A) 1%
\newpage%
\section*{Questão 22}

\section*{Solução}

\textbf{1. Retas que contêm os lados.} Os vértices são $O(0,0)$, $A(1,1)$, $B(0,2)$ e $C(-1,1)$.
\begin{itemize}
\item Lado $OA$: reta $y=x$, com $0\le x\le 1$. Lado $OC$: reta $y=-x$, com $-1\le x\le 0$. Juntos, eles formam o gráfico de $y=|x|$ para $|x|\le 1$ (a ``borda de baixo'').
\item Lado $AB$: reta pelos pontos $(1,1)$ e $(0,2)$, ou seja, $y=2-x$, com $0\le x\le 1$. Lado $CB$: $y=2+x$, com $-1\le x\le 0$. Juntos, eles formam o gráfico de $y=2-|x|$ para $|x|\le 1$ (a ``borda de cima'').
\end{itemize}

\textbf{2. Descrição da região.} Fixado $x$ com $|x|\le 1$, a reta vertical que passa por esse $x$ corta o quadrado no segmento que vai da borda de baixo à de cima. Logo,
\[
(x,y)\ \text{está na região} \iff |x|\le y\le 2-|x| .
\]
Não é preciso impor $|x|\le 1$ à parte. Da própria dupla desigualdade sai $|x|\le 2-|x|$, isto é, $|x|\le 1$. Ela também dá $y\ge |x|\ge 0$. Portanto o conjunto não tem pontos fora da faixa $-1\le x\le 1$.

\textbf{Verificação algébrica.} O quadrado tem centro $(0,1)$ e diagonais paralelas aos eixos, ambas de comprimento $2$ (semidiagonal $1$). Por isso ele é o conjunto $|x|+|y-1|\le 1$. Usando que $|t|\le a \iff -a\le t\le a$ (com $a\ge 0$):
\[
|y-1|\le 1-|x| \iff |x|-1\le y-1\le 1-|x| \iff |x|\le y\le 2-|x| .
\]
Os quatro vértices satisfazem a condição. Por exemplo, para $(1,1)$ temos $1\le 1\le 1$. O ponto interior $(0,1)$ também satisfaz, pois $0\le 1\le 2$.

\textbf{3. Por que as outras alternativas falham.}
\begin{itemize}
\item[(a)] $|x|\ge y\ge 2-|x|$ exige $|x|\ge 2-|x|$, ou seja, $|x|\ge 1$. O ponto $(0,1)$, que está dentro do quadrado, fica de fora. Já o ponto $(2,1)$, que está fora do quadrado, é aceito.
\item[(c)] $|x|\ge 2+|x|$ daria $0\ge 2$, o que é absurdo. O conjunto é vazio.
\item[(d)] $|x|\le y\le 2+|x|$ descreve uma região ilimitada. Por exemplo, o ponto $(5,6)$ satisfaz $5\le 6\le 7$ e não pertence ao quadrado.
\end{itemize}

Assim, a região é $\{(x,y)\in\mathbb{R}^2:\ |x|\le y\le 2-|x|\}$, que é a alternativa (b).

ANSWER: (B) $\{(x,y)\in\mathbb{R}^2: |x|\le y\le 2-|x|\}$%
\newpage%
\section*{Questão 23}

\section*{Solução}

\textbf{Ideia.} Uma função $f:A\to B$ é invertível se, e somente se, é bijetora: injetora (valores distintos de $x$ dão imagens distintas) e sobrejetora ($B=f(A)$). O contradomínio $B$ é livre, então basta tomar $B=f(A)$ e a sobrejetividade fica garantida. A questão se reduz a descobrir em qual intervalo $A$ a função $f$ é \emph{injetora}.

\textbf{Critério de injetividade.} Sejam $a,b\in A$ com $a\neq b$. Então
\[
f(a)-f(b)=(a^2-b^2)-(a-b)=(a-b)(a+b-1).
\]
Como $a-b\neq 0$, temos $f(a)=f(b)$ exatamente quando $a+b=1$. Logo $f$ é injetora em $A$ se, e somente se, $A$ não contém dois pontos distintos simétricos em relação a $x=\tfrac12$, que é a abscissa do vértice. Isso confere com a forma canônica
\[
f(x)=\left(x-\tfrac12\right)^2-\tfrac94 .
\]
O gráfico é uma parábola com concavidade para cima: $f$ é estritamente decrescente em $\left(-\infty,\tfrac12\right]$ e estritamente crescente em $\left[\tfrac12,+\infty\right)$. Assim, um intervalo serve se, e somente se, $\tfrac12$ não for ponto interior dele.

\textbf{Análise das alternativas.}
\begin{itemize}
\item[a)] $[0,2]$: contém $0$ e $1$, com $0+1=1$, e $f(0)=f(1)=-2$. Não é injetora.
\item[b)] $[-1,2]$: contém $-1$ e $2$, com $-1+2=1$, e $f(-1)=f(2)=0$. Não é injetora.
\item[c)] $[1,2]$: para $a,b\in[1,2]$ temos $a+b\ge 2>1$, logo $a+b-1\neq 0$. Então $f(a)=f(b)$ implica $a=b$, e $f$ é injetora. De fato, $[1,2]\subset\left[\tfrac12,+\infty\right)$, onde $f$ é estritamente crescente.
\item[d)] $[-2,1]$: contém $0$ e $1$, e $f(0)=f(1)=-2$. Não é injetora.
\end{itemize}

\textbf{Conferência.} Para $A=[1,2]$, como $f$ é contínua e crescente, temos $f(1)=1-1-2=-2$ e $f(2)=4-2-2=0$, logo $B=f(A)=[-2,0]$. Isolando $x$ em $y=\left(x-\tfrac12\right)^2-\tfrac94$, com $x-\tfrac12\ge 0$, obtemos a inversa
\[
f^{-1}:[-2,0]\to[1,2],\qquad f^{-1}(y)=\frac12+\sqrt{y+\frac94}.
\]
Ela está bem definida, pois $y+\tfrac94\ge \tfrac14>0$. Por exemplo, $f^{-1}(-2)=\tfrac12+\tfrac12=1$ e $f^{-1}(0)=\tfrac12+\tfrac32=2$.

Portanto, entre as alternativas, só $A=[1,2]$ torna $f$ invertível, com $B=[-2,0]$.

ANSWER: (C) $[1,2]$%
\newpage%
\section*{Questão 24}

\section*{Solução}

\textbf{Espaço amostral.} Como a retirada é sem reposição e a ordem não importa para o $\operatorname{mdc}$, os resultados possíveis são os pares $\{a,b\}$ com $a\neq b$ em $\{1,2,\dots,100\}$. Todos são equiprováveis:
\[
|\Omega|=\binom{100}{2}=\frac{100\cdot 99}{2}=4950 .
\]

\textbf{Quando o mdc é par?} Seja $d=\operatorname{mdc}(a,b)\geq 1$. Pela definição de mdc, um inteiro divide $d$ se, e somente se, divide $a$ e $b$ ao mesmo tempo. Com o inteiro $2$:
\[
d \text{ é par} \iff 2\mid d \iff 2\mid a \ \text{e}\ 2\mid b .
\]
Logo o mdc é par exatamente quando os dois números sorteados são pares. Basta que um deles seja ímpar para que o mdc seja ímpar. Por exemplo, $\operatorname{mdc}(3,9)=3$ e $\operatorname{mdc}(4,9)=1$ são ímpares, e $\operatorname{mdc}(4,6)=2$ é par.

\textbf{Evento complementar.} Entre $1$ e $100$ há $50$ números pares. Então
\[
P(\text{mdc par})=\frac{\binom{50}{2}}{\binom{100}{2}}=\frac{1225}{4950}=\frac{49}{198}.
\]
Isso coincide com o cálculo sequencial: a primeira bola é par com probabilidade $\tfrac{50}{100}$; dado isso, a segunda é par com probabilidade $\tfrac{49}{99}$. Assim, $\tfrac{50}{100}\cdot\tfrac{49}{99}=\tfrac{49}{198}$.

\textbf{Resposta.} Pelo complementar,
\[
P(\text{mdc ímpar})=1-\frac{49}{198}=\frac{149}{198}\approx 0{,}75 .
\]

\textbf{Conferência.} O resultado é coerente com uma estimativa rápida: a chance de os dois números serem pares é de cerca de $\tfrac12\cdot\tfrac12=\tfrac14$, e então o mdc é ímpar em cerca de $\tfrac34$ dos casos. Esse valor fica um pouco acima de $\tfrac34$, porque, sem reposição, a segunda bola par é um pouco menos provável ($\tfrac{49}{99}<\tfrac12$). Pela contagem direta, os pares favoráveis são $4950-1225=3725$, e $\tfrac{3725}{4950}=\tfrac{149}{198}$.

\textbf{Alternativas.} A alternativa (d) é a correta, pois é exatamente o valor encontrado. As demais estão erradas:
\begin{itemize}
\item (a) $\tfrac12$ supõe, sem justificativa, que mdc par e mdc ímpar têm a mesma chance.
\item (b) $\tfrac{50}{198}$ fica muito próximo de $\tfrac{49}{198}$, a probabilidade do evento complementar (mdc par).
\item (c) $\tfrac{100}{199}\approx 0{,}50$ fica muito abaixo de $0{,}75$.
\end{itemize}

ANSWER: (D) $\dfrac{149}{198}$%
\newpage%
\section*{Questão 25}

\section*{Solução}

\textbf{1. Limitação de $f(1/2)$.} Como $0<\tfrac12<1$, a hipótese $x^3<f(x)<x^2$ vale para $x=\tfrac12$:
\[
\left(\tfrac12\right)^3<f\!\left(\tfrac12\right)<\left(\tfrac12\right)^2
\quad\Longrightarrow\quad
\frac18<f\!\left(\tfrac12\right)<\frac14 .
\]
Em particular, $f(1/2)<1/2$. Logo $Q$ fica estritamente abaixo de $R$ e o triângulo não é degenerado.

\textbf{2. Área do triângulo.} Os pontos $Q=(1/2,\,f(1/2))$ e $R=(1/2,\,1/2)$ têm a mesma abscissa. Então o lado $QR$ está sobre a reta vertical $x=1/2$ e mede
\[
QR=\frac12-f\!\left(\tfrac12\right)>0 .
\]
Tomando $QR$ como base, a altura relativa a ela é a distância de $P=(0,0)$ à reta $x=1/2$, isto é, $h=\tfrac12$. Portanto
\[
A=\frac{QR\cdot h}{2}=\frac12\cdot\frac12\left(\frac12-f\!\left(\tfrac12\right)\right)=\frac14\left(\frac12-f\!\left(\tfrac12\right)\right).
\]
(Verificação pelo determinante: $A=\tfrac12\left|x_Q\,y_R-x_R\,y_Q\right|=\tfrac12\left|\tfrac14-\tfrac12 f(\tfrac12)\right|$, que dá a mesma expressão.)

\textbf{3. Limitação de $A$.} A partir de $\tfrac18<f(\tfrac12)<\tfrac14$, multiplicamos por $-1$, o que inverte as desigualdades, e somamos $\tfrac12$:
\[
\frac12-\frac14<\frac12-f\!\left(\tfrac12\right)<\frac12-\frac18
\quad\Longrightarrow\quad
\frac14<\frac12-f\!\left(\tfrac12\right)<\frac38 .
\]
Multiplicando por $\tfrac14>0$, as desigualdades se mantêm:
\[
\frac1{16}<A<\frac3{32}.
\]

\textbf{4. Conclusão.} Para qualquer $f$ que satisfaça a hipótese, $A$ está entre $\tfrac1{16}=0{,}0625$ e $\tfrac3{32}=0{,}09375$. Os dois extremos podem ser aproximados tanto quanto se queira, mas nunca são atingidos, porque as desigualdades da hipótese são estritas. Esse intervalo é exatamente o da alternativa (d). As alternativas (a), (b) e (c) dão intervalos contidos em $(0,\,1/16)$, disjuntos de $(1/16,\,3/32)$, e por isso são falsas.

ANSWER: (D) $1/16<A<3/32$%
\newpage%
\section*{Questão 26}

\section*{Solução}

\textbf{Dados.} $OA=OC=OP=r=1$ (raios), $BC=r=1$, $\angle ABO=15^\circ$; $P,O,B$ colineares (com $O$ entre $P$ e $B$) e $A,C,B$ colineares (com $C$ entre $A$ e $B$).

\textbf{1) Triângulo $OCB$.} Como $OC=CB=r$, o triângulo $OCB$ é isósceles de base $OB$, logo
\[
\angle COB=\angle CBO=15^\circ \quad\Longrightarrow\quad \angle OCB=180^\circ-2\cdot 15^\circ=150^\circ .
\]

\textbf{2) Triângulo $OAC$.} Como $A$, $C$, $B$ são colineares, $\angle OCA$ é o suplemento de $\angle OCB$:
\[
\angle OCA=180^\circ-150^\circ=30^\circ .
\]
Como $OA=OC=r$, o triângulo $OAC$ também é isósceles, então $\angle OAC=\angle OCA=30^\circ$.

\textbf{3) Ângulo $\angle AOP$.} No triângulo $AOB$ os ângulos em $A$ e em $B$ valem $30^\circ$ e $15^\circ$. Como $P$, $O$, $B$ são colineares, $\angle AOP$ é ângulo externo desse triângulo em $O$ e, pelo teorema do ângulo externo,
\[
\angle AOP=\angle OAB+\angle OBA=30^\circ+15^\circ=45^\circ .
\]
(É o resultado clássico: o ângulo central $\angle AOP$ vale o triplo de $\angle ABO$.)

\textbf{4) Lei dos cossenos no triângulo $AOP$.} Com $OA=OP=1$:
\[
AP^2=OA^2+OP^2-2\cdot OA\cdot OP\cos 45^\circ=1+1-2\cdot\frac{\sqrt2}{2}=2-\sqrt2,
\]
\[
AP=\sqrt{2-\sqrt2}\approx 0{,}765 .
\]

\textbf{Conferência.} Com $O=(0,0)$ temos $OB=2\cos 15^\circ\approx 1{,}932$ (no isósceles $OCB$), $A=(\cos 135^\circ,\ \operatorname{sen} 135^\circ)\approx(-0{,}707;\ 0{,}707)$ e $P=(-1,0)$. A reta por $B$ com inclinação de $15^\circ$ passa, em $x=-0{,}707$, por $y=\operatorname{tg}15^\circ\,(1{,}932+0{,}707)\approx 0{,}707$, ou seja, passa por $A$, como deveria. Além disso, $AP=\sqrt{0{,}293^2+0{,}707^2}\approx 0{,}765$. O valor é coerente com uma corda de ângulo central $45^\circ$: é menor que o raio, pois $45^\circ<60^\circ$.

\textbf{Alternativas.} Só a alternativa (a) vale $\sqrt{2-\sqrt2}\approx 0{,}765$. As outras dão aproximadamente $1{,}848$ (b), $1{,}612$ (c) e $2{,}298$ (d); este último valor é maior que o diâmetro $2$, o que é impossível para uma corda.

ANSWER: (a) $\sqrt{2-\sqrt{2}}$%
\newpage%
\section*{Questão 27}

\section*{Solução}

\textbf{1. O que os textos 1 e 2 mostram.} No texto 1, o folião está sem camisa e sem fantasia convencional: o que se destaca é a frase ``Não mate sua mina'', escrita nas costas. A legenda ``Quando o óbvio precisa ser escancarado'' deixa claro que se trata de uma mensagem \emph{real} e \emph{óbvia}: um alerta contra o feminicídio, que é uma norma moral elementar. No texto 2, o ilustrador repete a cena num tapete vermelho, entre confetes, com o título \emph{``Fantasia'' NOTA: 10}. Esse título retoma o jargão do julgamento dos desfiles de carnaval, em que o quesito ``fantasia'' é a roupa dos foliões.

\textbf{2. Qual acepção do texto 3 é usada.} O contexto do carnaval e a nota de quesito de desfile apontam para a acepção 4, ``traje que representa objetos, figuras históricas, imaginárias etc.''. As aspas em ``Fantasia'' mostram que o uso não é literal. O folião não usa um traje de verdade: a ``roupa'' dele é a própria frase escrita no corpo. Portanto, \emph{fantasia} tem aqui sentido figurado: o folião ``veste'' uma mensagem. Essa mensagem é uma ordem no imperativo (``Não mate''), com valor de dever moral. A nota 10 premia exatamente essa ``fantasia''.

\textbf{3. Análise das alternativas.}
\begin{itemize}
\item[a)] Incorreta. Usa a acepção 1 (capacidade de criar pela imaginação) e fala em ``situação utópica''. Não há utopia: a frase trata de uma violência real e de uma exigência mínima, ``o óbvio''.
\item[b)] Incorreta. Usa a acepção 2 (obra criada pela imaginação) e descreve o apelo como ``subjetivo e imaginário''. O apelo, porém, é concreto e de validade coletiva, não algo imaginado.
\item[c)] Incorreta. Usa a acepção 3 (coisa ideal ou ficcional) e fala em ``elemento ficcional''. A mensagem não é ficção; a graça da charge está justamente em tratar como fantasia algo que é real e literal.
\item[d)] Correta. Retoma a acepção 4 (traje). Indica o uso figurado (``em termos figurativos''): a frase nas costas faz as vezes de fantasia. Também identifica corretamente o conteúdo, um imperativo moral (``Não mate sua mina'').
\end{itemize}

\textbf{Conclusão.} O texto 2 chama de ``fantasia'' nota 10 a frase que o folião ``vestiu''. É a acepção de traje usada em sentido figurado, aplicada a uma mensagem que expressa um imperativo moral.

ANSWER: (D) trajado, em termos figurativos, algo que se apresenta como um imperativo moral.%
\newpage%
\section*{Questão 28}

\section*{Solução}

\textbf{O sentido de ``desengano''.} \emph{Desengano} é a perda de uma ilusão, isto é, a desilusão de quem percebe que aquilo em que acreditava não dura ou não era real. A canção liga essa palavra ao carnaval no início das duas primeiras estrofes (``Carnaval, desengano''). É preciso então ver, no próprio texto, que ilusão é desfeita e em que momento.

\textbf{Leitura das estrofes.}
\begin{itemize}
  \item 1ª estrofe: ``Deixei a dor em casa me esperando''. A dor não acaba; fica apenas adiada e espera o fim da festa. ``E brinquei e gritei e fui vestido de rei'': durante o carnaval o eu lírico vive uma fantasia, no sentido literal (a roupa de rei) e no figurado (deixa de ser quem é).
  \item O refrão ``Quarta-feira sempre desce o pano'' usa uma metáfora teatral. Na Quarta-feira de Cinzas, o pano (a cortina) desce: o espetáculo termina e a vida comum volta. O advérbio ``sempre'' mostra que esse fim é inevitável e se repete a cada ano.
  \item 2ª estrofe: o amor vivido na festa (``Mão na mão, pé no chão'') também se desfaz: ``hoje nem lembra não''. É mais um exemplo de ilusão desfeita quando o carnaval acaba.
\end{itemize}
O desengano, portanto, não está na festa. Está na passagem da fantasia de carnaval para a realidade da quarta-feira, quando a dor que ficou ``em casa'' volta a ser sentida.

\textbf{Análise das alternativas.}
\begin{itemize}
  \item[a)] Incorreta. A confraternização (``tomar a mão de cada irmão pela cidade'') aparece nas duas últimas estrofes e está ligada à \emph{esperança} (``No carnaval, esperança''), não ao desengano.
  \item[b)] Correta. A desilusão vem do fim da festa (``Quarta-feira sempre desce o pano''), que faz o eu lírico voltar à realidade: a dor que o esperava em casa e o amor que já não se lembra dele.
  \item[c)] Incorreta. O que a canção mostra é um encontro passageiro, que a morena já esqueceu, e não ``um grande amor''. Além disso, o desengano vem do fim desse encontro, e não do fato de vivê-lo.
  \item[d)] Incorreta. Durante a festa a dor fica suspensa (``Deixei a dor em casa''), e não aumentada. Ela só volta quando o carnaval termina.
\end{itemize}

\textbf{Conclusão.} Carnaval e desengano aparecem juntos porque a alegria e a fantasia da festa são passageiras. Na Quarta-feira de Cinzas a ilusão acaba e a realidade volta.

ANSWER: (B) ao retorno à realidade com o término da festa.%
\newpage%
\section*{Questão 29}

\section*{Solução}

\textbf{1. O que o texto de Módolo e Braga estabelece.} Segundo a leitura de Agostinho, o tempo é vivido interiormente em três modos: o \emph{passado} está na \emph{memória}, o \emph{presente} na \emph{atenção} e o \emph{futuro} na \emph{expectativa}. Por isso, os tempos verbais não são apenas etiquetas cronológicas. Eles mostram \emph{como o falante se posiciona} diante do que já foi, do que é e do que ainda não é. A questão pede, portanto, as formas verbais que exprimem um desejo ou uma esperança voltados para algo ainda não realizado, ou seja, a \emph{expectativa}.

\textbf{2. Onde está a expectativa na canção.} A última estrofe começa com ``No carnaval, \emph{esperança}'', palavra que anuncia o desejo. Em seguida vêm três orações dependentes desse substantivo (subentende-se ``esperança de que\ldots''):
\begin{itemize}
  \item ``Que gente longe \emph{viva} na lembrança'' (verso 12);
  \item ``Que gente triste \emph{possa} entrar na dança'' (verso 13);
  \item ``Que gente grande \emph{saiba} ser criança'' (verso 14).
\end{itemize}
As formas \emph{viva}, \emph{possa} e \emph{saiba} estão no \textbf{presente do subjuntivo}, que é o modo do desejo, da hipótese e do que é apenas possível. O eu lírico não afirma fatos. Ele projeta o que \emph{ainda não é} e deseja que venha a ser. Isso é exatamente a ``expectativa'' de Agostinho.

\textbf{3. Por que as outras alternativas estão erradas.}
\begin{itemize}
  \item[(a)] \emph{Brinquei}, \emph{gritei} e \emph{fui} estão no pretérito perfeito do indicativo. Narram ações já concluídas durante a festa, isto é, pertencem à \emph{memória} (``o que já foi''), e não à expectativa.
  \item[(b)] \emph{Esperando} e \emph{sonhando} são gerúndios, formas nominais que indicam ação em curso, sem tempo próprio. Apesar do sentido das palavras, em ``a dor [\ldots] me esperando'' e ``me deixou sonhando'' a ação acompanha fatos passados (\emph{deixei}, \emph{deixou}). O gerúndio não marca a projeção para o que ainda não é.
  \item[(c)] \emph{Desce} e \emph{lembra} estão no presente do indicativo, modo da constatação. ``Quarta-feira sempre desce o pano'' é um presente habitual, uma verdade que se repete. ``Hoje nem lembra não'' constata um fato atual. As duas formas ligam-se à \emph{atenção} ao que é, e não à expectativa.
\end{itemize}

\textbf{4. Conclusão.} Só o presente do subjuntivo de \emph{viva}, \emph{possa} e \emph{saiba} (versos 12 a 14) expressa o desejo voltado para o que ``ainda não é''. A estrofe é introduzida pela palavra \emph{esperança}, o que confirma essa leitura. Logo, a resposta é a alternativa (d).

ANSWER: (D) presente do subjuntivo (\emph{viva}, \emph{possa} e \emph{saiba}, nos versos 12, 13 e 14)%
\newpage%
\section*{Questão 30}

\section*{Solução}

\textbf{Conferência do enunciado.} A transcrição coincide com o original (página 9 do caderno): trecho adaptado de L. Karnal sobre a peça \emph{Para Sempre}, seguido da pergunta sobre o sentido de \emph{louvor-prisão}.

\textbf{1. Formação da palavra.} \emph{Louvor-prisão} é um substantivo composto por justaposição que une dois termos de sentidos opostos:
\begin{itemize}
  \item \emph{louvor}: elogio, enaltecimento (a forma do enunciado: ``Você passa bem café'');
  \item \emph{prisão}: confinamento, limitação (o efeito: prender Tony a um lugar inferior).
\end{itemize}
É, portanto, um neologismo de valor paradoxal, próximo do oxímoro: algo que se apresenta como elogio, mas atua como aprisionamento.

\textbf{2. O que o texto diz.} Logo após a expressão, o autor a explica com um aposto: ``um rebaixamento naquele universo onde o valor era dado pelos livros lidos/escritos''. Antes, já tinha dito que o elogio vinha ``com uma etiqueta de gênero (`você é o de menor importância, o doméstico, o que faz café')''. Ao final, resume a mesma ideia com fórmulas equivalentes: ``agredir com um elogio, preconceito com açúcar, ataque com punho de renda''. Assim, o \emph{louvor-prisão} é um preconceito (rebaixamento social, cultural e de gênero) disfarçado de elogio.

\textbf{3. Análise das alternativas.}
\begin{itemize}
  \item[a)] \textbf{Correta.} ``Postura discriminatória'' corresponde a ``rebaixamento'', ``preconceito'' e ``etiqueta de gênero''; ``disfarçada de enaltecimento'' corresponde a ``louvor'', ``com açúcar'' e ``com punho de renda''.
  \item[b)] Incorreta. A ofensa explícita é o ``veado'' dito pela vizinha da periferia, ou seja, o ``ataque direto e virulento''. O texto o opõe justamente ao louvor-prisão, que é sutil.
  \item[c)] Incorreta. Não há empatia no elogio: o texto diz que ele era ``mais agressivo'' que o insulto direto.
  \item[d)] Incorreta. A alternativa descreve a reação de quem recebe o elogio, e não o elogio em si. Além disso, o texto diz que diante do segundo tipo de ataque não cabe defesa: ``talvez, só possamos sorrir de forma amarela''. A resposta defensiva é a que se dá ao primeiro insulto.
\end{itemize}

\textbf{Conclusão.} \emph{Louvor-prisão} designa um elogio que, sob aparência de enaltecimento, prende o outro a uma posição inferior. É, portanto, uma postura discriminatória disfarçada de elogio.

ANSWER: (A) postura discriminatória disfarçada de enaltecimento.%
\newpage%
\section*{Questão 31}

\section*{Solução}

\textbf{O que se pede.} A questão pede uma característica que \emph{os dois} textos tenham em comum. Por isso, cada alternativa precisa valer para o Texto 1 e também para o Texto 2.

\textbf{Tese e recursos de cada texto.}
\begin{itemize}
  \item \textbf{Texto 1 (Carriello):} defende que ``o brasileiro trabalha menos do que seria esperado'' e atribui o desvio a ``uma preferência por maior quantidade de lazer''. Para isso usa dados estatísticos: uma pesquisa com 160 países, a média mundial de 42,7 horas semanais contra 40,1 horas no Brasil, e as relações entre horas trabalhadas, renda \emph{per capita}, impostos e transferências.
  \item \textbf{Texto 2 (Rangel):} defende que essa conclusão ``é um artefato da especificação econométrica'' e que, com controles institucionais, ``o Brasil passa a trabalhar mais do que o esperado''. Para isso também usa estatística: a mesma base de 160 países, o resíduo do Brasil em diferentes formas funcionais e métodos de ponderação, microdados da World Values Survey para 58 países e o teste de significância estatística.
\end{itemize}
Os dois sustentam um ponto de vista (opostos entre si) apoiados em dados e métodos estatísticos.

\textbf{Análise das alternativas.}
\begin{itemize}
  \item[a)] \textbf{Incorreta.} O Texto 1 não critica aspectos econômicos: ele usa a renda, os impostos e as transferências apenas como referência. Quanto ao aspecto cultural, ele o apresenta como explicação, não o critica. O Texto 2 também não critica aspectos culturais ou econômicos. O que ele critica é a \emph{metodologia} da análise.
  \item[b)] \textbf{Incorreta.} Só o Texto 2 questiona o argumento divulgado pela mídia, e ele vai além de relativizar: chega a inverter a conclusão. O Texto 1 é justamente a matéria de jornal que divulga e reforça esse argumento.
  \item[c)] \textbf{Correta.} Como mostrado acima, os dois textos usam informações estatísticas para defender sua tese: o Texto 1 usa médias de horas e comparações entre países; o Texto 2 usa resíduos de regressão, controles e significância estatística.
  \item[d)] \textbf{Incorreta.} A divergência existe \emph{entre} os dois textos, mas nenhum deles, isoladamente, apresenta divergências entre especialistas. O Texto 1 expõe uma única leitura dos dados. O Texto 2 contesta uma afirmação ``amplamente divulgada na mídia'', não uma disputa entre especialistas.
\end{itemize}

\textbf{Conclusão.} Os dois textos têm em comum o uso de informações estatísticas na defesa de um ponto de vista.

ANSWER: (C) recorrerem a informações estatísticas na defesa de um ponto de vista.%
\newpage%
\section*{Questão 32}

\section*{Solução}

\textbf{1. O que o texto 1 conclui.} Com base no banco de dados de Gethin e Saez (160 países), o texto 1 afirma que o brasileiro trabalha em média 40,1 horas semanais, contra 42,7 horas da média mundial. Daí tira duas conclusões:
\begin{itemize}
\item mesmo levando em conta as características demográficas, a renda \emph{per capita}, os impostos e as transferências, ``o brasileiro trabalha menos do que seria esperado'';
\item esse desvio provavelmente se explica por ``uma preferência por maior quantidade de lazer'', isto é, por uma questão cultural.
\end{itemize}

\textbf{2. O que a metodologia do texto 2 faz.} O autor do texto 2 usa \emph{o mesmo banco de dados} (``Usando os dados de Gethin e Saez (2025)'') e altera a análise:
\begin{itemize}
\item testa outras formas funcionais e outros métodos de ponderação e, sobretudo, inclui controles institucionais (tributação, informalidade e regulação de jornada). Com esses controles, ``o Brasil passa a trabalhar \emph{mais} do que o esperado'': o sinal do resíduo brasileiro se inverte;
\item testa a hipótese cultural com microdados da World Values Survey (58 países). Ela é \emph{rejeitada}, porque a variável que mede a disposição a priorizar trabalho sobre lazer deixa de ser estatisticamente significativa.
\end{itemize}
O texto 2 conclui, então, que o resultado do texto 1 é ``um artefato da especificação econométrica''.

\textbf{3. Análise das alternativas.}
\begin{itemize}
\item[a)] \emph{Falsa.} O texto 2 não acrescenta dados que confirmem ou ampliem os resultados do texto 1. Ele os contesta e mostra que o sinal do desvio brasileiro se inverte.
\item[b)] \emph{Verdadeira.} A conclusão do texto 1 é que o brasileiro trabalha menos do que o esperado. Com os controles institucionais, a metodologia do texto 2 mostra que ele trabalha mais do que o esperado. A conclusão é revertida, e a explicação cultural que a sustentava é rejeitada.
\item[c)] \emph{Falsa.} O banco de dados não é invalidado: o próprio texto 2 o utiliza. O que se critica é a especificação do modelo aplicado a esses dados, e não os dados em si.
\item[d)] \emph{Falsa.} A hipótese do texto 1 (preferência cultural por lazer) é testada e \emph{rejeitada}. Portanto, não é corroborada.
\end{itemize}

\textbf{Conclusão.} A metodologia do texto 2 usa os mesmos dados do texto 1 e, com a inclusão de controles institucionais, inverte o sinal do desvio brasileiro, além de rejeitar a hipótese cultural. Ela permite, portanto, reverter a conclusão do texto 1.

\medskip
ANSWER: (B) reverter a conclusão à qual chega o texto 1.%
\newpage%
\section*{Questão 33}

\section*{Solução}

\textbf{1. Quem fala e de onde fala.} O narrador é Brás Cubas, o ``defunto autor'', que conta a própria vida depois de morto. Por isso fala com franqueza cínica e sem remorso: relata o que fez sem pedir desculpas. O trecho pertence ao capítulo ``O menino é o pai do homem''. Nele, Brás descreve a infância de ``menino diabo'', mimado pelos pais e acostumado a ver todos os seus caprichos atendidos. O título já indica a tese do capítulo: o adulto (o homem da elite senhorial) forma-se nessa infância de mando.

\textbf{2. O que o trecho mostra.} Ao ler o texto, encontramos três marcas da relação senhor--escravizado:
\begin{itemize}
\item \emph{Animalização}: Prudêncio é ``o meu cavalo de todos os dias''. Põe ``as mãos no chão'', recebe um cordel ``à guisa de freio'' e é chamado de ``besta''.
\item \emph{Violência física}: o menino ``fustigava-o'' com a varinha, e o escravizado obedecia ``algumas vezes gemendo''.
\item \emph{Silenciamento}: Prudêncio ``obedecia sem dizer palavra''. O único lamento possível, ``ai, nhonhô!'', é logo calado pelo ``Cala a boca, besta!''. O próprio vocativo ``nhonhô'' (forma de tratamento do escravizado para o filho do senhor) confirma a hierarquia.
\end{itemize}
A ``brincadeira'' é, portanto, a forma que o poder senhorial assume no cotidiano. Uma criança já exerce sobre o subalternizado o domínio que aprendeu do seu meio. O narrador conta isso com naturalidade, e é a ironia machadiana que expõe ao leitor a brutalidade naturalizada da sociedade escravista. A obra reforça essa leitura no capítulo ``O vergalho''. Ali, Prudêncio, já liberto, chicoteia um escravizado que comprou e repete a fórmula ``Cala a boca, besta!''. A hierarquia, portanto, se \emph{reproduz}; não se inverte nem se atenua.

\textbf{3. Análise das alternativas.}
\begin{itemize}
\item[(a)] Falsa. O caráter lúdico existe só para Brás. Para Prudêncio há humilhação e dor (``gemendo''). O jogo não atenua a hierarquia: é o meio pelo qual ela se exerce.
\item[(b)] Falsa. O fato de o poder ser exercido por uma criança não enfraquece o quadro social. Ao contrário, mostra que a dominação é aprendida desde cedo e está naturalizada. Além disso, o narrador não denuncia nada: ele relata com indiferença.
\item[(c)] Verdadeira. O trecho \emph{caracteriza} as relações hierárquicas do Brasil escravista: o poder senhorial, ainda que nas mãos de um menino, exerce-se sobre o subalternizado, que é animalizado, agredido e silenciado.
\item[(d)] Falsa. O ``ai, nhonhô!'' é lamento de quem se submete, e é imediatamente reprimido. Não há inversão de papéis: Prudêncio ``obedecia''. Mais tarde, em ``O vergalho'', ele reproduz a mesma lógica, sem invertê-la.
\end{itemize}

\textbf{Conclusão.} Do ponto de vista do narrador, a cena revela o funcionamento ``normal'' da hierarquia senhorial sobre os subalternizados. A resposta é a alternativa (c).

ANSWER: (C) caracterização das relações hierárquicas, pois o poder senhorial se exerce sobre os subalternizados.%
\newpage%
\section*{Questão 34}

\section*{Solução}

\textbf{1. O que é o trecho.} A ``advertência'' abre \emph{Vida e morte de M. J. Gonzaga de Sá} (1919). Nela, quem fala se apresenta como colega de escola e de ofício de Augusto Machado e diz ter apenas revisado e providenciado a publicação de um livro escrito por ele. Fora do livro, porém, Augusto Machado não existe. Ele é o narrador do romance: o amigo de Gonzaga de Sá que conta a vida e a morte do protagonista e que aparece o tempo todo na narrativa.

\textbf{2. A função do recurso.} Trata-se do conhecido artifício do \emph{autor fictício} ou \emph{editor fictício}, em que o escritor real finge ser apenas o editor de um manuscrito alheio. Machado de Assis faz o mesmo com o Conselheiro Aires em \emph{Esaú e Jacó} e no \emph{Memorial de Aires}. Esse paratexto, portanto, já faz parte da ficção, por três razões:
\begin{itemize}
\item atribui a autoria a um personagem, Augusto Machado, que é o narrador-personagem da obra;
\item cria um efeito de verossimilhança, pois a história se apresenta como relato ``real'' de um amigo;
\item antecipa, com ironia, o caráter híbrido do livro. Os ``defeitos'' apontados (faltam datas e exatidão, e ``a pessoa do autor'' aparece e sobressai demais) são justamente o que faz da obra um romance, e não uma biografia: o narrador subjetivo, que ``trata mais dele do que do seu herói'', faz parte do projeto literário.
\end{itemize}

\textbf{3. Análise das alternativas.}
\begin{itemize}
\item[a)] \textbf{Falsa.} O trecho não é um paratexto editorial informativo. É um recurso ficcional, e o verdadeiro autor é Lima Barreto, não Augusto Machado, que é uma personagem.
\item[b)] \textbf{Falsa.} O ``editor'' não critica os defeitos: ele os chama de ``insignificantes'' e diz que ``eram de desprezar'' diante dos méritos do livro. Além disso, a obra não valoriza a discrição e a objetividade do narrador; o próprio trecho diz que o autor-narrador aparece e sobressai demais.
\item[c)] \textbf{Falsa.} O Modernismo não separou os gêneros literários; ao contrário, ele os misturou e rompeu suas fronteiras. Lima Barreto, autor pré-modernista, também não defende essa separação: o texto mostra justamente uma obra que fica entre a biografia e o romance.
\item[d)] \textbf{Verdadeira.} A advertência ajuda a construir a ficcionalidade da obra. Augusto Machado é o narrador-personagem e o autor fictício da narrativa, e Lima Barreto, fazendo-se de simples editor, atribui a ele a autoria.
\end{itemize}

\textbf{Conclusão.} O trecho é um paratexto ficcional: Lima Barreto finge ser o editor e atribui a autoria ao narrador-personagem Augusto Machado.

ANSWER: (D)%
\newpage%
\section*{Questão 35}

\section*{Solução}

\textbf{1. Leitura do poema.} O título ``À Bengala'' indica uma \emph{ode} (o livro chama-se \emph{Odes Mínimas}), gênero de exaltação dirigido a alguém ou a algo. O eu-lírico fala diretamente com o objeto (apóstrofe: ``\emph{Contigo} me faço''), e o poema inteiro é um único período, sintaticamente completo:
\begin{center}
``Contigo me faço / Pastor do rebanho / De meus próprios passos.''
\end{center}

\textbf{2. A metáfora central.} Os \emph{passos} do eu-lírico viram um \emph{rebanho}, e ele vira o \emph{pastor} desse rebanho. A palavra \emph{cajado} não aparece, mas está implícita: o pastor conduz o rebanho com o cajado, e a bengala tem a mesma forma (um bastão em que se apoia quem caminha). Assim, pela \textbf{aproximação semântica bengala $\leftrightarrow$ cajado}, o objeto ganha outro sentido. Em geral a bengala sugere velhice, limitação e dependência. Aqui ela vira instrumento de quem \emph{conduz}: com ela, o eu-lírico ``se faz'' (torna-se ativamente) senhor e guia dos próprios passos. O que poderia ser sinal de fragilidade aparece como \textbf{protagonismo}, ou seja, como autonomia para comandar o próprio caminhar.

\textbf{3. Análise das alternativas.}
\begin{itemize}
\item[a)] Incorreta. A comparação entre rebanho e passos não é ``difícil''. Ela é justamente a metáfora que sustenta o poema, e dela não sai nenhuma ``fantasia''. A bengala não é extensão da imaginação: ela dá ao eu-lírico o papel de condutor.
\item[b)] Incorreta. O texto não fala de trauma. Também não existem ``passos do rebanho'' comparados aos do eu-lírico: são os próprios passos do eu-lírico que \emph{formam} o rebanho.
\item[c)] \textbf{Correta.} A bengala expressa o protagonismo de quem conduz (o ``pastor''), e esse sentido só aparece porque bengala e cajado são semanticamente próximos (cajado é o bastão do pastor).
\item[d)] Incorreta. O tom é de exaltação (ode), não de lamento. Também não há ``dissolução sintática'': a frase é íntegra, e a ligação entre bengala e cajado acontece no plano do sentido, não da sintaxe.
\end{itemize}

\textbf{Conclusão.} A bengala, lida como um cajado, transforma o eu-lírico em pastor dos próprios passos, isto é, em protagonista do seu caminhar.

ANSWER: (C) expressam o protagonismo daquele que conduz, condição recuperada pela aproximação semântica entre as palavras bengala e cajado.%
\newpage%
\section*{Questão 36}

\section*{Solução}

\textbf{Conferência do enunciado.} A transcrição coincide com o original (página 11 do caderno, página 13 do PDF). A questão pede o que há em comum entre a fala da autora (texto 1), as memórias da narradora de \textit{Olhos d'água} (texto 2) e o narrador em terceira pessoa de \textit{Zaíta} (texto 3).

\textbf{Passo 1: o que diz o texto 1.} Conceição Evaristo afirma ser \emph{comprometida} e \emph{cumpliciada} com suas personagens. Diz também que escrever ``Zaíta'' é ``um processo doloroso''. A dor aparece, portanto, no \emph{processo de criação}. Ela é o modo como a autora vive o material que transforma em ficção, e não um assunto narrado nos contos. É a ideia de \emph{escrevivência}: a experiência vivida e sentida vira escrita.

\textbf{Passo 2: os contos (leitura integral).}
\begin{itemize}
\item \textit{Olhos d'água}: a narradora é a filha mais velha de sete e teve uma ``breve adolescência''. Ela relembra a infância pobre e a fome, que a mãe disfarçava com brincadeiras. Também lembra os silêncios da mãe nas horas difíceis. No fim, ela descobre que os olhos da mãe são ``olhos d'água'', isto é, de lágrimas. O conto trata da infância e da maternidade sob a pressão da miséria.
\item \textit{Zaíta}: Benícia é mãe de vários filhos na favela. Um dos filhos está envolvido com o tráfico, e a mãe ``recomenda o silêncio'' com a voz trêmula. No desfecho, a menina Zaíta morre atingida por uma bala num tiroteio. Novamente aparecem a infância e a maternidade sob pressão, agora a da violência urbana.
\end{itemize}

\textbf{Passo 3: análise das alternativas.}
\begin{itemize}
\item[(a)] Incorreta. Escrever sobre a infância difícil não ``diminui'' nada dentro das narrativas. Além disso, o tema não é o abandono familiar: nos dois contos a mãe está presente e luta pelos filhos.
\item[(b)] Incorreta. O silêncio familiar não diminui. Ele continua sendo marca das narrativas: os silêncios ``decifrados'' da mãe em \textit{Olhos d'água} e o silêncio recomendado por Benícia em \textit{Zaíta}. A cumplicidade da autora não altera esses fatos.
\item[(c)] Incorreta. A dor da escritora pertence ao ato de escrever (texto 1) e não é tema dos contos. Os temas são a infância, a maternidade, a pobreza e a violência. Além disso, \textit{Zaíta} não se limita à relação entre mãe e filhas, já que envolve também os filhos homens.
\item[(d)] Correta. A dor da autora diante da infância vulnerável e da maternidade pressionada pela miséria e pela violência é o que faz desses temas \emph{matéria literária}. Isso se vê nas memórias em primeira pessoa (texto 2) e na narração em terceira pessoa (texto 3). Essa leitura é coerente com a fala sobre o comprometimento e o ``processo doloroso'' de escrever.
\end{itemize}

\textbf{Conclusão.} Os três textos se articulam porque a experiência dolorosa da escritora diante da infância e da maternidade sob pressão é elaborada esteticamente e vira o assunto dos contos.

ANSWER: (D) a dor da escritora diante da infância e da maternidade sob pressão transforma esses temas em matéria literária dos contos.%
\newpage%
\section*{Questão 37}

\section*{Solução}

\textbf{1. Contexto do conto.} ``A historiadora obstinada'' narra três gerações de uma família igbo. A avó, Nwamgba, vê o filho ser levado pelos missionários: ele se converte, passa a se chamar Michael e rompe com os costumes do seu povo. Na neta, Nwamgba reconhece o espírito do marido morto e lhe dá o nome igbo \emph{Afamefuna}, ``meu nome não se perderá''. A menina, porém, é registrada e criada com o nome cristão \emph{Grace}. O trecho dado pertence ao desfecho, em que o narrador antecipa a vida adulta de Grace.

\textbf{2. Recursos do excerto.}
\begin{itemize}
\item \emph{Anáfora} ``Foi Grace quem\ldots'' combinada com o futuro do pretérito (\emph{largaria}, \emph{refletiria}, \emph{veria}). É uma prolepse: o narrador resume, de forma enfática, o destino da personagem, que se torna a historiadora ``obstinada'' do título.
\item \emph{Episódio do sr. Gboyega}: um africano formado pela tradição colonial acha absurdo que a história africana seja uma disciplina. Ao refletir sobre isso, Grace vê ``a ligação entre educação e dignidade''. A educação ocidental aparece, portanto, como algo que pode apagar a identidade africana. Grace troca a química pela história justamente para recuperar essa história.
\item \emph{Desfecho}: no fim da vida, já com ``seus prêmios'' e o ``jardim de rosas'', ela se sente ``distante de suas raízes'' e troca oficialmente \emph{Grace} (nome cristão, colonial) por \emph{Afamefuna}, o nome dado pela avó. Assim, realiza o desejo de Nwamgba de que seu nome e sua linhagem não se perdessem.
\end{itemize}

\textbf{3. Análise das alternativas.}

\textbf{(a) Falsa.} A anáfora de fato realça o protagonismo de Grace, mas esse protagonismo não acompanha a ocidentalização. Ele vai \emph{contra} ela: Grace dedica a vida a recuperar a história africana e termina rejeitando o nome cristão.

\textbf{(b) Verdadeira.} Ao adotar o nome \emph{Afamefuna}, que a avó lhe deu, Grace cumpre o que Nwamgba esperava (``meu nome não se perderá''). Com isso, ela se reconecta com os valores da tradição igbo anterior à colonização, dos quais o pai se afastara ao se converter.

\textbf{(c) Falsa.} Esta alternativa inverte o sentido do gesto. A troca do nome não exalta a educação ocidental nem trata o passado africano como ``atraso''. Ao contrário, o texto critica essa visão colonial, encarnada pelo sr. Gboyega, e mostra que o sucesso social não preencheu o vazio de estar ``distante de suas raízes''.

\textbf{(d) Falsa.} A anáfora não ``antecipa a tristeza''. Ela funciona como uma síntese enfática da trajetória da personagem. A tristeza aparece em um único episódio. Além disso, o conflito apresentado não opõe ``avanço da educação'' a ``ignorância'': o problema é a própria educação colonial, que desvaloriza a história africana.

\textbf{4. Conclusão.} Só a alternativa (b) interpreta corretamente o desfecho: a troca de nome é um retorno às raízes e aos valores da avó, ligados à tradição anterior à colonização.

ANSWER: (B)%
\newpage%
\section*{Questão 38}

\section*{Solução}

\textbf{1. Contexto do conto.} Em \emph{A viúva Montiel}, García Márquez narra a história de José Montiel, que enriqueceu durante a perseguição política promovida pelo alcaide: os opositores pobres eram assassinados e os ricos, expulsos. Montiel comprava as terras e o gado dos expulsos ``por um preço que ele mesmo se encarregava de fixar'' e assim se tornou, em poucos anos, o homem mais rico e mais odiado do lugarejo. A esposa, que depois fica viúva, vive alheia a tudo isso. Para ela, o marido era um homem bom, e, depois da morte dele, ela não entende o ódio da cidade nem a ruína que se segue.

\textbf{2. Análise do trecho.}
\begin{itemize}
\item A expressão ``Na verdade'' marca a voz do \emph{narrador}, que revela ao leitor o verdadeiro negócio de Montiel: tirar proveito da expulsão dos ricos e comprar seus bens a preço vil.
\item Logo em seguida vem a fala da mulher: ``Você vai arruinar-se ajudando-os [\ldots]''. Ela interpreta como \emph{caridade} uma ação que, na verdade, é \emph{exploração}. Diz que o marido vai se arruinar justamente com o que o enriquece.
\item O narrador coloca lado a lado o fato (enriquecimento à custa das vítimas) e a leitura ingênua da personagem (ajuda generosa que trará prejuízo). O leitor sabe o que a personagem ignora, e é dessa distância que nasce a \textbf{ironia}. Trata-se da ironia dramática ou situacional. Ela é construída pelo narrador, e não pela personagem, que fala com sinceridade.
\end{itemize}

\textbf{3. Exame das alternativas.}
\begin{itemize}
\item[(a)] Incorreta. Os pares ricos/pobres, morte/expulsão e tiros/compra apenas se opõem, sem formar uma afirmação contraditória em si. Isso seria, no máximo, \emph{antítese}, e não paradoxo. Além disso, o efeito central do trecho não está nessas oposições, e sim no contraste entre a verdade e a percepção da esposa.
\item[(b)] Incorreta. A esposa não constrói um paradoxo. Dizer que ajudar alguém pode levar à ruína e não trazer gratidão é uma previsão coerente dentro da lógica dela, não uma contradição aparente. O erro dela está na premissa (acredita que o marido ajuda), e não na forma do raciocínio.
\item[(c)] \textbf{Correta.} O narrador constrói a ironia ao mostrar a personagem numa situação cujo sentido real ela desconhece e o leitor já conhece: o marido não se arruína ajudando ninguém, mas lucra com a desgraça alheia.
\item[(d)] Incorreta. A mulher não sabe das reais intenções do marido. Ela acredita sinceramente na bondade dele, e o conto insiste em sua alienação. Não há, portanto, ironia intencional da personagem nem busca de efeito cômico.
\end{itemize}

\textbf{Conclusão.} O recurso do trecho é a ironia construída pelo narrador. Ela vem do contraste entre o que o leitor sabe (Montiel enriquece explorando os perseguidos) e o que a personagem ignora (ela pensa que ele se arruína por caridade).

ANSWER: (C) o narrador faz uso da ironia, que se constrói ao mostrar a personagem em uma situação na qual ela ignora algo que é do conhecimento do leitor.%
\newpage%
\section*{Questão 39}

\section*{Solução}

\textbf{Ideia central.} A questão pede a explicação \emph{principal} que os textos dão para a mudança na composição racial. Por isso, cada alternativa deve ser comparada com o que o Texto~2 afirma e com a tendência mostrada no gráfico do Texto~1.

\textbf{1. O que mostra o Texto 1 (gráfico).} De 1872 a 1940, a proporção de brancos sobe de cerca de $38\%$ para cerca de $63\%$, enquanto a de pardos cai para cerca de $21\%$ e a de pretos cai de cerca de $20\%$ para cerca de $15\%$. Esse período coincide com o auge da ideologia do \emph{embranquecimento}. Depois de 1940 a tendência se inverte: os pardos crescem até cerca de $45\%$ em 2022 e passam os brancos (cerca de $43\%$), o que combina com a valorização da \emph{mestiçagem}. Os pretos caem até cerca de $5\%$ a $6\%$ em 1980--1990 e voltam a cerca de $10\%$ em 2022, o que combina com a afirmação da \emph{negritude}. Essas oscilações são grandes e mudam de sentido. Por isso, não se explicam apenas por diferenças de natalidade e mortalidade entre os grupos.

\textbf{2. O que afirma o Texto 2.}
\begin{itemize}
  \item A mudança \emph{não} se explica só pelo crescimento vegetativo diferencial, nem pelas mudanças na pergunta do censo, nem pelo uso de hetero ou autoclassificação.
  \item O fator mais importante são as \emph{alterações nas ideologias raciais} (embranquecimento, mestiçagem, negritude). Elas pesam mais que os incentivos materiais criados pelas políticas públicas nos últimos vinte anos.
  \item Essas ideologias \emph{convivem e disputam entre si}, e nenhuma consegue suplantar as outras. Portanto, nenhuma é hegemônica de forma absoluta.
\end{itemize}

\textbf{3. Análise das alternativas.}
\begin{itemize}
  \item[(a)] Incorreta. Coloca o Estado como \emph{produtor} de ideologias \emph{hegemônicas} em cada período. O texto, porém, diz que as ideologias coexistem sem que uma suplante as outras e dá pouco peso relativo à ação estatal (as políticas públicas).
  \item[(b)] \textbf{Correta.} Traz a explicação do texto: há \emph{disputa entre classificações} ligadas a ideologias raciais que convivem. Essas ideologias moldam a forma como os indivíduos \emph{se percebem} e se declaram no censo. Também repercutem nas políticas públicas, como as de combate às desigualdades raciais. Assim, a alternativa trata as políticas como consequência das disputas, e não como sua causa principal.
  \item[(c)] Incorreta. Atribui a mudança a um cálculo racional voltado a benefícios sociais. O texto afirma justamente que esses incentivos materiais são \emph{menos} importantes que as ideologias.
  \item[(d)] Incorreta. Atribui a mudança ao instrumento e à metodologia de coleta. O texto descarta expressamente essa explicação como suficiente: ``nem pelas mudanças na questão censitária [\dots] ou pelo uso da hetero ou autoclassificação''.
\end{itemize}

\textbf{Conclusão.} Segundo os textos, a mudança se explica principalmente pela disputa entre classificações ligadas às ideologias raciais que convivem no Brasil. Essa disputa afeta a autopercepção dos indivíduos e repercute nas políticas públicas. Isso corresponde à alternativa (b).

ANSWER: (B) por disputas de classificações, abarcadas nas diferentes ideologias raciais brasileiras, com impacto na produção de políticas públicas e na forma como os indivíduos se percebem.%
\newpage%
\section*{Questão 40}

\section*{Solução}

\textbf{Tipo de questão.} A questão é de interpretação de texto (Sociologia/Ciência Política). Não há cálculo a fazer. Basta identificar a tese do texto de Przeworski e compará-la com cada alternativa.

\medskip
\textbf{1. Ideias centrais do texto.}
\begin{itemize}
  \item \emph{O que é a democracia:} ``um terreno no qual pessoas mais ou menos iguais e mais ou menos livres lutam pela realização de ideais, valores e interesses conflitantes''. Portanto, o conflito é parte normal da democracia, e não um sinal de que ela falhou.
  \item \emph{Qual é o seu mérito:} ``o único milagre da democracia é que esses conflitos podem ser administrados sem repressão e em paz''.
  \item \emph{O que vale sob polarização:} quando as pessoas discordam até sobre quais valores a democracia deve buscar, ela ``só pode ser defendida como o melhor método para administrar essas divergências''.
\end{itemize}
Conclusão: em contextos de polarização, o texto não diz que a democracia entra em crise nem que perde valor. Ao contrário, é justamente aí que ela deve ser defendida, como \emph{método} de gestão pacífica e sem repressão dos conflitos entre ideais, valores e interesses diferentes.

\medskip
\textbf{2. Análise das alternativas.}
\begin{itemize}
  \item[(a)] \textbf{Incorreta.} O texto não fala em crise nem em ``incapacidade de dialogar com a maioria''. O critério de ``maioria'' não aparece. Para o autor, o conflito entre valores faz parte da própria democracia.
  \item[(b)] \textbf{Incorreta.} Diz o oposto do texto. Para Przeworski, a democracia deve ser \emph{defendida}, e não questionada, como o melhor método de administrar conflitos.
  \item[(c)] \textbf{Incorreta.} Também contradiz o texto. O ``milagre'' da democracia é exatamente administrar os conflitos \emph{sem repressão}. A polarização não mostra que isso é impossível; é ela que torna esse método ainda mais necessário.
  \item[(d)] \textbf{Correta.} A alternativa repete a tese do autor: a democracia é o método para administrar em paz (``sem repressão e em paz'') as divergências entre ideias, valores e interesses existentes na sociedade. É assim que ela deve ser defendida diante da polarização.
\end{itemize}

\medskip
\textbf{Conclusão.} Segundo o texto, em contextos de polarização a democracia é o método eficaz para gerir de forma pacífica os conflitos entre diferentes ideias, valores e interesses.

\medskip
ANSWER: (D) apresenta-se como método eficaz para a gestão pacífica de conflitos entre diferentes ideias, valores e interesses existentes na sociedade.%
\newpage%
\section*{Questão 41}

\section*{Solução}

\textbf{1. O que o Texto 1 afirma.} O autor defende três ideias:
\begin{itemize}
  \item O que caracteriza os negacionismos é o \emph{investimento sistemático na fabricação de falsas controvérsias científicas}.
  \item Esse investimento fica nas \emph{franjas} das ciências estabelecidas, numa ``região de sobreposição'' ou ``região cinzenta''. Por isso \emph{``não se trata de um mero repúdio `pseudo' ou `anticientífico'''}.
  \item Nessa região, os negacionistas agem \emph{``apropriando-se deliberadamente de retóricas e credenciais científicas, com vistas a persuadir terceiros, criando dissensos e incertezas''}.
\end{itemize}

\textbf{2. O que o Texto 2 mostra.} No esquema, a fronteira entre ``ciências'' e ``negacionismos'' é tracejada e inclinada, ou seja, não há uma separação nítida. As elipses ``controvérsias'' (legítimas, próprias das ciências) e ``falsas controvérsias'' (próprias do negacionismo) se cruzam. A interseção sombreada é a ``região cinzenta'' do Texto 1. Nela, a falsa controvérsia se disfarça de debate científico real.

\textbf{3. Análise das alternativas.}
\begin{itemize}
  \item[(a)] \textbf{Incorreta.} Ela diz que os negacionistas buscam ``se afirmarem como forças anticientíficas''. O texto nega isso de forma explícita: não se trata de ``mero repúdio anticientífico''. Na verdade, eles imitam a ciência e se apresentam como científicos.
  \item[(b)] \textbf{Incorreta.} Ela descreve uma simples oposição às ciências estabelecidas. O texto e a figura mostram outra coisa: uma sobreposição com as ciências, nas suas franjas, e não um confronto externo. Além disso, o objetivo indicado no texto não é ``buscar legitimidade junto aos grupos sociais'', e sim persuadir terceiros e criar dissensos e incertezas.
  \item[(c)] \textbf{Correta.} Ela reproduz o mecanismo descrito no Texto 1: apropriação de retóricas e credenciais científicas, fabricação de falsas controvérsias e o objetivo de gerar dissensos e incertezas. Ela também corresponde à interseção sombreada do Texto 2, em que a falsa controvérsia ocupa o espaço das controvérsias legítimas.
  \item[(d)] \textbf{Incorreta.} Para o texto, os dissensos e incertezas são \emph{criados} pelos negacionistas. Não são algo já existente a que eles se ``sobrepõem''. A sobreposição descrita é com as ciências, e não com os dissensos. Por fim, os negacionistas não negam a existência das ciências: eles se apropriam das credenciais científicas.
\end{itemize}

\textbf{Conclusão.} Só a alternativa (c) reúne os três elementos dos textos: o meio (apropriação de retóricas e credenciais científicas), o produto (falsas controvérsias) e a finalidade (dissensos e incertezas).

ANSWER: (C) se apropriam de retóricas e credenciais científicas para produzir falsas controvérsias que visam gerar dissensos e incertezas na sociedade.%
\newpage%
\section*{Questão 42}

\section*{Solução}

\textbf{1. O que a tirinha mostra.} O ``chefe'' é uma molécula de DNA (dupla hélice) que dita uma mensagem. A ``secretária'' é uma fita simples que vai produzir uma cópia dessa mensagem para o ``Dr.\ Ribossomo''. A fita simples representa o RNA mensageiro, e o ribossomo é o lugar onde a mensagem será usada depois. Portanto, a cena mostra a passagem da informação do DNA para o RNA:
\[
\text{DNA} \xrightarrow{\text{transcrição}} \text{RNAm} \xrightarrow{\text{tradução}} \text{proteína}.
\]
A tirinha mostra somente a primeira etapa, que é a \textbf{transcrição}.

\textbf{2. Conferindo o pareamento de bases.} Na transcrição, a RNA polimerase usa uma fita do DNA como molde. Ela monta o RNA com bases complementares, e no RNA a uracila (U) ocupa o lugar da timina (T):
\[
\text{A} \rightarrow \text{U}, \qquad \text{T} \rightarrow \text{A}, \qquad \text{C} \rightarrow \text{G}, \qquad \text{G} \rightarrow \text{C}.
\]
Aplicando essas regras ao que o DNA ditou:
\[
\text{CACA}\ \text{CTCT}\ \text{ACT}\ \text{ACT} \;\longrightarrow\; \text{GUGU}\ \text{GAGA}\ \text{UGA}\ \text{UGA},
\]
que é exatamente o que a secretária escreveu. As duas piadas da tirinha correspondem a fatos biológicos:
\begin{itemize}
\item ``é o contrário do que eu falei'': a fita produzida é \emph{complementar} ao molde, e não uma cópia idêntica;
\item ``meu computador não tem o T'': o RNA não tem timina e usa uracila no lugar dela.
\end{itemize}

\textbf{3. Onde ocorre e para onde vai a mensagem.} Nos eucariotos, o DNA fica no núcleo. Por isso a transcrição ocorre no \textbf{núcleo}. O RNAm formado sai pelos poros nucleares e vai para o \textbf{citoplasma}, onde os ribossomos fazem a tradução. Isso explica a mensagem ``enviada ao Dr.\ Ribossomo''.

\textbf{4. Análise das alternativas.}
\begin{itemize}
\item[(a)] \textbf{Correta.} A alternativa diz que é transcrição, no núcleo, com o DNA como molde e um RNAm de bases complementares que leva a informação ao citoplasma. Tudo isso confere com os itens 1 a 3.
\item[(b)] Incorreta. A tirinha não mostra a tradução, porque o DNA participa da cena. Além disso, a tradução não produz uma fita de RNAt: ela produz uma cadeia de aminoácidos. O RNAt apenas leva aminoácidos e reconhece os códons por meio do anticódon.
\item[(c)] Incorreta. A transcrição não ocorre no ribossomo, e o molde não é o RNAt. O molde é o DNA.
\item[(d)] Incorreta. O processo não é a tradução, e a tradução não ocorre no núcleo. A molécula formada a partir do DNA na tirinha é o RNAm, que é quem leva a mensagem ao ribossomo, e não o RNAt.
\end{itemize}

\textbf{Conclusão.} A tirinha mostra a transcrição: o DNA serve de molde, o RNAm formado tem bases complementares com U no lugar de T, o processo ocorre no núcleo e o RNAm leva a informação para o citoplasma.

ANSWER: (A) transcrição, que ocorre no núcleo; o DNA serve de molde para a síntese de um RNAm complementar, que leva a informação ao citoplasma.%
\newpage%
\section*{Questão 43}

\section*{Solução}

\textbf{Ideia central.} O texto diz que o veneno da jararaca-ilhoa é \emph{mais letal para aves}, sua principal presa, que isso se deve a \emph{modificações genéticas} (como a duplicação de genes) e que as toxinas sofreram \emph{seleção positiva}. Precisamos separar três processos evolutivos diferentes:
\begin{itemize}
  \item \textbf{Mutação} (por exemplo, a duplicação gênica): é a fonte da variação. Ocorre ao acaso, sem ``direção'' e sem ser causada pelo ambiente.
  \item \textbf{Seleção natural}: quando uma variante aumenta a sobrevivência e/ou a reprodução de quem a carrega, sua frequência sobe de geração em geração. É isso que significa \emph{seleção positiva}.
  \item \textbf{Deriva genética} e \textbf{fluxo gênico}: a deriva é a variação aleatória das frequências alélicas, mais forte em populações pequenas. O fluxo gênico é a troca de alelos entre populações por migração.
\end{itemize}

\textbf{Raciocínio.} Na ilha, as aves são a principal fonte de alimento. Surgiu por mutação a duplicação de genes de toxinas, que aumenta a quantidade ou a diversidade de toxinas eficazes contra aves. As serpentes com essa variante capturam presas com mais sucesso, sobrevivem e se reproduzem mais, e por isso deixam mais descendentes. Assim, ao longo de muitas gerações, a variante duplicada ficou mais frequente na população. O resultado é a adaptação do veneno à presa.

\textbf{Análise das alternativas.}
\begin{itemize}
  \item[(a)] \textbf{Incorreta.} A adaptação à presa existe, mas não vem de fluxo gênico. Além disso, a seleção natural não ``causa'' fluxo gênico. Fluxo gênico é a troca de alelos entre populações. Uma população isolada numa ilha, como esta, tem pouco ou nenhum fluxo gênico com o continente. É justamente esse isolamento que permite a diferenciação.
  \item[(b)] \textbf{Incorreta.} A duplicação gênica é de fato um tipo de mutação. Porém mutações não são ``motivadas'' por uma necessidade. Além disso, a especiação não foi súbita: a diferenciação foi gradual, ao longo de cerca de 100 mil anos, com isolamento geográfico, mutação e seleção agindo em conjunto.
  \item[(c)] \textbf{Correta.} A alternativa descreve exatamente a seleção positiva: a variante duplicada deu vantagem de sobrevivência e/ou reprodução às serpentes que a carregavam e, por isso, ficou mais frequente com o tempo.
  \item[(d)] \textbf{Incorreta.} O isolamento na ilha de fato favoreceu a especiação, que é alopátrica. Mas a deriva genética apenas muda frequências alélicas ao acaso. Ela não \emph{causa} mutações como a duplicação gênica. Além disso, o texto atribui a evolução das toxinas à seleção positiva, e não à deriva.
\end{itemize}

\textbf{Conclusão.} A única alternativa coerente com o texto e com a teoria sintética da evolução é a (c). A mutação (duplicação) gera a variação, e a seleção natural aumenta a frequência da variante vantajosa.

ANSWER: (C)%
\newpage%
\section*{Questão 44}

\section*{Solução}

Conferi o enunciado com a prova original (página 15 do PDF). A questão pede cinco lacunas: duas sobre a evolução do sistema nervoso e três sobre a fisiologia do neurônio.

\textbf{(i) Primeira rede nervosa difusa: cnidários.} Águas-vivas, anêmonas, hidras e corais são cnidários. Têm simetria radial e são os animais mais simples com células nervosas. Esses neurônios formam uma \emph{rede nervosa difusa}, sem centralização. Essa rede responde por respostas rápidas de retração (por exemplo, o pólipo que se contrai ao ser tocado) e pelas contrações rítmicas da umbrela da medusa durante a natação. Os equinodermos (estrelas-do-mar, ouriços) não possuem cérebro, mas não são o grupo em que a rede difusa surge. Eles são deuterostômios, que surgiram depois na evolução. Seu sistema nervoso, formado por um anel nervoso e nervos radiais, deriva de ancestrais bilaterais, e sua simetria radial é secundária.

\textbf{(ii) Cordões longitudinais e gânglio anterior: platelmintos.} Os platelmintos (planárias, tênias) são os primeiros animais com simetria bilateral e \emph{cefalização}. Na região anterior há um par de gânglios cerebroides, que é o agrupamento de células nervosas citado no texto. Desses gânglios partem cordões nervosos longitudinais ligados por nervos transversais, formando o sistema nervoso \emph{ganglionar} ou \emph{em escada}. Por isso a lacuna (ii) não pode ser preenchida por cnidários nem por equinodermos.

\textbf{(iii) Local do potencial de ação: membrana plasmática.} O potencial de ação é uma variação rápida da diferença de potencial elétrico através da \emph{membrana plasmática} do neurônio, sobretudo do axônio. Ele depende de canais iônicos dependentes de voltagem que existem nessa membrana. O envoltório nuclear apenas separa o núcleo do citoplasma e não participa da propagação do impulso nervoso.

\textbf{(iv) e (v) Ordem dos fluxos iônicos.} No repouso, a face interna da membrana fica negativa, com cerca de $-70\ \text{mV}$. A bomba de sódio e potássio mantém $[\text{Na}^+]$ maior fora da célula e $[\text{K}^+]$ maior dentro dela, transportando $3\,\text{Na}^+$ para fora e $2\,\text{K}^+$ para dentro.
\begin{itemize}
\item \emph{Despolarização:} ao ser atingido o limiar, abrem-se os canais de $\text{Na}^+$. O $\text{Na}^+$ \textbf{entra} a favor do seu gradiente eletroquímico, e o potencial sobe até cerca de $+30\ \text{mV}$. Logo, (iv) é \textbf{sódio}.
\item \emph{Repolarização:} os canais de $\text{Na}^+$ se inativam e abrem-se os canais de $\text{K}^+$. O $\text{K}^+$ \textbf{sai} da célula, e a face interna volta a ficar negativa. Logo, (v) é \textbf{potássio}.
\end{itemize}
A saída de $\text{K}^+$ devolve a membrana ao potencial de repouso. Em seguida, a bomba $\text{Na}^+/\text{K}^+$-ATPase restabelece as concentrações iniciais dos íons. Uma ordem invertida (entrada de $\text{K}^+$ e depois saída de $\text{Na}^+$) contrariaria os gradientes de concentração e não geraria o potencial de ação.

\textbf{Conferência das alternativas.}
\begin{itemize}
\item (a) Está errada: inverte a evolução, pondo equinodermos em (i) e cnidários em (ii), e cita o envoltório nuclear.
\item (c) Está errada: põe equinodermos em (i) e inverte os íons.
\item (d) Está errada: põe equinodermos em (ii), cita o envoltório nuclear e inverte os íons.
\item (b) É a única alternativa que acerta as cinco lacunas: cnidários, platelmintos, membrana plasmática, sódio e potássio.
\end{itemize}

ANSWER: (B) (i) cnidários; (ii) platelmintos; (iii) da membrana plasmática; (iv) sódio; (v) potássio.%
\newpage%
\section*{Questão 45}

\section*{Solução}

\textbf{Leitura da figura.} O Sol está no alto, à esquerda, e o caule se curva para a esquerda, em direção à luz. As setas 1 e 3 ficam no lado esquerdo do caule, o lado \emph{iluminado}. A seta 2 fica no lado direito da parte superior do caule, o lado \emph{sombreado}. As setas 4, 5 e 6 estão na raiz: 4 e 5 nos dois lados e 6 no ápice.

\textbf{Fundamento: hipótese de Cholodny--Went.}
\begin{enumerate}
\item A auxina (AIA) é produzida principalmente no ápice do caule e se desloca para baixo por transporte polar.
\item \emph{Fototropismo:} quando a luz incide de um só lado, os fotorreceptores (fototropinas) fazem a auxina migrar lateralmente para o lado \emph{sombreado}. No caule, a auxina em concentração adequada estimula o alongamento celular, pois acidifica e afrouxa a parede celular. Por isso, as células do lado sombreado se alongam mais que as do lado iluminado, e o caule se curva para a luz: \textbf{fototropismo positivo do caule}. A curvatura acontece na \emph{zona de alongamento}, logo abaixo do ápice, e não por divisão celular no meristema.
\item \emph{Geotropismo (gravitropismo):} a coifa percebe a gravidade por meio de amiloplastos (estatólitos) que se depositam no fundo das células. Com isso, a \emph{auxina} se acumula no lado inferior da raiz. A raiz é muito mais sensível à auxina que o caule, então essa concentração mais alta \emph{inibe} o alongamento: o lado de cima cresce mais e a raiz se curva para baixo (\textbf{geotropismo positivo da raiz}). O hormônio responsável pelos tropismos é a auxina, não a citocinina. A citocinina atua sobretudo na divisão celular e no retardo da senescência.
\end{enumerate}

\textbf{Análise das alternativas.}
\begin{itemize}
\item[a)] \textbf{Falsa.} A seta 1 indica o lado \emph{iluminado}, onde a concentração de auxina \emph{diminui}, porque a auxina migra para o lado oposto. Além disso, a curvatura se deve ao alongamento diferencial na zona de alongamento, não a um aumento de auxina no meristema apical do lado da luz.
\item[b)] \textbf{Falsa.} O geotropismo positivo da raiz é controlado pela redistribuição de \emph{auxina}, percebida na coifa por estatólitos. Ele não é causado por aumento de citocinina.
\item[c)] \textbf{Verdadeira.} A curvatura em direção ao Sol mostra fototropismo positivo na parte aérea. A seta 2 está no lado sombreado, para onde a auxina migra. A seta 3 está no lado iluminado, que perde auxina. Logo,
\[
[\text{auxina}]_{2} > [\text{auxina}]_{3},
\]
e é exatamente essa diferença que faz o lado 2 alongar mais e o caule se curvar para a luz.
\item[d)] \textbf{Falsa.} O geotropismo positivo ocorre de fato na raiz, mas a assimetria hormonal que o explica é de \emph{auxina}, não de citocinina. Além disso, a raiz da figura cresce verticalmente, e nada no enunciado indica diferença de citocinina entre os lados 4 e 5.
\end{itemize}

\textbf{Conclusão.} A única afirmação coerente com o mecanismo do fototropismo é a (c): a auxina se concentra no lado sombreado (seta 2), mais do que no lado iluminado (seta 3), e o caule se curva para a luz.

ANSWER: (C)%
\newpage%
\section*{Questão 46}

\section*{Solução}

\textbf{1. Leitura da figura.} Os quatro níveis representam, de baixo para cima, produtores (branco), consumidores primários (cinza-claro), secundários (cinza-escuro) e terciários (preto). Comparando as larguras nos dois tempos:

\begin{center}
\begin{tabular}{lcc}
\hline
Nível trófico & $t_1$ & $t_2$ \\
\hline
Produtores (branco) & muito largo & muito estreito (forte redução) \\
Consumidores primários (cinza-claro) & largo & estreito (forte redução) \\
Consumidores secundários (cinza-escuro) & médio & praticamente igual \\
Consumidores terciários (preto) & estreito & praticamente igual \\
\hline
\end{tabular}
\end{center}

Em $t_1$ a pirâmide é \emph{regular}. Em $t_2$ ela fica \emph{parcialmente invertida}: os níveis de baixo encolheram, mas os de cima não mudaram, e o nível dos consumidores secundários ficou maior que o dos níveis que o sustentam.

\textbf{2. Não pode ser uma pirâmide de energia (eliminamos a e c).} A energia entra pelos produtores. A cada passagem de um nível para o seguinte, a maior parte dela se perde como calor na respiração, nas fezes e nos restos não consumidos. Só cerca de $10\%$ passa adiante. Pela conservação da energia, o fluxo que chega a um nível nunca pode ser maior que o do nível anterior:
\[
E_1 > E_2 > E_3 > E_4 .
\]
Por isso a pirâmide de energia é \emph{sempre} regular e nunca se inverte. Em $t_2$, porém, o nível 3 é maior que os níveis 1 e 2, o que seria impossível para fluxo de energia. Logo, as alternativas (a) e (c) estão erradas.

\textbf{3. É uma pirâmide de biomassa.} Uma pirâmide de números também pode se inverter, mas as alternativas só oferecem energia ou biomassa. A biomassa é a quantidade de matéria orgânica presente num dado momento, e ela \emph{pode} aparecer invertida. Isso acontece quando a base é destruída e os consumidores ainda não diminuíram, ou quando os produtores se renovam muito depressa, como o fitoplâncton. Portanto, a figura representa pirâmides de biomassa.

\textbf{4. Comparando (b) e (d).}
\begin{itemize}
\item A tabela mostra que \emph{duas} bases foram reduzidas: os produtores e também os consumidores primários. Os níveis superiores ficaram iguais.
\item Os consumidores dos níveis de cima mantiveram sua biomassa, mas o alimento que os sustenta diminuiu muito. Eles ainda não sentiram a falta de alimento, porque a resposta de cada nível vem com atraso. Com o tempo, vão declinar por falta de recurso. Assim, $t_2$ é um estado \emph{transitório}, em que a teia ainda está se reorganizando, e não um novo equilíbrio. Isso é reforçado porque a mesma teia, em equilíbrio ($t_1$), tinha pirâmide regular.
\end{itemize}

A alternativa (d) diz que o equilíbrio foi restabelecido em $t_2$, o que contradiz o item acima. Ela também cita apenas os produtores, embora os consumidores primários também tenham caído. Portanto, (d) é falsa.

A alternativa (b) afirma justamente que se trata de pirâmides de biomassa, que produtores e consumidores primários foram atingidos por um desequilíbrio e que a teia ainda está se reestruturando em $t_2$. Tudo isso confere com a análise.

\textbf{Conclusão.} São pirâmides de biomassa. Em $t_2$, a redução de produtores e consumidores primários, com os níveis superiores ainda inalterados, mostra uma teia em desequilíbrio, ainda se reestruturando.

ANSWER: (B) biomassa; produtores e consumidores primários foram impactados por desequilíbrio, com a teia ainda se reestruturando em $t_2$.%
\newpage%
\section*{Questão 47}

\section*{Solução}

\textbf{Conceitos-base.}
\begin{itemize}
\item O HIV é um retrovírus. Ele infecta principalmente os linfócitos T auxiliares $\text{CD4}^{+}$ (além de macrófagos). Por meio da transcriptase reversa, o RNA viral vira DNA, que a integrase insere no genoma da célula hospedeira. Por isso se formam \emph{reservatórios latentes} que o sistema imunológico não consegue eliminar: a infecção dura a vida toda e ainda não tem cura.
\item A AIDS (síndrome da imunodeficiência adquirida) não é uma doença separada. Ela é o \emph{estágio mais avançado} da infecção pelo HIV. Sem tratamento, o número de linfócitos $\text{CD4}^{+}$ cai progressivamente. Fala-se em AIDS quando essa contagem fica abaixo de cerca de $200$ células/mm$^3$ ou quando surgem doenças oportunistas definidoras (pneumocistose, tuberculose, sarcoma de Kaposi, candidíase esofágica etc.).
\item A terapia antirretroviral (TARV), feita de forma contínua, impede a replicação do vírus. Assim, a carga viral no sangue pode ficar \emph{indetectável}, a imunidade se recupera e a evolução para a AIDS é evitada. Além disso, a pessoa com carga viral indetectável não transmite o vírus por via sexual (princípio ``indetectável = intransmissível'').
\end{itemize}

\textbf{Análise das alternativas.}

\textbf{a) Falsa.} O sistema imunológico \emph{não} elimina o HIV, porque o vírus se integra ao DNA das células e permanece latente. Sem tratamento, a maioria das pessoas infectadas evolui para a AIDS, depois de um período assintomático que pode durar anos.

\textbf{b) Falsa.} A primeira frase está correta, mas a conclusão não. Mesmo sem cura, os antirretrovirais \emph{previnem} a transmissão (carga viral indetectável) e a progressão da doença. Além disso, as doenças oportunistas não ``causam a AIDS'': elas são consequência da imunodeficiência.

\textbf{c) Falsa.} É verdade que o HIV permanece no organismo por toda a vida e ataca o sistema imunológico. O erro está na relação de causa e efeito: não existem ``doenças que causam a AIDS''. É o HIV que causa a AIDS, e a imunodeficiência é que deixa a pessoa suscetível às doenças oportunistas.

\textbf{d) Verdadeira.} A AIDS é de fato o estágio mais grave da infecção pelo HIV. A TARV mantém a replicação viral suprimida, a ponto de a carga viral ficar indetectável. Com isso, a contagem de linfócitos $\text{CD4}^{+}$ se preserva e a evolução para a doença (AIDS) é evitada. As duas afirmações são corretas e a relação de causa entre elas também.

\textbf{Conclusão.} A única alternativa em que a definição de AIDS e o papel do tratamento antirretroviral estão corretos é a \textbf{d}.

\medskip
ANSWER: (D)%
\newpage%
\section*{Questão 48}

\section*{Solução}

\textbf{1. O que a questão pede.} O enunciado diz que a lesão medular rompe os \emph{axônios}. A pergunta é em que região da medula espinhal esses axônios ficam concentrados. Para responder, é preciso lembrar como os componentes do neurônio se distribuem no corte transversal da medula.

\textbf{2. Organização da medula espinhal (corte transversal).}
\begin{itemize}
  \item \textbf{Substância cinzenta:} fica na \emph{região central} e tem forma de ``H'' ou de borboleta. Nela se concentram os \emph{corpos celulares} dos neurônios (por exemplo, os neurônios motores dos cornos anteriores), os dendritos, os interneurônios e as sinapses. A cor acinzentada se deve à pouca mielina.
  \item \textbf{Substância branca:} fica na \emph{região periférica}, envolvendo a cinzenta. É formada por \emph{fibras nervosas}, isto é, axônios, em grande parte mielinizados. A cor esbranquiçada vem justamente da mielina, que é rica em lipídios. Esses axônios se organizam em feixes (tratos):
  \begin{itemize}
    \item \emph{ascendentes} (sensitivos), que levam ao encéfalo as informações vindas do corpo;
    \item \emph{descendentes} (motores), que levam às regiões inferiores os comandos gerados no encéfalo.
  \end{itemize}
\end{itemize}
No encéfalo (córtex) a disposição é a inversa: a substância cinzenta fica por fora e a branca por dentro. Muitas alternativas exploram essa confusão.

\textbf{3. Ligação com o enunciado.} A lesão interrompe esses tratos de axônios na substância branca. Por isso há perda motora (vias descendentes) e sensorial (vias ascendentes) \emph{abaixo} do ponto lesado, como diz o texto. Isso confirma que os axônios relevantes estão na substância branca periférica.

\textbf{4. Análise das alternativas.}
\begin{itemize}
  \item[(a)] Incorreta. A descrição das fibras ascendentes e descendentes na periferia está certa, mas essa região é a substância \emph{branca}, e não a cinzenta.
  \item[(b)] Incorreta. A substância branca não é central, e os corpos celulares dos neurônios motores ficam na substância cinzenta.
  \item[(c)] \textbf{Correta.} A substância branca é periférica e formada por fibras nervosas (axônios) que conduzem impulsos ascendentes ao encéfalo e descendentes gerados por ele.
  \item[(d)] Incorreta. A descrição da substância cinzenta está correta (central, com corpos celulares), mas ali se concentram os \emph{corpos celulares}, e não os axônios, que é o que a questão pede.
\end{itemize}

\textbf{Conclusão.} Na medula espinhal, os axônios ficam concentrados na substância branca. Ela é periférica e formada pelos tratos de fibras ascendentes e descendentes.

ANSWER: (C) substância branca, periférica, com fibras que conduzem impulsos ascendentes e descendentes%
\newpage%
\section*{Questão 49}

\section*{Solução}

\textbf{Ideia central do texto.} Para Milton Santos (\emph{Por uma outra globalização}, 2000), só se entende o mundo atual quando se consideram ao mesmo tempo três dimensões:
\begin{enumerate}
\item \textbf{a globalização como fábula}: o mundo como nos fazem vê-lo, sustentado por mitos como a ``aldeia global'', a ``morte do Estado'' e a ideia de que tudo está ao alcance de todos;
\item \textbf{a globalização como perversidade}: o mundo como ele é, comandado pela tirania do dinheiro e da informação. Essa modernização seletiva aprofunda o desemprego, a pobreza e a desigualdade;
\item \textbf{uma outra globalização}: o mundo como pode ser. Ela já aparece como \emph{tendência} no presente, por meio da humanidade misturada e concentrada nas grandes cidades e da cultura popular, que funciona como contrarracionalidade diante da ordem vigente.
\end{enumerate}
Para o autor, essas dimensões não são etapas separadas. Elas são aspectos de \emph{uma mesma realidade}, que é contraditória e está em movimento: o mundo é uma totalidade em processo, aberta à transformação. Por isso a terceira dimensão não é utopia. Ela se apoia em processos que já existem, sobretudo nas metrópoles.

\textbf{Análise das alternativas.}

\textbf{a) Incorreta.} A ``aldeia global'' conectada pelas técnicas da informação, em que o consumo reduziria as desigualdades, é justamente a globalização como \emph{fábula} que o autor critica. Na realidade (a dimensão da perversidade), as desigualdades aumentam.

\textbf{b) Incorreta.} As dimensões não estão dissociadas. Além disso, a terceira dimensão não é ``fábula'': é uma possibilidade concreta, apreendida como tendência no presente. A alternativa confunde a primeira dimensão com a terceira.

\textbf{c) Incorreta.} A perversidade é só uma das dimensões e não esgota o mundo concreto. A globalização também não ``homogeneizou'' o mundo: ela é seletiva e aprofunda as diferenças entre lugares. Por fim, para Santos as grandes cidades periféricas, com sua escassez, são o lugar onde nasce a contrarracionalidade (a cultura popular e a ``outra globalização''), e não apenas vítimas passivas.

\textbf{d) Correta.} Fábula, perversidade e possibilidade formam uma unidade. O mundo aparece como uma totalidade aberta e em movimento, e as transformações que apontam para ``uma outra globalização'' já estão em curso nas grandes cidades, onde vive a humanidade misturada e se desenvolve a cultura popular, como diz o próprio texto.

\textbf{Conclusão.} A única alternativa que articula as três dimensões e mostra a terceira como tendência já presente no espaço urbano é a (d).

ANSWER: (D) As três dimensões formam uma unidade: o mundo é uma totalidade aberta, em movimento, com transformações já em curso nas grandes cidades.%
\newpage%
\section*{Questão 50}

\section*{Solução}

\textbf{1. O que o texto mostra.} A notícia apresenta dois atores centrais na corrida pela inteligência artificial (IA) militar:
\begin{itemize}
  \item a \textbf{China}, que alerta para a ``militarização desenfreada'' da IA;
  \item os \textbf{EUA}, destinatários do alerta e ``envolvidos em um intenso debate sobre as aplicações militares'' da IA.
\end{itemize}
O texto também traz dois tipos de dilema:
\begin{itemize}
  \item \emph{ético}: algoritmos com ``decisão de vida e morte sobre seres humanos'';
  \item \emph{existencial}: a distopia de ``O Exterminador do Futuro'', em que as máquinas tomam o poder.
\end{itemize}
Tudo isso ocorre num cenário de ``forte tensão geopolítica'' e ``escalada de conflitos entre as Nações''.

\textbf{2. Contexto geopolítico.}
\begin{itemize}
  \item Hoje a infraestrutura da IA (chips avançados, \emph{data centers}, nuvem, grandes modelos e terras raras) concentra-se sobretudo nos EUA e na China. São as duas potências que disputam a liderança tecnológica.
  \item Essas economias são \emph{concorrentes} (guerra comercial e restrições à exportação de semicondutores), mas também \emph{interdependentes}: comércio bilateral intenso, cadeias produtivas integradas e dependência mútua de insumos e mercados.
  \item O sistema internacional atual é descrito como \emph{multipolar} e em rápida transformação. Há outros polos de poder (União Europeia, Rússia, Índia, BRICS), e a hegemonia tecnológica sino-americana se exerce dentro desse quadro, sem formar blocos fechados como na Guerra Fria.
\end{itemize}

\textbf{3. Análise das alternativas.}
\begin{itemize}
  \item[(a)] \textbf{Falsa.} Falar em ordem \emph{unipolar} dos EUA contradiz o próprio texto. Nele, a China aparece como potência que confronta e adverte os EUA, ou seja, o poder não está concentrado em um só país.
  \item[(b)] \textbf{Verdadeira.} A alternativa reúne todos os elementos do texto e do contexto:
  \begin{itemize}
    \item EUA e China como os dois países que dominam a infraestrutura de IA;
    \item economias interdependentes e, ao mesmo tempo, concorrentes;
    \item hegemonia exercida num mundo multipolar em transformação acelerada;
    \item dilemas ético (máquinas decidindo sobre a vida humana) e existencial (o risco distópico das máquinas no poder) candentes.
  \end{itemize}
  \item[(c)] \textbf{Falsa.} O texto não diz que a China seguiu tratados de não proliferação ao desenvolver sua IA. Também não existe tratado internacional consolidado sobre armas autônomas com IA. Além disso, o alerta chinês é um discurso diplomático contra o rival, e não prova de que a China lidere uma ordem ``ética e não distópica''.
  \item[(d)] \textbf{Falsa.} A alternativa tem três erros:
  \begin{itemize}
    \item fala num mundo \emph{bipolar} ``recriado'', mas a ordem atual é multipolar;
    \item diz que as economias se desenvolvem ``de forma independente'', mas elas são profundamente interdependentes;
    \item afirma, de forma genérica, que tudo ocorre ``sem critérios éticos'', mas o próprio texto registra o debate ético nos EUA e o alerta chinês.
  \end{itemize}
\end{itemize}

\textbf{Conclusão.} Na era da IA, a geopolítica fortalece o poder dos EUA e da China, potências interdependentes e rivais que dominam a infraestrutura tecnológica. Elas exercem hegemonia num mundo multipolar e em rápida mudança, marcado por dilemas éticos e existenciais sobre o uso militar da IA. Por isso a resposta é a alternativa (b).

ANSWER: (B) EUA e China, economias interdependentes e concorrentes que dominam a infraestrutura de IA, num mundo multipolar com dilemas ético e existencial%
\newpage%
\section*{Questão 51}

\section*{Solução}

\textbf{1. Leitura do gráfico.} Em Mato Grosso, a participação da cana-de-açúcar na produção de etanol caiu de cerca de $93\%$ (2017) para cerca de $52\%$ (2024). No mesmo período, a do milho subiu de cerca de $7\%$ para cerca de $48\%$. Em 2024 as duas matérias-primas já têm pesos quase iguais. Esse avanço do etanol de milho é o que a questão chama de transformação produtiva e socioespacial.

\textbf{2. Diferença entre as duas matérias-primas.}
\begin{itemize}
\item \emph{Cana-de-açúcar}: é perecível. Depois do corte, precisa ser moída em pouco tempo (cerca de 24 a 48 horas), senão perde sacarose. Por isso o canavial fica perto da usina. A cana também é uma cultura semiperene: depois de plantada, rebrota por vários anos e tem safra definida, com entressafra em que a usina fica parada.
\item \emph{Milho}: é um grão seco e pode ser guardado em silos e armazéns por meses. Pode ser transportado de longe e processado o ano todo. Em Mato Grosso, ele é cultivado sobretudo como ``milho safrinha'', a segunda safra plantada depois da soja.
\end{itemize}

\textbf{3. Análise das alternativas.}

\textbf{a) Correta.} Surgiram as chamadas \emph{usinas flex}: unidades adaptadas que processam cana durante a safra e milho armazenado na entressafra da cana. A usina deixa de ficar parada parte do ano e ganha flexibilidade para escolher a matéria-prima conforme os preços e a disponibilidade. Há também as usinas \emph{full}, que usam só milho. É exatamente o que a alternativa descreve.

\textbf{b) Incorreta.} A estocagem e a logística das duas não se correspondem. A cana não pode ser estocada e exige proximidade da usina. O milho pode ser guardado em silos e transportado por longas distâncias.

\textbf{c) Incorreta.} Por ser semiperene, a cana ocupa a mesma terra por vários anos seguidos. Não existe rotação com milho na entressafra da cana dentro do mesmo talhão. O milho usado para etanol em Mato Grosso vem principalmente da sucessão soja--milho safrinha, e não de rotação com a cana.

\textbf{d) Incorreta.} A colheita de milho, sobretudo nas grandes propriedades de Mato Grosso, é amplamente mecanizada, feita com colheitadeiras. Ela não absorve mais mão de obra que a cana. Assim, a expansão do etanol de milho não amplia o emprego no campo.

\textbf{4. Conclusão.} A única alternativa coerente com as características das duas matérias-primas é a (a). A possibilidade de armazenar o milho permite que as usinas processem as duas matérias-primas e mantenham a produção de biocombustível ao longo do ano, com maior flexibilidade.

ANSWER: (A)%
\newpage%
\section*{Questão 52}

\section*{Solução}

\textbf{1. Leitura da tabela.} As taxas conferem com
\[
\text{taxa}=\frac{P_{2022}-P_{2010}}{P_{2010}}\cdot 100\%,
\]
por exemplo, para Senador Canedo: $\dfrac{155635-84443}{84443}\cdot 100\% \approx 84{,}31\%$. Em doze anos, todos os municípios cresceram entre $50\%$ e $84\%$. O que importa, porém, é identificar \emph{por que} cada um cresceu. Para isso, agrupamos os municípios pela função que exercem no território:

\begin{itemize}
\item \textbf{Fronteira agrícola (agronegócio de grãos no Cerrado e na Amazônia Legal):} Luís Eduardo Magalhães (BA), no oeste baiano, e Sinop e Sorriso (MT), no eixo da BR-163. São cidades ligadas à soja, ao milho e ao algodão, que oferecem serviços, armazenagem e comércio ao agronegócio.
\item \textbf{Fronteira mineral:} Parauapebas (PA), sede da mineração de ferro na Serra dos Carajás.
\item \textbf{Combustíveis fósseis:} Maricá (RJ), cujo crescimento está ligado aos \emph{royalties} do petróleo do pré-sal na Bacia de Santos, além da expansão a partir da metrópole fluminense.
\item \textbf{Turismo e litoral:} Camboriú (SC), vizinha de Balneário Camboriú, polo turístico e imobiliário.
\item \textbf{Dispersão metropolitana (periurbanização):} Senador Canedo (RM de Goiânia), Fazenda Rio Grande (RM de Curitiba), Palhoça (Grande Florianópolis) e São José de Ribamar (Grande São Luís). Esses municípios recebem a população que deixa o núcleo metropolitano em busca de terra mais barata, com loteamentos e condomínios em antigas áreas semirrurais, ou seja, com expansão do mercado imobiliário.
\end{itemize}

\textbf{2. Análise das alternativas.}

\textbf{a) Correta.} A alternativa reúne exatamente esses vetores: o \emph{front} agrícola (LEM, Sinop, Sorriso) e mineral (Parauapebas), os combustíveis fósseis (Maricá), o turismo (Camboriú) e a dispersão metropolitana para áreas semirrurais, com expansão imobiliária (Senador Canedo, Fazenda Rio Grande, Palhoça, São José de Ribamar). Ela cobre os dez casos da tabela.

\textbf{b) Incorreta.} A indústria de base (siderurgia, petroquímica etc.) não explica o crescimento recente dessas cidades. Nenhuma delas é um polo portuário que concentre sedes corporativas. Na dinâmica recente, a indústria se desconcentra, mas não é esse o vetor dominante; e as cidades do agronegócio crescem por causa da produção primária e dos serviços ligados a ela.

\textbf{c) Incorreta.} A premissa é falsa. O crime organizado e a violência se \emph{interiorizaram} nas últimas décadas, inclusive em áreas de fronteira agrícola e mineral. Além disso, a segurança pública não é um fator que explique migrações dessa escala.

\textbf{d) Incorreta.} Programas como o Minha Casa Minha Vida atuam sobretudo \emph{dentro} das próprias regiões metropolitanas e cidades, muitas vezes nas periferias. Eles não promovem uma transferência planejada de população de áreas adensadas para o interior do país. Também não explicam o crescimento de cidades do agronegócio, da mineração ou do petróleo.

\textbf{3. Conclusão.} O crescimento das cidades médias resulta da combinação entre a expansão das fronteiras econômicas (agropecuária, mineração, petróleo, turismo) e a dispersão das metrópoles para municípios vizinhos, com expansão do mercado imobiliário. Essa é a alternativa (a).

ANSWER: (A)%
\newpage%
\section*{Questão 53}

\section*{Solução}

\textbf{Leitura do mapa.} O mapa sobrepõe duas camadas: as \emph{cores} indicam os biomas e as \emph{letras} indicam as zonas de uso da terra, que são A (relativamente natural), B (uso extensivo) e C (uso intensivo). As zonas não coincidem com os limites dos biomas, então um mesmo bioma pode receber letras diferentes. Algumas letras ficam fora do território e apontam para uma área por meio de uma linha de chamada. Nesses casos, vale a cor do ponto indicado pela linha, e não a do lugar onde a letra está escrita. A questão pede o cruzamento das duas camadas para dois biomas: o deserto e a savana.

\textbf{Deserto (amarelo-claro).} É a maior mancha do mapa e cobre o centro e o oeste do continente.
\begin{itemize}
\item Há um único rótulo A, mas ele marca o enorme bloco central e ocidental, que corresponde à maior parte da área do bioma.
\item Há quatro rótulos C, todos em bordas do bioma: o litoral oeste (Carnarvon, na região de Gascoyne), o litoral da Grande Baía Australiana (a linha de chamada atravessa a estreita faixa costeira e termina no amarelo), o lobo leste no sudoeste de Queensland e o oeste de Nova Gales do Sul (região de Broken Hill). O lobo leste não é pequeno, mas, somado aos trechos litorâneos, fica bem menor que o bloco A.
\item Não há rótulo B sobre o deserto. Os B vizinhos, como os da Península de Eyre e do litoral sudoeste, apontam para a vegetação mediterrânea ou para a floresta.
\end{itemize}
A predominância de A tem explicação física e humana: clima árido, solos pobres e povoamento muito rarefeito. \emph{Conclusão:} em área, o deserto está \emph{majoritariamente} na zona relativamente natural, e só duas das três zonas aparecem nele.

\textbf{Savana (laranja).} Ocupa o norte tropical (Kimberley, Top End e Golfo de Carpentária), Pilbara, uma mancha isolada no centro, o interior de Queensland e uma faixa no sul da Austrália Ocidental. Aparecem nela as três letras:
\begin{itemize}
\item \textbf{A} em Kimberley, na Terra de Arnhem, em Pilbara, na mancha central, no interior de Queensland e na faixa do sul da Austrália Ocidental;
\item \textbf{C} no interior do Top End, no interior do norte de Queensland e, por linhas de chamada, em trechos do litoral de Kimberley e da Terra de Arnhem;
\item \textbf{B} no interior do sudeste de Queensland.
\end{itemize}
Os rótulos B e C do litoral de Queensland, perto de Cairns e de Mackay, apontam para as estreitas faixas de \emph{floresta} e por isso não entram nessa contagem. A zona A ainda é a mais extensa da savana, mas as zonas C e B ocupam porções grandes do bioma. \emph{Conclusão:} a savana está distribuída pelas três zonas, de forma \emph{mais equilibrada} que o deserto.

\textbf{Termos relativos.} A alternativa (a) é comparativa. Ela não afirma que o deserto está \emph{inteiramente} em A, e sim \emph{majoritariamente} em A, o que os C periféricos não contrariam. Também não afirma que a savana se divide em partes iguais, e sim que sua distribuição é \emph{mais equilibrada} que a do deserto. Isso se confirma: a savana tem A, B e C, e o deserto, quase todo em A, não tem nenhum B.

\textbf{Análise das alternativas.}
\begin{itemize}
\item[(a)] \textbf{Correta.} Descreve exatamente o que se concluiu acima.
\item[(b)] \textbf{Incorreta.} A alternativa inverte os biomas. A savana não está ``quase totalmente'' em A, porque tem trechos C e B. E o deserto não pode estar equilibrado entre as zonas extensiva, intensiva e natural, já que não há nenhum B sobre ele e o A domina sua área.
\item[(c)] \textbf{Incorreta.} O interior desértico é, em sua maior parte, zona A. Os C do deserto ficam em bordas. As principais áreas agrícolas e urbanas estão na periferia mais úmida do leste, do sudeste e do sudoeste, com floresta, vegetação mediterrânea e campos, onde ficam Sydney, Melbourne, Brisbane, Adelaide e Perth. Nessa periferia aparecem os rótulos B e C.
\item[(d)] \textbf{Incorreta.} O mapa mostra rótulos C sobre a savana, no Top End, no litoral de Kimberley e de Arnhem e no interior do norte de Queensland. Além disso, a justificativa é determinista: o clima tropical não \emph{impede} a agropecuária nem a infraestrutura.
\end{itemize}

\textbf{Os dados numéricos.} Segundo o enunciado, $86{,}5\%$ do continente mantém vegetação nativa e $13{,}2\%$ está desmatado, coberto por vegetação não nativa ou edificado. Esses números servem de contexto e não definem a alternativa. A vegetação nativa não equivale à zona A: ela também existe nas zonas de uso extensivo, como no pastoreio sobre a vegetação natural. Ainda assim, os números são coerentes com o mapa, em que a maior parte do território está nas zonas A e B, e não na C.

ANSWER: (A) Os desertos concentram-se majoritariamente na zona relativamente natural, enquanto as savanas apresentam distribuição mais equilibrada entre as zonas natural, extensiva e intensiva.%
\newpage%
\section*{Questão 54}

\section*{Solução}

\textbf{Critério.} Sejam $P$ a precipitação anual (entrada de água, Mapa 1) e $ETP$ a evapotranspiração potencial anual (demanda atmosférica, Mapa 2), ambas em mm/ano. No balanço hídrico climatológico, $P - ETP > 0$ indica excedente anual: sobra água para escoar, recarregar aquíferos e encher reservatórios. Já $P - ETP < 0$ indica déficit. Os mapas mostram classes (intervalos), e não valores exatos. Por isso, só se pode afirmar $P > ETP$ quando a classe de $P$ fica inteiramente acima da de $ETP$. Se as duas caem na mesma classe, os valores são apenas \emph{semelhantes} e a comparação fica indeterminada.

A média anual, porém, não mostra como a chuva se distribui nos meses. No Sudeste, de clima tropical com verão chuvoso e inverno seco, pode haver excedente no ano e mesmo assim déficit no inverno. O abastecimento depende então de armazenar o excedente do verão. Esses são justamente os fatores que o enunciado manda considerar: sazonalidade, armazenamento e demanda.

\textbf{Leitura dos mapas} (figura original ampliada; valores em mm):
\begin{center}
\small
\begin{tabular}{|l|c|p{2.9cm}|p{3.1cm}|}
\hline
RM & $P$ (Mapa 1) & $ETP$ (Mapa 2) & Comparação \\ \hline
1 -- Belo Horizonte & 1400--1600 & 1400--1600 & mesma classe: $P \approx ETP$ \\ \hline
2 -- Rio de Janeiro & 1600--1800 & oeste: 1400--1600; leste: 1600--1800 & oeste: $P > ETP$; leste: mesma classe \\ \hline
3 -- São Paulo & 1600--1800 & oeste: 1400--1600; leste: $< 1400$ & $P > ETP$ em toda a RM \\ \hline
4 -- Vitória & 1600--1800 & 1600--1800 & mesma classe: $P \approx ETP$ \\ \hline
\end{tabular}
\end{center}

\textbf{(a) Falsa.} A premissa é verdadeira: em BH, $P$ e $ETP$ estão na mesma classe (1400--1600 mm). A conclusão é que está errada. Valores semelhantes indicam balanço anual próximo de zero, sem folga de excedente, e nada dizem sobre a distribuição mensal. Com inverno seco, há déficit sazonal. Não se conclui ``disponibilidade hídrica ao longo do ano'', muito menos que os alertas possam diminuir. A margem pequena pede, ao contrário, mais atenção.

\textbf{(b) Falsa.} A premissa é só parcialmente verdadeira. A RM do Rio está dividida entre duas classes de $ETP$: no oeste as faixas são distintas ($P > ETP$), no leste coincidem. O erro decisivo está na conclusão:
(i) dados anuais não informam nada sobre o que ocorre ``ao longo do ano'';
(ii) nenhum balanço médio ``assegura'' o atendimento de uma metrópole nem ``dispensa'' alertas, porque a sazonalidade, a variabilidade entre anos e a alta demanda continuam existindo.

\textbf{(c) Verdadeira.} \emph{Premissa:} toda a RM de São Paulo está na classe de 1600--1800 mm de $P$. Sua $ETP$ está em 1400--1600 mm na porção oeste e abaixo de 1400 mm na porção leste. As classes não se sobrepõem, logo
\[
P > 1600\ \text{mm} \geq ETP \quad\Longrightarrow\quad P - ETP > 0
\]
em toda a RM. Há excedente anual: entre 0 e 400 mm no oeste e acima de 200 mm no leste.

\emph{Conclusão:} a alternativa não confunde esse excedente médio com garantia de abastecimento. Ela pondera exatamente os fatores do enunciado:
(i) a \emph{sazonalidade}, com chuvas concentradas no verão e inverno seco, quando a vazão cai;
(ii) o \emph{armazenamento} limitado dos reservatórios (sistemas como o Cantareira);
(iii) a \emph{demanda} muito alta de uma metrópole de cerca de 20 milhões de habitantes.
Daí a necessidade de monitorar os reservatórios constantemente e emitir alertas nos períodos de menor disponibilidade. A crise hídrica de 2014--2015 ilustra isso. Apesar da normal climatológica favorável, uma estiagem anômala, com chuvas de verão muito abaixo da média, somada à alta demanda e ao armazenamento limitado, levou o Cantareira ao volume morto.

\textbf{(d) Falsa.} No Mapa 1, a RM de Vitória está na mancha mais úmida do litoral capixaba (1600--1800 mm), mais chuvosa que o entorno (1200--1400 mm). No Mapa 2, está na classe de 1600--1800 mm. Como $P$ e $ETP$ caem na mesma classe, os mapas não permitem afirmar que $ETP > P$. A premissa não se sustenta, e a conclusão de ``maior vulnerabilidade'' tirada dela também não.

\textbf{Conclusão.} A (c) é a única alternativa com premissa confirmada pelos mapas ($P > ETP$ sem ambiguidade na RM de São Paulo). É também a única cuja conclusão reconhece que o balanço anual positivo não basta, por causa da sazonalidade, do armazenamento e da demanda.

ANSWER: (C)%
\newpage%
\section*{Questão 55}

\section*{Solução}

\textbf{Ideia central.} A questão pede que se diferencie uma \emph{representação simbólica} (logotipo, marca visual) de uma \emph{representação cartográfica} (mapa).

\textbf{1) O que torna uma representação um mapa.} Um mapa precisa, no mínimo, dos elementos cartográficos que permitem ler e medir o espaço:
\begin{itemize}
  \item \emph{escala}: relação de proporção entre as distâncias no papel e as distâncias no terreno;
  \item \emph{orientação}: a indicação do Norte (presente na foto aérea, ausente no logotipo);
  \item \emph{legenda}: o significado de cada cor ou símbolo;
  \item \emph{correspondência posicional}: cada elemento desenhado ocupa a posição real do objeto que representa (vias, edifícios, áreas verdes).
\end{itemize}
Só com esses elementos é possível localizar lugares, medir distâncias e planejar trajetos.

\textbf{2) Análise do logotipo.} O logotipo não tem escala, nem orientação, nem legenda. Ele faz uma \emph{estilização}: guarda apenas a ideia geral da forma do \emph{campus}, com o anel viário circular, os setores radiais e o centro. Os segmentos pretos não correspondem um a um aos edifícios e às vias reais, e os três círculos vermelhos são elementos gráficos da marca, não lugares localizáveis. A função do logotipo é \emph{identificar e simbolizar} a instituição em documentos, placas e \emph{sites}. A semelhança com a foto aérea é uma referência visual (um ``diálogo'') com o desenho urbano do \emph{campus}, não uma reprodução fiel dele.

\textbf{3) Avaliação das alternativas.}
\begin{itemize}
  \item[(a)] \textbf{Correta.} Reconhece a função simbólica e identitária do logotipo e admite a semelhança com a organização espacial do \emph{campus}, mas o classifica corretamente como uma representação gráfica simplificada, e não como um mapa.
  \item[(b)] Incorreta. O logotipo não é mapa e não permite reconhecer trajetos de circulação reais: não há escala, orientação nem correspondência entre o desenho e as vias.
  \item[(c)] Incorreta. A primeira parte está certa, mas o logotipo \emph{não} reproduz a distribuição dos edifícios e das vias e, portanto, não funciona como representação cartográfica.
  \item[(d)] Incorreta. O logotipo não é mapa e não serve para orientar deslocamentos entre as áreas do \emph{campus}.
\end{itemize}

\textbf{Conclusão.} O logotipo é um símbolo institucional inspirado na forma do \emph{campus}, e não um documento cartográfico.

ANSWER: (A)%
\newpage%
\section*{Questão 56}

\section*{Solução}

\textbf{Leitura do gráfico.} A barra superior mostra a distribuição de \emph{todos} os entrevistados: $32\%$ acham que \emph{todos} (\emph{all}) os imigrantes em situação irregular devem ser deportados, $51\%$ acham que \emph{alguns} (\emph{some}) devem ser e $16\%$ acham que \emph{nenhum} (\emph{none}) deve ser. A soma dá $99\%$ por causa de arredondamentos ou de quem não respondeu.

As barras inferiores valem \emph{apenas para o subgrupo} que apoia \emph{algumas} deportações, ou seja, os $51\%$. Cada valor é a porcentagem desse subgrupo que apoia a deportação em cada situação. Para passar um valor $p$ dessas barras para o total de entrevistados, é preciso multiplicar: $p \cdot 0{,}51$.

\textbf{Análise das alternativas.}

\begin{itemize}
\item[a)] \textbf{Falsa.} A deportação irrestrita (\emph{all}) tem o apoio de $32\%$ dos entrevistados, menos da metade e próximo de um terço.

\item[b)] \textbf{Verdadeira.} No subgrupo que apoia deportações condicionadas, apenas $14\%$ deportariam pais de crianças nascidas nos EUA e apenas $5\%$ deportariam quem é casado com cidadão americano. Logo, cerca de $86\%$ e $95\%$ desse subgrupo, respectivamente, \emph{não} deportariam essas pessoas. Nos critérios de crime, o apoio à deportação é bem maior ($97\%$ para crimes violentos e $52\%$ para crimes não violentos). Assim, esse grupo é claramente mais tolerante com a permanência de quem tem vínculos familiares no país.

\item[c)] \textbf{Falsa.} Os $52\%$ que deportariam autores de crimes não violentos (\emph{brandos}) se referem ao subgrupo dos $51\%$. Em relação ao total de entrevistados, isso equivale a
\[
0{,}52 \cdot 51\% \approx 26{,}5\%,
\]
ou seja, cerca de um quarto, e não um terço. Se forem somados também os $32\%$ que querem deportar todos, o resultado passa a ser $32\% + 26{,}5\% \approx 58{,}5\%$, que também não é um terço. O valor próximo de um terço ($32\%$) corresponde à deportação \emph{irrestrita}, e não à deportação condicionada a crimes brandos.

\item[d)] \textbf{Falsa.} Os $97\%$ também se referem apenas ao subgrupo dos $51\%$. Em relação ao total, apoiam a deportação de autores de crimes violentos, no máximo,
\[
32\% + 0{,}97 \cdot 51\% \approx 32\% + 49{,}5\% \approx 81{,}5\%.
\]
Além disso, $16\%$ dos entrevistados rejeitam qualquer deportação. Portanto, esse grupo não chega à ``quase totalidade'' dos respondentes. O erro da alternativa está em tratar a porcentagem do subgrupo como se fosse do total.
\end{itemize}

\textbf{Conclusão.} A única afirmação apoiada pelos dados é a (b): entre os que defendem deportações condicionadas, só $14\%$ e $5\%$ deportariam imigrantes com vínculos familiares (filhos nascidos nos EUA e cônjuge americano). Isso mostra grande tolerância à permanência dessas pessoas.

ANSWER: (B)%
\newpage%
\section*{Questão 57}

\section*{Solução}

\textbf{Ideia central de cada depoimento.}

\begin{itemize}
\item \textbf{Wagner Moura:} ``Any sort of art is political''; quem manda o artista ``ficar na sua'' ignora que os artistas sempre ajudaram a moldar a política. ``I think they're afraid of us'': essas pessoas não têm medo do que ele \emph{diz} (``They're not afraid of what I'm saying [\ldots]''), mas do que os artistas \emph{fazem} (``with our films, with our music, with our poetry''). Conclusão: ``art is revolutionary''. Moura não diz quem são ``these people'': elas não são identificadas.
\item \textbf{Toni Morrison:} ``art is dangerous''. A história da arte ``has always been bloody'', porque artistas foram assassinados e presos. Quem os ataca são ``dictators and people in office and people who want to control and deceive'', ou seja, ditadores, ocupantes de cargos de poder e quem quer controlar e enganar. Esses grupos sabem que os artistas vão atrapalhar seus planos, porque ``They're the ones that tell the truth''.
\end{itemize}

\textbf{Ponto em comum.} Para os dois, a arte incomoda e ameaça o poder: ela é ``revolutionary'' para Moura e ``dangerous'' para Morrison, e os poderosos têm medo dos artistas. Essa é a \emph{natureza transgressora da arte}.

\textbf{Análise das alternativas.}

\begin{itemize}
\item[a)] \textbf{Correta.} As duas partes da alternativa são verdadeiras. Os dois textos concordam que a arte é transgressora. Só Morrison diz que tipos de pessoa atacam artistas (ditadores, ocupantes de cargos públicos, quem quer controlar e enganar). Moura fala apenas em ``these people'' e ``they'', sem dizer quem são.
\item[b)] \textbf{Incorreta.} Moura afirma o \emph{contrário}: essas pessoas não temem o que ele diz no microfone, e sim o que os artistas fazem com filmes, música e poesia.
\item[c)] \textbf{Incorreta.} A alternativa usa um falso cognato: \emph{deceive} quer dizer ``enganar'', e não ``decepcionar''. Morrison fala de pessoas que querem controlar e enganar. Ela não fala de uma ``decepção'' dos artistas por terem sido perseguidos.
\item[d)] \textbf{Incorreta.} Nenhum dos dois textos trata de uma ``luta contra a punição de artistas''. Só Morrison menciona punições (assassinatos e prisões), e ela não propõe nenhuma luta contra elas. É verdade que só Moura cita o cinema (``our films''), mas a primeira parte da alternativa já a invalida.
\end{itemize}

Portanto, os depoimentos convergem quanto à natureza transgressora da arte, e só Morrison especifica quem ataca os artistas.

ANSWER: (A)%
\newpage%
\section*{Questão 58}

\section*{Solução}

\textbf{1. O que o texto descreve.} Trechos-chave do texto e suas traduções:
\begin{itemize}
  \item \emph{``plant originating in the Amazon region''}: planta nativa da Amazônia;
  \item \emph{``popular name comes from a Tupi word which alludes to its red dye''}: o nome popular vem de uma palavra tupi que se refere ao seu corante vermelho. É o \textbf{urucum} (do tupi \emph{uru-ku}, ``vermelho''), \emph{Bixa orellana};
  \item \emph{``small, teardrop-shaped capsules''}: cápsulas pequenas em forma de gota;
  \item \emph{``covered by a non-smooth coating and malleable bristles that give them a fuzzy appearance''}: casca áspera, coberta de cerdas flexíveis, com aspecto ``peludo'';
  \item \emph{``seeds can be used for body painting and protection of the skin against insect bites''}: as sementes servem para pintura corporal e para proteger a pele contra picadas de insetos, um uso tradicional do urucum entre povos indígenas.
\end{itemize}
A ilustração correta precisa mostrar, ao mesmo tempo, \emph{cor vermelha}, \emph{forma de gota} e \emph{cerdas}.

\textbf{2. Análise das alternativas.}
\begin{itemize}
  \item[a)] Fruto grande e alongado, de casca \emph{lisa} com sulcos longitudinais, parecido com o cacau. Não tem cerdas nem aspecto peludo. Falsa.
  \item[b)] Frutos \emph{redondos} e alaranjados, de casca lisa, com polpa amarela e caroço (tipo tucumã/pupunha). Não têm forma de gota nem cerdas. Falsa.
  \item[c)] Frutos \emph{vermelhos}, pequenos, em forma de \emph{gota}, cobertos de \emph{cerdas flexíveis} que dão aspecto peludo: são as cápsulas de urucum. Atendem a todos os elementos da descrição. \textbf{Verdadeira.}
  \item[d)] Cachos de frutos vermelhos abertos, com sementes pretas envoltas em arilo branco (aspecto de ``olho''), típicos do guaraná. A casca não tem cerdas. Falsa.
\end{itemize}

\textbf{3. Conclusão.} A única ilustração que reúne cor vermelha, forma de gota e cobertura de cerdas maleáveis, como diz o texto, é a da alternativa c), que representa o fruto do urucum.

ANSWER: (C)%
\newpage%
\section*{Questão 59}

\section*{Solução}

\textbf{Ideia central do texto.} Aleksic descreve um ciclo de circulação das gírias. Grupos marginalizados (pessoas negras e LGBTQIA+, a partir da \textit{ballroom scene} do fim dos anos 1980) criam palavras como \textit{slay}, \textit{serve}, \textit{mother} e \textit{gagged}. Depois, essas palavras viralizam no TikTok e são adotadas por grupos dominantes (\textit{straight white girls}, ``everybody''), que se interessam pelo seu \textit{perceived coolness}. Os trechos que respondem à questão são estes:
\begin{itemize}
  \item ``The original appeal of these words was that they \textbf{signaled in-group status}'': as gírias marcavam quem pertence ao grupo.
  \item ``marginalized groups creating slang \textbf{as a tool of power to build community and shared identity}'': as gírias são instrumento de construção de comunidade e de identidade compartilhada.
\end{itemize}
Assim, o autor explica a \emph{função} dessas gírias: delimitar um espaço próprio do grupo (território) e afirmar uma identidade cultural.

\textbf{Análise das alternativas.}
\begin{itemize}
  \item[a)] Incorreta. O autor não \emph{defende} as redes sociais. Ele mostra que elas espalham a linguagem por apropriação (``by straight white girls adopting what gay and black people are saying''), o que faz as gírias perderem a exclusividade (``they've lost that exclusivity''). O tom é descritivo e crítico, e não uma celebração da diversidade sociolinguística.
  \item[b)] \textbf{Correta.} A alternativa corresponde exatamente a ``signaled in-group status'' e ``tool of power to build community and shared identity''.
  \item[c)] Incorreta. O texto não fala de um \emph{direito} de recusar a apropriação. Ao contrário, o autor prevê que ela vai acontecer de novo (``it's going to be taken too'') dentro de um ``never-ending cycle''.
  \item[d)] Incorreta. A alternativa inverte os papéis. Quem cria gírias com valor de exclusividade são os grupos \emph{marginalizados*}. Os grupos dominantes apenas tomam essas gírias para ``capitalize on its perceived coolness'' e, com isso, acabam com a exclusividade.
\end{itemize}

\textbf{Conclusão.} Aleksic explica que algumas gírias servem para demarcar pertencimento (território) e identidade cultural de grupos marginalizados. Por isso, a alternativa correta é a (b).

ANSWER: (B) esclarece a função de algumas gírias na demarcação de território e identidade cultural.%
\newpage%
\section*{Questão 60}

\section*{Solução}

\textbf{Ideia central do texto.} O resumo diz, logo na primeira frase, que o artigo teoriza como uma nova geração de produtos de ``companheiros'' de IA \emph{``can commodify intimacy through emotionally manipulative design and racialised and gendered aesthetics''}. Em português: esses produtos \textbf{podem transformar a intimidade em mercadoria} por meio de um \emph{design} emocionalmente manipulador e de uma estética racializada e generificada. Em seguida, os autores propõem o conceito de \emph{cruel companionship} (``companheirismo cruel''): o usuário cria um apego profundo que promete intimidade, mas a própria estrutura do produto \emph{``forecloses''} (impede) uma relação de fato recíproca.

\textbf{Análise das alternativas.}

\begin{itemize}
\item[a)] \textbf{Incorreta.} Economia política e estudos culturais são as \emph{lentes de análise} usadas pelos autores (\emph{``By analysing AI companions through the lens of political economy and cultural studies''}), e não o assunto dos estereótipos. Os estereótipos reproduzidos pelos assistentes são \emph{raciais e de gênero}, ligados à servidão e à docilidade.

\item[b)] \textbf{Incorreta.} O texto afirma o contrário: a dinâmica desses produtos \emph{``structurally foreclose[s] the possibility of genuinely reciprocal relationships''}, ou seja, impede a reciprocidade. Além disso, o ``companheirismo cruel'' é um conceito criado pelos autores para descrever essas dinâmicas, não algo que os assistentes ``reforçam''.

\item[c)] \textbf{Correta.} É uma paráfrase direta da primeira frase: \emph{``AI companion products can commodify intimacy through emotionally manipulative design''}. ``A forma como são desenvolvidos'' corresponde a \emph{design}, ``tem o potencial'' corresponde a \emph{can}, e ``transformar a intimidade em mercadoria'' corresponde a \emph{commodify intimacy}.

\item[d)] \textbf{Incorreta.} A primeira parte (marcas de raça e gênero) está de acordo com o texto, mas a conclusão o contradiz. Essas identidades \emph{``draw on longstanding stereotypes of servitude and docility''}, isto é, \emph{recorrem} a estereótipos de servidão e docilidade, que os produtos reproduzem e reembalam (\emph{``repackage''}) no meio digital, em vez de combatê-los.
\end{itemize}

\textbf{Conclusão.} A única alternativa fiel ao texto é a (c): segundo o resumo, o \emph{design} dos companheiros de IA pode transformar a intimidade em mercadoria.

\medskip
ANSWER: (C) A forma como assistentes de IA são desenvolvidos tem o potencial de transformar a intimidade em mercadoria.%
\newpage%
\section*{Questão 61}

\section*{Solução}

\textbf{Ideia central do poema.} O texto opõe dois planos: o \emph{passado colonial de violência} contra a África (exploração, apagamento, mortes) e o \emph{presente de vida e resistência} (a África que respira, gera, dança, e a eu lírica que carrega sua memória e ``se ergue''). Essa oposição aparece sobretudo nos \textbf{tempos verbais}.

\textbf{1) Passado e particípio $\to$ trauma.}
\begin{itemize}
  \item \textit{``a continent who's \textbf{carried} the cross of the world''} (present perfect: um sofrimento acumulado ao longo da história);
  \item \textit{``Gold \textbf{stripped} from her rivers, / names \textbf{erased} from her maps. / Bones \textbf{buried} beneath railroads and sugarcane fields''} (particípios: saque, apagamento, mortes no trabalho forçado);
  \item \textit{``for those who \textbf{burned}. For those who \textbf{bled}, for those who \textbf{built} and were \textbf{buried}''} (passado simples e voz passiva: as vítimas).
\end{itemize}

\textbf{2) Presente simples $\to$ resistência.}
\begin{itemize}
  \item \textit{``\textbf{Yet} she \textbf{breathes} / she \textbf{births} / she \textbf{dances} with dust in her ankles and fire in her lungs''} (o conectivo \textit{yet} marca a virada: apesar de tudo, ela vive);
  \item \textit{``I \textbf{carry} her in my marrow'', ``I \textbf{am} her memory'', ``So I \textbf{wear} it for her'', ``It \textbf{pulses} in my blood''};
  \item \textit{``And still / we \textbf{rise} / crowned''} (eco de \textit{Still I Rise}, de Maya Angelou, poema de resistência).
\end{itemize}
O contraste fica explícito no verso \textit{``She's \textbf{carried} the cross of the world but \textbf{has not worn} its crown. / So I \textbf{wear} it for her''}: o passado é de dor; o presente, de reivindicação.

\textbf{Análise das alternativas.}
\begin{itemize}
  \item[(a)] Falsa. O sofrimento é representado pela \emph{cruz} (\textit{the cross of the world}). A \emph{coroa} é o oposto: dignidade e legado que não podem ser tirados (\textit{``Not gold but grace. / Not jewels but justice. / A crown that cannot be taken''}).
  \item[(b)] Falsa. O presente simples não expressa a persistência da opressão colonial. Ele expressa a vitalidade e a resistência (\textit{breathes, dances, wear, rise}). As violências coloniais aparecem no passado e no particípio.
  \item[(c)] Falsa. Há, sim, personificação (\textit{she breathes, she births, her ankles, her lungs, the womb of the world}), mas a África é um corpo vivo e materno que \emph{não} perdeu suas raízes. A própria eu lírica as preserva: \textit{``I carry her in my marrow''}, \textit{``I am her memory''}.
  \item[(d)] Verdadeira. É exatamente a alternância mostrada nos itens 1 e 2: as formas de passado e particípio (\textit{carried, stripped, erased, buried, burned, bled}) registram os traumas, e as de presente simples (\textit{breathes, births, dances, carry, wear, rise}) mostram a resistência.
\end{itemize}

ANSWER: (D)%
\newpage%
\section*{Questão 62}

\section*{Solução}

\textbf{Ideia central do texto.} Francione e Charlton apresentam o princípio da \emph{Abordagem Abolicionista}: todo ser senciente, \emph{human or nonhuman}, tem o direito básico de não ser tratado como propriedade de outros. Os autores sustentam essa tese com três afirmações:
\begin{enumerate}
  \item ``There is absolutely no doubt that the animals we routinely exploit are sentient'', ou seja, não há dúvida de que os animais explorados são sencientes;
  \item ``All sentient beings have at least two interests: the interest in not suffering and the interest in not dying'', ou seja, todo ser senciente tem, no mínimo, os interesses de não sofrer e de não morrer;
  \item ``the use of animals as property for food, clothing, and other purposes implicates at least two interests that animals have'', ou seja, usar animais como propriedade atinge esses dois interesses.
\end{enumerate}
No fim do trecho, os autores dizem que a exploração é moralmente injustificável \emph{however ``humane''} (por mais ``humanitária'' que seja) e que se devem julgar \emph{ações}, não \emph{pessoas}.

\textbf{Análise das alternativas.}

\textbf{a) Falsa.} O texto não admite animais explorados que não sintam. Ao contrário, afirma que eles são, sem nenhuma dúvida, sencientes. É justamente da senciência que vem o direito de não ser propriedade (``all sentient beings \ldots\ have the basic right''). Por isso, a concessão ``mesmo que alguns deles não sejam capazes de sentir'' contradiz o fundamento do argumento.

\textbf{b) Falsa.} O texto não fala em \emph{criminalizar}. A obrigação citada é a de ``make clear'' (deixar claro) que a exploração é \emph{moralmente} injustificável. Além disso, a justificativa ``porque essa é uma prática desumana'' inverte o argumento: para os autores, a exploração é injustificável \emph{mesmo quando} é ``humane''. O problema está no próprio uso dos animais como propriedade, e não na falta de humanidade com que eles são tratados.

\textbf{c) Falsa.} O texto diz o contrário: ``it's not a matter of judging others, but judging actions. The former is never a good idea; the latter is essential''. Julgar pessoas ``nunca é uma boa ideia''. O que é essencial para o movimento por justiça social para os animais é julgar as \emph{ações}.

\textbf{d) Verdadeira.} A alternativa junta corretamente as afirmações 2 e 3. Como os interesses de não sofrer e de não morrer pertencem a \emph{todos} os seres sencientes, humanos ou não, eles são comuns a humanos e animais. Segundo o texto, o uso de animais como propriedade (``como produto'', para comida, vestuário etc.) ``implicates'' (atinge, fere) esses dois interesses.

\textbf{Conclusão.} Só a alternativa (d) está de acordo com o texto: o uso dos animais como propriedade fere os dois interesses que todo ser senciente, humano ou não, tem em comum, o de viver e o de não sofrer.

ANSWER: (D) O uso de animais como produto fere dois interesses que humanos e animais têm em comum: o interesse em viver e em não sofrer.%
\newpage%
\section*{Questão 63}

\section*{Solução}

\textbf{Hipóteses.} O enunciado diz que a onda gravitacional se propaga com velocidade muito próxima à da luz. Vamos então tratar a propagação como movimento uniforme com velocidade $v \approx c$ ao longo do caminho até a Terra. O tempo de viagem é $\Delta t = 1{,}3$ bilhão de anos, e $1$ bilhão $= 10^9$. Nesse nível de modelo, desprezamos a expansão do Universo: a distância obtida é a distância percorrida pela onda no tempo de viagem.

\textbf{1) Conversão do tempo para segundos.} Usamos $1\ \text{ano}\approx 3{,}2\times10^{7}\ \text{s}$:
\[
\Delta t = 1{,}3\times10^{9}\ \text{anos}\times 3{,}2\times10^{7}\ \frac{\text{s}}{\text{ano}}
= 4{,}16\times10^{16}\ \text{s}.
\]

\textbf{2) Movimento uniforme.} Com velocidade constante, $d = v\,\Delta t$:
\[
d = c\,\Delta t = \left(3{,}0\times10^{8}\ \frac{\text{m}}{\text{s}}\right)\left(4{,}16\times10^{16}\ \text{s}\right)
= 12{,}48\times10^{24}\ \text{m}\approx 1{,}2\times10^{25}\ \text{m}.
\]
Quanto às unidades, $\dfrac{\text{m}}{\text{s}}\cdot\text{s}=\text{m}$, como esperado.

\textbf{Verificação da ordem de grandeza.} Com os dados da questão, um ano-luz vale
$3{,}0\times10^{8}\cdot 3{,}2\times10^{7}\approx 9{,}6\times10^{15}\ \text{m}\sim10^{16}\ \text{m}$.
Então $1{,}3\times10^{9}$ anos-luz correspondem a
$1{,}3\times10^{9}\cdot 9{,}6\times10^{15}\ \text{m}\approx 1{,}2\times10^{25}\ \text{m}$, o que confirma o resultado.

\textbf{Análise das alternativas.}
\begin{itemize}
\item (b) $3{,}9\times10^{17}$ m vem de $3{,}0\times10^{8}\cdot1{,}3\times10^{9}$, ou seja, do esquecimento da conversão de anos para segundos. O resultado tem unidade incoerente.
\item (d) $3{,}9\times10^{14}$ m vem de tomar ``bilhão'' como $10^{6}$ e também esquecer a conversão de anos para segundos.
\item (c) $1{,}2\times10^{22}$ m vem de fazer a conversão para segundos corretamente, mas tomar ``bilhão'' como $10^{6}$ em vez de $10^{9}$.
\item (a) $1{,}2\times10^{25}$ m é o único valor que usa $1$ bilhão $=10^{9}$ e converte o tempo para segundos. Por isso é a alternativa correta.
\end{itemize}

ANSWER: (a) $1{,}2\times10^{25}$ m%
\newpage%
\section*{Questão 64}

\section*{Solução}

\textbf{Hipóteses.} Tratamos a Terra como uma esfera de massa $M$ e raio $R$ com distribuição de massa esfericamente simétrica. Pelo teorema das cascas de Newton, para um ponto externo ou na superfície, ela atrai como se toda a massa estivesse concentrada no centro. Desprezamos a rotação da Terra, como faz o enunciado, que pede apenas a parte gravitacional.

\textbf{1) Lei da gravitação universal.} Um corpo de massa $m$ na superfície fica à distância $R$ do centro da Terra. O módulo da força que a Terra exerce sobre ele é
\[
F = \frac{G\,M\,m}{R^{2}},
\]
em que $G$ é a constante da gravitação universal.

\textbf{2) Segunda lei de Newton.} Se essa é a única força considerada, a aceleração adquirida pelo corpo, que é a aceleração da gravidade, vale
\[
g = \frac{F}{m} = \frac{G\,M\,m}{m\,R^{2}} \quad\Longrightarrow\quad g = \frac{G\,M}{R^{2}}.
\]
A massa $m$ aparece tanto na força gravitacional quanto na inércia, e por isso se cancela. É a igualdade entre massa gravitacional e massa inercial: todos os corpos caem com a mesma aceleração.

\textbf{3) Dependências de $g$.}
\begin{itemize}
\item $g$ não depende de $m$;
\item $g \propto M$: dobrando $M$, $g$ dobra;
\item $g \propto \dfrac{1}{R^{2}}$: dobrando $R$, $g$ cai para $\dfrac{1}{4}$ do valor.
\end{itemize}

\textbf{Verificações.}
\begin{itemize}
\item Unidades: $[G\,M/R^{2}] = \dfrac{\mathrm{N\,m^{2}\,kg^{-2}}\cdot \mathrm{kg}}{\mathrm{m^{2}}} = \mathrm{N/kg} = \mathrm{m/s^{2}}$, que é unidade de aceleração.
\item Ordem de grandeza: com $G \approx 6,67\times 10^{-11}\ \mathrm{N\,m^{2}/kg^{2}}$, $M \approx 5,97\times 10^{24}\ \mathrm{kg}$ e $R \approx 6,37\times 10^{6}\ \mathrm{m}$, obtém-se
\[
g \approx \frac{6,67\times 10^{-11}\cdot 5,97\times 10^{24}}{\left(6,37\times 10^{6}\right)^{2}} \approx 9,8\ \mathrm{m/s^{2}},
\]
que é o valor conhecido.
\end{itemize}

\textbf{Análise das alternativas.}
\begin{itemize}
\item (a) Errada: diz que $g$ é proporcional a $m$ e a $R^{2}$. Isso é falso.
\item (b) Correta: independente de $m$, proporcional a $M$ e inversamente proporcional a $R^{2}$.
\item (c) Errada: diz que $g$ é proporcional a $m$. Isso vale para a \emph{força} (o peso), não para a aceleração.
\item (d) Errada: diz que $g$ é diretamente proporcional a $R^{2}$, o que é falso.
\end{itemize}

ANSWER: (B) independente de $m$, diretamente proporcional a $M$ e inversamente proporcional ao quadrado do raio da Terra ($g = GM/R^{2}$)%
\newpage%
\section*{Questão 65}

\section*{Solução}

\textbf{Hipóteses.} Fios ideais (inextensíveis e de massa desprezível) e polias ideais (sem massa e sem atrito). Assim, a tração é a mesma em todos os pontos de um mesmo fio, e as polias fixas só mudam a direção da força. Na iminência de o bloco ser levantado, o sistema ainda está em repouso, então vale o equilíbrio estático, $\sum \vec F = \vec 0$, para cada corpo.

\textbf{1) Bloco (iminência de subir).} Sobre o bloco atuam o peso $mg$ (para baixo), a tração $T$ do fio (para cima) e a normal $N$ do piso. Na iminência de o bloco perder o contato com o piso, $N = 0$. Logo,
\[
T = mg = 12 \cdot 10 = 120\ \text{N}.
\]

\textbf{2) Polia móvel.} O segundo fio passa por cima da polia móvel e tem tração $T$ nos dois ramos: um ramo está preso à base e o outro sustenta o bloco. Como esse fio é ideal, a tração é a mesma ($T$) nos dois ramos, como indicam as setas $\vec T$ da figura. Ambos os ramos puxam a polia para baixo. Do outro lado, o primeiro fio, ligado ao eixo da polia, a puxa para cima com tração $T_1$. Como a polia tem massa desprezível e está em equilíbrio,
\[
T_1 = 2T = 2 \cdot 120 = 240\ \text{N}.
\]

\textbf{3) Polias fixas e manopla.} O primeiro fio vai do eixo da polia móvel até a manopla passando pelas duas polias fixas. Como fio e polias são ideais, a tração ao longo dele permanece $T_1$. Na extremidade livre, a pessoa aplica $F$ e o fio está em equilíbrio, portanto
\[
F = T_1 = 2mg = 240\ \text{N}.
\]

\textbf{Verificação.} As unidades conferem: $\text{kg}\cdot\text{m/s}^2 = \text{N}$. Quanto à ordem de grandeza, $240$ N equivale ao peso de cerca de $24$ kg, valor plausível para um aparelho de musculação. O resultado também faz sentido físico. Aqui a carga está pendurada num dos ramos do fio que passa pela polia móvel, e não presa ao eixo dela. Já o fio da manopla puxa justamente o eixo. Por isso, em vez de reduzir o esforço à metade, a montagem o \emph{dobra}: é uma desvantagem mecânica proposital, o que é conveniente para exercitar a musculatura.

\textbf{Análise das alternativas.} A alternativa (b), $60$ N, aplica sem critério a ideia de que a polia móvel reduz a força à metade. A (c), $120$ N, ignora a polia móvel e iguala $F$ ao peso. A (a), $12$ N, confunde massa com peso. A única compatível com o equilíbrio é a (d).

ANSWER: (D) 240 N%
\newpage%
\section*{Questão 66}

\section*{Solução}

\textbf{1) Tipo de atrito.} O enunciado pede a força de atrito \emph{durante o deslocamento} do disco. Se há movimento relativo entre o disco e o gelo, o atrito é \textbf{cinético}. O coeficiente estático $\mu_E$ só vale quando não há escorregamento (disco em repouso ou prestes a deslizar), e então determina apenas o valor máximo do atrito estático. Logo, $\mu_E = 0{,}1$ é um dado que não entra no cálculo.

\textbf{2) Força normal.} A pista é plana e horizontal, e as únicas forças verticais sobre o disco são o peso e a normal. Como não há aceleração vertical, pela 2ª lei de Newton nessa direção:
\[
N = P = m g = 20 \cdot 10 = 200 \ \text{N}.
\]

\textbf{3) Atrito cinético.} Pelo modelo de Coulomb, o módulo do atrito cinético é
\[
f_{at} = \mu_C \, N = 0{,}02 \cdot 200 = 4{,}0 \ \text{N}.
\]
O atrito tem sentido oposto ao da velocidade do disco em relação ao gelo e não depende da velocidade nem da área de contato.

\textbf{4) Conferência e análise das alternativas.}
\begin{itemize}
\item Unidade: $[\mu_C]\cdot[N] = 1 \cdot \text{N} = \text{N}$. O valor $4\ \text{N} \ll P = 200\ \text{N}$ é coerente com o gelo, uma superfície muito escorregadia: o atrito é apenas $2\%$ do peso. A desaceleração vale $a = \mu_C g = 0{,}2\ \text{m/s}^2$, uma ordem de grandeza razoável para um disco que desliza dezenas de metros.
\item (a) $0{,}4$ N $= \mu_C\, m$: esquece a gravidade, e o resultado sai em unidade errada.
\item (b) $2{,}0$ N $= \mu_E\, m$: usa o coeficiente errado e também esquece $g$.
\item (d) $20$ N $= \mu_E\, m g$: seria o atrito estático \emph{máximo}, que só vale na iminência de movimento, e não durante o deslizamento.
\item (c) $4{,}0$ N $= \mu_C\, m g$: é a única que usa o coeficiente cinético, adequado a um disco em movimento, com a normal correta. Portanto, é a alternativa correta.
\end{itemize}

ANSWER: (C) 4,0 N%
\newpage%
\section*{Questão 67}

\section*{Solução}

\textbf{1. O que o gráfico precisa mostrar.} A amostra recebe calor do \emph{chip}, portanto sua temperatura não pode diminuir. Isso já elimina a alternativa (d), cuja temperatura cai de $20$ para $0\ ^\circ\mathrm{C}$. O processo tem três etapas:
\begin{itemize}
\item aquecimento do sólido de $T_1 = 20\ ^\circ\mathrm{C}$ até $T_2 = 30\ ^\circ\mathrm{C}$ (calor sensível, a temperatura sobe);
\item fusão a $T_2 = 30\ ^\circ\mathrm{C}$. Para uma substância pura sob pressão constante, a mudança de fase ocorre a temperatura constante, e todo o calor recebido é latente: $Q = m\,L$. No gráfico, isso aparece como um \emph{patamar horizontal}. O valor de cerca de $30\ ^\circ\mathrm{C}$ é coerente com o ponto de fusão real do gálio, de aproximadamente $29{,}8\ ^\circ\mathrm{C}$;
\item aquecimento do líquido de $T_2$ até $T_3$ (calor sensível), terminando em $t = 230\ \mathrm{s}$.
\end{itemize}
A alternativa (c) é uma reta sem patamar, ou seja, ignora a fusão. Além disso, ela chega a $50\ ^\circ\mathrm{C}$ em $t \approx 220\ \mathrm{s}$, e não em $230\ \mathrm{s}$. Por isso também é descartada.

\textbf{2. Cálculo de $T_3$.} Na fase líquida não há mudança de fase, e vale a equação fundamental da calorimetria:
\[
Q = m_{Ga}\,c_{liq}\,(T_3 - T_2)
\quad\Longrightarrow\quad
T_3 - T_2 = \frac{Q}{m_{Ga}\,c_{liq}} = \frac{80\ \mathrm{J}}{(10\ \mathrm{g})\cdot\left(0{,}4\ \mathrm{J/(g\cdot {}^\circ C)}\right)} = \frac{80}{4}\ ^\circ\mathrm{C} = 20\ ^\circ\mathrm{C}.
\]
Conferência de unidades: $\mathrm{J} / \left(\mathrm{g}\cdot \mathrm{J/(g\cdot{}^\circ C)}\right) = {}^\circ\mathrm{C}$. Logo
\[
T_3 = 30 + 20 = 50\ ^\circ\mathrm{C}.
\]

\textbf{3. Escolha entre (a) e (b).} As duas alternativas têm a forma correta: aquecimento de $20$ a $30\ ^\circ\mathrm{C}$ entre $0$ e $10\ \mathrm{s}$, patamar em $30\ ^\circ\mathrm{C}$ até $210\ \mathrm{s}$ e aquecimento do líquido até $t = 230\ \mathrm{s}$. A diferença está na temperatura final.
\begin{itemize}
\item Em (a), $T_3 = 120\ ^\circ\mathrm{C}$. Isso exigiria $\Delta T = 90\ ^\circ\mathrm{C}$, ou seja, $Q = 10\cdot 0{,}4\cdot 90 = 360\ \mathrm{J} \neq 80\ \mathrm{J}$. Está incorreta.
\item Em (b), $T_3 = 50\ ^\circ\mathrm{C}$ em $t = 230\ \mathrm{s}$, exatamente o valor calculado. Está correta.
\end{itemize}
Como a potência do \emph{chip} é variável, as inclinações dos trechos não precisam ser comparadas entre si. Para decidir, bastam a forma qualitativa (subida, patamar, subida) e os valores $T_1$, $T_2$, $T_3$ e $t = 230\ \mathrm{s}$.

\textbf{Conclusão.} O gráfico que representa o processo é o da alternativa (b). Ele mostra aquecimento do sólido, patamar de fusão em $30\ ^\circ\mathrm{C}$ e aquecimento do líquido até $T_3 = 50\ ^\circ\mathrm{C}$ em $t = 230\ \mathrm{s}$.

ANSWER: (B) gráfico com patamar de fusão em $30\ ^\circ\mathrm{C}$ e $T_3 = 50\ ^\circ\mathrm{C}$ em $t = 230\ \mathrm{s}$%
\newpage%
\section*{Questão 68}

\section*{Solução}

\textbf{Dados:} potência consumida pelo roteador ocioso $P = 3{,}0\ \text{W}$; tensão de alimentação $U = 12\ \text{V}$ (contínua, fornecida pela fonte do roteador).

\textbf{Lei utilizada.} A potência elétrica entregue a qualquer elemento de circuito é o produto da diferença de potencial entre seus terminais pela corrente que o atravessa:
\[
P = U\,i .
\]
Essa relação vem da definição de potência como energia por unidade de tempo: a carga $\Delta q$ que atravessa a ddp $U$ recebe energia $U\,\Delta q$, então $P = U\,\Delta q/\Delta t = U\,i$. Ela vale para qualquer dispositivo, ôhmico ou não. Por isso, não é preciso conhecer a resistência interna do roteador nem supor que ele seja ôhmico.

\textbf{Cálculo.} Isolando a corrente:
\[
i = \frac{P}{U} = \frac{3{,}0\ \text{W}}{12\ \text{V}} = 0{,}25\ \text{A}.
\]

\textbf{Unidades:} $\dfrac{\text{W}}{\text{V}} = \dfrac{\text{J/s}}{\text{J/C}} = \dfrac{\text{C}}{\text{s}} = \text{A}$, como esperado.

\textbf{Verificação:} $U\,i = 12 \cdot 0{,}25 = 3{,}0\ \text{W}$, que é a potência dada. A ordem de grandeza também faz sentido: uma corrente de alguns décimos de ampère é típica de pequenos aparelhos eletrônicos alimentados por fonte de 12 V.

\textbf{Análise das alternativas:}
\begin{itemize}
\item a) $0{,}03\ \text{A}$: é cerca de 8 vezes menor que o valor correto e daria $P = 12 \cdot 0{,}03 = 0{,}36\ \text{W} \neq 3{,}0\ \text{W}$.
\item b) $0{,}25\ \text{A}$: é igual a $P/U$. \textbf{Correta.}
\item c) $4{,}0\ \text{A}$: é a razão invertida, $U/P = 12/3$, e daria $P = 48\ \text{W}$.
\item d) $36\ \text{A}$: é o produto $P\cdot U = 3\cdot 12$, que nem tem unidade de corrente.
\end{itemize}

Portanto, a corrente no roteador ocioso é $i = 0{,}25\ \text{A}$.

ANSWER: (B) 0,25 A%
\newpage%
\section*{Questão 69}

\section*{Solução}

\textbf{Hipóteses.} A bateria é ideal (sem resistência interna) e os fios e a chave têm resistência desprezível. Os resistores são ôhmicos e iguais:
\[
R_1=R_2=R_3=R_4=R=1{,}2\ \text{k}\Omega = 1200\ \Omega .
\]

\textbf{1. Configuração do circuito com o alarme ativado.} Pela figura, os quatro ramos ligam os mesmos dois nós, que estão ligados diretamente aos terminais da bateria. Logo, os ramos estão em paralelo e cada um fica submetido à tensão $U_b$. Com o dispositivo intacto, a chave do ramo de $R_4$ está aberta e só $R_1$, $R_2$ e $R_3$ conduzem. Quando o alarme é ativado, a chave se fecha e $R_4$ também passa a conduzir: são \emph{quatro} resistores em paralelo.

\textbf{2. Resistência equivalente.} Para resistores em paralelo,
\[
\frac{1}{R_{eq}}=\frac{1}{R_1}+\frac{1}{R_2}+\frac{1}{R_3}+\frac{1}{R_4}=\frac{4}{R}
\quad\Longrightarrow\quad
R_{eq}=\frac{R}{4}=\frac{1200\ \Omega}{4}=300\ \Omega .
\]

\textbf{3. Corrente total (lei de Ohm).}
\[
i=\frac{U_b}{R_{eq}}=\frac{9{,}0\ \text{V}}{300\ \Omega}=3{,}0\times 10^{-2}\ \text{A}=30\ \text{mA}.
\]

\textbf{Conferência.} Cada ramo recebe $9{,}0$ V, portanto
\[
i_k=\frac{U_b}{R}=\frac{9{,}0\ \text{V}}{1200\ \Omega}=7{,}5\ \text{mA}.
\]
Pela lei dos nós, $i=4\cdot 7{,}5\ \text{mA}=30\ \text{mA}$, o que confirma o resultado. A unidade está correta, pois $\text{V}/\Omega=\text{A}$. A ordem de grandeza, dezenas de mA, é coerente com uma resistência equivalente de centenas de ohms sob 9 V.

\textbf{Por que as outras alternativas estão erradas.}
\begin{itemize}
\item (a) $1{,}9$ mA $\approx 9{,}0/4800$ A: corresponde a ligar os quatro resistores em série, o que não é o caso.
\item (b) $7{,}5$ mA: é a corrente em um único ramo, e não a corrente total.
\item (c) $22$ mA $\approx 9{,}0/400$ A: corresponde a apenas três ramos em paralelo, ou seja, à situação \emph{sem} violação, com a chave aberta.
\end{itemize}
A alternativa correta é a (d), pois com a chave fechada os quatro ramos conduzem e $i=30$ mA.

ANSWER: (D) 30 mA%
\newpage%
\section*{Questão 70}

\section*{Solução}

\textbf{1. A tese do trecho.} Habermas define os dois princípios da ética do discurso:
\begin{itemize}
\item \emph{justiça}: postula \emph{direitos iguais para cada indivíduo} e diz respeito à liberdade subjetiva de indivíduos inalienáveis;
\item \emph{solidariedade}: reclama empatia e cuidado com o bem-estar do próximo e se refere às partes \emph{irmanadas numa forma de vida partilhada intersubjetivamente}.
\end{itemize}
A seguir ele afirma que os dois princípios \emph{remontam à mesma raiz da moral}, que é a vulnerabilidade dos seres vivos. Esses seres \emph{só se particularizam em indivíduos por ação da socialização}. Daí ele conclui que a moral, ao defender os direitos do indivíduo, \emph{é obrigada a defender igualmente o bem-estar da comunidade}. Portanto os dois princípios são complementares e estão ligados entre si: um não se sustenta sem o outro.

\textbf{2. Análise das alternativas.}

\textbf{a) Falsa.} O texto diz o contrário: a ética do discurso \emph{explica} por que justiça e solidariedade têm uma raiz comum. Não há ``contradição insolúvel'' assumida. O que existe é uma articulação necessária entre os dois princípios.

\textbf{b) Falsa.} Habermas não diz que a modernidade impede a vida ética. Ele apenas caracteriza a justiça ``em sentido moderno'' pela liberdade subjetiva e, no mesmo período, põe ao lado dela a solidariedade. A liberdade individual não exclui a solidariedade; as duas se exigem mutuamente.

\textbf{c) Falsa.} A alternativa inverte o texto de duas maneiras. Primeiro, a ética do discurso não se opõe à vulnerabilidade humana: parte dela, como raiz da moral. Segundo, a particularização do indivíduo não é independente da socialização. Ela ocorre \emph{somente} ``por ação da socialização''.

\textbf{d) Verdadeira.} A relação entre os princípios é estreita, pois eles têm a mesma raiz e defender um obriga a defender o outro. A divisão de funções da alternativa coincide com o texto. A justiça garante direitos iguais e a liberdade subjetiva de cada indivíduo. A solidariedade cuida do bem-estar das partes ligadas por uma forma de vida partilhada intersubjetivamente, ou seja, das relações intersubjetivas.

\textbf{3. Conclusão.} Só a alternativa (d) traduz a tese central do trecho. Para Habermas, como o indivíduo se forma na socialização, proteger o indivíduo (justiça) e proteger a comunidade e os vínculos intersubjetivos (solidariedade) são duas faces da mesma exigência moral.

ANSWER: (D) Na ética do discurso há uma relação estreita entre o princípio da justiça (preserva os direitos dos indivíduos) e o da solidariedade (preserva as relações intersubjetivas).%
\newpage%
\section*{Questão 71}

\section*{Solução}

\textbf{Ideia central dos textos.} Para Simone Weil, a \emph{força} é ``aquilo que transforma quem quer que lhe seja submetido em uma coisa''. Esse efeito atinge os dois lados da relação de força:
\begin{itemize}
  \item \textbf{quem a sofre}: a força que mata faz do homem um cadáver, uma coisa no sentido literal. A força que não mata transforma o homem ``que continua vivo em pedra'': ele ``tem uma alma'', mas ``é uma coisa'';
  \item \textbf{quem a exerce}: ``com a mesma dureza com que a força esmaga os vencidos, embriaga aquele que a possui, ou julga possuí-la''. Embriagado, o agente também perde o domínio de si e se torna instrumento da força. Por isso Weil conclui: ``Ninguém a possui verdadeiramente''.
\end{itemize}
No Texto 2, a fala de Aquiles aparece como exceção (``palavras razoáveis''). Ela reconhece que bens materiais podem ser obtidos de novo, mas que a vida humana perdida é irrecuperável. É um momento de lucidez diante da ``idolatria da força''.

\textbf{Análise das alternativas.}

\begin{enumerate}
  \item[a)] \textbf{Incorreta.} A alternativa supõe que os vencedores ``detêm'' a força. Weil afirma o contrário: o vencedor apenas ``julga possuí-la'', e ``ninguém a possui verdadeiramente''. Além disso, o vencedor não fica de fora do processo, porque também é ``embriagado'' pela força.
  \item[b)] \textbf{Incorreta.} Weil chama as palavras de Aquiles de ``razoáveis''. Elas mostram justamente consciência e humanidade, pois reconhecem o valor insubstituível da vida. Não mostram ``falta de humanidade''.
  \item[c)] \textbf{Correta.} A força reifica quem a sofre, que vira cadáver ou ``pedra'' viva. Ela também desumaniza quem a usa, que é ``embriagado'' por ela e vira seu instrumento, sem jamais possuí-la de fato. É exatamente a tese do Texto 1: a força desumaniza os dois polos.
  \item[d)] \textbf{Incorreta.} Bois e ovelhas aparecem na fala de Aquiles como \emph{bens} que podem ser obtidos ou recuperados, em oposição à vida humana, que, ``uma vez perdida, não pode ser recuperada''. O contraste trata da irreversibilidade da perda da vida e da recusa ao culto da força. Não há uma hierarquia entre o ser humano e os outros animais.
\end{enumerate}

\textbf{Conclusão.} A única alternativa fiel aos dois textos é a (c). Para Weil, a força transforma em coisa tanto quem a sofre quanto quem a exerce.

ANSWER: (C) A força desumaniza aqueles sobre os quais recai e aqueles que dela se utilizam para transformar seres humanos em coisas.%
\newpage%
\section*{Questão 72}

\section*{Solução}

\textbf{Ideia central do texto.} Para Dussel, a Modernidade foi \emph{gestada} nas cidades medievais europeias, mas só \emph{nasceu} em 1492, quando a Europa ``pôde se confrontar com o seu Outro e controlá-lo, vencê-lo, violentá-lo, colonizá-lo''. Nesse momento, ela se definiu ``como um \emph{ego} descobridor, conquistador, colonizador''. Logo, a subjetividade moderna não nasce só de um ato de pensamento. Ela nasce também de uma relação de dominação sobre a Alteridade (o indígena, o africano, o colonizado), que por isso é chamada de ``constitutiva da própria Modernidade''. O título fala em \emph{encobrimento}, e não em ``descobrimento'', do outro: o Outro não foi reconhecido, foi negado e dominado. A proposta final de Dussel, uma filosofia do diálogo dentro de uma filosofia da libertação do oprimido, aparece como \emph{tarefa} a cumprir, e não como algo que a Modernidade já tenha feito.

\textbf{Análise das alternativas.}

\begin{itemize}
\item[(a)] \textbf{Correta.} Na filosofia, o marco tradicional do sujeito moderno é o \textit{cogito} cartesiano (\textit{Penso, logo existo}). Dussel acrescenta que o \emph{ego} moderno também se forma como \emph{ego} ``descobridor, conquistador, colonizador''. Em sua obra, ele chama isso de \textit{ego conquiro}, que vem antes do \textit{ego cogito}. Dizer que o \emph{ego} moderno se constitui ``não apenas'' pelo \textit{cogito}, ``mas também'' pela conquista e colonização do Outro resume fielmente a tese do texto.

\item[(b)] \textbf{Incorreta.} O texto não diz que a filosofia moderna foi um diálogo com o Outro. A relação descrita é de controle, violência e colonização. O diálogo é justamente o que \emph{ainda falta} construir (``a tarefa, portanto, é desenvolver uma filosofia do diálogo''). Além disso, Dussel critica a ideia de ``descoberta'': o Outro foi \emph{encoberto}.

\item[(c)] \textbf{Incorreta.} Segundo o texto, a Modernidade não surge ``dos escombros da Idade Média''. Ela se origina nas ``cidades europeias medievais, livres, centros de enorme criatividade''. O par ``progresso e retrocesso'' e a ideia de uma contradição ``por definição'' não estão no texto.

\item[(d)] \textbf{Incorreta.} A filosofia da libertação não é apresentada como uma mediação neutra entre europeus e colonizados para ``dirimir'' (resolver, apaziguar) a violência colonial. Ela é a filosofia ``da libertação do oprimido, do incomunicado, do excluído, do Outro''. Ela parte, portanto, do ponto de vista da vítima e busca sua emancipação. O diálogo é \emph{parte} desse projeto, e não sua finalidade de conciliação.
\end{itemize}

\textbf{Conclusão.} Só a alternativa (a) expressa a tese de Dussel: o sujeito moderno se constitui também, e de forma decisiva, pela conquista e colonização do Outro, e não apenas pelo \textit{cogito} cartesiano.

\medskip
ANSWER: (A) O ego moderno não se constitui apenas pelo \textit{Penso, logo existo}, de Descartes, mas também pela conquista e colonização do Outro.%
\newpage%
\section*{Apêndice: parecer do avaliador}%
\subsection*{Questão 1: nota 9,0 (aprovada, rodada 2 de 2)}%
\textbf{Resposta da resolução às cegas do avaliador:} Alternativa B. Pelo trecho, a Igreja primeiro se beneficiou do apoio imperial e depois ajudou a preservar o legado de Roma. As duas instituições se fortaleceram uma à outra, numa relação "por vezes simbiótica", e isso favoreceu grupos como o clero e a aristocracia senatorial. A alternativa A fala em uma separação entre Igreja e Estado, que não houve no século V. A alternativa C diz que a expansão ocorreu sem apoio imperial, o que contradiz o texto. A alternativa D confunde "preservar" o legado romano com "substituir" o Estado e não vale para o Oriente.\par %
Conferi o enunciado na página 3 do PDF e a transcrição está fiel. O candidato marca (B), igual à minha referência, e a resposta está correta. A resolução é sólida e completa.\newline%
\newline%
**Seção 1:** separa bem as duas fases descritas no trecho, o apoio imperial e depois a preservação do legado romano. Apoia cada fase em marcos corretos: Galério em 311, Milão em 313, isenções do clero, Niceia em 325, Tessalônica em 380 com o termo 'católico', e 476.\newline%
\newline%
**Alternativas:**\newline%
{-} (a) Mostra o anacronismo da separação entre Igreja e Estado, que no Oriente o vínculo ficou mais estreito e que as 'disputas entre classes sociais' não têm base.\newline%
{-} (b) Explica os dois qualificadores. Para o 'por vezes', cita Constâncio II e Atanásio, Juliano e Ambrósio contra Teodósio em 390. Para os 'determinados grupos sociais', cita o alto clero e a aristocracia senatorial (Ambrósio, Sidônio Apolinário, episcopalis audientia). Também trata a imprecisão do nome 'Igreja Católica Apostólica Romana' sem descartar a alternativa por isso.\newline%
{-} (c) Mostra que ela contradiz o texto de forma literal.\newline%
{-} (d) Trata como o distrator mais forte, distingue 'preservar' de 'substituir', mostra a inversão entre causa e consequência e lembra que o Império do Oriente continuou existindo.\newline%
\newline%
Todas as armadilhas da minha referência foram tratadas. Os problemas são pequenos imprecisões historiográficas: Calcedônia como 'base' do cesaropapismo, o fim da perseguição simplificado em 311 e Juliano classificado como 'atrito' e não como interrupção. Também falta a menção aos reinos germânicos em (d). Nenhum desses pontos afeta a conclusão. A resolução tem qualidade de gabarito comentado, com pequenos ajustes.%
\begin{itemize}%
\item%
{[}minor{]} Na alternativa (a), o candidato diz que o Concílio de Calcedônia (451), convocado por Marciano, foi a 'base do chamado cesaropapismo bizantino'. É exagero. É verdade que Marciano convocou o concílio e que isso mostra o vínculo estreito entre o imperador e a Igreja no Oriente. Mas a historiografia não trata Calcedônia como 'base' do cesaropapismo: o modelo costuma ser ligado à tradição que vem de Constantino e chega ao auge com Justiniano. O certo seria dizer apenas que a convocação imperial dos concílios ecumênicos (Niceia, Constantinopla, Calcedônia) mostra que, no Oriente, o vínculo ficou mais estreito.%
\item%
{[}minor{]} Na seção 1, o candidato afirma que o édito de Galério (311) 'havia encerrado a Grande Perseguição'. É uma simplificação: Maximino Daia retomou a perseguição no Oriente até 312--313. O correto seria escrever que o édito marcou o fim oficial da perseguição, retomada depois de forma localizada por Maximino Daia.%
\item%
{[}minor{]} Em (b), Juliano (361--363) aparece como exemplo de 'atrito dentro dessa aliança'. Na verdade, a tentativa de restaurar o paganismo interrompeu a aliança, não foi um atrito dentro dela. O exemplo ainda serve para justificar o 'por vezes', mas a classificação deveria ser 'interrupção' e não 'atrito interno'. Os outros exemplos, Constâncio II contra Atanásio e Ambrósio contra Teodósio, estão bem colocados.%
\item%
{[}minor{]} Em (d), ficou faltando lembrar que, no Ocidente, quem herdou boa parte da estrutura política e administrativa romana foram os reinos germânicos, e não a Igreja. Esse argumento reforçaria que não houve 'substituição' das instituições pela Igreja. O argumento do Oriente, já apresentado, basta para descartar a alternativa, então a falta é pequena.%
\end{itemize}%
\subsection*{Questão 2: nota 9,5 (aprovada, rodada 1 de 1)}%
\textbf{Resposta da resolução às cegas do avaliador:} Alternativa C. Para alguns povos ameríndios, o mundo não se organizava pela oposição entre humanos e não humanos. A distinção entre seres familiarizados e selvagens sustentava formas de organização social diferentes das europeias.\par %
Conferi o PDF (página 3) e a transcrição do enunciado e das alternativas está fiel ao original. O candidato marcou C, a mesma resposta da minha referência, e essa resposta está correta.\newline%
\newline%
A resolução começa resumindo o texto em três ideias:\newline%
{-} não há separação rígida entre humanos e não humanos;\newline%
{-} a distinção central é entre seres selvagens e familiarizados, sendo familiarizados os integrados ao convívio;\newline%
{-} os animais europeus foram incorporados segundo referências próprias, diferentes da pecuária e da propriedade animal europeias.\newline%
\newline%
Depois, cada alternativa errada é eliminada com base no trecho do texto que ela contradiz:\newline%
{-} a) Inverte a definição de seres familiarizados. O candidato também nota que a segunda metade da alternativa é plausível, mas não salva a definição errada.\newline%
{-} b) O texto nega a oposição humano/não humano como eixo, e o candidato aproveita bem o advérbio 'sobretudo'. Também critica o termo 'simbólica', porque o texto descreve uma integração concreta ao convívio.\newline%
{-} d) A primeira parte está certa, mas a conclusão contradiz o texto, que diz 'diferentes das concepções europeias'.\newline%
\newline%
A alternativa (c) é justificada ponto por ponto, inclusive com a observação de que o quantificador 'alguns povos' está de acordo com 'muitos povos' e 'diversas comunidades'. Todas as armadilhas que eu tinha identificado foram tratadas, incluindo a diferença entre familiarização e domesticação ou propriedade.\newline%
\newline%
Só achei detalhes menores. Em (c), o candidato troca 'contribuía' por 'gerava'. Ele também não comenta a pequena extrapolação 'o mundo... Por isso' presente em (c). Fora isso, a resolução está pronta para um gabarito comentado.%
\begin{itemize}%
\item%
{[}minor{]} Em (c), o candidato reescreve 'contribuía para formas de organização social' como 'gerava formas de organização social'. 'Gerava' afirma uma causalidade mais forte do que a alternativa diz. O certo é manter o verbo original ('contribuía'), que é justamente o que torna (c) defensável.%
\item%
{[}minor{]} O candidato não comenta que (c) generaliza um pouco o texto: diz que 'o mundo não era organizado pela oposição' e liga as ideias com 'Por isso', enquanto o texto fala apenas em 'não separar rigidamente' humanos e não humanos na relação com os animais. Bastaria uma frase dizendo que essa extrapolação é pequena e que (c) continua sendo a única alternativa compatível, porque as outras contradizem o texto de forma explícita.%
\end{itemize}%
\subsection*{Questão 3: nota 9,5 (aprovada, rodada 1 de 1)}%
\textbf{Resposta da resolução às cegas do avaliador:} Alternativa C: o voto era censitário. A Constituição de 1824 condicionava a participação eleitoral a uma renda líquida anual mínima: 100 mil réis para votar nas assembleias paroquiais e 200 mil réis para ser eleitor, como no Art. 94 citado. Havia ainda outras restrições legais, como a exclusão dos libertos e dos criminosos pronunciados.\par %
Conferi a transcrição com a página 3 do PDF e ela está fiel. A resposta final do candidato é a alternativa C (voto censitário), a mesma da minha referência, e está correta. A resolução é sólida e historicamente correta:\newline%
{-} Distingue votantes, que precisavam de 100 mil réis pelo art. 92, de eleitores, que precisavam de 200 mil réis pelo art. 94, e mostra que as eleições eram indiretas.\newline%
{-} Explica a origem do termo 'censitário'.\newline%
{-} Descarta a alternativa a com as exclusões do próprio trecho, além das mulheres e dos escravizados.\newline%
{-} Descarta a alternativa b ao mostrar que não havia exigência de instrução em 1824 e que os analfabetos só foram excluídos pela Lei Saraiva, de 1881.\newline%
{-} Descarta a alternativa d porque o critério era a condição jurídica de liberto, e não a cor.\newline%
{-} Liga a expressão 'entre outros temas legais' aos incisos II e III.\newline%
Todas as armadilhas que identifiquei foram tratadas. Ficam três imprecisões menores: o exagero de dizer que a participação ficava restrita às 'camadas proprietárias', quando a renda podia vir de emprego e o piso de 1º grau era baixo; a frase em que o candidato fala só do 'requisito de renda' e deixa de fora os demais requisitos do art. 92; e um tratamento apenas implícito do anacronismo. Nada disso afeta a resposta nem a validade dos argumentos. A resolução tem qualidade de gabarito comentado.%
\begin{itemize}%
\item%
{[}minor{]} Na conclusão, o candidato diz que a exigência de renda 'restringia a participação política às camadas proprietárias'. É um exagero. O próprio inciso I aceita renda vinda de 'emprego', ou seja, não só de propriedade. Além disso, o piso de 100 mil réis para votar no 1º grau era baixo, e a historiografia registra participação ampla de homens livres nas assembleias paroquiais. Deveria dizer que a exigência de renda excluía os mais pobres e ficava mais restritiva nos graus superiores, onde se exigiam 200, 400 e 800 mil réis.%
\item%
{[}minor{]} Na alternativa d, o candidato afirma que 'um homem livre de nascimento, qualquer que fosse sua cor, podia ser eleitor se cumprisse o requisito de renda'. Faltam os demais requisitos do art. 92, que o art. 94 incorpora: idade mínima de 25 anos (com exceções), não ser filho{-}família nem criado de servir, entre outros. Bastaria escrever 'se cumprisse os requisitos legais, sem distinção de cor'.%
\item%
{[}minor{]} O ponto do anacronismo, ligado à expressão 'de acordo com a concepção {[}...{]} presente na Constituição', aparece só de forma implícita, na oposição entre o critério legal explícito e o efeito social. Também falta dizer que os libertos nascidos no Brasil podiam votar no 1º grau, informação que reforçaria o argumento contra a alternativa d. Nenhuma das duas omissões compromete a resolução.%
\end{itemize}%
\subsection*{Questão 4: nota 9,0 (aprovada, rodada 1 de 1)}%
\textbf{Resposta da resolução às cegas do avaliador:} Alternativa C: o projeto fascista de modernização agrícola enfrentou resistência em certas regiões da Itália (o sul semiárido e pobre), e isso levou o Estado a financiar estratégias alternativas adaptadas às condições locais ("adaptação profunda", que valorizava a agrobiodiversidade local).\par %
Conferi a transcrição com as páginas 3 e 4 do PDF e ela está fiel. O candidato marcou a alternativa C, a mesma da minha referência. Refiz a análise e confirmo que C está correta: é uma paráfrase direta do texto, com os três elementos que ele traz. Primeiro, houve resistência regional, no sul semiárido e pobre. Segundo, essa resistência levou a abordagens alternativas. Terceiro, essas abordagens foram de \textbackslash{}"adaptação profunda\textbackslash{}", financiadas pelo Estado e baseadas na agrobiodiversidade local.\newline%
\newline%
A resolução tratou bem todas as armadilhas:\newline%
{-} **(a):** descartada porque o texto diz que o caso fascista \textbackslash{}"não foi um bloco homogêneo\textbackslash{}" e foi \textbackslash{}"desigual e contraditório\textbackslash{}". O candidato reconhece que a padronização vale para o pós{-}guerra.\newline%
{-} **(b):** descartada com conhecimento histórico externo ao texto. A Revolução Verde do pós{-}guerra se apoiou justamente em pesquisa e melhoramento genético, o que é uma semelhança com a Batalha do Trigo, não uma diferença.\newline%
{-} **(d):** o candidato mostra que a alternativa inverte a causalidade. O Estado tentou impor as novas variedades no sul, houve resistência e depois vieram as alternativas, sem abandono da região.\newline%
{-} **(c):** aceita sem excesso de rigor com a diferença entre \textbackslash{}"financiar\textbackslash{}" e \textbackslash{}"desenvolver\textbackslash{}".\newline%
{-} **Revolução Verde:** a resolução distingue a \textbackslash{}"Revolução Verde\textbackslash{}" fascista da clássica e contextualiza o pacote tecnológico do pós{-}guerra.\newline%
\newline%
Os defeitos são todos de redação ou de profundidade, sem nenhum erro de conteúdo:\newline%
{-} Uma frase ambígua na refutação de (d).\newline%
{-} Em (d), falta retomar a causa que o texto dá para a resistência (condições semiáridas, pobreza e apego às sementes tradicionais).\newline%
{-} O contexto histórico é um pouco genérico, sem Borlaug, CIMMYT, IRRI ou Strampelli.\newline%
{-} Em (a), não se diz explicitamente que basta uma metade falsa para invalidar a afirmação.\newline%
\newline%
A resolução é clara, bem organizada e publicável como gabarito comentado, com pequenos ajustes.%
\begin{itemize}%
\item%
{[}minor{]} Na alternativa (d), a frase "O Estado não abandonou o sul, e foi isso que provocou a resistência" é ambígua. Pode ser lida como "o não abandono (a imposição) provocou a resistência", que é o sentido correto, ou como uma negação mal construída da tese da alternativa. A frase seguinte desfaz a dúvida, mas um gabarito comentado deveria dizer direto: "O Estado não abandonou o sul; ao contrário, tentou impor ali as novas variedades, e foi essa imposição que gerou a resistência."%
\item%
{[}minor{]} Em (d), a resolução não retoma a causa que o próprio texto dá para a resistência: as condições semiáridas e a pobreza, somadas ao apego dos agricultores às próprias sementes e práticas. Esses fatores aparecem só na seção 'Ideia central'. Citá{-}los na refutação de (d) deixaria o descarte mais completo.%
\item%
{[}minor{]} O contexto histórico é correto, mas genérico. Faltam referências que tornariam o descarte de (b) mais sólido e o comentário mais rico: Borlaug, CIMMYT e IRRI, as variedades semianãs de alto rendimento e, do lado italiano, a autarquia e os trigos de Strampelli. Não é obrigatório numa questão objetiva, mas um bom gabarito comentado incluiria.%
\item%
{[}minor{]} Em (a), a resolução mostra que o caso fascista não foi homogêneo, mas não diz explicitamente que a afirmação é uma conjunção e que basta uma das metades ser falsa para invalidá{-}la. A ideia fica implícita.%
\end{itemize}%
\subsection*{Questão 5: nota 9,0 (aprovada, rodada 1 de 1)}%
\textbf{Resposta da resolução às cegas do avaliador:} Alternativa b): as bolsas de mandinga eram recursos espirituais de proteção (para "fechar o corpo") amplamente usados e difundidos no Atlântico escravista. Elas combinavam elementos africanos, católicos e americanos em práticas religiosas novas, híbridas.\par %
Conferi a página 4 do PDF, e a transcrição bate com o original. O candidato respondeu (B), a mesma resposta da minha referência, e ela está correta. A resolução é bem organizada. Primeiro, ela levanta o que o texto diz sobre função (fechar o corpo, proteção espiritual), composição ('fruto do contato entre diferentes tradições culturais', com itens americanos, católicos e corporais) e difusão (redes entre Minas, Recife, Rio, Lisboa, Costa da Mina e Angola, com crescimento no século XVIII). Depois, ela usa esses pontos para julgar cada alternativa. As armadilhas principais foram tratadas. Em (a), o candidato mostra que a ideia de 'imunidade' contradiz os elementos europeus e americanos. Em (c), ele separa incorporação ou ressignificação de substituição e usa a perseguição da Inquisição como prova de heterodoxia, que é o argumento esperado. Em (d), ele aponta que o texto não fala em Islã nem em Alcorão e que a alternativa reduz o fenômeno a uma só matriz. Ele também justifica as três partes de (b) com o texto, sem se limitar a dizer que é a 'mais completa'. Os problemas que encontrei são pequenos: (1) chama de cristãos os 'outros símbolos gráficos', o que o texto não diz; (2) atribui as partes do corpo a uma matriz africana, o que o excerto não afirma; (3) chama a prática de 'não cristã' de forma categórica demais; (4) trata pouco a armadilha da etimologia de 'mandinga' e não menciona que os usuários eram de origem africana em geral, e não só mandinga. Nenhum desses pontos afeta a conclusão. É uma resolução publicável como gabarito comentado, com pequenos ajustes de precisão.%
\begin{itemize}%
\item%
{[}minor{]} Na análise de (d), o candidato diz que 'os escritos citados são orações e símbolos cristãos'. O texto fala em 'orações, cruzes e outros símbolos gráficos', e esses 'outros símbolos gráficos' não são necessariamente cristãos. A frase vai um pouco além do texto e, sem querer, deixa espaço para quem defende (d). Deveria dizer apenas que o excerto não menciona Islã nem Alcorão e lista elementos católicos explícitos (pedra d'ara, hóstias, cruzes).%
\item%
{[}minor{]} No Passo 1, o candidato atribui as partes do corpo humano (unhas, cabelos, ossos, dentes) a 'práticas mágicas de matriz africana'. O excerto não diz isso, e esse tipo de elemento também aparece na magia popular europeia. Seria melhor apresentar isso como inferência ('provavelmente ligadas a\ldots{}') ou não atribuir origem nenhuma.%
\item%
{[}minor{]} Em (c), chamar a prática de 'mágica, que não era cristã' é categórico demais. Ela era híbrida e heterodoxa, e incorporava elementos cristãos. Dizer 'não ortodoxa' ou 'não reconhecida pela Igreja' seria mais preciso e mais coerente com o argumento do próprio candidato.%
\item%
{[}minor{]} A resolução trata a etimologia de 'mandinga' (povo mandinka, islamizado) só com a frase 'O nome não autoriza essa conclusão', sem explicar por que (d) é uma armadilha. Também não aponta que o texto fala em usuários 'sobretudo de origem africana' em geral, e não só em 'escravizados de origem mandinga'. O argumento continua certo, só que mais raso do que poderia ser.%
\item%
{[}minor{]} Para justificar 'amplamente utilizados', a resolução cita as redes atlânticas e o crescimento no século XVIII, mas não menciona que o texto chama as bolsas de 'célebres', o que reforçaria a ideia de uso amplo. A omissão não prejudica a conclusão.%
\end{itemize}%
\subsection*{Questão 6: nota 9,0 (aprovada, rodada 1 de 1)}%
\textbf{Resposta da resolução às cegas do avaliador:} Alternativa b): os projetos de ocupação territorial da ditadura expressavam uma concepção de desenvolvimento que subordinava a permanência dos povos indígenas em seus territórios aos interesses do Estado. É uma paráfrase direta de "A sobrevivência dos povos indígenas e a existência de suas culturas foram subordinadas aos grandes projetos de colonização do país".\par %
Conferi a transcrição com a página 4 do PDF e ela bate com o original. As duas resoluções dão a mesma resposta: b). Ela está correta, porque parafraseia a frase central do texto: a sobrevivência dos povos indígenas e suas culturas foram 'subordinadas aos grandes projetos de colonização do país'.\newline%
\newline%
O candidato organiza bem o texto em três eixos: falta de controle social, Doutrina de Segurança Nacional 'etnocêntrica, racista e desenvolvimentista' e a conclusão explícita de subordinação, com o exemplo da BR{-}80. Também situa o trecho no contexto histórico (PIN de 1970, Transamazônica, lema 'integrar para não entregar'). Refuta cada distrator com elementos do próprio texto:\newline%
{-} a) a palavra 'subordinadas' e o corte do Xingu;\newline%
{-} c) censura, decisões em Brasília e a Funai justificando a obra;\newline%
{-} d) fragmentação do Xingu, doenças e conflitos.\newline%
\newline%
As armadilhas principais foram tratadas: o vocabulário oficial em (a), a falta de participação em (c) e a inversão de causa e efeito em (d). Faltam detalhes menores: o anacronismo de 'autonomia' e a tutela do Estatuto do Índio de 1973. Há também uma leve extrapolação ao falar de 'territórios' e 'conflitos agravados', que não está escrita no texto, embora seja historicamente correta.\newline%
\newline%
É uma resolução clara, correta e bem fundamentada, de nível próximo a um gabarito comentado.%
\begin{itemize}%
\item%
{[}minor{]} Na alternativa (a), o candidato não aponta que 'autonomia indígena' é anacrônica. A política do regime era tutelar e integracionista (Funai de 1967, Estatuto do Índio de 1973, que tratava o indígena como 'relativamente incapaz'), e a autonomia só passa a ser reconhecida, em parte, com a CF/1988. A refutação feita só com o texto basta para eliminar (a). Faltou, porém, esse argumento histórico, que tornaria a refutação completa.%
\item%
{[}minor{]} Na alternativa (c), faltou citar a tutela jurídica do Estatuto do Índio. Ela excluía de forma institucional a participação indígena e reforça o argumento além do que o texto mostra (censura, decisões em Brasília, falta de controle social).%
\item%
{[}minor{]} O candidato diz que as rodovias 'agravaram os conflitos com colonos e fazendeiros' e, na conclusão, que os 'territórios' foram subordinados. É uma pequena extrapolação: o texto fala em 'sobrevivência' e 'existência de suas culturas' e não liga diretamente as arbitrariedades de colonos às rodovias. A leitura é historicamente correta e o exemplo da BR{-}80 a sustenta, mas seria mais rigoroso deixar claro o que vem do texto e o que vem do contexto histórico.%
\item%
{[}minor{]} Ao tratar o Xingu como 'vitrine', o candidato só diz que 'nem uma área protegida como o Xingu foi poupada'. Não aponta explicitamente a armadilha de confundir a imagem protecionista da Funai com proteção real, nem comenta a ironia de o próprio presidente do órgão indigenista ter respondido apenas com um comentário. A ideia está implícita, mas não foi desenvolvida.%
\end{itemize}%
\subsection*{Questão 7: nota 9,5 (aprovada, rodada 1 de 1)}%
\textbf{Resposta da resolução às cegas do avaliador:} Alternativa A. O console criou uma cultura material própria e transformou a relação com a tela da TV, que deixou de ser só expositiva e passou a ser interativa. Isso redefiniu a sociabilidade (o jogo saiu dos arcades e foi para dentro de casa), o consumo (Atari, Nintendo, Sega) e a memória (games expostos em museu e tratados como fonte histórica).\par %
Conferi o enunciado no PDF (páginas 4 e 5) e a transcrição está correta. A resposta do candidato, (A), coincide com a minha referência e está correta.\newline%
\newline%
A resolução cruza as duas fontes como a questão pede:\newline%
{-} **Imagem:** a vitrine de museu da exposição 'Game Masters' sustenta 'cultura material' e 'memória'. O candidato também usa o título do artigo ('videogames como fontes históricas') para apoiar a ideia de memória.\newline%
{-} **Excerto:** a passagem da tela 'expositiva' para a tela 'interativa e responsiva' sustenta a 'relação com as telas da TV'. A ida do jogo dos arcades e espaços públicos para dentro de casa sustenta a sociabilidade, e a trajetória comercial (Odyssey, Atari, Nintendo, Sega) sustenta o consumo.\newline%
\newline%
Todas as armadilhas foram tratadas:\newline%
{-} **A:** cada termo da alternativa é justificado.\newline%
{-} **B:** o candidato identifica que o erro é atribuir ao console um impacto sobre a produção audiovisual ('filmar e exibir'), algo que o texto não diz.\newline%
{-} **C:** a relação com Hollywood e o aprofundamento das diferenças entre cinema e TV não aparecem nas fontes.\newline%
{-} **D:** a alternativa contradiz o texto (o console reforçou o uso doméstico da TV) e inventa uma 'recusa' da linguagem dos games.\newline%
\newline%
Os problemas são apenas de redação: um 'pelo contrário' retoricamente forte demais na (c), o termo 'passivo' em vez de 'expositiva' e o fato de não admitir explicitamente que a primeira metade da (b) é compatível com o texto. Nada disso compromete o raciocínio. A resolução está no nível de um gabarito comentado publicável.%
\begin{itemize}%
\item%
{[}minor{]} Na alternativa (c), o candidato escreve 'Pelo contrário' ao dizer que os games 'articulam visualidade, som, regra e performance'. O contraste é exagerado. O fato de os games articularem várias linguagens não refuta logicamente a tese de que as diferenças entre cinema e TV se aprofundaram, apenas mostra que essa tese não tem apoio no texto. Bastava dizer que essa ideia é estranha às fontes, sem apresentar a frase como contraprova.%
\item%
{[}minor{]} No Passo 1, o candidato afirma que 'o espectador, antes passivo, passa a agir sobre ela'. O excerto diz que a tela era 'predominantemente expositiva', e 'passivo' é uma leve extrapolação. Convém manter a formulação da fonte (de expositiva para interativa e responsiva).%
\item%
{[}minor{]} Na alternativa (b), o candidato não reconhece explicitamente que a primeira oração ('radicalizar a relação entre ver e agir') é compatível com o texto e que o erro está só na conclusão, sobre a indústria do cinema e da TV e as 'novas formas de filmar'. Essa leitura aparece de forma implícita ('do lado do jogador'). Dizer isso com clareza deixaria mais didático o tratamento do distrator mais forte.%
\end{itemize}%
\subsection*{Questão 8: nota 9,2 (aprovada, rodada 1 de 1)}%
\textbf{Resposta da resolução às cegas do avaliador:} Alternativa a). O gás é CO$_{2}$. Quando ele entra no frasco fechado, a pressão salta ($\approx$ 1,08 atm). Depois cai aos poucos ($\to$ $\approx$ 1,01 atm), porque o CO$_{2}$ se dissolve na água. Ao mesmo tempo, a condutividade sobe continuamente de $\approx$ 0,1 para dezenas de $\mu$S/cm, porque se forma H$_{2}$CO$_{3}$ ? H? + HCO$_{3}$$^{-}$. Com He, a pressão subiria e ficaria constante, e a condutividade não mudaria; nenhum gráfico mostra isso.\par %
Conferi a transcrição com a p. 5 do PDF e ela está correta. Os gráficos são de dois tipos: a e c são iguais, b e d são iguais. O candidato marcou a alternativa a (CO2, pressão que salta e depois cai, condutividade crescente), a mesma da minha referência, e a resposta está correta. O raciocínio cobre as armadilhas principais:\newline%
{-} Com pressão constante (b e d), o gráfico é impossível para qualquer gás, porque o frasco fechado recebe gás a mais (pV = nRT).\newline%
{-} Com He, a pressão subiria e ficaria num patamar, e a condutividade não mudaria. Por isso c sai também pela curva de pressão.\newline%
{-} O CO2 dissolvido forma H? + HCO3$^{-}$, e é isso que explica o aumento da condutividade.\newline%
{-} Com CO2, a condutividade cresce de forma monotônica, sem pico seguido de queda como em b.\newline%
{-} A queda de pressão e o aumento de condutividade são ligados como dois lados da absorção do CO2.\newline%
\newline%
A conta de ordem de grandeza confere: {[}H?{]} $\approx$ 1,24$\times$10$^{-}$$^{4}$ mol/L, pH $\approx$ 3,9 e $\kappa$ $\approx$ 49 $\mu$S/cm. Os problemas são pequenos:\newline%
{-} A conta usa p\_CO2 = 1 atm, mas no frasco p\_CO2 ? 0,08 atm, o que daria $\approx$ 14 $\mu$S/cm. Faltou dizer que o valor calculado é um limite superior.\newline%
{-} A conversão de mol/L para mol/cm$^3$ não está escrita.\newline%
{-} Não se explica por que a pressão fica em $\approx$ 1,01 atm, acima de 1,00 atm.\newline%
{-} A relação entre a autoionização e 0,1 $\mu$S/cm é imprecisa.\newline%
\newline%
É uma resolução clara e completa, de nível gabarito comentado.%
\begin{itemize}%
\item%
{[}minor{]} A estimativa de ordem de grandeza usa p\_CO2 = 1 atm. Pelo próprio gráfico, porém, o frasco já tem ar a 1,00 atm e o CO2 entra só até $\approx$ 1,08 atm, então p\_CO2 no frasco é no máximo $\approx$ 0,08 atm. Com esse valor, o equilíbrio daria $\kappa$ $\approx$ 14 $\mu$S/cm, contra $\approx$ 49 $\mu$S/cm a 1 atm (conferido em Python). Dizer que o resultado é 'coerente com os $\approx$ 60 $\mu$S/cm' sem avisar que a conta é um limite superior é um pouco otimista. A conclusão qualitativa (aumento de duas ordens de grandeza na escala log) continua válida. O candidato deveria ter feito a conta também para p $\approx$ 0,1 atm, ou dito que 1 atm é cota superior.%
\item%
{[}minor{]} Notação de unidades: {[}H?{]} = 1,2$\times$10$^{-}$$^{4}$ mol/L aparece como 1,2$\times$10$^{-}$$^{7}$ em $\kappa$ = ?$\cdot$c, sem dizer que houve conversão para mol/cm$^3$. A conta está certa ($\approx$ 4,7$\times$10$^{-}$$^{5}$ S/cm), mas a conversão deveria estar escrita.%
\item%
{[}minor{]} O candidato não diz por que a pressão fica em $\approx$ 1,01 atm em vez de voltar a 1,00 atm: a dissolução tende a um equilíbrio (lei de Henry) e sobra CO2 na fase gasosa. A ideia de equilíbrio só aparece na parte de condutividade.%
\item%
{[}minor{]} Dizer que a autoionização com {[}H?{]} = {[}OH$^{-}${]} = 10$^{-}$$^{7}$ mol/L dá 'cerca de 0,1 $\mu$S/cm' é impreciso. O valor teórico da água pura é $\approx$ 0,055 $\mu$S/cm; 0,1 $\mu$S/cm é o típico de água deionizada exposta ao ar. Não afeta a análise.%
\end{itemize}%
\subsection*{Questão 9: nota 9,5 (aprovada, rodada 1 de 1)}%
\textbf{Resposta da resolução às cegas do avaliador:} Alternativa b): o íon é o Au$^3$?. Como E$^\circ$(Au$^3$?/Au) = +0,93 V é maior que E$^\circ$(Ag?/Ag) = +0,80 V, temos $\Delta$E$^\circ$ = +0,13 V e a reação 3 Ag + Au$^3$? $\to$ 3 Ag? + Au é espontânea. Como 3 Ag saem para cada Au que entra, com volumes molares quase iguais, o nanocubo maciço vira uma estrutura de liga Au--Ag oca e porosa (tipo gaiola).\par %
Conferi a transcrição com o original (pág. 6 do PDF) e ela confere: a) e c) mostram cubo maciço $\to$ cubo maciço idêntico; b) e d) mostram cubo maciço $\to$ cubo oco e poroso. A resposta final do candidato, (b), coincide com a minha referência e está correta.\newline%
\newline%
**Critério termodinâmico.** O candidato usa o sentido certo: a Ag é oxidada, e o íon precisa ter E$^\circ$red maior que +0,80 V. Ele escreve as semirreações e a equação global 3 Ag + Au$^3$? $\to$ 3 Ag? + Au, com $\Delta$E$^\circ$ = 0,93 - 0,80 = +0,13 V, sem multiplicar o potencial pelo coeficiente. $\Delta$G$^\circ$ $\approx$ -38 kJ/mol; conferi e dá -37,6 kJ. Para o Cu$^2$?, obtém $\Delta$E$^\circ$ = -0,46 V e reconhece que a reação espontânea é a inversa, o que descarta c) e d). Comenta ainda Pt$^2$? (serviria) e Zn$^2$? (não serviria).\newline%
\newline%
**Morfologia.** A justificativa vem da estequiometria 3 Ag : 1 Au com volumes atômicos praticamente iguais, de modo que só cerca de 1/3 do volume oxidado é reposto. Disso resulta a casca de liga Au--Ag oca e porosa (nanogaiola). O candidato também cita o dado do enunciado de que houve \textbackslash{}"alteração de composição e morfologia\textbackslash{}" para descartar a) e c).\newline%
\newline%
**Armadilhas.** Quase todas foram tratadas: sentido da reação, eliminação do Cu$^2$?, armadilha da alternativa a), citação do enunciado, potenciais sem multiplicar e os distratores Pt$^2$? e Zn$^2$?. Ficaram de fora só a observação sobre a escala das figuras e o erro tipográfico, ambos irrelevantes.\newline%
\newline%
**Ressalvas.** São pequenas: o dado de parâmetro de rede é externo à prova e o mecanismo de formação dos poros é descrito de forma um pouco vaga.\newline%
\newline%
A resolução é clara, correta e completa, com nível de gabarito comentado.%
\begin{itemize}%
\item%
{[}minor{]} Para dizer que Ag e Au têm volumes atômicos iguais, a resolução usa o parâmetro de rede ($\approx$4,08 Å), um dado que não está na prova. O valor está correto (Ag 4,086 Å, Au 4,078 Å), mas convinha dizer que é conhecimento externo. Outra opção seria mostrar que a conclusão qualitativa vale só pela estequiometria 3:1, já que sai muito mais átomo do que entra. Ou ainda usar massas molares e densidades: Ag $\approx$ 10,3 e Au $\approx$ 10,2 cm$^3$/mol.%
\item%
{[}minor{]} O mecanismo da morfologia ficou um pouco vago ("a prata é consumida pelo interior"). Faltou dizer que a oxidação começa num ponto de defeito da superfície e segue pelo interior. Também faltou dizer que a remoção da prata da liga (desliga) é o que abre os furos da parede, levando de nanocaixa a nanogaiola. Para uma questão de múltipla escolha, o nível qualitativo basta.%
\item%
{[}minor{]} A resolução não comenta a instrução de que as figuras são ilustrativas e fora de escala. Ela também não menciona o erro tipográfico "Cu$^2$+" nas alternativas c) e d). Nenhum dos dois pontos afeta a resposta.%
\end{itemize}%
\subsection*{Questão 10: nota 9,5 (aprovada, rodada 1 de 1)}%
\textbf{Resposta da resolução às cegas do avaliador:} Alternativa d. Na figura, K diminui quando T aumenta, então a reação é exotérmica e o aumento de temperatura desfavorece a formação de NH$_{3}$. Como $\Delta$n(gás) = 1 - 2 = -1, o aumento de pressão desloca o equilíbrio para o produto e favorece sua formação.\par %
Conferi o enunciado na página 6 do PDF. A reação é $\frac12$N$_{2}$ + 3/2H$_{2}$ = NH$_{3}$, a curva de K em escala logarítmica de 10$^{-}$$^{6}$ a 10$^{6}$ é decrescente entre 200 e 900 K e traz o rótulo 300 atm. As alternativas coincidem com a transcrição.\newline%
\newline%
O candidato e a referência chegam à mesma resposta, alternativa d, e ela está correta.\newline%
\newline%
**Temperatura.** A leitura do gráfico está certa: K vai de cerca de 10$^{5}$ em 250 K a cerca de 10$^{-}$$^{4}$ em 880 K. Daí o candidato conclui que o aumento de T desfavorece o produto e justifica com van 't Hoff ($\Delta$H$^\circ$ < 0) e com Le Chatelier, que desloca o equilíbrio no sentido endotérmico.\newline%
\newline%
**Pressão.** A dedução é rigorosa. Pela estequiometria, $\Delta$n(gás) = 1 - ($\frac12$ + 3/2) = -1. Escrevendo p? = x?P, obtém{-}se K$\cdot$P = x\_NH$_{3}$/(x\_N$_{2}$\^{}$\frac12$$\cdot$x\_H$_{2}$\^{}\{3/2\}). O candidato deixa claro que K não muda com P e que o que muda é a composição. Com isso evita o erro clássico de dizer que a pressão aumenta K. As contas e o sinal do expoente de P estão corretos.\newline%
\newline%
**Armadilhas.** Todas foram tratadas:\newline%
{-} A menção do texto à 'alta temperatura' na indústria é explicada como escolha cinética.\newline%
{-} A escala logarítmica é lida corretamente.\newline%
{-} O sinal de $\Delta$H é deduzido a partir da figura.\newline%
{-} Fica claro que o gráfico é a pressão fixa e que o efeito da pressão sai da estequiometria.\newline%
{-} Os coeficientes fracionários são somados corretamente.\newline%
{-} As alternativas c e d são separadas pelo efeito da pressão.\newline%
\newline%
**Ressalvas menores:**\newline%
{-} Falta uma frase ligando o aumento da razão de frações molares ao aumento de x\_NH$_{3}$.\newline%
{-} Falta a ressalva de que a independência de K em relação a P vale para gás ideal.\newline%
{-} Falta dizer de forma explícita que o catalisador não desloca o equilíbrio.\newline%
\newline%
A resolução tem qualidade de gabarito comentado.%
\begin{itemize}%
\item%
{[}minor{]} A conclusão 'aumentar P aumenta a razão x\_NH3/(x\_N2\^{}\{1/2\} x\_H2\^{}\{3/2\}), ou seja, aumenta a fração de NH3' tem um pequeno salto. Falta dizer que, para uma composição inicial fixa, essa razão cresce de forma monotônica com o avanço da reação. Logo, se a razão aumenta, o avanço aumenta e x\_NH3 também. Bastava uma frase para fechar o argumento.%
\item%
{[}minor{]} A afirmação 'K não depende da pressão' só vale de forma estrita para gases ideais, com K definido em relação ao estado{-}padrão. O próprio gráfico traz o rótulo '300 atm', o que sugere uma possível correção por fugacidade. Uma ressalva curta (gás ideal ou Kp termodinâmico) deixaria o texto mais rigoroso. Isso não afeta a resposta.%
\item%
{[}minor{]} Na observação sobre a indústria, o texto cita o catalisador ('junto com o efeito do catalisador') sem dizer de forma explícita que ele não desloca o equilíbrio, só acelera a chegada até ele. Não há erro, mas a armadilha do catalisador ficou tratada apenas de forma implícita.%
\end{itemize}%
\subsection*{Questão 11: nota 9,5 (aprovada, rodada 1 de 1)}%
\textbf{Resposta da resolução às cegas do avaliador:} Alternativa **c**: 1ª opção **A**, 2ª opção **B**.\newline%
\newline%
A ordem de hidrofilicidade dos ligantes é A > B > C.\newline%
{-} **A** é o mais hidrofílico: além dos 2 --COOH, tem o grupo --SO$_{3}$H, que é muito polar e ácido forte, e faz várias ligações de hidrogênio e interações íon--dipolo com a água.\newline%
{-} **B** vem em seguida: tem 2 --COOH, que ainda fazem ligações de hidrogênio.\newline%
{-} **C** é o pior: tem poucos sítios polares e um --CH$_{3}$ apolar. No ZIF{-}8, os seus N ficam coordenados ao Zn, e esse material é hidrofóbico.\par %
Conferi o enunciado no PDF (págs. 4--5 do caderno): A é o ácido 2{-}sulfotereftálico, B o ácido tereftálico e C o 2{-}metilimidazol. A pergunta pede a 1ª e a 2ª opção de ligante para captar água do ar. O candidato marcou (c), A e B, a mesma resposta da minha referência, e ela está correta.\newline%
\newline%
A resolução tem quatro pontos fortes:\newline%
{-} Usa o critério certo: a hidrofilicidade e a natureza das interações com a água (ligação de hidrogênio e íon--dipolo, contra dispersão). Não se deixa levar pelo contexto de acidez, CO$_{2}$ ou catálise.\newline%
{-} Justifica bem A em 1º lugar: o --SO$_{3}$H tem 3 O aceptores e 1 O--H doador, pode se ionizar a --SO$_{3}$$^{-}$ e fica livre no poro, porque os carboxilatos é que estão coordenados ao metal.\newline%
{-} Explica por que C é o pior: N--H desprotonado, os dois N coordenados ao Zn$^2$?, poro revestido de --CH$_{3}$ e C--H, e o ZIF{-}8 é hidrofóbico.\newline%
{-} Dá a ordem A > B > C, discute as quatro alternativas e trata a questão da ordem (1ª e 2ª opções).\newline%
\newline%
Todas as armadilhas que listei foram tratadas.\newline%
\newline%
As ressalvas são pequenas:\newline%
{-} Afirma, sem fundamento firme, que os O de carboxilato coordenados ainda aceitam ligação de hidrogênio.\newline%
{-} Mistura o critério do ligante livre (em B) com o do ligante dentro da MOF (em C), sem dizer que os dois dão a mesma ordem.\newline%
{-} Diz que não sobra 'nenhum' sítio livre no ZIF{-}8, o que é um pouco forte.\newline%
\newline%
É uma resolução praticamente de gabarito comentado.%
\begin{itemize}%
\item%
{[}minor{]} Na análise de B, o candidato afirma que 'mesmo coordenados ao metal, os oxigênios dos carboxilatos ainda podem aceitar ligações de hidrogênio'. A afirmação não tem sustentação. Num carboxilato em ponte entre centros de Zr, os pares livres do O já estão envolvidos na coordenação, e esses oxigênios são, no máximo, aceptores fracos. No UiO{-}66 real, a hidrofilicidade vem principalmente dos grupos $\mu$3{-}OH/$\mu$3{-}O dos clusters de Zr. Bastaria dizer que os carboxilatos (ou os --COOH do ligante livre) oferecem mais sítios polares que o 2{-}metilimidazol, sem afirmar o que acontece com os O coordenados.%
\item%
{[}minor{]} O critério muda de um ligante para outro. Para C, o candidato usa o estado dentro da MOF (N coordenados ao Zn, sem sítio livre). Para B, conta sobretudo os grupos do ligante livre (4 O e 2 O--H, doadores e aceptores). A ordem A > B > C sai igual pelos dois critérios, mas o texto deveria dizer isso explicitamente para ficar coerente. Por exemplo: no ligante livre, B tem 4 O e 2 O--H contra 1 N e 1 N--H de C; na MOF, C fica sem sítios polares expostos.%
\item%
{[}minor{]} A frase 'não sobra nenhum sítio livre para ligação de hidrogênio' no ZIF{-}8 é forte demais. Vale para a estrutura idealizada, mas convinha dizer 'praticamente nenhum' (há C--H do anel, que são doadores muito fracos, e há defeitos). Isso não afeta a conclusão.%
\end{itemize}%
\subsection*{Questão 12: nota 9,5 (aprovada, rodada 1 de 1)}%
\textbf{Resposta da resolução às cegas do avaliador:} Alternativa C: {[}CH3OH{]} > 50 mg/dL e acidose metabólica, com HCO3$^{-}$ < 15 mmol/L e pH arterial < 7,30. O ácido fórmico vindo do metanol consome o bicarbonato e baixa o pH.\par %
Conferi a transcrição com o original (p. 7 do PDF, p. 5 do caderno) e ela está fiel, inclusive no termo 'acidez metabólica'. O candidato marca (C), a mesma resposta da minha referência, e a resposta está correta.\newline%
\newline%
A justificativa está certa e bem encadeada. O candidato percorre a rota metabólica CH3OH $\to$ HCHO $\to$ HCOOH e mostra que o HCOOH fica quase todo ionizado no sangue (pKa 3,75, Ka $\approx$ 1,8$\times$10$^{-}$$^{4}$; conferi). Explica que o tampão consome o bicarbonato (H? + HCO3$^{-}$ $\to$ CO2 + H2O) e usa Henderson--Hasselbalch para chegar ao pH normal de 7,40. Nota que o problema vem do metabolismo, por isso a acidose é metabólica, e registra que o enunciado diz 'acidez' em vez de 'acidose'. Usa as faixas normais do sangue (pH 7,35--7,45 e HCO3$^{-}$ 22--26 mmol/L), o que evita comparar o pH com 7. Justifica o limite de metanol (o tratamento vale acima dele) e converte 50 mg/dL $\approx$ 16 mmol/L (conferi: 15,6 mmol/L). Usa o próprio tratamento para confirmar o quadro: o etanol e o fomepizol bloqueiam a ADH, e a hemodiálise remove o metanol e o formiato.\newline%
\newline%
As armadilhas principais foram tratadas: acidose $\times$ alcalose, sentido do bicarbonato, sentido do limite de metanol, pH < 7,30 ser ácido no sangue e impropriedade 'acidez'. Ficaram omissões pequenas na análise das alternativas: a contradição interna de (b) e a de (d) quanto ao bicarbonato não foram apontadas, e faltou a conta de pH com HCO3$^{-}$ = 15. Também faltou a ressalva de que os critérios clínicos costumam ser alternativos. É uma resolução publicável como gabarito comentado.%
\begin{itemize}%
\item%
{[}minor{]} Na alternativa (b), a resolução não aponta a contradição interna: a alternativa chama de 'acidez metabólica' um quadro com pH > 7,30 e HCO3$^{-}$ > 15 mmol/L. Ela aponta essa incoerência em (a), mas não aqui. Uma frase a mais deixaria a análise simétrica.%
\item%
{[}minor{]} Na alternativa (d), falta dizer que HCO3$^{-}$ < 15 mmol/L é marca de acidose metabólica, o que contradiz o rótulo 'alcalose'. A resolução descarta (d) pelo metanol baixo e pelo pH > 7,30, o que basta para eliminá{-}la, mas a incoerência entre bicarbonato baixo e alcalose ficou sem comentário.%
\item%
{[}minor{]} A resolução não mostra com uma conta que HCO3$^{-}$ < 15 é compatível com pH < 7,30. Com Henderson--Hasselbalch e pCO2 = 40 mmHg, HCO3$^{-}$ = 15 mmol/L dá pH $\approx$ 7,20. Ela só calcula o caso normal (pH $\approx$ 7,40), embora a fórmula já estivesse escrita.%
\item%
{[}minor{]} A resolução não comenta que, nos protocolos clínicos, os critérios costumam ser alternativos ('ou'), enquanto a alternativa usa 'e'. Isso não muda o gabarito, mas é uma ressalva útil num gabarito comentado.%
\end{itemize}%
\subsection*{Questão 13: nota 9,5 (aprovada, rodada 1 de 1)}%
\textbf{Resposta da resolução às cegas do avaliador:} Alternativa D. Os gases são H$_{2}$ (no catodo, pela redução da água) e Cl$_{2}$ (no anodo, pela oxidação do Cl$^{-}$). A equação é Cl$_{2}$(g) + H$_{2}$O(l) ? HClO(aq) + H?(aq) + Cl$^{-}$(aq), balanceada em átomos e em carga.\par %
Conferi o enunciado na página 7 do PDF e a transcrição está correta. A resposta final do candidato é (D), a mesma da minha referência, e está correta. Os gases da eletrólise da solução aquosa de NaCl são H$_{2}$ no catodo, pela redução da água (E$^\circ$ do Na?/Na = -2,71 V, muito negativo), e Cl$_{2}$ no anodo, porque a sobretensão do O$_{2}$ favorece a oxidação do Cl$^{-}$. A equação Cl$_{2}$ + H$_{2}$O ? HClO + H? + Cl$^{-}$ é um desproporcionamento, com o Cl indo de 0 para +1 e para -1. Verifiquei com Python o balanço de átomos e de carga das quatro alternativas, e ele coincide exatamente com o que o candidato afirma: em (a) H 2 $\to$ 1 e carga 0 $\to$ -1; em (b) H 2 $\to$ 3 e carga -1 $\to$ 0; em (c) O 3 $\to$ 1 e carga -1 $\to$ +1; (d) está balanceada.\newline%
\newline%
Todas as armadilhas foram tratadas:\newline%
{-} **(a) versus (d):** o candidato as distingue pelo H? que falta em (a).\newline%
{-} **Balanço de carga:** verificado além do de átomos, com tabela.\newline%
{-} **Catodo:** explica que o Na? não se reduz.\newline%
{-} **Anodo:** explica por que o O$_{2}$ não se forma, pela sobretensão.\newline%
{-} **Desproporcionamento:** identificado, inclusive com a conservação de elétrons.\newline%
\newline%
A observação final, de que o HClO se regenera em Cl$^{-}$ e por isso o sal não é consumido, está correta e é pertinente.\newline%
\newline%
Há só detalhes menores. O argumento auxiliar de que o produto precisa conter cloro é logicamente frágil. A afirmação de que O$_{2}$ não se forma é absoluta demais. E o meio do potencial de 1,23 V não foi indicado. Nenhum deles afeta a conclusão. A resolução tem qualidade de gabarito comentado.%
\begin{itemize}%
\item%
{[}minor{]} O argumento auxiliar 'esse produto precisa conter cloro, e o único que contém é o Cl$_{2}$' não vale como está escrito. Um produto sem cloro, como um oxidante tipo O$_{2}$, poderia em tese oxidar o Cl$^{-}$ já presente na água, e é exatamente isso que a alternativa (c) propõe. Por isso o argumento não descarta (c) sozinho. Quem descarta (c) é a sobretensão do O$_{2}$ e o fato de a equação de (c) não estar balanceada (O: 3 $\to$ 1; carga: -1 $\to$ +1), e a resolução apresenta as duas coisas. Melhor seria escrever que o produto que gera HClO é o Cl$_{2}$, por desproporcionamento, ou apenas tratar o enunciado como um reforço.%
\item%
{[}minor{]} A afirmação 'O$_{2}$ não é produto da eletrólise', feita em (b) e (c), é categórica demais. Em salmoura diluída, como a de piscina (cerca de 3 a 4 g/L), parte do O$_{2}$ é desprendida junto com o Cl$_{2}$. O produto principal é o Cl$_{2}$. Isso não muda a resposta, mas o ideal seria dizer 'não é o produto principal ou relevante'.%
\item%
{[}minor{]} O potencial de 1,23 V para O$_{2}$/H$_{2}$O é o valor padrão em meio ácido. Em pH próximo de 7 ele vale cerca de 0,82 V, o que reforça ainda mais o argumento da sobretensão. A conclusão continua certa. Para coerência com o -0,83 V citado em meio básico no catodo, valeria indicar o meio.%
\end{itemize}%
\subsection*{Questão 14: nota 9,5 (aprovada, rodada 1 de 1)}%
\textbf{Resposta da resolução às cegas do avaliador:} Alternativa b. A absorção é mais rápida que a metabolização: a concentração sobe em cerca de 5--6 h e cai ao longo de cerca de 18--19 h. Pelo gráfico, a dose diária é de 35 mg (0,7 mg/kg/dia). A dose acumulada mínima de 120 mg/kg $\times$ 50 kg = 6000 mg exige 6000/35 $\approx$ 171 dias, ou seja, cerca de 170 dias.\par %
Conferi o enunciado e a figura na página 7 do PDF original. A transcrição está correta: a dose acumulada sobe em 10 degraus de cerca de 35 mg até $\approx$ 350 mg em 240 h, e a concentração tem forma de dente de serra, com subida rápida, queda lenta e acúmulo até o estado estacionário. A resposta do candidato, (B), coincide com a minha referência e está correta.\newline%
\newline%
Contas conferidas:\newline%
{-} dose diária: 35 mg/dia = 0,7 mg/kg/dia, dentro da faixa da bula;\newline%
{-} faixa de dose acumulada: 120 $\times$ 50 = 6000 mg e 150 $\times$ 50 = 7500 mg;\newline%
{-} tempo mínimo: 6000/35 $\approx$ 171,4 dias, ou seja, cerca de 170 dias;\newline%
{-} tempo para o limite superior: 7500/35 $\approx$ 214,3 dias, ou seja, cerca de 215 dias.\newline%
\newline%
Armadilhas, todas tratadas:\newline%
{-} a dose diária foi lida do gráfico, e não tirada dos extremos da bula;\newline%
{-} multiplicou pela massa de 50 kg;\newline%
{-} associou 'no mínimo' ao limite inferior de 120 mg/kg e mostrou explicitamente que os 215 dias correspondem ao limite superior (é a armadilha que separa a da b);\newline%
{-} leu a dose no eixo da direita, sem confundir com o da concentração;\newline%
{-} concluiu absorção > metabolização a partir do balanço dC/dt = v\_abs - v\_met durante a subida;\newline%
{-} não arredondou a dose diária para 50 mg.\newline%
\newline%
As ressalvas são pequenas e de rigor:\newline%
{-} disse que a queda é 'só eliminação' sem justificar;\newline%
{-} usou o acúmulo como evidência direta da comparação entre as velocidades;\newline%
{-} não comentou que no estado estacionário as quantidades por ciclo se igualam;\newline%
{-} deu estimativas de taxa um pouco vagas.\newline%
\newline%
É uma resolução clara, correta e publicável como gabarito comentado.%
\begin{itemize}%
\item%
{[}minor{]} O candidato chama a fase de queda de 'só metabolização/eliminação' sem justificar. A absorção pode continuar, mais fraca, durante a queda. O correto é dizer que na queda a eliminação supera a entrada (dC/dt < 0). A conclusão não muda, mas a hipótese deveria vir explícita ou ser escrita com mais cuidado.%
\item%
{[}minor{]} O acúmulo de picos e vales foi usado como 'outra evidência' de que a absorção é maior que a metabolização. Na verdade, o acúmulo mostra que a eliminação é lenta em comparação com o intervalo de 24 h entre doses, e isso não prova sozinho a comparação entre as velocidades. Faltou também dizer explicitamente que no estado estacionário as quantidades absorvida e eliminada por ciclo se igualam, e que a pergunta compara velocidades (inclinações), não quantidades por ciclo.%
\item%
{[}minor{]} As taxas foram estimadas só no 1º ciclo e de forma vaga ('dezenas' e 'alguns' ng/mL/h). Ficaria mais claro dar os valores ($\approx$ 60 contra $\approx$ 6 ng$\cdot$mL$^{-}$$^{1}$$\cdot$h$^{-}$$^{1}$) ou usar também o ciclo estacionário ($\approx$ 250 ng/mL em $\approx$ 5 h contra $\approx$ 250 ng/mL em $\approx$ 19 h).%
\end{itemize}%
\subsection*{Questão 15: nota 10,0 (aprovada, rodada 1 de 1)}%
\textbf{Resposta da resolução às cegas do avaliador:} Foram vendidas 95 pizzas de brigadeiro (e 110 de calabresa): alternativa a) 95.\par %
Conferi o enunciado na página 8 do PDF e a transcrição está fiel: calabresa custa R\$ 45, brigadeiro custa R\$ 40, foram 205 pizzas e R\$ 8.750 arrecadados. A pergunta é sobre as pizzas de brigadeiro.\newline%
\newline%
O candidato chega ao mesmo resultado que a referência: b = 95, c = 110, alternativa a). Não há divergência.\newline%
\newline%
O candidato montou corretamente o sistema c + b = 205 e 45c + 40b = 8750, com os preços nos sabores certos. Resolveu por substituição sem erro de conta: 45$\cdot$205 = 9225, 9225 - 8750 = 475 e 475/5 = 95.\newline%
\newline%
Ele também apresenta o método da suposição e usa corretamente a variação de R\$ 5 por troca. Na verificação, testa as duas equações: 4950 + 3800 = 8750 e 110 + 95 = 205.\newline%
\newline%
Testei em Python as outras alternativas e os valores que ele dá estão certos: b = 100 dá 8725, b = 105 dá 8700 e b = 110 dá 8675.\newline%
\newline%
A unicidade fica justificada porque a arrecadação cai R\$ 5 a cada pizza de brigadeiro a mais, com o total fixo em 205. Isso equivale ao argumento do determinante não nulo.\newline%
\newline%
Ele trata todas as armadilhas, inclusive o distrator d) = 110, que é a quantidade de calabresa. Também verifica que a solução é inteira e não negativa.\newline%
\newline%
A resolução é clara, completa e rigorosa, com qualidade de gabarito comentado. Não encontrei nenhum problema.%
\subsection*{Questão 16: nota 9,5 (aprovada, rodada 1 de 1)}%
\textbf{Resposta da resolução às cegas do avaliador:} Alternativa **d**: o vermelho venceu com 2 pontos. A pedra vermelha (1/2, 1/5) é a mais próxima do centro (r $\approx$ 0,54). A melhor pedra azul é (-3/2, -1), então d = $\surd$13/2 $\approx$ 1,80. Ficam a menos de d só as vermelhas (1/2, 1/5) e (-1, 1), com r = $\surd$2 $\approx$ 1,41. A pedra (2, 0) tem r = 2 > d e não conta.\par %
Respostas finais: as duas dão a alternativa (d), Vermelho venceu com 2 pontos, e a resposta está correta. Conferi o enunciado na p. 8 do PDF (p. 6 do caderno) e ele bate com a transcrição. Refiz todas as contas com frações exatas em Python: os r$^2$ da azul são \{5; 13/4; 13; 17\} e os da vermelha são \{29/100; 4; 2; 8\}, logo d$^2$ = 13/4. As vermelhas com r$^2$ < 13/4 são 29/100 e 2, ou seja, 2 pontos. A ordem crescente das oito pedras é V, V, A, V, A, V, A, A, o que confirma o resultado. A resolução do candidato coincide passo a passo com a referência.\newline%
\newline%
Armadilhas, uma a uma:\newline%
{-} Compara r$^2$ com frações exatas e justifica por que pode fazer isso (r $\geq$ 0).\newline%
{-} Diz explicitamente que os raios dos círculos não entram na regra, então não cai na contagem por anel.\newline%
{-} Trata a pedra (2,0) do jeito certo: r = 2 > d $\approx$ 1,80, então ela não conta.\newline%
{-} Toma d como a distância da pedra adversária (azul) mais próxima do centro.\newline%
{-} Calcula 1/4 + 1/25 = 29/100 e 9/4 + 1 = 13/4 sem erro.\newline%
{-} Verifica que nenhuma pedra fica exatamente à distância d, então não há empate a tratar.\newline%
{-} Elimina as alternativas a, b e c com justificativa.\newline%
{-} A pedra (4,-1), fora do alvo, fica coberta pela observação de que os raios do alvo não importam.\newline%
\newline%
Só há falhas de redação, que não afetam a correção. A resolução está completa, rigorosa e clara, e serve como gabarito comentado.%
\begin{itemize}%
\item%
{[}minor{]} Na seção 'Vencedora da rodada', o texto diz 'A pedra azul mais próxima tem r$^2$=3,25' sem dizer explicitamente que esse é o mínimo de \{5; 13/4; 13; 17\}. A conclusão está correta e dá para tirá{-}la da lista anterior, mas valeria escrever d$^2$ = min\{5; 13/4; 13; 17\} = 13/4 logo ali, e não só na seção de pontuação.%
\item%
{[}minor{]} Detalhe de redação: 'Ela está menos de 1 unidade do centro' deveria ser 'a menos de 1 unidade'. A frase sobre o botão também é dispensável, porque a regra não usa os anéis. Não afeta a correção.%
\end{itemize}%
\subsection*{Questão 17: nota 9,5 (aprovada, rodada 1 de 1)}%
\textbf{Resposta da resolução às cegas do avaliador:} det A(x) = -x$^3$ + 27x + 10 = -(x + 5)(x$^2$ - 5x - 2). É uma cúbica com coeficiente líder negativo e y(0) = 10 > 0. Tem mínimo local em (-3, -44), máximo local em (3, 64) e raízes -5, (5 $\mp$ $\surd$33)/2 $\approx$ -0,37 e 5,37. Alternativa **B**.\par %
Confronto com o PDF original (p. 8 do arquivo, Questão 17): o enunciado e os quatro gráficos batem com a transcrição.\newline%
\newline%
A resolução do candidato está correta. A resposta final, alternativa (B) com y = -x$^3$ + 27x + 10, coincide com a minha referência.\newline%
\newline%
Conferi o cálculo passo a passo:\newline%
{-} Na expansão de Laplace pela 1ª coluna, o cofator de a21 = -7 é -(10 - x), e o candidato chega certo a +7(10 - x).\newline%
{-} (-3 - x)(1 - x) - 27 = x$^2$ + 2x - 30.\newline%
{-} (2 - x)(x$^2$ + 2x - 30) = -x$^3$ + 34x - 60.\newline%
{-} Somando: -x$^3$ + 27x + 10.\newline%
\newline%
Também verifiquei numericamente em Python, para x = -5, -3, -1, 0, 1, 3, 5 e 7,5. O determinante bate com o polinômio e com a forma fatorada -(x + 5)(x$^2$ - 5x - 2). As raízes são -5 e (5 $\mp$ $\surd$33)/2 $\approx$ -0,372 e 5,372.\newline%
\newline%
Todas as armadilhas foram tratadas explicitamente:\newline%
{-} O grau é 3, porque x$^3$ só aparece no produto da diagonal. Isso elimina (c) e (d).\newline%
{-} O coeficiente líder é (-1)$^3$ = -1. Isso elimina (a).\newline%
{-} Não há termo em x$^2$.\newline%
{-} O termo independente é 10 > 0: a curva corta o eixo y um pouco acima da origem, e 10 é pequeno perto dos extremos -44 e 64.\newline%
{-} O estudo de sinal de y? dá um mínimo em (-3, -44), à esquerda do eixo y e abaixo do eixo x, e um máximo em (3, 64), à direita do eixo y e acima do eixo x.\newline%
{-} Há três raízes reais, e a posição de cada uma em relação aos extremos e à origem está correta.\newline%
\newline%
A comparação final com a figura (b) é precisa. Há apenas detalhes menores: a conferência só cobre o termo independente e a divisão polinomial não foi mostrada. A resolução tem nível de gabarito comentado.%
\begin{itemize}%
\item%
{[}minor{]} A 'conferência' em x = 0 só valida o termo independente (10). Erros nos coeficientes de x ou de x$^3$ passariam por ela. Uma conferência completa poderia usar traço = 0 (ausência de x$^2$) e soma dos menores principais = -27, ou então avaliar em um segundo ponto, por exemplo y(1) = 36 = det A(1). Mesmo assim, ela é independente da expansão principal (feita pela 1ª linha, enquanto a expansão foi pela 1ª coluna), e isso tem algum valor.%
\item%
{[}minor{]} A divisão por (x + 5) não aparece: o resultado -(x + 5)(x$^2$ - 5x - 2) é dado direto. Seria bom mostrar Briot{-}Ruffini ou a multiplicação de volta. Conferi e o resultado está correto.%
\end{itemize}%
\subsection*{Questão 18: nota 9,5 (aprovada, rodada 1 de 1)}%
\textbf{Resposta da resolução às cegas do avaliador:} O maior valor possível é b = 10, alternativa c. Ele é atingido com p(x) = (x + 1)$^2$(x + 8) = x$^3$ + 10x$^2$ + 17x + 8, de raízes -1 (dupla) e -8.\par %
Conferi o enunciado na página 9 do PDF (Questão 18): a transcrição está fiel. A resposta do candidato, (C) b = 10, coincide com a minha referência e está correta. Uma força bruta em Python sobre todas as ternas de divisores de 8 com produto -8 deu os valores possíveis de b: \{-8, -5, -3, -2, 1, 6, 7, 10\}. O máximo é 10, com as raízes (-1, -1, -8).\newline%
\newline%
A resolução trata todas as armadilhas:\newline%
{-} **Hipótese:** interpreta corretamente que as 3 raízes complexas, contadas com multiplicidade, são inteiras. Também justifica essa leitura com um contraexemplo válido: a raiz -n e duas raízes conjugadas $\pm$i$\surd$(8/n) dão produto -8 e b = n, então b não teria limite.\newline%
{-} **Sinais de Girard:** estão corretos (soma = -b, produto = -8), e o problema passa a ser minimizar a soma das raízes.\newline%
{-} **Raízes repetidas:** admite raízes repetidas, então não cai na resposta 7.\newline%
{-} **Enumeração:** lista completa no caso de três raízes negativas, e no caso de uma negativa dá a cota u + v - m $\geq$ -6, que está correta e é atingida em (-8, 1, 1).\newline%
{-} **Valor atingido:** verifica (x+1)$^2$(x+8) = x$^3$ + 10x$^2$ + 17x + 8.\newline%
{-} **Alternativas:** descarta 13 (a soma mínima é -10) e mostra que 6 e 7 são atingíveis, mas não são o máximo.\newline%
\newline%
A exposição é clara e organizada, no nível de um gabarito comentado. Só há duas observações pequenas de rigor, ambas fáceis de completar: justificar m $\leq$ 8 por m = 8/(uv), e dizer explicitamente que as raízes são não nulas.%
\begin{itemize}%
\item%
{[}minor{]} No caso de uma raiz negativa, a resolução usa m $\leq$ 8 sem dizer de onde vem. Bastaria escrever m = 8/(uv) $\leq$ 8, já que u, v $\geq$ 1. A cota está correta, mas faltou meia linha de justificativa.%
\item%
{[}minor{]} Não diz explicitamente que nenhuma raiz é zero (o produto é -8 $\neq$ 0). Isso garante que cada raiz seja só positiva ou só negativa, e por isso a divisão em dois casos cobre todas as possibilidades. Fica implícito em 'o número de raízes negativas é ímpar', mas um gabarito comentado registraria isso em uma frase.%
\end{itemize}%
\subsection*{Questão 19: nota 9,5 (aprovada, rodada 1 de 1)}%
\textbf{Resposta da resolução às cegas do avaliador:} A maré é mínima (0,6 m) quando sen{[}($\pi$/12)(2t+3){]} = -1. Isso dá t = 7,5 + 12k e, dentro de {[}0, 24{]}, t = 7,5 h e t = 19,5 h, ou seja, às 07h30 e às 19h30. Alternativa (d).\par %
A resposta final do candidato é (d), 07h30 e 19h30, a mesma da minha referência, e está correta. Conferi as contas com Python:\newline%
{-} $\theta$(0) = $\pi$/4 e $\theta$(24) = 51$\pi$/12 = 17$\pi$/4.\newline%
{-} No intervalo, só as soluções 3$\pi$/2 e 7$\pi$/2 servem. Os candidatos vizinhos ficam fora: 11$\pi$/2 > 17$\pi$/4 e -$\pi$/2 < $\pi$/4.\newline%
{-} Daí saem t = 7,5 h e t = 19,5 h, com A = 0,6 m em ambos.\newline%
{-} Em t = 7: sen(17$\pi$/12) $\approx$ -0,966 e A(7) $\approx$ 0,624 m, como o candidato afirma.\newline%
\newline%
O método de restringir $\theta$ ao intervalo imagem é elegante e mais rigoroso que só listar valores de k. Ele também cobre os extremos do domínio sem esforço extra, porque o seno vale -1 apenas nesses dois pontos.\newline%
\newline%
O candidato tratou todas as armadilhas:\newline%
{-} usou a amplitude positiva para concluir que o mínimo ocorre com sen = -1;\newline%
{-} igualou o argumento a 3$\pi$/2 + 2k$\pi$, e não a $\pi$/2;\newline%
{-} obteve o período correto de 12 h;\newline%
{-} restringiu as soluções ao domínio;\newline%
{-} converteu 0,5 h em 30 min;\newline%
{-} descartou os distratores próximos com um valor numérico.\newline%
\newline%
Restam só detalhes de exposição: não justifica que $\theta$ é crescente, não mostra a forma ($\pi$/6)t + $\pi$/4 e não dá valores numéricos para t = 6 e t = 6,5. A resolução tem qualidade de gabarito comentado.%
\begin{itemize}%
\item%
{[}minor{]} A afirmação de que $\theta$ é crescente em t não é justificada. A justificativa é curta: o coeficiente de t em $\theta$(t) = ($\pi$/6)t + $\pi$/4 é positivo. Basta uma frase.%
\item%
{[}minor{]} Na verificação do período, o candidato escreve $\omega$ = 2$\pi$/12 sem mostrar que esse valor vem de ($\pi$/12)$\cdot$2. É melhor escrever explicitamente $\theta$(t) = ($\pi$/6)t + $\pi$/4, porque isso também deixa claro o deslocamento de fase de 1,5 h. Essa é a armadilha das alternativas (a) e (c).%
\item%
{[}minor{]} Para t = 6 e t = 6,5, o candidato só afirma que sen $\theta$ > -1, sem dar valores. Seria mais completo informar A(6) $\approx$ 0,805 m e A(6,5) $\approx$ 0,694 m, como fez com A(7) $\approx$ 0,62 m. Essa omissão não afeta a correção, porque o mínimo global já está estabelecido.%
\end{itemize}%
\subsection*{Questão 20: nota 9,5 (aprovada, rodada 1 de 1)}%
\textbf{Resposta da resolução às cegas do avaliador:} A área de B é $\Sigma$\_\{n$\geq$0\} (1/3)\^{}n = 1/(1 - 1/3) = 3/2. Alternativa c.\par %
Conferi o enunciado no original (p. 9 do PDF) e a transcrição está correta. A resposta do candidato é (c) 3/2, a mesma da minha referência, e está certa.\newline%
\newline%
A resolução está completa e correta:\newline%
{-} Calcula a área de cada retângulo, 1/3\^{}n.\newline%
{-} Justifica por que a área da união é a soma das áreas: os retângulos só se encostam em lados verticais, de área nula.\newline%
{-} Identifica a PG infinita com a$_{0}$ = 1 e q = 1/3, mostra que converge porque |q| < 1 e aplica a fórmula correta a$_{0}$/(1-q) = 3/2.\newline%
{-} Verifica que A ? B porque f é decrescente.\newline%
{-} Confere a ordem de grandeza: Área(A) = 1/ln 3 $\approx$ 0,910 < 1,5.\newline%
{-} Mostra as somas parciais 13/9 $\approx$ 1,444 e 40/27 $\approx$ 1,481. Conferi tudo isso numericamente.\newline%
\newline%
Todas as armadilhas foram tratadas:\newline%
{-} (d) 40/27 é a soma só dos quatro retângulos desenhados.\newline%
{-} A soma começa em n = 0, então o primeiro retângulo, de altura 1, está incluído.\newline%
{-} A fórmula da PG está certa e a convergência foi justificada.\newline%
{-} (a) 0,9 < Área(A), então não pode ser estimativa por excesso.\newline%
{-} (b) foi descartada.\newline%
\newline%
Só há ressalvas pequenas de redação: (b) também é a razão da PG; somas parciais não 'confirmam' um limite; a cobertura de x $\geq$ 0 pelos intervalos ficou implícita. A resolução tem qualidade de gabarito comentado.%
\begin{itemize}%
\item%
{[}minor{]} Na alternativa (b), o candidato diz que 3/9 = 1/3 é só a área do segundo retângulo. O valor também é a razão q da PG, que é a confusão mais provável de quem marca (b). Vale citar as duas leituras.%
\item%
{[}minor{]} A frase 'as somas parciais crescem e se aproximam de 1,5, o que confirma o limite' exagera um pouco. Três ou quatro somas parciais não confirmam um limite, apenas mostram que o resultado é coerente. Quem prova a convergência é o critério |q| < 1, que o candidato já citou. O certo seria chamar isso de checagem numérica.%
\item%
{[}minor{]} Para mostrar que A ? B, o candidato usa que f é decrescente. Ficou implícito que os intervalos {[}n, n+1{]} com n $\geq$ 0 cobrem todo o semieixo x $\geq$ 0, e isso poderia ser dito numa linha. A colagem {[}{-}3\^{}\{{-}x\}/ln 3{]}\_0\^{}$\infty$ também usa um limite no infinito, que é notação usual para integral imprópria. Detalhes sem impacto.%
\end{itemize}%
\subsection*{Questão 21: nota 10,0 (aprovada, rodada 1 de 1)}%
\textbf{Resposta da resolução às cegas do avaliador:} a = 1, alternativa (a). Com a razão r, os termos são b = a + r, c = a + 2r e d = a + 3r. Então det M = -2r$^2$ = -8 dá r = -2 (porque r < 0), e o traço 2a + 3r = -4 dá a = 1. A matriz fica M = {[}{[}1, -1{]},{[}-3, -5{]}{]}.\par %
Conferi o enunciado na página 9 do PDF (Questão 21) e a transcrição está correta. A resposta do candidato, a = 1 (alternativa a), é a mesma da minha referência e está certa.\newline%
\newline%
A resolução está correta, completa e bem escrita:\newline%
{-} **PA:** a PA foi montada na ordem certa, lendo por linhas: b = a + r, c = a + 2r, d = a + 3r.\newline%
{-} **Traço:** o traço foi tomado como a soma da diagonal principal, a + d = 2a + 3r = -4.\newline%
{-} **Determinante:** usou ad - bc e expandiu sem erro de sinal ou coeficiente. Os termos a$^2$ e 3ar se cancelam e sobra -2r$^2$. O candidato ainda observa que o determinante não depende de a.\newline%
{-} **Raiz r = +2:** de r$^2$ = 4 sai r = $\pm$2, e r = +2 é descartado explicitamente pela hipótese r < 0. O candidato também mostra que r = +2 levaria a a = -5, a alternativa (d), que é o distrator. Conferi essa conta: 2a + 6 = -4 dá a = -5.\newline%
{-} **Verificação:** a PA fica (1, -1, -3, -5), a matriz fica M = {[}{[}1, -1{]},{[}-3, -5{]}{]}, e de fato det M = -8 e tr M = -4.\newline%
\newline%
Todas as armadilhas que listei foram tratadas: o descarte de r = +2, a expansão do determinante, a definição de traço, a ordem da PA, a fórmula do determinante 2$\times$2 e a verificação final. Não encontrei nenhum erro, salto lógico ou omissão. A resolução serve como gabarito comentado.%
\subsection*{Questão 22: nota 9,5 (aprovada, rodada 1 de 1)}%
\textbf{Resposta da resolução às cegas do avaliador:} Alternativa **B**: \{(x,y) $\in$ $\mathbb{R}$$^2$: |x| $\leq$ y $\leq$ 2 {-} |x|\}. A fronteira inferior do quadrado é y = |x| e a superior é y = 2 {-} |x|. A restrição |x| $\leq$ 1 já decorre das duas desigualdades (|x| $\leq$ 2 {-} |x|).\par %
Conferi o enunciado no PDF (página 9, questão 22). A transcrição está correta, só falta uma vírgula entre (-1,1) e (0,2). A resposta do candidato, (B), coincide com a minha referência e está correta. Também conferi com uma malha racional exata de passo 1/20, comparando com o quadrado |x|+|y-1| $\leq$ 1: (B) coincide em todos os pontos, (C) é vazia, e (A) e (D) divergem em milhares de pontos. Os contraexemplos do candidato estão certos: (2,1) pertence a (A) e fica fora do quadrado; (0,1) está no quadrado e não pertence a (A); (5,6) pertence a (D) e fica fora do quadrado.\newline%
\newline%
A resolução é completa e bem organizada:\newline%
{-} Deduz as quatro retas dos lados e as junta nas fronteiras y = |x| (embaixo) e y = 2-|x| (em cima).\newline%
{-} Descreve a região por cortes verticais.\newline%
{-} Mostra que a restrição |x| $\leq$ 1 já decorre da dupla desigualdade, uma armadilha importante.\newline%
{-} Confirma tudo por um segundo caminho, a "bola" da norma da soma |x|+|y-1| $\leq$ 1, centrada em (0,1).\newline%
{-} Descarta cada alternativa com um argumento próprio: (C) é vazia por exigir 0 $\geq$ 2; (A) exige |x| $\geq$ 1 e tem contraexemplos dos dois tipos (ponto do quadrado que fica de fora e ponto de fora que entra); (D) é ilimitada.\newline%
\newline%
Todas as armadilhas que listei foram tratadas: a ordem das fronteiras, o sinal em 2-|x|, a restrição implícita |x| $\leq$ 1 e o fato de (D) conter o quadrado sem ser igual a ele. Restam só detalhes: a equivalência com |t| $\leq$ a foi citada apenas para a $\geq$ 0 (o caso |x| > 1 fica implícito), a conferência dos vértices mostra só um deles, e a vírgula faltando no enunciado não foi comentada. Do jeito que está, a resolução poderia ser publicada como gabarito comentado.%
\begin{itemize}%
\item%
{[}minor{]} A equivalência |y-1| $\leq$ 1-|x| ? |x| $\leq$ y $\leq$ 2-|x| foi justificada com a regra |t| $\leq$ a ? -a $\leq$ t $\leq$ a, citada só para a $\geq$ 0. O caso 1-|x| < 0 (|x| > 1) não foi tratado nesse ponto. A equivalência continua válida, porque os dois lados são falsos, e isso decorre da observação anterior de que a dupla desigualdade força |x| $\leq$ 1. Mesmo assim, valeria uma linha para fechar o caso ali mesmo.%
\item%
{[}minor{]} O texto diz que 'os quatro vértices satisfazem a condição', mas só mostra a conta para (1,1). Num gabarito comentado, convém listar os quatro de forma explícita: (0,0): 0$\leq$0$\leq$2; ($\pm$1,1): 1$\leq$1$\leq$1; (0,2): 0$\leq$2$\leq$2.%
\item%
{[}minor{]} Falta uma vírgula entre (-1,1) e (0,2) no enunciado original, e a resolução não comenta isso. A figura também não é usada para confirmar que o interior está sombreado. Não muda nada no resultado, é só um detalhe de leitura atenta do enunciado.%
\end{itemize}%
\subsection*{Questão 23: nota 9,5 (aprovada, rodada 1 de 1)}%
\textbf{Resposta da resolução às cegas do avaliador:} c) {[}1,2{]}. É o único intervalo que não contém o vértice x = 1/2 da parábola no seu interior. Nele f é estritamente crescente e, portanto, injetora. Tomando B = {[}-2, 0{]}, a inversa é f$^{-}$$^{1}$(y) = 1/2 + $\surd$(y + 9/4).\par %
Conferi a transcrição com o PDF original (p. 9, Questão 23) e ela está correta. O candidato responde C) {[}1,2{]}, a mesma resposta da minha referência, e a resposta está certa.\newline%
\newline%
A resolução é de gabarito. Ela começa dizendo explicitamente que invertível equivale a bijetora e que, como B é livre, a sobrejetividade se resolve tomando B = f(A); assim a condição que restringe A é a injetividade. Em seguida o candidato fatora f(a) - f(b) = (a - b)(a + b - 1). Esse é um critério algébrico mais forte que o argumento só pelo gráfico: ele mostra que f(a) = f(b) com a $\neq$ b exatamente quando a + b = 1, ou seja, quando os dois pontos são simétricos em relação a x = 1/2. O candidato confere isso com a forma canônica (x - 1/2)$^2$ - 9/4.\newline%
\newline%
As alternativas erradas são descartadas com contraexemplos concretos e corretos:\newline%
{-} a) e d): f(0) = f(1) = -2;\newline%
{-} b): f(-1) = f(2) = 0.\newline%
\newline%
Em c), ele prova a injetividade de {[}1,2{]} por a + b $\geq$ 2 > 1 e também pela monotonicidade em {[}1/2, $\infty$). Por fim, dá B = {[}-2, 0{]} e a inversa f$^{-}$$^{1}$(y) = 1/2 + $\surd$(y + 9/4), verificada nos extremos. Conferi tudo numericamente: só em {[}1,2{]} a função é monótona, a imagem é {[}-2, 0{]} e a fórmula da inversa confere.\newline%
\newline%
Todas as armadilhas foram tratadas: a posição do vértice, a insuficiência de comparar só os extremos (resolvida com contraexemplos internos), a escolha de B, o critério de ponto interior e a justificativa de c). Restam apenas detalhes mínimos de redação, como a ida do critério do intervalo, que é óbvia mas não foi escrita.%
\begin{itemize}%
\item%
{[}minor{]} A frase 'um intervalo serve se, e somente se, 1/2 não for ponto interior dele' é afirmada sem demonstrar a ida (se 1/2 for interior, os pontos 1/2 $\pm$ t pertencem ao intervalo e somam 1). Isso sai direto do critério a + b = 1 que o próprio candidato estabeleceu, e as alternativas são todas descartadas com contraexemplos explícitos, então a lacuna não afeta a resposta. Bastaria uma linha a mais para fechar o argumento.%
\item%
{[}minor{]} No item d), o candidato não repete a observação 0 + 1 = 1, como fez em a). O contraexemplo f(0) = f(1) = -2 está correto e basta; a falta é só de uniformidade na redação.%
\end{itemize}%
\subsection*{Questão 24: nota 9,5 (aprovada, rodada 1 de 1)}%
\textbf{Resposta da resolução às cegas do avaliador:} P(mdc ímpar) = 1 - C(50,2)/C(100,2) = 1 - 49/198 = 149/198. Alternativa **d**.\par %
Conferi a transcrição com o PDF original (p. 10), e ela confere: sorteio sem reposição de 1 a 100, alternativas a) 1/2, b) 50/198, c) 100/199, d) 149/198. O candidato e a minha referência chegam à mesma resposta, (d) 149/198. A força bruta em Python confirma: o mdc é ímpar em 3725 dos 4950 pares não ordenados, e 3725/4950 = 149/198 $\approx$ 0,7525.\newline%
\newline%
A resolução está correta e bem organizada. Ela define o espaço amostral não ordenado, C(100,2) = 4950, e justifica a equivalência central: o mdc é par $\Leftrightarrow$ 2 | a e 2 | b. Também dá exemplos que separam 'mdc ímpar' de 'mdc = 1', como mdc(3,9) = 3. Em seguida, calcula o complementar de forma coerente de duas maneiras: C(50,2)/C(100,2) e 50/100 $\cdot$ 49/99, ambas dando 49/198. Por fim, chega a 1 - 49/198 e ainda explica por que o valor fica um pouco acima de 3/4, o que trata bem a armadilha do sorteio sem reposição.\newline%
\newline%
Todas as armadilhas foram cobertas: a equivalência justificada, mdc ímpar $\neq$ mdc 1, sorteio sem reposição, coerência entre numerador e denominador, não responder com o complementar e os 50 pares/50 ímpares. Os problemas que encontrei são menores: chamar de 'definição' o que é uma propriedade do mdc, a explicação fraca dos distratores (b) e (c) e uma conferência que só repete a conta do complementar. A resolução serve como gabarito comentado.%
\begin{itemize}%
\item%
{[}minor{]} O candidato diz que 'um inteiro divide d se e somente se divide a e b' vem 'pela definição de mdc'. Na definição escolar, porém, o mdc é o maior divisor comum, e essa equivalência é uma propriedade dele (sai de Bézout ou da fatoração em primos). Nesta questão, só a ida é realmente necessária: se 2 | d, então 2 | a e 2 | b. Essa parte segue da definição, porque d divide a e b. A volta também é imediata: se 2 é divisor comum, 2 $\leq$ d e, mais precisamente, 2 | d. Mesmo assim, chamar isso de 'definição' é impreciso. O certo seria dizer 'propriedade do mdc' ou justificar em uma linha.%
\item%
{[}minor{]} A explicação do distrator (b) 50/198 é superficial. O candidato só diz que o valor 'fica muito próximo' de 49/198, sem mostrar um erro de raciocínio que leve a ele. Um mecanismo plausível é contar 50$\cdot$50 casos sobre 100$\cdot$99 pares ordenados, misturando critérios. Para (c), o argumento também é só numérico. Como a análise dos distratores não era obrigatória, o peso disso é pequeno.%
\item%
{[}minor{]} Na conferência, os 3725 casos favoráveis foram obtidos por diferença (4950 - 1225), o que repete a conta do complementar. Uma contagem direta independente, 50$\cdot$50 (um par e um ímpar) + C(50,2) (dois ímpares) = 2500 + 1225 = 3725, valeria como verificação de verdade. A estimativa heurística $\approx$ 3/4 compensa em parte.%
\end{itemize}%
\subsection*{Questão 25: nota 9,5 (aprovada, rodada 1 de 1)}%
\textbf{Resposta da resolução às cegas do avaliador:} A resposta é a alternativa **d**: 1/16 < A < 3/32.\newline%
\newline%
Como 1/8 < f(1/2) < 1/4, a área é A = (1/4)(1/2 - f(1/2)). Isso dá 1/16 < A < 3/32.\par %
Conferi o enunciado no original (p. 10 do PDF, Questão 25) e a transcrição está fiel. O candidato marcou a alternativa d (1/16 < A < 3/32), a mesma da minha referência. Refiz as contas e a resposta está certa.\newline%
\newline%
Os passos do candidato:\newline%
{-} Aplica a hipótese em x = 1/2 e obtém 1/8 < f(1/2) < 1/4.\newline%
{-} Mostra que f(1/2) < 1/2, então Q fica abaixo de R e o triângulo não é degenerado. Isso também justifica tirar o módulo.\newline%
{-} Toma o lado vertical QR = 1/2 - f(1/2) como base, com altura 1/2 (distância de P à reta x = 1/2), e chega a A = (1/4)(1/2 - f(1/2)).\newline%
{-} Confere pela fórmula do determinante com P na origem: $\frac12$|x\_Q$\cdot$y\_R - x\_R$\cdot$y\_Q| = $\frac12$|1/4 - f/2|. É a mesma expressão.\newline%
{-} Inverte as desigualdades de forma explícita ao multiplicar por -1 e mantém o sentido ao multiplicar por 1/4 > 0. Obtém 1/16 < A < 3/32 com os extremos na ordem certa (f = 1/8 dá o máximo 3/32 e f = 1/4 dá o mínimo 1/16).\newline%
{-} Descarta as alternativas a, b e c porque estão contidas em (0, 1/16).\newline%
\newline%
Todas as armadilhas da minha lista foram tratadas: a base e a altura corretas (sem cair no distrator c, que vem de A = f(1/2)/4), a inversão por A ser decrescente em f, os valores (1/2)$^3$ = 1/8 e (1/2)$^2$ = 1/4, a não degenerescência, o fato de só f(1/2) importar, e as desigualdades estritas.\newline%
\newline%
Só há ressalvas pequenas: a afirmação de que os extremos podem ser aproximados tanto quanto se queira não vem com um exemplo de f, e a letra da alternativa está maiúscula. É uma resolução clara, rigorosa e completa, no nível de um gabarito comentado.%
\begin{itemize}%
\item%
{[}minor{]} Na conclusão, o candidato diz que os extremos 1/16 e 3/32 'podem ser aproximados tanto quanto se queira', mas não justifica. A afirmação é verdadeira e bastaria um exemplo: f(x) = x$^3$ + t(x$^2$ - x$^3$), com t $\in$ (0,1), satisfaz a hipótese e dá f(1/2) = 1/8 + t/8, que percorre todo o intervalo (1/8, 1/4). Para escolher a alternativa isso não é necessário, mas a frase fica sem demonstração. Ou se inclui o exemplo, ou se retira a afirmação.%
\item%
{[}minor{]} Detalhe de notação: as alternativas da prova são minúsculas (d), mas o candidato escreveu 'ANSWER: (D)'. Não muda a resposta.%
\end{itemize}%
\subsection*{Questão 26: nota 9,0 (aprovada, rodada 1 de 1)}%
\textbf{Resposta da resolução às cegas do avaliador:} AP = $\surd$(2 - $\surd$2) $\approx$ 0,765. Alternativa a). Como BC = OC = r, o ângulo COB vale 15$^\circ$. Daí OCA = OAC = 30$^\circ$, e o ângulo externo dá AOP = 30$^\circ$ + 15$^\circ$ = 45$^\circ$. Pela lei dos cossenos, AP$^2$ = 2 - 2cos45$^\circ$ = 2 - $\surd$2.\par %
Conferi o enunciado na página 10 do PDF e a transcrição está fiel. A resposta final do candidato é a alternativa (a), AP = $\surd$(2-$\surd$2) $\approx$ 0,765, a mesma da minha referência. Está correta.\newline%
\newline%
A cadeia de raciocínio é a mesma da referência e está toda certa:\newline%
{-} OC = CB = r torna o triângulo OCB isósceles, então $\angle$COB = $\angle$CBO = 15$^\circ$ e $\angle$OCB = 150$^\circ$.\newline%
{-} O suplemento dá $\angle$OCA = 30$^\circ$, e o isósceles OAC dá $\angle$OAC = 30$^\circ$.\newline%
{-} Pelo ângulo externo do triângulo AOB, $\angle$AOP = 30$^\circ$ + 15$^\circ$ = 45$^\circ$.\newline%
{-} A lei dos cossenos com os raios OA = OP = 1 dá AP$^2$ = 2 - $\surd$2.\newline%
\newline%
O candidato tratou as armadilhas principais:\newline%
{-} usou OC = r;\newline%
{-} identificou corretamente qual ponto da secante é A e qual é C;\newline%
{-} obteve $\angle$OCA = 30$^\circ$, e não 15$^\circ$;\newline%
{-} usou $\angle$AOP = 45$^\circ$, e não o ângulo central de 135$^\circ$, que levaria à distratora (b);\newline%
{-} usou os raios na lei dos cossenos, e não o diâmetro.\newline%
\newline%
Ele também fez uma conferência por coordenadas, que está correta embora seja em parte circular, e um teste de plausibilidade: a corda de 45$^\circ$ é menor que o raio.\newline%
\newline%
Problemas encontrados:\newline%
{-} Os valores numéricos das alternativas (c) e (d) estão errados: ele escreve 1,612 e 2,298, quando o certo é 1,608 e 2,327. A conclusão não muda.\newline%
{-} Não diz que as posições "entre" (O entre P e B, C entre A e B) foram lidas da figura.\newline%
{-} Não justifica a unicidade da configuração, descartando o caso de 165$^\circ$, que é degenerado.\newline%
\newline%
São falhas menores. A resolução é clara e rigorosa no essencial e serve como gabarito depois que esses números forem corrigidos.%
\begin{itemize}%
\item%
{[}minor{]} Dois dos valores numéricos das alternativas descartadas estão errados. O candidato dá $\approx$1,612 para (c) e $\approx$2,298 para (d). Os valores corretos são $\surd$(4-$\surd$2) $\approx$ 1,608 e $\surd$(4+$\surd$2) $\approx$ 2,327, conferidos em python. Isso não muda a conclusão, já que (d) continua maior que o diâmetro 2. Mesmo assim, números errados não deveriam aparecer num gabarito comentado, e as contas devem ser refeitas.%
\item%
{[}minor{]} O candidato põe nos 'Dados' que O fica entre P e B e que C fica entre A e B. O texto do enunciado só fala em colinearidade, e essas posições vêm da figura. Faltou dizer que foram lidas da figura e explicar por que a configuração é única: pela lei dos senos, $\angle$COB = 15$^\circ$ ou 165$^\circ$, e o caso de 165$^\circ$ degenera o triângulo OCB. Com isso, $\angle$COB = $\angle$CBO = 15$^\circ$ sai demonstrado, e não apenas afirmado a partir do isósceles.%
\item%
{[}minor{]} A conferência é em parte circular. Ela já parte de A em 135$^\circ$, que é o resultado obtido, e verifica apenas que a reta por B passa por esse ponto. O teste é válido como verificação de consistência, mas seria mais forte calcular as interseções da reta com a circunferência e checar que BC = 1 e BA = 1+$\surd$3, por exemplo pela potência de ponto BA$\cdot$BC = OB$^2$ - 1.%
\end{itemize}%
\subsection*{Questão 27: nota 9,5 (aprovada, rodada 1 de 1)}%
\textbf{Resposta da resolução às cegas do avaliador:} Alternativa D. No texto 2, "fantasia" tem o sentido 4 do dicionário (traje de carnaval), usado em sentido figurado. O folião não vestia um traje de verdade. Sua "fantasia" é a frase "Não mate sua mina" escrita no corpo, uma ordem moral contra o feminicídio. As aspas e o "NOTA: 10", que lembra o quesito Fantasia dos desfiles, marcam esse uso figurado e elogioso.\par %
Conferi a página 10 do PDF original. A transcrição bate com o original: título \textbackslash{}"'FANTASIA' NOTA:10\textbackslash{}" com FANTASIA entre aspas, tapete vermelho e confetes. As duas respostas finais coincidem: alternativa D, que está correta. A palavra 'fantasia' aparece no sentido 4 (traje de carnaval), usado de forma figurada. A 'fantasia' do folião é a frase \textbackslash{}"Não mate sua mina\textbackslash{}" escrita nas costas, um imperativo moral contra o feminicídio.\textbackslash{}n\textbackslash{}nA resolução trata todas as armadilhas da referência:\textbackslash{}n{-} o contexto do carnaval, que leva ao sentido 4;\textbackslash{}n{-} as aspas como marca de uso não literal;\textbackslash{}n{-} o 'NOTA: 10' ligado ao quesito Fantasia dos desfiles;\textbackslash{}n{-} a legenda 'Quando o óbvio precisa ser escancarado', que mostra que a mensagem é real e não utópica nem ficcional;\textbackslash{}n{-} o imperativo negativo 'Não mate', que justifica a expressão 'imperativo moral';\textbackslash{}n{-} a ligação de cada distrator com as acepções 1, 2 e 3 do dicionário.\textbackslash{}n\textbackslash{}nAs eliminações de a, b e c estão corretas e bem justificadas. Os problemas são só de detalhe: a grafia do título, o uso do termo 'charge', a falta de comentário sobre 'obra' e 'apelo emotivo' nas alternativas b e c, e a gíria 'mina' sem explicação. A resolução tem nível de gabarito comentado.%
\begin{itemize}%
\item%
{[}minor{]} Detalhe da transcrição: o título no original é "'FANTASIA' NOTA:10", em maiúsculas, e o candidato escreve "Fantasia". Também chama o texto 2 de "charge", quando ele é uma ilustração comemorativa sem intenção crítica ou humorística. Nada disso afeta o raciocínio, mas o melhor é citar o título como ele aparece e chamar o texto 2 de ilustração.%
\item%
{[}minor{]} A refutação de (b) e (c) se apoia só nos termos 'imaginário' e 'ficcional'. Fica de fora outra falha das duas: o uso de 'fantasia' no texto 2 não vem de uma 'obra' elaborada pelo folião (b) nem de um 'apelo emotivo' idealizado (c), e sim do traje figurado. Uma frase sobre isso deixaria a eliminação mais completa.%
\item%
{[}minor{]} O candidato não explica a gíria 'mina' (namorada, mulher). Ele chega ao tema do feminicídio, então entendeu o sentido, mas um gabarito comentado deveria dizer isso de forma explícita.%
\end{itemize}%
\subsection*{Questão 28: nota 9,5 (aprovada, rodada 1 de 1)}%
\textbf{Resposta da resolução às cegas do avaliador:} Alternativa B: no texto, o carnaval aparece ligado ao desengano porque a fantasia acaba e se volta à realidade quando a festa termina. O refrão "Quarta{-}feira sempre desce o pano" marca esse fim. A dor ficou "em casa me esperando" e volta depois da festa. O "rei" era só fantasia, e a "morena" "hoje nem lembra não".\par %
Conferi a transcrição com o original (págs. 10--11 do PDF, Q28) e ela está correta. A resposta do candidato é B, a mesma da minha referência, e B está certa. A resolução define bem 'desengano' como perda de ilusão, que é o ponto que derruba quem lê a palavra só como 'tristeza' e marca a D. Ela também se apoia em vários trechos da letra: a dor 'em casa me esperando', que fica adiada e não aumenta; o 'vestido de rei' como fantasia no sentido literal e no figurado; a imagem teatral de 'desce o pano', com o 'sempre' indicando um fim certo; e a morena que 'hoje nem lembra não'. O candidato tratou todas as armadilhas. Na (a), mostra que a confraternização aparece ligada a 'esperança', o lado oposto ao desengano. Na (c), mostra que foi um encontro passageiro, não um grande amor, e que o desengano vem do fim dele. Na (d), mostra que a dor fica suspensa durante a festa, não ampliada, e que a alternativa inverte a relação de causa. Os problemas que encontrei são só de acabamento: o 'rei' não volta a ser citado na análise da (b), há uma frase um pouco forte ('não está na festa') e a oposição entre as estrofes podia ser mais explícita. A resolução pode ser publicada como gabarito comentado.%
\begin{itemize}%
\item%
{[}minor{]} Na análise da alternativa (b), a resolução cita só a dor que esperava em casa e o amor esquecido. Deixou de fora o 'vestido de rei', que ela mesma tinha apresentado antes como fantasia. Ficaria mais completo dizer que essa fantasia também acaba quando 'desce o pano', o que fecha as três ilusões desfeitas.%
\item%
{[}minor{]} A frase 'O desengano, portanto, não está na festa' está um pouco forte. A canção liga as duas coisas já no verso 'Carnaval, desengano', e o próprio texto do candidato diz logo depois que o desengano está na passagem da festa para a quarta{-}feira. Seria mais preciso dizer que o desengano não vem do que se vive durante a festa, mas do fim dela.%
\item%
{[}minor{]} A oposição entre as estrofes (as duas primeiras ligam carnaval e desengano, a última liga carnaval e esperança) só aparece dentro da análise da alternativa (a). Valeria apresentá{-}la como chave de leitura do texto todo. A conclusão está certa assim mesmo.%
\end{itemize}%
\subsection*{Questão 29: nota 9,0 (aprovada, rodada 1 de 1)}%
\textbf{Resposta da resolução às cegas do avaliador:} A resposta é a alternativa d), presente do subjuntivo (viva, possa, saiba, nos versos 12, 13 e 14). Essas formas completam o substantivo "esperança" do verso 11 e exprimem desejos ainda não realizados, ou seja, a expectativa por "o que ainda não é". Mesmo com o nome de "presente", elas apontam para o futuro, o que confirma a ideia do texto de que os tempos verbais não são etiquetas rígidas.\par %
Conferi a transcrição com o original (p. 9 do caderno, p. 11 do PDF): ela confere, inclusive na numeração dos versos sem o título (12 viva, 13 possa, 14 saiba). O candidato marca (d), a mesma resposta da minha referência, e está correta. As formas viva, possa e saiba, no presente do subjuntivo, completam o substantivo 'esperança' do v. 11 ('esperança de que...') e exprimem desejos ainda não realizados, isto é, a expectativa agostiniana pelo 'que ainda não é'.\newline%
\newline%
Pontos fortes:\newline%
{-} Resume com fidelidade a tese do texto de apoio: memória, atenção e expectativa como modos de experiência, e tempos verbais que mostram a posição do falante.\newline%
{-} Liga corretamente as três orações com 'Que' ao substantivo 'esperança'.\newline%
{-} Descarta cada distrator com o argumento certo: (a) pretérito perfeito ligado à memória; (b) gerúndio, que indica ação em curso presa a fatos passados, resistindo à tentação do sentido de 'esperar/sonhar'; (c) presente habitual ('sempre desce o pano') e presente atual ('hoje nem lembra'), que são constatação e atenção, não desejo.\newline%
{-} Associa cada alternativa a uma das três dimensões de Agostinho.\newline%
\newline%
Lacunas, todas menores:\newline%
{-} Não usa o caso do 'presente' do subjuntivo que aponta para o futuro como exemplo explícito da tese das 'etiquetas não rígidas'.\newline%
{-} Não observa que quem espera, no v. 2, é a dor (o desengano).\newline%
{-} Há uma pequena imprecisão terminológica entre tempo e modo.\newline%
{-} Não comenta 'lembrança' (v. 12).\newline%
\newline%
Mesmo assim, a resolução é correta, completa e clara, no nível de um gabarito comentado.%
\begin{itemize}%
\item%
{[}minor{]} A resolução não diz explicitamente que as formas se chamam presente do subjuntivo e mesmo assim apontam para o futuro desejado. Esse é o melhor exemplo da tese de Módolo e Braga de que os tempos verbais 'não funcionam como etiquetas rígidas'. O candidato cita a tese de forma genérica na seção 1, mas não a aplica à alternativa correta. Deveria acrescentar uma frase como: 'embora nomeado presente, o subjuntivo aqui projeta o que ainda não é, o que confirma que o nome do tempo não é etiqueta rígida'.%
\item%
{[}minor{]} Na alternativa (b), a justificativa está certa: o gerúndio indica ação em curso e está preso a fatos passados ('deixei', 'deixou'). Faltou reforçar que quem 'espera' é a dor, ou seja, o desengano da volta à realidade, e não uma esperança do eu lírico, e que o sonho provocado pela morena já se desfez ('hoje nem lembra não'). Isso desmontaria por completo a armadilha do significado da palavra.%
\item%
{[}minor{]} A frase '{[}as formas{]} estão no presente do subjuntivo, que é o modo do desejo' está um pouco frouxa na terminologia: o subjuntivo é o modo, e o presente é o tempo. O mais exato seria 'no presente do subjuntivo; o subjuntivo é o modo do desejo, do virtual, aqui com valor optativo'.%
\item%
{[}minor{]} Não comenta a possível distração de 'lembrança' (v. 12), palavra que remete à memória. O ponto é que a forma verbal 'viva' projeta um desejo. É uma omissão pequena, porque a pergunta trata das formas verbais.%
\end{itemize}%
\subsection*{Questão 30: nota 9,5 (aprovada, rodada 1 de 1)}%
\textbf{Resposta da resolução às cegas do avaliador:} Alternativa A: postura discriminatória disfarçada de enaltecimento. "Louvor{-}prisão" é o elogio ("você passa bem café") que, na verdade, rebaixa Tony e o prende a um papel doméstico e subalterno. O texto resume isso como "preconceito com açúcar".\par %
Conferi a página 11 do PDF (página 9 do caderno). A transcrição está fiel e a questão não tem figura. A resposta do candidato, alternativa A ("postura discriminatória disfarçada de enaltecimento"), é a mesma da minha referência e está correta.\newline%
\newline%
A resolução cobre todas as armadilhas que eu tinha listado:\newline%
(1) Diz que o termo é um neologismo formado por justaposição de termos opostos, próximo do oxímoro.\newline%
(2) Apoia a resposta nas paráfrases do próprio texto: o aposto "um rebaixamento naquele universo...", a "etiqueta de gênero" e o fechamento "agredir com um elogio, preconceito com açúcar, ataque com punho de renda".\newline%
(3) Liga cada parte da alternativa A a trechos do texto: "discriminatória" corresponde a rebaixamento, preconceito e etiqueta de gênero, e "enaltecimento" corresponde a louvor, açúcar e punho de renda.\newline%
(4) Na alternativa B, separa a ofensa explícita ("veado", o ataque direto e virulento) do louvor{-}prisão, que é sutil.\newline%
(5) Na alternativa D, a pegadinha principal, mostra que a alternativa troca quem age por quem reage. Também lembra que a conduta defensiva é a resposta ao primeiro insulto e que, diante do elogio, a reação é o sorriso amarelo.\newline%
(6) Menciona o preconceito social, cultural e de gênero.\newline%
\newline%
Só há duas falhas pequenas. A refutação da alternativa C não aponta que o louvor{-}prisão é a própria agressão, e não uma resposta a ela. Na alternativa D, "não cabe defesa" exagera um pouco o "talvez" do texto. Nenhuma das duas compromete o raciocínio. A resolução é clara, bem fundamentada e serve como gabarito comentado.%
\begin{itemize}%
\item%
{[}minor{]} Na alternativa c), o candidato só refuta a ideia de 'empatia' citando que o elogio era 'mais agressivo'. Faltou dizer que a alternativa também erra ao colocar o louvor{-}prisão como resposta 'frente a uma agressão verbal', quando ele é a própria agressão. A refutação está certa, mas fica incompleta. Bastaria uma frase a mais.%
\item%
{[}minor{]} Na alternativa d), a frase 'o texto diz que diante do segundo tipo de ataque não cabe defesa' exagera um pouco: o texto usa 'talvez, só possamos sorrir de forma amarela'. O correto seria escrever 'o texto sugere que talvez só reste o sorriso amarelo'. O argumento principal continua válido: a alternativa descreve a reação de quem recebe o elogio, e não o elogio em si.%
\end{itemize}%
\subsection*{Questão 31: nota 9,5 (aprovada, rodada 1 de 1)}%
\textbf{Resposta da resolução às cegas do avaliador:} Alternativa C. Os dois textos usam dados estatísticos para defender seu ponto de vista. O Texto 1 compara as médias de horas semanais (42,7 h no mundo e 40,1 h no Brasil) e usa a base de Gethin e Saez para concluir que o brasileiro trabalha menos que o esperado, por preferir lazer. O Texto 2 usa a mesma base, com especificação econométrica e controles institucionais, e microdados da World Values Survey com teste de significância estatística, para concluir o oposto.\par %
Conferi a transcrição com o original (págs. 9--10 do caderno, páginas 11--12 do PDF) e ela está fiel. A resposta final do candidato é (C), a mesma da minha referência, e está certa. Os dois textos usam dados estatísticos para defender pontos de vista opostos. O Texto 1 usa 160 países, as médias de 42,7 h e 40,1 h e a base de Gethin e Saez. O Texto 2 usa especificação econométrica, resíduo, controles, a World Values Survey com 58 países e teste de significância.\newline%
\newline%
A resolução trata as armadilhas principais:\newline%
{-} Começa explicando que \textbackslash{}"em comum\textbackslash{}" exige que a característica valha para os dois textos.\newline%
{-} Reconhece que o Texto 1 também defende uma tese e não é um relato neutro, citando os trechos.\newline%
{-} Descarta (b) porque só o Texto 2 contesta o argumento da mídia, e o Texto 1 é a própria matéria que o divulga.\newline%
{-} Descarta (d) separando a divergência entre os textos da divergência dentro de cada texto.\newline%
{-} Descarta (a) mostrando que atribuir uma causa cultural não é criticar a cultura e que o Texto 2 critica o método.\newline%
\newline%
As falhas são pequenas. Em (a), falta tratar a frase avaliativa \textbackslash{}"nem pode ser considerado particularmente esforçado\textbackslash{}" e a expressão \textbackslash{}"ao descreverem\textbackslash{}". Em (b), a ideia de que o Texto 2 \textbackslash{}"vai além de relativizar\textbackslash{}" é discutível. Há também uma grafia errada no nome do autor. É uma resolução clara, organizada e publicável como gabarito comentado.%
\begin{itemize}%
\item%
{[}minor{]} Na alternativa (a), o candidato não trata a frase do Texto 1 "Nem pode ser considerado particularmente esforçado". Ela é o que mais faz a alternativa (a) parecer certa, porque soa como crítica cultural. A alternativa continua eliminada, já que o Texto 2 não critica aspectos culturais nem econômicos, mas a análise ficaria completa com uma linha mostrando que a frase é uma constatação tirada dos dados e não uma crítica. O trecho "ao descreverem uma situação" também não é comentado: os dois textos argumentam em vez de só descrever.%
\item%
{[}minor{]} Na alternativa (b), dizer que o Texto 2 "vai além de relativizar" é discutível. O texto afirma que o sinal do resíduo muda conforme a especificação, e isso também é uma forma de relativizar. O que realmente elimina a alternativa é o Texto 1, que é o próprio argumento divulgado pela mídia e não o relativiza. O candidato cita esse ponto, então a eliminação está correta. Só seria bom deixar claro qual é o critério decisivo.%
\item%
{[}minor{]} Grafia do autor: a fonte do Texto 1 no original é "CARRIELO", e o candidato escreve "Carriello". A expressão "resíduos de regressão" é uma inferência razoável a partir de "especificação econométrica" e "resíduo do Brasil", mas o texto não usa a palavra "regressão".%
\end{itemize}%
\subsection*{Questão 32: nota 9,5 (aprovada, rodada 1 de 1)}%
\textbf{Resposta da resolução às cegas do avaliador:} Alternativa B: a metodologia do texto 2 permite reverter a conclusão à qual chega o texto 1. Com os controles institucionais (tributação, informalidade e regulação de jornada), o resíduo do Brasil muda de sinal e o país passa a trabalhar mais do que o esperado. Além disso, a hipótese da preferência cultural por lazer é rejeitada.\par %
Conferi o enunciado nas páginas 11 e 12 do PDF, e a transcrição está fiel. O candidato marcou B, a mesma alternativa da minha referência, e a resposta está correta. A resolução separa bem as duas afirmações do texto 1: o brasileiro trabalha menos do que o esperado, e isso se deveria a uma preferência cultural por lazer. Também descreve com fidelidade a metodologia do texto 2. Ela usa o mesmo banco de Gethin e Saez, testa outras formas funcionais e ponderações e acrescenta controles institucionais. Com isso, o resíduo do Brasil muda de sinal e o país passa a trabalhar 'mais do que o esperado'. Por fim, a hipótese cultural é rejeitada com microdados da World Values Survey. A resolução cita literalmente as passagens que justificam 'reverter'.\newline%
\newline%
Cobertura das armadilhas:\newline%
{-} C, o distrator principal: tratado corretamente. O texto 2 usa os dados e critica a especificação, não os dados.\newline%
{-} A: refutado, porque o texto 2 contesta o texto 1 em vez de complementá{-}lo.\newline%
{-} D: refutado com o teste e a rejeição da hipótese.\newline%
{-} Comando da questão: tratado como pergunta sobre a metodologia.\newline%
{-} Inversão literal do sinal: explicitada.\newline%
{-} Dado bruto x conclusão: a resolução separa os dois só de forma implícita.\newline%
\newline%
Os problemas que encontrei são de redação e de pouco peso. A resolução não diz explicitamente que as médias brutas não são contestadas. A expressão 'a explicação cultural que a sustentava' está um pouco imprecisa. A ressalva em (a) depende da qualificação 'que confirmem ou ampliem'. No conjunto, é uma resolução clara, completa e bem fundamentada no texto, em nível de gabarito comentado.%
\begin{itemize}%
\item%
{[}minor{]} A resolução não diz explicitamente que o texto 2 aceita os dados brutos (40,1 h contra 42,7 h) e que a reversão vale para a conclusão condicional ('menos do que seria esperado'). Na prática ela já separa as duas coisas, porque em (b) define a conclusão do texto 1 como 'trabalha menos do que o esperado'. Ainda assim, uma frase dizendo que as médias descritivas continuam valendo e que muda apenas a comparação com o valor previsto pelo modelo deixaria esse ponto completo.%
\item%
{[}minor{]} Em (b), a frase 'a explicação cultural que a sustentava' inverte um pouco a relação lógica. No texto 1, a preferência por lazer não sustenta a conclusão. Ela é a explicação proposta para o desvio. O certo seria algo como 'e a explicação cultural proposta para o desvio é rejeitada'. Minha referência tem a mesma imprecisão.%
\item%
{[}minor{]} Em (a), a frase 'não acrescenta dados que confirmem ou ampliem' está correta só por causa da ressalva. O texto 2 traz dados novos, os microdados da World Values Survey, mas usa esses dados para refutar o texto 1, não para complementá{-}lo. Valeria deixar isso explícito para fechar a leitura 'traz dados novos, logo complementa'.%
\end{itemize}%
\subsection*{Questão 33: nota 9,0 (aprovada, rodada 1 de 1)}%
\textbf{Resposta da resolução às cegas do avaliador:} Alternativa C: o trecho caracteriza as relações hierárquicas da sociedade escravista. O poder senhorial, já presente no menino Brás, é exercido sobre o escravizado Prudêncio, que é animalizado ("cavalo", "besta"), castigado e silenciado.\par %
Conferi o enunciado na página 12 do PDF, e a transcrição está fiel. O candidato marca a alternativa C, a mesma da minha referência, e a resposta está correta. A justificativa é sólida e cobre todas as armadilhas. Primeiro, o candidato separa o ponto de vista de Brás, que narra com cinismo e naturalidade, da ironia machadiana que expõe a brutalidade da cena. Também mostra três marcas no texto: animalização ('cavalo', 'freio', 'besta'), violência ('fustigava{-}o', 'gemendo') e silenciamento ('sem dizer palavra', 'Cala a boca, besta!'). Interpreta bem o título do capítulo como a formação do homem da elite senhorial. Refuta cada alternativa pela justificativa que vem depois do 'pois': o lúdico existe só para Brás; o poder nas mãos de uma criança mostra que ele está naturalizado e não o enfraquece; o lamento é submissão, e não resistência. Usa corretamente 'O vergalho' como reprodução da hierarquia, e não como inversão. Encontrei só problemas pequenos: um erro de LaTeX (`\textbackslash{}emph\{reproduz*`), a falta dos números dos capítulos e da origem do título, e uma refutação da alternativa (d) que poderia ser um pouco mais explícita. Serve como gabarito comentado.%
\begin{itemize}%
\item%
{[}minor{]} Erro de LaTeX: em `\textbackslash{}emph\{reproduz*;` a chave aberta é fechada com asterisco. Isso impede a compilação ou deixa em itálico todo o texto seguinte. Deveria ser `\textbackslash{}emph\{reproduz\};`.%
\item%
{[}minor{]} Faltam os números dos capítulos citados: XI para 'O menino é o pai do homem' e LXVIII para 'O vergalho'. Também não há referência à origem do título (Wordsworth). Não é erro, mas um gabarito comentado ficaria mais preciso com esses dados.%
\item%
{[}minor{]} Na alternativa (d), a análise poderia usar de novo a frase 'obedecia sem dizer palavra' e o fato de o 'ai, nhonhô!' manter o tratamento de senhor. Os dois pontos aparecem só no item 2, não na refutação da alternativa. Na alternativa (b), poderia também separar a ironia do autor, que de certo modo faz o leitor estranhar a cena, da justificativa errada da alternativa ('o poder exercido por uma criança enfraquece o quadro social').%
\end{itemize}%
\subsection*{Questão 34: nota 9,5 (aprovada, rodada 1 de 1)}%
\textbf{Resposta da resolução às cegas do avaliador:} Alternativa D. A "Advertência" faz parte da ficção: é o recurso do editor fictício. Augusto Machado é o personagem{-}narrador e o autor fictício da narrativa, e Lima Barreto, fazendo{-}se de editor, atribui a ele a autoria.\par %
Conferi o enunciado no original (página 12 do PDF), e ele bate com a transcrição. O candidato marcou D, a mesma resposta da minha referência, e a resposta está correta. A resolução reconhece o recurso do editor fictício (manuscrito alheio) e o compara com as advertências de Machado de Assis em Esaú e Jacó e no Memorial de Aires. Mostra que Augusto Machado é ao mesmo tempo narrador{-}personagem e autor fictício, usando as próprias pistas do trecho ('trata mais dele do que do seu herói', 'sobressair demasiado, a pessoa do autor'). Explica também que a discussão 'biografia ou romance' faz parte do jogo ficcional, porque cria efeito de verossimilhança e marca o caráter de romance. Todas as alternativas são refutadas corretamente. Na a, o texto é ficção e não paratexto informativo, e o autor real é Lima Barreto. Na b, a resolução aponta que o texto foi lido ao contrário: os defeitos são desprezados, não criticados, e o narrador aparece demais em vez de ser discreto. Na c, o Modernismo não separou os gêneros, e Lima Barreto é pré{-}modernista. Quase todas as armadilhas foram tratadas. Faltaram dois detalhes pequenos: dizer explicitamente, na b, que a voz é do editor fictício e não de Lima Barreto; e usar, na c, o argumento da data (1919, antes de 1922). O texto é claro, organizado e serviria como gabarito comentado.%
\begin{itemize}%
\item%
{[}minor{]} Na alternativa b, a resolução refuta o conteúdo (os defeitos são relativizados, não criticados; e o narrador aparece demais, o oposto de discrição e objetividade), mas não diz que quem fala é o editor fictício. Por isso a advertência não pode ser lida como expressão direta das concepções literárias de Lima Barreto. Bastaria uma frase para fechar esse ponto, que a própria resolução já preparou no item 2.%
\item%
{[}minor{]} Na alternativa c, falta o argumento cronológico explícito: a obra é de 1919, antes da Semana de 1922, e por isso falar em 'Modernismo que se iniciava' como referência das concepções do autor é anacrônico. A resolução só diz que Lima Barreto é 'pré{-}modernista'. A data de 1919 aparece no item 1, mas não é usada aqui. A afirmação de que o Modernismo 'misturou' os gêneros também é uma generalização, embora aceitável para o nível do vestibular.%
\end{itemize}%
\subsection*{Questão 35: nota 9,5 (aprovada, rodada 1 de 1)}%
\textbf{Resposta da resolução às cegas do avaliador:} Alternativa C: "expressam o protagonismo daquele que conduz, condição recuperada pela aproximação semântica entre as palavras bengala e cajado".\par %
Conferi a transcrição com o original (p. 12 do PDF, que corresponde à p. 10 da prova): a questão 35 confere integralmente. O candidato marca a alternativa C, a mesma da minha referência, e essa é a resposta correta. A bengala, lida como cajado pela metáfora pastoril (passos = rebanho, eu{-}lírico = pastor), converte a dependência em condução ativa ('me faço', verbo reflexivo e ativo). Por isso a alternativa fala em protagonismo de quem conduz, recuperado pela aproximação semântica entre bengala e cajado.\newline%
\newline%
A resolução trata quase todas as armadilhas:\newline%
{-} **Cajado implícito:** mostra que 'cajado' não aparece no texto e que é inferido pela semelhança de forma e de função com a bengala.\newline%
{-} **Rebanho:** lê corretamente que o rebanho é formado pelos próprios passos, e usa isso para derrubar (B).\newline%
{-} **Tom:** identifica o tom de exaltação da ode contra o 'lamento' de (D).\newline%
{-} **Sintaxe:** nega a 'dissolução sintática', já que o poema é um único período íntegro e a ligação entre bengala e cajado é semântica.\newline%
{-} **Alternativa (A):** refuta a 'dificuldade de comparação'.\newline%
{-} **Senso comum:** aborda a leitura da bengala como sinal de velhice ou fragilidade.\newline%
{-} **Diálogo:** explica a apóstrofe ('Contigo').\newline%
\newline%
As falhas são pequenas: (a) não comenta 'apoio material'; (b) não comenta 'apoio psicológico' nem desmonta explicitamente a armadilha biográfica; (c) não nomeia a personificação. O texto é claro, rigoroso e conciso, e descarta cada distrator por um defeito concreto. É publicável como gabarito comentado.%
\begin{itemize}%
\item%
{[}minor{]} Na alternativa (a), o candidato não comenta 'apoio material'. Essa parte até é aceitável, porque a bengala é de fato um apoio físico, e o defeito está em 'extensão da imaginação' e em 'dificuldade de comparação'. Uma frase dizendo que o apoio material existe, mas não ganha sentido de fantasia, deixaria claro que a alternativa erra pela segunda parte, e não pela primeira.%
\item%
{[}minor{]} Na alternativa (b), o candidato refuta 'trauma' e a leitura errada dos 'passos do rebanho', mas não comenta 'apoio psicológico'. Também não diz que a ideia de trauma só poderia vir de fora do texto, por exemplo da biografia do autor (o livro também traz 'À minha perna esquerda'). Faltou desmontar explicitamente a armadilha biográfica. O candidato trata só a leitura de senso comum da bengala como velhice e dependência.%
\item%
{[}minor{]} O 'diálogo' citado no enunciado aparece como apóstrofe ('Contigo'), mas a personificação do objeto, que dá sentido à expressão 'com quem dialoga', não é nomeada. É uma lacuna pequena de vocabulário analítico.%
\item%
{[}minor{]} O candidato poderia dizer que (C) é a única alternativa que junta um efeito de sentido sustentado pelo texto ('me faço pastor', ou seja, conduzir) a um mecanismo correto (a aproximação semântica, não sintática). Isso deixaria a escolha mais claramente motivada por eliminação. Mesmo assim, cada distrator já foi descartado por um defeito concreto.%
\end{itemize}%
\subsection*{Questão 36: nota 9,0 (aprovada, rodada 1 de 1)}%
\textbf{Resposta da resolução às cegas do avaliador:} Alternativa D. A dor que a autora expressa no Texto 1, diante de uma infância exposta à violência e de uma maternidade sob pressão, é transformada em matéria literária nos contos, como mostram os Textos 2 e 3. A alternativa A é falsa porque as mães estão presentes e escrever não "diminui" o que a ficção retrata. A B é falsa porque o silêncio continua presente nos excertos. A C é falsa porque confunde a dor da autora ao escrever, que está fora dos contos, com o tema deles.\par %
Conferi o enunciado no PDF (p. 13, impressa 11), e a transcrição está correta. O candidato marca D, a mesma resposta da minha referência, e D está correta. Ele analisa as quatro alternativas e mostra o erro de cada distrator. Na A, não há abandono, e escrever não 'diminui' nada. Na B, o silêncio continua presente nos dois excertos. Na C, que é o principal distrator, ele separa bem a dor do ato de escrever, que aparece só no Texto 1, dos temas dos contos, e ainda observa que 'Zaíta' não trata só de mães e filhas. A leitura integral dos contos está certa: a fome disfarçada com brincadeiras e os 'olhos d'água' da mãe no primeiro conto, o filho envolvido com o tráfico e a morte de Zaíta num tiroteio no segundo. A distinção entre o processo de criação e o tema dos contos, ligada à escrevivência, é o ponto central da questão e foi bem resolvida. Os problemas são pequenos: citações levemente alteradas entre aspas; a dor ligada a 'Zaíta' estendida a 'Olhos d'água' sem justificativa; nenhuma ligação entre a indignação da leitora e a morte da menina; a isca 'cumplicidade' na B não foi apontada; e uma pequena falta de cuidado entre autora e narradora na conclusão. A resolução é clara e correta e serve de gabarito comentado, com pequenos ajustes de precisão.%
\begin{itemize}%
\item%
{[}minor{]} Há citações levemente alteradas entre aspas. O candidato escreve 'recomenda o silêncio', mas o original diz 'recomendou então o silêncio'. Escreve também silêncios 'decifrados', enquanto o original traz 'Decifrava os seus silêncios'. Em prova de Literatura, o que vai entre aspas deve reproduzir o texto exatamente, ou então a passagem deve ser parafraseada sem aspas.%
\item%
{[}minor{]} O Texto 1 fala da dor da autora só em relação a 'Zaíta'. Para chegar à alternativa D, essa dor precisa ser estendida também a 'Olhos d'água', e o candidato faz essa extensão sem justificá{-}la. Bastava uma frase dizendo que o 'comprometimento' e a 'cumplicidade' declarados valem para as personagens em geral, o que caracteriza a escrevivência, e que isso autoriza estender a afirmação aos dois contos.%
\item%
{[}minor{]} O candidato não relaciona a indignação da leitora no Texto 1 com o desfecho de 'Zaíta', a morte da menina. Essa relação mostra por que a escrita é 'dolorosa' e reforça a D com a leitura integral exigida pelo enunciado. O desfecho é mencionado no Passo 2, mas não é ligado ao Texto 1.%
\item%
{[}minor{]} Na alternativa B, o candidato não aponta de forma explícita a isca lexical 'cumplicidade'/'cumpliciada' nem diz que a primeira parte da alternativa é plausível e que o erro está no 'diminui o silêncio'. A refutação está correta, mas poderia ser mais didática.%
\item%
{[}minor{]} Na conclusão ('a experiência dolorosa da escritora {[}...{]} vira o assunto dos contos'), a autora e a narradora quase se confundem. Seria bom deixar claro que a narradora de 'Olhos d'água' é uma personagem, e não a própria autora, mesmo que a escrevivência aproxime vida e ficção.%
\end{itemize}%
\subsection*{Questão 37: nota 9,0 (aprovada, rodada 1 de 1)}%
\textbf{Resposta da resolução às cegas do avaliador:} Alternativa B. Ao trocar o nome cristão Grace por Afamefuna, nome igbo que lhe foi dado pela avó Nwamgba e que significa "meu nome não se perderá", a personagem se reconecta com as expectativas e os valores da avó, ligados à tradição anterior à colonização.\par %
O enunciado foi conferido na p. 13 do PDF (p. 11 do caderno) e a transcrição está fiel. O candidato marcou (B), a mesma resposta da referência, e ela está correta. Ao trocar o nome cristão Grace por Afamefuna, nome igbo dado pela avó Nwamgba e que significa 'meu nome não se perderá', a personagem retoma os valores da avó, ligados à tradição anterior à colonização. A resolução apresenta bem o contexto do conto inteiro: as três gerações, a conversão de Michael, a origem e o significado do nome. Ela identifica a prolepse marcada pelo futuro do pretérito e pela anáfora 'Foi Grace quem', interpreta o episódio do sr. Gboyega como crítica à visão colonial e analisa cada alternativa com argumento próprio. Na (c), mostra que a alternativa inverte o sentido do gesto e que o sucesso não preencheu o vazio de estar 'distante de suas raízes'. Na (d), mostra que a anáfora não antecipa tristeza, que a tristeza aparece só no episódio Gboyega e que o problema é a própria educação colonial, não a ignorância. Todas as armadilhas da referência foram tratadas. A única lacuna relevante está na (a): o candidato aceita o 'protagonismo' sem notar que ele não vale 'ao longo da narrativa', já que a protagonista é Nwamgba. Há também pequenas imprecisões de formulação. A resolução está clara, correta e completa, no nível de um gabarito comentado, com pequenos ajustes possíveis.%
\begin{itemize}%
\item%
{[}minor{]} Na alternativa (a), o candidato concede que 'a anáfora de fato realça o protagonismo de Grace' e refuta só a parte da ocidentalização. Faltou rebater o trecho 'ao longo da narrativa'. Grace só ganha destaque no desfecho, num resumo do seu futuro (prolepse), e a protagonista da maior parte do conto é Nwamgba. A refutação seria mais completa se apontasse que as duas partes da alternativa estão erradas.%
\item%
{[}minor{]} No item 2, o 'portanto' em 'A educação ocidental aparece, portanto, como algo que pode apagar a identidade africana' é um pequeno salto. O excerto fala da 'ligação entre educação e dignidade' e não diz isso de forma direta. A conclusão se apoia no exemplo do sr. Gboyega e no conto inteiro, e o texto deveria deixar claro que é uma inferência.%
\item%
{[}minor{]} Algumas expressões do contexto são um pouco fortes ou supõem coisas fora do excerto: Michael 'rompe com os costumes', e a menina é 'registrada' como Grace (no conto, ela é batizada Grace). Isso não afeta a análise, mas vale ser mais preciso.%
\end{itemize}%
\subsection*{Questão 38: nota 9,1 (aprovada, rodada 1 de 1)}%
\textbf{Resposta da resolução às cegas do avaliador:} Alternativa C. O narrador conta ao leitor o verdadeiro negócio de Montiel: ele lucra com a expulsão dos ricos comprando os bens deles pelo preço que ele mesmo fixa. Logo depois mostra a mulher, que acredita sinceramente que o marido "ajuda" os expulsos e corre o risco de "arruinar{-}se". É ironia dramática: a personagem ignora o que o leitor sabe. Isso é coerente com o conto inteiro, em que só a viúva não se sente "vingada" com a morte dele.\par %
Conferi o enunciado na página 13 do PDF (página 11 da prova): a transcrição está fiel. O candidato marcou C, a mesma alternativa da minha referência, e a resposta está correta. Na passagem, o narrador primeiro revela ao leitor o negócio real de Montiel: ele lucra com a expulsão dos ricos, comprando os bens deles pelo preço que ele mesmo fixa. Logo depois vem a fala sincera da mulher, que vê nisso caridade e risco de ruína. É ironia dramática construída pelo narrador: a personagem ignora o que o leitor sabe.\newline%
\newline%
A resolução trata bem as armadilhas principais:\newline%
{-} **(a):** distingue antítese ou contraste de paradoxo.\newline%
{-} **(b):** mostra que o raciocínio da esposa é coerente e que o erro dela está na premissa, não na forma do argumento.\newline%
{-} **(c):** atribui a ironia ao narrador, e não à personagem.\newline%
{-} **(d):** diferencia ironia verbal intencional de ironia dramática, notando que não há consciência nem efeito cômico.\newline%
{-} **Contexto:** usa o conto inteiro --- o conluio com o alcaide, a viúva alheia e a ruína depois da morte de Montiel.\newline%
\newline%
A frase \textbackslash{}"diz que o marido vai se arruinar justamente com o que o enriquece\textbackslash{}" resume bem o mecanismo. Há só lacunas pequenas, e nenhuma compromete a resposta: falta o duplo sentido de \textbackslash{}"nunca irão agradecer{-}lhe\textbackslash{}", a função de \textbackslash{}"Na verdade\textbackslash{}" está descrita de forma imprecisa e falta a citação \textbackslash{}"menos a viúva\textbackslash{}". A resolução é clara, completa e publicável como gabarito comentado, com pequenos acréscimos.%
\begin{itemize}%
\item%
{[}minor{]} A resolução não explora o duplo sentido de "eles nunca irão agradecer{-}lhe". A frase é verdadeira, mas pelo motivo oposto ao que a mulher imagina: as vítimas não agradecem a quem as espoliou. Esse é o ponto em que a ironia aparece na própria fala da personagem, e mencioná{-}lo daria mais força à alternativa (c). Bastaria uma frase sobre isso na análise do trecho.%
\item%
{[}minor{]} A resolução diz que "'Na verdade' marca a voz do narrador", mas a expressão não serve para isso: o primeiro parágrafo inteiro já é narração. O papel de "Na verdade\ldots{} não era X, mas sim Y" é corrigir e revelar, ou seja, o narrador conta ao leitor o mecanismo real do negócio. Dizer isso também afastaria de vez a leitura paradoxal da alternativa (a). Convém reescrever: "o narrador, com 'Na verdade\ldots{} não era\ldots{} mas sim\ldots{}', corrige e revela\ldots{}".%
\item%
{[}minor{]} O contexto do conto está correto, mas falta citar o dado textual mais forte, a abertura "todo mundo se sentiu vingado, menos a viúva". Ela comprova que só a viúva desconhece quem o marido era, e o enunciado pede leitura "no contexto geral do conto". A expressão "o homem mais rico e mais odiado" é uma paráfrase aceitável, mas não está ancorada em citação.%
\end{itemize}%
\subsection*{Questão 39: nota 9,0 (aprovada, rodada 1 de 1)}%
\textbf{Resposta da resolução às cegas do avaliador:} Alternativa B: segundo o texto, a explicação principal está nas disputas de classificação dentro das ideologias raciais (embranquecimento, mestiçagem, negritude), que convivem sem que uma suplante as outras. Essas disputas afetam a forma como os indivíduos se percebem e se declaram, e também as políticas públicas.\par %
Conferi a transcrição com a página 14 do PDF e ela está correta. A resposta do candidato, alternativa (B), coincide com a da referência e está certa.\newline%
\newline%
A resolução está bem organizada. Ela descreve as curvas do gráfico com valores coerentes com o original (brancos com pico de cerca de 63\% em 1940, pardos passando os brancos em 2022, pretos caindo a 5--6\% e voltando a cerca de 10\%). Liga cada fase a uma das ideologias: embranquecimento, mestiçagem e negritude. Também lê com cuidado a expressão 'não pode ser explicado somente' e descreve a metodologia como 'descartada como suficiente', sem dizer que ela é irrelevante.\newline%
\newline%
Cada alternativa errada é eliminada por contradição com o texto:\newline%
{-} (a): 'hegemônicas' contradiz o trecho 'incapazes de suplantar umas às outras', e o texto não diz que o Estado produz as ideologias.\newline%
{-} (c): os incentivos materiais são declarados menos importantes que as ideologias.\newline%
{-} (d): o texto rejeita expressamente a explicação pela pergunta do censo e pela forma de classificação.\newline%
\newline%
As armadilhas principais foram tratadas. Ficaram só lacunas pequenas:\newline%
{-} não reconhece que 'impacto na produção de políticas públicas' é uma inferência e não uma frase do texto;\newline%
{-} não aponta o problema de 'mais precisos' em (d);\newline%
{-} não comenta a pegadinha weberiana e o título do artigo em (c).\newline%
\newline%
É uma resolução de nível de gabarito comentado, com pequenos pontos a melhorar.%
\begin{itemize}%
\item%
{[}minor{]} Na alternativa (b), o trecho 'impacto na produção de políticas públicas' não está escrito no Texto 2. O candidato justifica esse trecho com uma leitura interpretativa ('a alternativa trata as políticas como consequência das disputas') sem dizer que é uma inferência. Deveria declarar que esse trecho não vem literalmente do texto e que (b) é escolhida porque é a única que põe as ideologias raciais e a autopercepção no centro da explicação, enquanto as outras contradizem o texto.%
\item%
{[}minor{]} Na alternativa (d), faltou apontar um segundo erro: o texto nunca diz que o instrumento censitário 'foi aprimorado e se tornou mais preciso'. Essa ideia também supõe uma raça 'objetiva' a ser medida, o que vai contra a visão do texto de que a classificação é uma construção social.%
\item%
{[}minor{]} A alternativa (c) foi descartada corretamente, mas sem identificar a origem da pegadinha. 'Ação racional com relação a fins' é um conceito weberiano, e o título do artigo ('oportunismo político') pode levar o leitor a essa opção. Citar isso deixaria a análise mais completa.%
\item%
{[}minor{]} Na parte 1, a conclusão 'Por isso, não se explicam apenas por diferenças de natalidade e mortalidade' é tirada só do gráfico, o que é um pequeno salto. Seria mais rigoroso apoiá{-}la diretamente na afirmação do Texto 2, ou lembrar que outros fatores demográficos, como a imigração, também existiram, e que mesmo assim o texto põe as ideologias como fator principal.%
\end{itemize}%
\subsection*{Questão 40: nota 9,0 (aprovada, rodada 1 de 1)}%
\textbf{Resposta da resolução às cegas do avaliador:} Alternativa D. Em contextos de polarização, a democracia se apresenta como método eficaz, "o melhor método" segundo o texto, para administrar de forma pacífica, sem repressão, os conflitos entre ideais, valores e interesses diferentes.\par %
Conferi o enunciado na página 14 do PDF (página impressa 12): a transcrição está correta. As respostas coincidem: tanto o candidato quanto a minha referência marcam (D), e a resposta está correta. A resolução identifica bem as três ideias do texto de Przeworski: a democracia como terreno de conflito, seu 'milagre' de administrar os conflitos sem repressão e, sob polarização, a defesa da democracia como método. O candidato também mostra que o conflito faz parte da democracia e não é sinal de falha. Quase todas as armadilhas foram tratadas. A leitura de 'crise' como senso comum e a ausência de 'maioria' no texto aparecem em (a). A inversão 'defendida'/'questionada' aparece em (b), e a inversão de 'sem repressão' em (c). A tese procedimental de democracia como 'método' é explicitada. Ficaram pequenas lacunas: 'maioria' não é apontada em (b); em (c), 'ainda mais necessário' vai um pouco além do texto; e em (d) as equivalências não estão completas ('melhor' corresponde a 'eficaz'; a troca 'ideais'/'ideias' não muda o sentido). A resolução é clara, bem organizada e serviria como gabarito comentado, com pequenos ajustes de precisão textual.%
\begin{itemize}%
\item%
{[}minor{]} Na alternativa (b), o candidato só aponta a troca de 'defendida' por 'questionada'. Não comenta a segunda falha, a menção à 'maioria', que não está no texto. Ele trata disso em (a), mas uma análise completa deveria apontar o mesmo erro também em (b).%
\item%
{[}minor{]} Na alternativa (c), a frase 'é ela que torna esse método ainda mais necessário' vai um pouco além do texto. O autor diz que, sob polarização, a democracia 'só pode ser defendida como o melhor método', e não que ela fica 'mais necessária'. O candidato também não contrapõe explicitamente 'perde importância' a 'só pode ser defendida'. Seria melhor citar o trecho literal.%
\item%
{[}minor{]} Na alternativa (d), faltam as equivalências lexicais completas. 'Melhor método para administrar essas divergências' corresponde a 'método eficaz para a gestão de conflitos', mas o candidato só faz a ligação entre 'sem repressão e em paz' e 'pacífica'. Também não observa que a troca 'ideais' por 'ideias' é apenas de redação. A pequena diferença entre 'apresenta{-}se' (tom descritivo) e 'pode ser defendida' (tom normativo) também não é comentada. Isso não muda a resposta, mas deixaria a justificativa mais firme.%
\end{itemize}%
\subsection*{Questão 41: nota 9,0 (aprovada, rodada 1 de 1)}%
\textbf{Resposta da resolução às cegas do avaliador:} Alternativa C: grupos e discursos negacionistas se apropriam de retóricas e credenciais científicas para produzir falsas controvérsias que visam gerar dissensos e incertezas na sociedade.\par %
Conferi o enunciado no original (p. 14 do PDF, p. 12 do caderno): a transcrição, a figura e as alternativas estão corretas. O candidato responde (C), igual à minha referência. A resposta está correta, porque (c) é uma paráfrase quase literal do Texto 1: apropriação de retóricas e credenciais científicas, fabricação de falsas controvérsias e o objetivo de criar dissensos e incertezas.\newline%
\newline%
A resolução é bem organizada. Primeiro resume a tese do Texto 1. Depois interpreta corretamente o diagrama: a fronteira tracejada e inclinada indica que não há separação nítida, e a interseção sombreada é a 'região cinzenta'. Por fim, refuta cada distrator com base no texto:\newline%
{-} (a): contradiz o trecho 'não se trata de um mero repúdio... anticientífico'.\newline%
{-} (b): reduz o fenômeno a uma oposição frontal, e a finalidade apontada ('buscar legitimidade') não é a do texto.\newline%
{-} (d): inverte a causalidade, pois os negacionistas criam os dissensos em vez de se sobreporem a eles; além disso, não negam a existência das ciências.\newline%
\newline%
A conclusão sintetiza bem os três elementos: meio, produto e finalidade. Todas as armadilhas principais foram tratadas. As falhas são pequenas: não aponta que 'desinformação' é um termo externo ao texto em (a) e (d), comprime uma citação dentro de aspas e não discute o 'se digladiam' que dá força ao distrator (b). É uma resolução de nível de gabarito comentado, com pequenas lacunas de completude.%
\begin{itemize}%
\item%
{[}minor{]} Na análise de (a), o candidato não aponta que 'produzir desinformação' é um termo que não está no texto e vai além dele. Só refuta 'se afirmarem como forças anticientíficas'. A refutação feita basta para eliminar a alternativa, mas a análise ficaria completa se registrasse também o termo externo. O mesmo vale para (d), onde 'produzir desinformação' também aparece e não é comentado.%
\item%
{[}minor{]} A citação 'mero repúdio anticientífico' aparece entre aspas como se fosse literal, mas o original diz "mero repúdio 'pseudo' ou 'anticientífico'". Num gabarito comentado, a citação deveria ser fiel ou vir marcada como paráfrase.%
\item%
{[}minor{]} Em (b), o candidato não trata o trecho que torna esse distrator atraente: 'se digladiam para estender suas fronteiras em face das ciências'. O texto sugere algum embate, mas ele acontece na zona de sobreposição e por meio de apropriação, não como oposição frontal. Esse ponto deixaria a eliminação de (b), o distrator mais forte, mais robusta.%
\end{itemize}%
\subsection*{Questão 42: nota 9,5 (aprovada, rodada 1 de 1)}%
\textbf{Resposta da resolução às cegas do avaliador:} Alternativa A: a tirinha mostra a transcrição. Ela ocorre no núcleo; o DNA serve de molde para um RNAm com bases complementares (U no lugar de T), e esse RNAm leva a informação ao citoplasma, onde ficam os ribossomos.\par %
Conferi o enunciado no original (p. 15 do PDF): a transcrição está fiel. Refiz o pareamento com python3: CACA$\to$GUGU, CTCT$\to$GAGA, ACT$\to$UGA. O resultado bate com o que a secretária escreveu.\newline%
\newline%
O candidato marca a alternativa A, a mesma da minha referência, e a resposta está correta. A tirinha mostra a transcrição: o DNA serve de molde para um RNAm complementar, com U no lugar de T. Isso ocorre no núcleo, e o RNAm vai depois para o citoplasma, onde ficam os ribossomos.\newline%
\newline%
A resolução trata quase todas as armadilhas:\newline%
{-} distingue transcrição de tradução e mostra o fluxo DNA$\to$RNAm$\to$proteína;\newline%
{-} interpreta 'o contrário' como complementaridade, e não como sequência invertida;\newline%
{-} explica a falta do 'T' pela uracila;\newline%
{-} identifica o produto como RNAm, e não como RNAt;\newline%
{-} localiza a transcrição no núcleo e a tradução no citoplasma;\newline%
{-} refuta cada alternativa errada: c põe a transcrição no ribossomo com o RNAt como molde; d põe a tradução no núcleo e diz que o RNAt leva a mensagem; b diz que a tradução produz RNAt.\newline%
\newline%
Ela também não se apoia no detalhe de UGA ser códon de parada.\newline%
\newline%
Restam só ressalvas menores:\newline%
{-} a ordem do argumento no item 1 é um pouco circular;\newline%
{-} falta comentar explicitamente procariotos e organelas;\newline%
{-} a refutação da b tem um argumento visual mais fraco;\newline%
{-} não cita a orientação antiparalela nem o processamento do pré{-}RNAm.\newline%
\newline%
É uma resolução clara e correta, com qualidade de gabarito comentado.%
\begin{itemize}%
\item%
{[}minor{]} O item 1 afirma que a fita simples é o RNA mensageiro antes de apresentar as evidências. A justificativa só aparece nos itens 2 e 3: as bases complementares, o U no lugar do T e o destino 'Dr. Ribossomo'. A conclusão está certa, mas a ordem do argumento fica um pouco circular. Deveria primeiro conferir as bases e o U e só depois concluir que é RNAm, e não RNAt ou RNAr, pelo fato de levar a mensagem ao ribossomo.%
\item%
{[}minor{]} A resolução diz 'Nos eucariotos, o DNA fica no núcleo', mas não comenta que em procariotos, e no DNA de mitocôndrias e cloroplastos, a transcrição ocorre fora do núcleo. Também não explica por que o contexto da questão permite supor uma célula eucariótica. Uma frase sobre isso daria mais segurança ao 'ocorre no núcleo' da alternativa A.%
\item%
{[}minor{]} Na alternativa (b), o primeiro argumento ('a tirinha não mostra a tradução, porque o DNA participa da cena') se apoia na interpretação da imagem, não em um erro biológico. Ele está certo, mas é fraco. O texto também não diz explicitamente que o RNAm não serve de molde para o RNAt, que é transcrito do DNA. A segunda frase compensa em parte, ao dizer que a tradução produz aminoácidos, e não RNAt.%
\item%
{[}minor{]} A resolução não comenta que a fita de RNA é antiparalela ao molde. Também não comenta que o RNAm de eucariotos é processado (capa 5', cauda poli{-}A, splicing) antes de sair do núcleo. Nenhum dos dois pontos é necessário para escolher a alternativa, mas uma nota de rodapé deixaria o gabarito comentado mais completo.%
\end{itemize}%
\subsection*{Questão 43: nota 9,0 (aprovada, rodada 1 de 1)}%
\textbf{Resposta da resolução às cegas do avaliador:} Alternativa C. Pelo texto, as toxinas sofreram "seleção positiva": a variante com o gene duplicado aumentou a sobrevivência e/ou a reprodução das serpentes que a carregavam e, por isso, ficou mais frequente na população com o tempo. As outras alternativas erram pelos seguintes motivos: (a) o isolamento corta o fluxo gênico, e a seleção natural não causa fluxo gênico; (b) a especiação não foi súbita e a mutação não surge com uma finalidade; (d) a deriva genética não cria duplicação gênica, quem cria é a mutação.\par %
A resposta do candidato é C, a mesma da minha referência, e está correta. Conferi o enunciado na página 15 do PDF e a transcrição está fiel. O candidato separa bem mutação, seleção natural, deriva e fluxo gênico, e todas as armadilhas foram tratadas:\newline%
{-} (a): a seleção natural não causa fluxo gênico, e o isolamento corta esse fluxo.\newline%
{-} (b): a especiação não foi súbita.\newline%
{-} (c): "seleção positiva" corresponde ao aumento de frequência de uma variante que dá mais aptidão; o "e/ou" foi bem aproveitado.\newline%
{-} (d): a deriva não cria duplicação gênica e o texto fala em seleção, não em deriva.\newline%
\newline%
Também evita o raciocínio finalista, ao dizer que a mutação ocorre ao acaso e que o ambiente apenas filtra as variantes. As falhas são pequenas: em (b), uma objeção responde a uma leitura que não está na frase ("motivada por necessidade" em vez de "motivou a especiação"); o "há cerca de 100 mil anos" do texto foi lido como duração; e há uma extrapolação sobre o efeito da duplicação que não vem marcada como inferência. A resolução é clara, completa e serve como gabarito comentado.%
\begin{itemize}%
\item%
{[}minor{]} Na alternativa (b), o candidato refuta a ideia de que 'mutações não são motivadas por uma necessidade'. Mas o enunciado diz que a duplicação 'motivou a especiação', ou seja, a mutação como causa da especiação, e não como algo provocado por necessidade. A refutação responde a uma leitura que não está na frase. O erro de fato é atribuir a especiação a uma única mutação e dizer que ela foi súbita, quando o motor foi o isolamento geográfico. O argumento sobre a especiação não ter sido 'súbita' basta para descartar a alternativa, mas a primeira objeção deveria ser reescrita.%
\item%
{[}minor{]} O candidato escreve que a diferenciação foi gradual 'ao longo de cerca de 100 mil anos'. O texto diz que a espécie 'se diferenciou ... há cerca de 100 mil anos', o que é uma data e não uma duração. A conclusão biológica (especiação gradual e alopátrica) continua correta, mas a citação do texto é imprecisa. Minha referência tem a mesma imprecisão.%
\item%
{[}minor{]} No 'Raciocínio', a afirmação de que a duplicação 'aumenta a quantidade ou a diversidade de toxinas' é uma extrapolação razoável (dosagem gênica/neofuncionalização), mas o texto não diz isso. Poderia vir marcada como inferência ('possivelmente').%
\end{itemize}%
\subsection*{Questão 44: nota 9,5 (aprovada, rodada 1 de 1)}%
\textbf{Resposta da resolução às cegas do avaliador:} Alternativa B --- (i) cnidários; (ii) platelmintos; (iii) da membrana plasmática; (iv) sódio; (v) potássio.\par %
Conferi o enunciado na página 15 do PDF (página 13 do caderno) e a transcrição está correta. A única diferença é tipográfica: '(III)' na alternativa d. A resposta final do candidato, alternativa B (cnidários; platelmintos; membrana plasmática; sódio; potássio), coincide com a minha referência em todas as cinco lacunas e está correta. Cada lacuna tem justificativa sólida:\newline%
{-} (i) Rede difusa nos cnidários. O candidato desarma a armadilha do texto de apoio ao explicar por que os equinodermos, apesar de não terem cérebro, não servem: são deuterostômios, têm anel nervoso com nervos radiais e simetria radial secundária.\newline%
{-} (ii) Platelmintos. O candidato aponta a cefalização, os gânglios cerebroides anteriores e os cordões longitudinais com comissuras (sistema 'em escada').\newline%
{-} (iii) Membrana plasmática, por causa dos canais dependentes de voltagem. O envoltório nuclear é descartado corretamente.\newline%
{-} (iv)/(v) Entrada de Na? na despolarização e saída de K? na repolarização, com valores típicos (-70 mV e +30 mV) e a estequiometria da bomba (3 Na? para 2 K?). O candidato ainda argumenta que a ordem invertida contrariaria os gradientes.\newline%
{-} A conferência das alternativas identifica corretamente todos os erros de a, c e d.\newline%
\newline%
Os problemas são só pontuais: a bomba descrita como etapa posterior ao potencial de ação, e as afirmações de que cnidários e platelmintos são os 'primeiros' grupos, que são simplificações de livro didático. Nenhum deles afeta a resposta. A resolução tem qualidade de gabarito comentado.%
\begin{itemize}%
\item%
{[}minor{]} A frase 'Em seguida, a bomba Na?/K?{-}ATPase restabelece as concentrações iniciais dos íons' sugere que a bomba entra em ação depois de cada potencial de ação. Na verdade, ela funciona o tempo todo, e um único potencial de ação muda as concentrações de forma desprezível. O certo seria dizer que a bomba mantém os gradientes a longo prazo. É uma simplificação comum nos livros didáticos e não afeta a resposta.%
\item%
{[}minor{]} Dizer que os cnidários são 'os animais mais simples com células nervosas' é uma simplificação: os ctenóforos também têm rede nervosa, e há debate sobre a posição deles na base da árvore dos animais. Do mesmo modo, chamar os platelmintos de 'os primeiros animais com simetria bilateral e cefalização' ignora os acelomorfos. Nos dois casos é a versão dos livros didáticos, e bastaria acrescentar 'segundo a abordagem clássica'. Além disso, a tênia é um exemplo fraco de cefalização, porque é parasita e tem o escólex reduzido. A planária basta como exemplo.%
\item%
{[}minor{]} Faltou destacar explicitamente que só a lacuna (iii) e a ordem dos íons já eliminam a, c e d, chegando à B sem depender da zoologia. A conferência das alternativas lista todos os erros de cada uma, então isso está implícito, mas a observação deixaria a resposta mais robusta.%
\end{itemize}%
\subsection*{Questão 45: nota 9,0 (aprovada, rodada 1 de 1)}%
\textbf{Resposta da resolução às cegas do avaliador:} A alternativa correta é a C. O fototropismo positivo ocorre na parte aérea, e há mais auxina na seta 2 (lado sombreado, perto do ápice) do que na seta 3 (lado iluminado). A auxina migra para o lado sem luz, que se alonga mais, e o caule se curva para o Sol.\par %
A figura original (p. 16 do PDF) confirma a leitura: o Sol está no canto superior esquerdo, as setas 1 e 3 ficam no lado iluminado do caule, a seta 2 no lado sombreado perto do ápice, as setas 4 e 5 nos lados da raiz e a seta 6 na ponta da raiz. A resposta final do candidato é (C), igual à minha referência, e está correta.\newline%
\newline%
A resolução é sólida. Expõe o modelo de Cholodny--Went de forma correta:\newline%
{-} produção da auxina no ápice e transporte polar;\newline%
{-} fototropinas e migração lateral da auxina para o lado sombreado;\newline%
{-} crescimento ácido e curvatura na zona de alongamento, não no meristema;\newline%
{-} estatólitos na coifa e a maior sensibilidade da raiz à auxina, que por isso inibe o alongamento do lado de baixo;\newline%
{-} distinção explícita entre auxina e citocinina.\newline%
\newline%
Ela trata as armadilhas principais:\newline%
{-} a inversão na (a): a auxina diminui no lado iluminado (seta 1), e a curvatura não ocorre no meristema;\newline%
{-} a troca do hormônio por citocinina na (b) e na (d);\newline%
{-} o fato de a primeira parte da (d) estar certa e o erro estar só no hormônio;\newline%
{-} o efeito oposto da auxina em caule e raiz.\newline%
\newline%
Faltaram só pontos menores:\newline%
{-} na (c), não usou o segundo argumento para comparar 2 com 3 (a proximidade do ápice);\newline%
{-} na (b), não separou o local de percepção da gravidade (seta 6, coifa) do local da curvatura da raiz (zona de alongamento).\newline%
\newline%
É uma resolução que poderia ser publicada como gabarito comentado com pequenos ajustes.%
\begin{itemize}%
\item%
{[}minor{]} Na (c), a comparação entre as setas 2 e 3 se apoia só na migração lateral da auxina para o lado sombreado. Como as setas estão em alturas diferentes, faltou o segundo argumento, que reforça a conclusão: a seta 2 fica mais perto do ápice, onde a auxina é produzida e de onde desce por transporte polar, e a seta 3 fica mais abaixo. Deveria ter citado esse gradiente apical{-}basal, ou ao menos justificado por que dá para comparar regiões de alturas diferentes.%
\item%
{[}minor{]} Na (b), a resolução derruba a alternativa só pela troca de auxina por citocinina. Não comenta que a seta 6 (coifa/ponta) é o lugar onde a gravidade é percebida, e não onde a raiz se curva, que é a zona de alongamento, pouco acima. A resolução até menciona a zona de alongamento, mas só ao tratar do fototropismo do caule, e não volta a ela ao analisar a raiz. Bastaria uma frase para separar 'percepção na coifa' de 'curvatura na zona de alongamento' também na raiz.%
\item%
{[}minor{]} Na (d), o argumento extra de que 'a raiz da figura cresce verticalmente' é dispensável e um pouco frágil, porque no desenho ela é só aproximadamente vertical. O que basta é o argumento central: o hormônio que explica o geotropismo é a auxina, não a citocinina. O comentário não está errado, mas não acrescenta rigor.%
\end{itemize}%
\subsection*{Questão 46: nota 9,0 (aprovada, rodada 1 de 1)}%
\textbf{Resposta da resolução às cegas do avaliador:} Alternativa B. São pirâmides de biomassa. Uma pirâmide de fluxo de energia nunca fica invertida, por isso as alternativas a e c caem. Na figura, produtores e consumidores primários encolheram muito de t1 para t2, enquanto consumidores secundários e terciários continuaram do mesmo tamanho. Assim, a biomassa dos consumidores secundários fica maior que a dos níveis que os sustentam: isso indica desequilíbrio, com a teia ainda se reestruturando em t2, e não equilíbrio restabelecido (o que elimina a d).\par %
Conferi a figura no original (p. 16 do PDF, impressa como 14). Em t1 a pirâmide é regular. Em t2 a base branca e a faixa cinza{-}clara encolheram muito, enquanto as faixas cinza{-}escura e preta ficaram praticamente iguais, o que inverte a pirâmide até o 3º nível. O candidato leu a figura corretamente.\newline%
\newline%
A resposta final é B, a mesma da minha referência, e está correta. A justificativa é sólida. Uma pirâmide de fluxo de energia nunca se inverte, por causa das perdas a cada nível e da transferência de cerca de 10\%, e isso elimina a e c. Biomassa pode aparecer invertida, e o candidato cita tanto o caso transitório quanto o do fitoplâncton. Assim ele trata a armadilha de que inversão sozinha não prova desequilíbrio, ao basear o desequilíbrio na mudança de t1 para t2 e no atraso de resposta dos predadores. Também distingue b de d pelo 'ainda se reestruturando' contra 'equilíbrio restabelecido' e confere que produtores e consumidores primários foram os níveis atingidos.\newline%
\newline%
Não há erros de conteúdo. As falhas são de completude e precisão: faltam os erros secundários de a e c (invasão faria a faixa cinza{-}clara crescer; os níveis superiores não mudaram); o argumento 'd cita só produtores' é fraco se tomado sozinho; o equilíbrio em t1 é suposto sem dizer que é suposição; e há pequenas imprecisões de redação ('duas bases'; 'conservação da energia' onde caberia citar as perdas, isto é, a 2ª lei). Mesmo assim, a resolução está no nível de um gabarito comentado.%
\begin{itemize}%
\item%
{[}minor{]} As alternativas a) e c) são eliminadas só pelo tipo de pirâmide. Isso basta para a resposta, mas faltou apontar os outros erros delas. Na c), uma invasão por consumidor primário faria a faixa cinza{-}clara crescer, e ela diminuiu. Na a), as faixas dos consumidores secundários e terciários não mudaram; a mudança está na base. Um gabarito comentado deveria trazer essas observações.%
\item%
{[}minor{]} Um dos argumentos contra a d) é fraco: 'ela cita apenas os produtores'. Uma limitação que atinge os produtores pode afetar os consumidores primários em cascata, então essa omissão sozinha não torna a d) falsa. O argumento decisivo, que o candidato também dá, é que o equilíbrio não foi 'restabelecido em t2'. Convinha deixar claro que esse é o motivo principal e o outro é só acessório.%
\item%
{[}minor{]} O candidato diz que a mesma teia estava 'em equilíbrio (t1)', mas o enunciado não afirma isso; é uma suposição razoável a partir da pirâmide regular e deveria ser apresentada como tal. A expressão 'duas bases foram reduzidas' também é imprecisa: quem foi reduzido são os dois níveis inferiores.%
\item%
{[}minor{]} A frase 'pela conservação da energia, o fluxo nunca pode ser maior' sozinha só garante E\_\{n+1\} $\leq$ E\_n. A desigualdade estrita vem das perdas por respiração, calor e excretas, que o candidato cita logo antes. A justificativa está correta, mas a atribuição ao princípio físico ficou imprecisa: as perdas vêm da 2ª lei da termodinâmica.%
\end{itemize}%
\subsection*{Questão 47: nota 9,5 (aprovada, rodada 1 de 1)}%
\textbf{Resposta da resolução às cegas do avaliador:} Alternativa D: a AIDS é o estágio mais grave da infecção pelo HIV, e a terapia antirretroviral, ao tornar a carga viral indetectável, impede a evolução para a doença.\par %
Conferi a transcrição com o original (p. 16 do PDF) e ela está fiel. A resposta do candidato é D, a mesma da minha referência, e está correta: a AIDS é o estágio mais avançado da infecção pelo HIV, e a terapia antirretroviral, ao tornar a carga viral indetectável, impede a evolução para a doença. A resolução trata as armadilhas principais. Separa HIV (o vírus) de AIDS (a síndrome, com CD4 < 200/mm$^3$ ou doenças definidoras). Em (a), refuta a eliminação natural do vírus com o argumento do provírus latente. Em (b), aponta que não ter cura não torna o tratamento ineficaz e que as doenças oportunistas são consequência da imunodeficiência, não a causa da AIDS. Em (c), a distratora principal, reconhece que as duas primeiras afirmações estão certas e localiza o erro na inversão de causa e efeito. Em (d), confere cada afirmação e a relação de causa entre elas e não confunde indetectável com cura, porque já disse que a infecção não tem cura. Os conceitos estão corretos e o texto é claro e enxuto. As falhas são só de profundidade: a refutação de (b) poderia ser mais detalhada (I = I por via sexual, transmissão vertical, prevenção das doenças oportunistas), faltou citar os controladores de elite em (a), faltou a ressalva sobre as curas raras por transplante CCR5$\Delta$32 e não há menção ao contexto das vacinas. Um professor publicaria como gabarito comentado.%
\begin{itemize}%
\item%
{[}minor{]} Na alternativa (b), a resolução diz que os antirretrovirais previnem a transmissão e a progressão, mas não detalha o argumento. Faltou dizer que o I = I vale para a via sexual (isso aparece só nos conceitos{-}base) e citar a prevenção da transmissão vertical e a PrEP/PEP. Faltou também dizer que, ao preservar os CD4, a terapia evita as doenças oportunistas, que é o que a alternativa nega. Uma frase a mais deixaria a refutação completa.%
\item%
{[}minor{]} Em (a), a resolução não diz que a primeira frase ("a AIDS é causada pela infecção com HIV") está correta. Também não comenta os controladores de elite (menos de 1\%), que controlam a replicação mas não eliminam o vírus. Esse detalhe deixaria o 'sem tratamento, a maioria evolui' mais preciso, mas não muda a conclusão.%
\item%
{[}minor{]} A resolução diz que a infecção 'ainda não tem cura' sem ressalva. Poderia citar os casos raros de transplante de medula com doador CCR5$\Delta$32 e explicar que eles não invalidam 'permanece por toda a vida' como regra geral em (c). É um refinamento e não afeta a resposta.%
\item%
{[}minor{]} A resolução não comenta o texto{-}base (vacinas contra Covid{-}19 não contêm HIV nem causam imunodeficiência). Isso não é necessário para a resposta, mas um gabarito comentado normalmente liga o enunciado ao contexto em uma frase.%
\end{itemize}%
\subsection*{Questão 48: nota 9,5 (aprovada, rodada 1 de 1)}%
\textbf{Resposta da resolução às cegas do avaliador:} Alternativa C: os axônios concentram{-}se na substância branca, que fica na região periférica da medula e é formada por fibras nervosas (axônios mielinizados, organizados em tratos) que conduzem impulsos ascendentes ao encéfalo e descendentes gerados pelo encéfalo.\par %
Conferi a transcrição com a página 16 do PDF (página 14 impressa) e ela está fiel; a questão não tem figura. O candidato marcou a alternativa C, a mesma da minha resolução de referência, e a resposta está correta. Na medula espinhal, a substância cinzenta é central, tem forma de 'H' e concentra os corpos celulares. A substância branca é periférica e é formada por axônios mielinizados organizados em tratos ascendentes (sensitivos) e descendentes (motores). A resolução é correta, completa e bem organizada. Explica por que a substância branca é branca: a mielina é rica em lipídios. Liga a resposta ao texto{-}base: a lesão corta os tratos e causa perda motora e sensorial abaixo do nível lesado. Analisa as quatro alternativas apontando o elemento errado de cada uma: em (a), o tipo de substância; em (b), a posição e a composição. Explica por que a (d) não responde à pergunta, mesmo descrevendo bem a substância cinzenta. Também aponta a armadilha da inversão entre encéfalo e medula. Das armadilhas que listei, só não tratou de forma explícita a leitura de 'região periférica' como sistema nervoso periférico, e não comentou que a substância cinzenta também tem algumas fibras. Os problemas encontrados são pequenos: a generalização 'no encéfalo a disposição é a inversa', que deveria se limitar ao cérebro e ao cerebelo, e a descrição da (d) chamada de 'correta' sem a ressalva de que é incompleta. Um professor publicaria o texto como gabarito comentado com retoques mínimos.%
\begin{itemize}%
\item%
{[}minor{]} Ao dizer 'No encéfalo (córtex) a disposição é a inversa', o candidato generaliza demais. A inversão vale para o córtex cerebral e o cerebelar. O encéfalo também tem substância cinzenta profunda (núcleos da base, tálamo, núcleos do tronco), e o tronco encefálico tem organização mista. O correto seria escrever 'No cérebro e no cerebelo, o córtex (cinzento) fica por fora e a substância branca por dentro'.%
\item%
{[}minor{]} Faltou deixar claro que 'região periférica' indica a posição no corte transversal da medula, que faz parte do SNC, e não o sistema nervoso periférico. A expressão 'envolvendo a cinzenta' já esclarece isso de forma implícita, mas uma frase explícita resolveria a ambiguidade.%
\item%
{[}minor{]} Na alternativa (d), o candidato diz que a descrição da substância cinzenta está 'correta' sem apontar que ela é parcial: a cinzenta também abriga interneurônios e recebe aferências sensitivas no corno dorsal. Ele também não ressalva que a cinzenta tem algumas fibras (segmentos axonais e fibras locais amielínicas), embora não chegue a dizer o contrário. São ajustes de precisão e não mudam a conclusão.%
\end{itemize}%
\subsection*{Questão 49: nota 9,5 (aprovada, rodada 1 de 1)}%
\textbf{Resposta da resolução às cegas do avaliador:} Alternativa D. As três dimensões (a globalização como fábula, como perversidade e como possibilidade) formam uma unidade. O mundo é uma totalidade aberta e em movimento, e a "outra globalização" já pode ser vista em processos concretos que estão em curso nas grandes cidades.\par %
Conferi o enunciado nas páginas 16 e 17 do PDF (págs. 14 e 15 do caderno). A transcrição está fiel. O candidato marcou D, a mesma letra da minha referência, e essa resposta está correta: só a alternativa D trata fábula, perversidade e possibilidade como uma unidade, isto é, uma totalidade aberta e em movimento. Ela também apresenta a 'outra globalização' como tendência que já está em curso nas grandes cidades, como diz o último período do texto.\newline%
\newline%
Todas as alternativas foram analisadas e todas as armadilhas foram tratadas:\newline%
{-} **(a):** a 'aldeia global' e a redução das desigualdades pelo consumo pertencem à fábula, enquanto a realidade aprofunda a desigualdade.\newline%
{-} **(b):** as dimensões não estão dissociadas, e a terceira dimensão não é fábula. A resolução nota que a alternativa troca a primeira dimensão pela terceira.\newline%
{-} **(c):** são apontados três erros. A perversidade não esgota o mundo concreto. A globalização é seletiva e não 'homogeneizou' o mundo. E as cidades periféricas são o lugar onde surge a contrarracionalidade, não apenas vítimas passivas.\newline%
{-} **(d):** a justificativa liga a alternativa ao texto.\newline%
\newline%
Os acréscimos tirados da obra estão corretos: a 'morte do Estado' como uma das fábulas e a totalidade em processo. O texto é claro e bem organizado.\newline%
\newline%
Os únicos reparos são de detalhe:\newline%
{-} Em (b), o termo 'projeção de um futuro' poderia ser rebatido de forma explícita.\newline%
{-} Poderia ficar mais claro o que vem do trecho citado e o que vem da obra como um todo.\newline%
\newline%
Com o nível de um gabarito comentado publicável, a nota é 9,5.%
\begin{itemize}%
\item%
{[}minor{]} Em (b), a resolução rebate bem o 'dissociadas' e a troca da terceira dimensão pela 'fábula'. Faltou, porém, comentar explicitamente a expressão 'projeção de um futuro possível'. A tendência descrita no texto já se apreende no presente ('através da presença de uma humanidade misturada'), e não é só uma projeção. A ideia aparece de forma implícita ('apreendida como tendência no presente'), mas o contraste poderia estar explícito.%
\item%
{[}minor{]} As afirmações de que a 'homogeneização' faz parte do discurso do pensamento único e de que a escassez é a base da contrarracionalidade vêm da obra de Santos, não do trecho citado. Estão corretas, mas um gabarito comentado poderia apontar com mais clareza o que se tira do texto e o que se tira da obra.%
\end{itemize}%
\subsection*{Questão 50: nota 9,0 (aprovada, rodada 1 de 1)}%
\textbf{Resposta da resolução às cegas do avaliador:} B: a geopolítica na era da IA fortalece o poder dos EUA e da China, economias interdependentes e concorrentes que dominam a infraestrutura de IA, num mundo multipolar em transformação acelerada, com dilemas ético e existencial em aberto.\par %
Conferi o enunciado na página 17 do PDF (Questão 50, página 15 do caderno), e ele confere com a transcrição. A resposta do candidato é a alternativa (b), igual à minha referência e correta.\newline%
\newline%
A resolução é bem organizada. Primeiro o candidato separa os dados do texto:\newline%
{-} os atores são a China, que faz o alerta, e os EUA, que recebem o alerta e vivem um intenso debate sobre o tema;\newline%
{-} o dilema ético é algoritmos decidirem sobre vida e morte;\newline%
{-} o dilema existencial é a distopia de 'O Exterminador do Futuro';\newline%
{-} o pano de fundo é de forte tensão geopolítica.\newline%
\newline%
Depois traz o contexto de Geografia necessário: EUA e China são concorrentes e interdependentes, a ordem mundial é multipolar com outros polos de poder e não há blocos fechados como na Guerra Fria.\newline%
\newline%
O candidato tratou todas as armadilhas centrais:\newline%
{-} unipolaridade em (a), contrariada pela China, que aparece como contrapoder;\newline%
{-} bipolaridade 'recriada' em (d);\newline%
{-} 'de forma independente' contra a interdependência real, em (d);\newline%
{-} generalização 'sem critérios éticos', que contradiz o debate citado no texto, em (d);\newline%
{-} tomar o discurso diplomático chinês como fato e a inexistência de tratado sobre armas autônomas com IA, em (c);\newline%
{-} leitura correta dos dilemas ético e existencial que sustentam (b).\newline%
\newline%
As falhas são pequenas:\newline%
{-} não reconhece a parte verdadeira de (a);\newline%
{-} não menciona que a China também militariza a IA;\newline%
{-} afirma, de modo impreciso, que chips avançados e terras raras 'se concentram' nos EUA e na China;\newline%
{-} trata de leve a tensão entre 'hegemonia' e 'multipolar' em (b).\newline%
\newline%
Nenhuma dessas falhas compromete a resposta. É uma resolução publicável como gabarito comentado, com pequenos retoques.%
\begin{itemize}%
\item%
{[}minor{]} Na alternativa (a), o candidato só aponta o erro da 'ordem unipolar'. Não comenta que a parte final ('objetos inteligentes voltados à guerra que estão sob discussão ética') está de acordo com o texto, nem que o 'tecnopoder compartilhado com grandes empresas' não basta para salvar a alternativa. Seria melhor dizer que essa parte é verdadeira, mas que o núcleo da afirmação é falso.%
\item%
{[}minor{]} Na alternativa (c), falta lembrar que a própria China também desenvolve aplicações militares de IA. Isso derruba de vez a ideia de uma liderança 'ética e não distópica'. Também falta comentar a afirmação de que a China 'assumiu a liderança na construção da ordem multipolar'. Os argumentos que o candidato dá já bastam para descartar a alternativa, mas a refutação ficaria completa com esses dois pontos.%
\item%
{[}minor{]} Há uma imprecisão factual no contexto. O candidato diz que a infraestrutura de IA, incluindo 'chips avançados' e 'terras raras', 'concentra{-}se sobretudo nos EUA e na China'. Na verdade, a fabricação de chips de ponta está concentrada em Taiwan e na Coreia do Sul (TSMC, Samsung), e o refino de terras raras está concentrado na China, não nos EUA. O certo seria falar em 'controle/liderança' das duas potências sobre a cadeia (projeto de chips, modelos, data centers, nuvem, insumos), e não em concentração física de todos esses itens nos dois países. Isso não afeta a escolha da alternativa.%
\item%
{[}minor{]} A tensão entre 'hegemonia' e 'mundo multipolar' na alternativa (b) aparece só de passagem ('a hegemonia tecnológica sino{-}americana se exerce dentro desse quadro'). Valeria explicitar que 'hegemonia' aqui quer dizer liderança tecnológica e econômica, e não domínio exclusivo. Isso evitaria uma objeção natural do leitor.%
\end{itemize}%
\subsection*{Questão 51: nota 9,0 (aprovada, rodada 1 de 1)}%
\textbf{Resposta da resolução às cegas do avaliador:} Alternativa A. Com a expansão do etanol de milho em Mato Grosso surgem as usinas flex, unidades adaptadas para processar tanto a cana quanto o milho. Elas usam o milho, que pode ser armazenado, durante a entressafra da cana e assim operam o ano todo, com maior flexibilidade na produção do biocombustível.\par %
Conferi o enunciado na p. 17 do PDF e a transcrição está fiel. O candidato marcou (A), a mesma resposta da minha referência, e eu a considero correta: a expansão do etanol de milho em MT fez surgir as usinas flex, que processam cana na safra e milho estocado na entressafra, ganham flexibilidade e operam o ano todo.\newline%
\newline%
A resolução trata as armadilhas principais:\newline%
{-} explica a perecibilidade da cana frente à armazenabilidade do milho e usa isso para refutar (b);\newline%
{-} usa o caráter semiperene da cana e a sucessão soja--milho safrinha para refutar (c), separando a entressafra industrial de uma suposta rotação no campo;\newline%
{-} usa a mecanização da colheita do milho em grandes propriedades para refutar (d);\newline%
{-} cita usinas flex e full em (a).\newline%
\newline%
O texto é claro, conciso e bem organizado. Faltam alguns refinamentos:\newline%
{-} Em (a), uma frase ambígua dá a entender que a alternativa descreveria as usinas full, e falta comentar o 'poderá'.\newline%
{-} Em (b), não reconhece a parte verdadeira do distrator, o compartilhamento de parte da planta industrial.\newline%
{-} Não ressalta que o gráfico mostra participações relativas.\newline%
{-} Em (d), faz afirmações categóricas sobre emprego sem apoio.\newline%
\newline%
São problemas menores, e a resolução tem nível de gabarito comentado publicável.%
\begin{itemize}%
\item%
{[}minor{]} Em (a), a frase 'Há também as usinas full, que usam só milho. É exatamente o que a alternativa descreve.' fica ambígua: pode parecer que a alternativa descreve as usinas full, quando ela descreve as usinas flex. Deveria dizer explicitamente que o 'poderá' da alternativa indica uma possibilidade (usinas flex) e que a existência de usinas full não contradiz a alternativa.%
\item%
{[}minor{]} Em (b), a refutação está correta, mas o candidato não reconhece a parte verdadeira que torna o distrator atraente: há compartilhamento de parte da planta industrial (destilaria, tanques, caldeira a bagaço) nas usinas flex. Deveria dizer que o erro está na justificativa (estocagem equivalente e proximidade obrigatória para as duas), não na ideia de aproveitar infraestrutura.%
\item%
{[}minor{]} Na leitura do gráfico, não se observa que os valores são participações relativas (\%). A queda da cana de $\approx$93\% para $\approx$52\% não significa necessariamente queda absoluta da produção de etanol de cana. O candidato não comete o erro, mas deixa de fazer a ressalva.%
\item%
{[}minor{]} Em (d), as afirmações 'Ela não absorve mais mão de obra que a cana' e 'não amplia o emprego no campo' são categóricas e não têm apoio. Seria melhor dizer que a colheita da cana também foi amplamente mecanizada (o 'ao contrário da cana' da alternativa é falso) e que a cadeia grão--etanol emprega relativamente pouca mão de obra no campo, o que invalida a premissa da 'pouca mecanização'.%
\end{itemize}%
\subsection*{Questão 52: nota 9,0 (aprovada, rodada 1 de 1)}%
\textbf{Resposta da resolução às cegas do avaliador:} Alternativa A. O crescimento das cidades médias da tabela se explica pela combinação de três tipos de vetor. O primeiro é o avanço da fronteira agrícola (Sinop, Sorriso e Luís Eduardo Magalhães) e da fronteira mineral (Parauapebas/Carajás). O segundo são os combustíveis fósseis, com os royalties do pré{-}sal em Maricá, e o turismo, em Camboriú e Palhoça. O terceiro é a dispersão metropolitana com expansão imobiliária em áreas periurbanas e semirrurais (Senador Canedo, Fazenda Rio Grande, Palhoça, São José de Ribamar e Maricá). A alternativa B está errada porque não houve industrialização de base nem crescimento ligado a polos portuários e corporativos. A C está errada porque a violência se interiorizou, em vez de diminuir. A D está errada porque os programas habitacionais não transferem população para o interior e não explicam o crescimento das cidades do agronegócio e da mineração.\par %
A resposta final (A) está correta e coincide com a minha referência. Conferi o enunciado na página 18 do PDF e a transcrição está fiel. Refiz as dez taxas com Python e todas conferem com (P2022/P2010 - 1)$\cdot$100.\newline%
\newline%
A resolução faz o mais importante. Ela não justifica a alternativa A só por eliminação: associa cada vetor da alternativa a municípios concretos da tabela:\newline%
{-} fronteira agrícola: LEM no MATOPIBA, Sinop e Sorriso na BR{-}163;\newline%
{-} fronteira mineral: Parauapebas, em Carajás;\newline%
{-} combustíveis fósseis: Maricá, com os royalties do pré{-}sal da Bacia de Santos, o que está correto;\newline%
{-} turismo: Camboriú;\newline%
{-} dispersão metropolitana com expansão imobiliária em áreas semirrurais: Senador Canedo, Fazenda Rio Grande, Palhoça e São José de Ribamar.\newline%
Os dez casos ficam cobertos.\newline%
\newline%
As refutações tratam as armadilhas principais:\newline%
{-} B: a indústria de base não é o vetor e não há polos portuários com funções corporativas;\newline%
{-} C: o fenômeno real é a interiorização da violência, e segurança pública não é vetor demográfico;\newline%
{-} D: o MCMV atua dentro das próprias regiões metropolitanas, não transfere população para o interior e não explica as cidades do agro, da mineração e do petróleo.\newline%
\newline%
As falhas são todas menores:\newline%
{-} a verificação numérica foi feita em um único exemplo e não diz explicitamente que a taxa é acumulada;\newline%
{-} falta a comparação com o crescimento nacional;\newline%
{-} não aponta os municípios que 'ascenderam' a partir de menos de 100 mil habitantes;\newline%
{-} a frase sobre desindustrialização na B é confusa;\newline%
{-} Palhoça aparece sem o vetor turístico.\newline%
\newline%
É uma resolução que pode ser publicada como gabarito comentado, com pequenos ajustes de completude.%
\begin{itemize}%
\item%
{[}minor{]} O candidato afirma que 'as taxas conferem', mas só faz a conta para Senador Canedo. Refiz as dez com Python e todas batem (84,31; 82,27; 79,53; 74,02; 73,58; 66,32; 65,29; 62,09; 54,77; 50,01). Portanto a afirmação é verdadeira, mas não foi demonstrada. Ele também não diz explicitamente que a taxa é acumulada no período, e não anual (cerca de 5,2 \% ao ano no máximo). Bastava uma frase dizendo isso.%
\item%
{[}minor{]} Falta comparar com o crescimento nacional (cerca de 6,5 \% entre 2010 e 2022, ou 0,52 \% ao ano). Essa comparação mostra que as cidades cresceram de 8 a 13 vezes acima da média, o que só se explica por migração. Ela reforçaria a ideia de 'vetores' que atraem população.%
\item%
{[}minor{]} O candidato não aponta que cinco municípios tinham menos de 100 mil habitantes em 2010 (Senador Canedo, Fazenda Rio Grande, LEM, Sorriso e Camboriú). Eles 'ascenderam' à condição de cidade média, como diz o enunciado. Isso não afeta a resposta, mas deixa a leitura da tabela menos completa.%
\item%
{[}minor{]} Na refutação da B, a frase 'Na dinâmica recente, a indústria se desconcentra, mas não é esse o vetor dominante' é confusa e não usa o conceito esperado de desindustrialização relativa. O argumento central está correto: a indústria de base não é o motor, não há polos portuários concentradores de sedes corporativas e essas funções seguem nas metrópoles. Faltou só formulá{-}lo com mais precisão.%
\item%
{[}minor{]} A associação de Palhoça ficou incompleta. O candidato a coloca apenas na dispersão metropolitana, sem citar o vetor turístico e imobiliário litorâneo, que também atua ali, na Grande Florianópolis. Maricá, por sua vez, recebe as duas leituras (royalties e expansão a partir da metrópole), o que é correto. Isso não compromete a conclusão.%
\item%
{[}minor{]} Na D, o candidato não comenta a parte da alternativa que diz que os programas 'ajudam a reduzir o déficit habitacional', e o déficit continua elevado. O argumento principal, porém, está correto: não há transferência de população para o interior e esse vetor não explica as cidades do agro, da mineração e do petróleo.%
\end{itemize}%
\subsection*{Questão 53: nota 9,5 (aprovada, rodada 3 de 3)}%
\textbf{Resposta da resolução às cegas do avaliador:} Alternativa (a): os desertos ficam majoritariamente na zona relativamente natural (A), enquanto as savanas aparecem nas três zonas (A, B e C), com distribuição mais equilibrada.\par %
A resposta final é a alternativa (a), a mesma da minha referência, e está correta.\newline%
\newline%
Conferi o mapa no PDF (p. 18), ampliando o recorte, e a leitura do candidato está mais precisa que a minha. Minha referência dizia que não há rótulos B ou C no interior desértico. O mapa mostra o contrário: há quatro rótulos C sobre o amarelo{-}claro do deserto, que o candidato identificou e localizou corretamente:\newline%
{-} o litoral de Carnarvon/Gascoyne;\newline%
{-} a Grande Baía Australiana, onde a linha de chamada termina no amarelo;\newline%
{-} o lobo leste no sudoeste de Queensland;\newline%
{-} a região de Broken Hill, a cerca de 141,5$^\circ$E e 32$^\circ$S.\newline%
\newline%
O candidato também viu que os rótulos B e C do litoral de Queensland apontam para as faixas de floresta, e não para a savana. As linhas de chamada terminam, com boa precisão, perto de Cairns (\textasciitilde{}146$^\circ$E, 17$^\circ$S) e de Mackay (\textasciitilde{}149$^\circ$E, 21$^\circ$S). Nesse ponto minha referência errava ao contar esses rótulos como savana. A conclusão não muda, porque continua havendo C sobre a savana (interior do Top End, interior do norte de Queensland e as linhas de chamada no litoral norte) e B no sudeste de Queensland. Também está certo que nenhum B fica sobre o deserto: os B vizinhos apontam para a vegetação mediterrânea ou para a floresta.\newline%
\newline%
Com essa leitura, o deserto fica majoritariamente em A, com alguns C nas bordas e nenhum B, e a savana aparece nas três zonas. A alternativa (a) se confirma. O candidato tratou todas as armadilhas:\newline%
{-} cruzou as duas camadas do mapa e explicou o critério para as linhas de chamada;\newline%
{-} mostrou que (b) inverte os biomas;\newline%
{-} explicou que 'majoritariamente' e 'mais equilibrada' são termos comparativos;\newline%
{-} em (c), admitiu de forma honesta que existem C no deserto, mostrou que ficam na periferia e lembrou que a agricultura e as cidades estão no litoral úmido;\newline%
{-} em (d), apontou os C sobre a savana e o caráter determinista da justificativa;\newline%
{-} tratou os dados numéricos como contexto.\newline%
\newline%
Os problemas são poucos e menores: uma afirmação sem base sobre qual zona é a mais extensa na savana e um comentário pouco verificado sobre a coerência dos números. A resolução é clara, bem organizada e publicável como gabarito comentado.%
\begin{itemize}%
\item%
{[}minor{]} O candidato afirma que 'A zona A ainda é a mais extensa da savana', mas não mostra nenhuma base para isso. As zonas C também são grandes: o bloco do norte de Queensland, do Golfo até Cape York, e o interior do Top End podem ter área comparável à das zonas A. A afirmação não é necessária para o argumento. Bastaria dizer que a savana tem porções expressivas nas três zonas.%
\item%
{[}minor{]} Na parte dos dados numéricos, o candidato diz que os números 'são coerentes com o mapa, em que a maior parte do território está nas zonas A e B', mas não verifica isso. Também deixou de lado duas checagens simples: 664 Mha ÷ 0,865 $\approx$ 768 Mha, que bate com a área da Austrália, e a soma 86,5\% + 13,2\% = 99,7\%. Ele faz bem em separar 'vegetação nativa' de 'zona A'. Mesmo assim, o parágrafo poderia ser mais cuidadoso ou mais curto, porque esses números não decidem a questão.%
\item%
{[}minor{]} A descrição da extensão da savana não cita Cape York, embora a península inteira esteja em laranja. Isso não afeta a conclusão.%
\end{itemize}%
\subsection*{Questão 54: nota 9,0 (aprovada, rodada 2 de 2)}%
\textbf{Resposta da resolução às cegas do avaliador:} Alternativa C. Pelo Mapa 1, a precipitação anual da RM de São Paulo fica toda em 1600--1800 mm. Pelo Mapa 2, a evapotranspiração potencial (ETP) fica abaixo de 1400 mm no leste da RM e entre 1400 e 1600 mm no oeste. Logo, P > ETP no balanço anual. Mesmo assim, as chuvas se concentram no verão, a reserva depende de reservatórios e a demanda urbana é muito alta, o que exige monitoramento e alertas nos períodos secos. As alternativas A e B erram nas conclusões. A D erra na leitura do mapa: em Vitória, P e ETP estão na mesma classe, 1600--1800 mm.\par %
A resposta do candidato, (C), coincide com a minha referência e com o que se lê no original ampliado.\newline%
\newline%
**Leitura dos mapas.** Conferi cada RM nos recortes do PDF.\newline%
{-} **Belo Horizonte:** precipitação (P) de 1400--1600 e evapotranspiração potencial (ETP) de 1400--1600.\newline%
{-} **Rio de Janeiro:** P de 1600--1800; ETP dividida entre 1400--1600 (oeste) e 1600--1800 (leste).\newline%
{-} **São Paulo:** P inteira em 1600--1800; ETP de 1400--1600 no oeste e < 1400 no leste. Como P > 1600 $\geq$ ETP, vale P > ETP sem ambiguidade.\newline%
{-} **Vitória:** P de 1600--1800, na mancha costeira, e ETP de 1600--1800.\newline%
\newline%
**Armadilhas tratadas.**\newline%
{-} Leu o original, e não a transcrição errada que levaria à (d).\newline%
{-} Usou a regra correta de comparação entre classes: só afirma 'maior' quando as classes não se sobrepõem.\newline%
{-} Na (a), separou a premissa verdadeira da inferência falsa (P $\approx$ ETP indica margem pequena, e média anual não descreve o que ocorre 'ao longo do ano').\newline%
{-} Na (b), atacou a linguagem absoluta ('assegura', 'dispensa') e a ausência de sazonalidade.\newline%
{-} Na (c), justificou com os três fatores do enunciado (sazonalidade, armazenamento, demanda), com uma conta de cota correta: P - ETP entre 0 e 400 mm no oeste e acima de 200 mm no leste. Citou de forma correta a crise do Cantareira de 2014--2015.\newline%
{-} Na (d), derrubou a premissa por P e ETP estarem na mesma classe.\newline%
\newline%
**Ressalvas, todas menores.**\newline%
{-} Faltou ressalvar as bordas dos polígonos: o extremo leste do Rio na precipitação e a ponta norte de Vitória.\newline%
{-} Não citou a incoerência 'faixas distintas $\to$ equilíbrio' na (b).\newline%
{-} Não registrou que a transcrição diverge do original.\newline%
{-} Igualou o excedente a P - ETP, quando P - ETP é só uma cota inferior.\newline%
\newline%
A resolução é clara, rigorosa e completa, com qualidade de gabarito comentado.%
\begin{itemize}%
\item%
{[}minor{]} Na tabela, a precipitação da RM do Rio aparece como 1600--1800 sem ressalva. No original ampliado, o extremo leste do polígono parece tocar a faixa de 1400--1600. Isso não muda nada na análise da (b), mas a leitura deveria vir com a ressalva 'quase toda em 1600--1800', como foi feito para a ETP.%
\item%
{[}minor{]} Em Vitória, o candidato diz que a RM 'está na mancha' de 1600--1800 sem citar que a pequena ponta norte do polígono pode tocar a classe 1200--1400. Também chama o entorno de '1200--1400', mas a oeste há trecho em 1400--1600. Vale registrar a ressalva e explicar que, como a área é dominante, P e ETP ficam na mesma classe e a premissa da (d) não se sustenta.%
\item%
{[}minor{]} Na (b), faltou apontar a incoerência interna da alternativa: 'faixas distintas' com P > ETP indicariam excedente, e não 'equilíbrio hídrico'. O candidato derruba a alternativa pelo caráter absoluto ('assegura', 'dispensa'), o que basta, mas o argumento ficaria mais completo.%
\item%
{[}minor{]} A resolução não avisa que a transcrição do enunciado diverge do original: a transcrição põe Vitória em 1200--1400 e SP em $\approx$1400--1600 de precipitação. O candidato leu corretamente o original ampliado, mas um gabarito comentado deveria registrar a divergência, porque é exatamente ela que levaria à (d).%
\item%
{[}minor{]} O candidato iguala o 'excedente anual' a P - ETP ('entre 0 e 400 mm'). A rigor, no balanço de Thornthwaite o excedente é P - ETR $\geq$ P - ETP, pois ETR $\leq$ ETP. A conclusão (P - ETP > 0 $\Rightarrow$ há excedente) está certa, mas os números são uma cota inferior, não o valor do excedente.%
\end{itemize}%
\subsection*{Questão 55: nota 9,0 (aprovada, rodada 1 de 1)}%
\textbf{Resposta da resolução às cegas do avaliador:} Alternativa A. O logotipo identifica e simboliza a universidade. Ele lembra a organização circular{-}radial do campus, mas é uma representação gráfica simplificada de identidade institucional, e não um mapa.\par %
O enunciado foi conferido no PDF original (página 19 do arquivo, p. 17 do caderno), e a transcrição está correta. A resposta do candidato é (A), a mesma da minha referência, e está correta. A resolução separa bem representação simbólica de cartográfica e usa os elementos do mapa (escala, orientação, legenda e correspondência posicional) para mostrar que o logotipo não é mapa. Também observa que o Norte aparece na foto aérea, não no logotipo. Todas as armadilhas foram tratadas. A semelhança visual é tratada como 'diálogo' ou estilização, e não como cartografia. Em (c), a resolução aponta o erro da justificativa ('não reproduz a distribuição dos edifícios e das vias'), e não só a parte certa. Em (b) e (d), recusa as funções de reconhecer trajetos e orientar deslocamentos. Em (a), preserva a ressalva 'dialoga visualmente' sem negar a relação entre o logotipo e o campus. Os problemas são pequenos: (1) a afirmação de que todo mapa precisa 'no mínimo' de escala, orientação e legenda é exagerada; (2) a análise de (d) não refuta explicitamente 'reproduz sua organização espacial'; (3) a resolução atribui um significado aos círculos vermelhos sem base no enunciado; (4) a comparação visual concreta é pouco usada (o logotipo é um anel aberto, o anel viário é fechado). A resolução é clara, correta e completa o bastante para servir de gabarito comentado, com pequenos ajustes de precisão. Nota 9.%
\begin{itemize}%
\item%
{[}minor{]} A frase 'Um mapa precisa, no mínimo, dos elementos...' e o fecho 'Só com esses elementos é possível localizar lugares, medir distâncias e planejar trajetos' exageram. Croquis e mapas esquemáticos (como os de metrô) não têm escala, ou só mantêm relações topológicas, e mesmo assim orientam trajetos. O certo seria dizer que esses elementos caracterizam a representação cartográfica convencional e permitem uma leitura métrica e orientada do espaço. A lista também omite o título e as coordenadas/projeção, que costumam constar entre os elementos do mapa.%
\item%
{[}minor{]} Em (d), a refutação é curta demais: diz só que o logotipo 'não é mapa e não serve para orientar deslocamentos'. Não responde explicitamente à afirmação de que ele 'reproduz sua organização espacial'. Deveria remeter à seção 2: o logotipo estiliza, não reproduz, porque não tem escala, orientação nem legenda e não guarda correspondência posicional.%
\item%
{[}minor{]} A resolução afirma, sem base no enunciado, que os três círculos vermelhos 'são elementos gráficos da marca, não lugares localizáveis'. A prova não diz o que eles representam. Seria melhor dizer que, sem legenda, não é possível associá{-}los a lugares do campus. Isso também reforça o argumento.%
\item%
{[}minor{]} Faltou um argumento visual direto que reforçaria a tese de que o logotipo não reproduz fielmente o campus: ele é um anel aberto, em forma de 'C', com a abertura onde estão os círculos vermelhos, enquanto o anel viário da foto é fechado. A comparação concreta entre as duas imagens fica um pouco genérica.%
\end{itemize}%
\subsection*{Questão 56: nota 9,0 (aprovada, rodada 1 de 1)}%
\textbf{Resposta da resolução às cegas do avaliador:} Alternativa B. Entre os que apoiam deportações condicionadas, só 14\% deportariam pais/mães de crianças nascidas nos EUA e só 5\% deportariam quem é casado com cidadão americano. São índices muito menores que os de crimes violentos (97\%), crimes não violentos (52\%) e chegada recente (44\%), o que mostra mais tolerância com quem tem vínculos familiares no país.\par %
Conferi o enunciado na página 20 do PDF original e a transcrição está fiel. O candidato marcou a alternativa B, a mesma da minha referência, e a resposta está certa.\newline%
\newline%
A leitura do gráfico está correta. O ponto central é que os percentuais das barras de baixo valem só para o subgrupo 'Some' (51\% do total), e o candidato explica isso e mostra como levar cada valor ao total multiplicando por 0,51. As contas conferem: 0,52$\cdot$51 $\approx$ 26,5\%, 0,97$\cdot$51 $\approx$ 49,5\%, 32 + 26,5 $\approx$ 58,5\% e 32 + 49,5 $\approx$ 81,5\%.\newline%
\newline%
Armadilhas tratadas:\newline%
{-} **C:** ele mostra que os 52\% equivalem a cerca de um quarto do total, e não a um terço, e que o 32\% próximo de 1/3 é o grupo da deportação irrestrita.\newline%
{-} **D:** ele mostra que os 97\% valem só para o subgrupo e que os 16\% do 'None' impedem a 'quase totalidade'.\newline%
{-} **B:** ele identifica corretamente as categorias de vínculo familiar (14\% e 5\%) e as traduz em tolerância à permanência (86\% e 95\%).\newline%
\newline%
Os defeitos são pequenos e de precisão ou completude:\newline%
{-} Na (d), chama de cota máxima ('no máximo') o que é a estimativa direta de 81,5\%; o teto é 84\%.\newline%
{-} Na (a), não explicita a confusão entre 'Some' (51\%) e deportação irrestrita.\newline%
{-} Na (b), poderia comparar também com 'chegada recente' (44\%) e 'ter emprego' (15\%).\newline%
\newline%
É uma resolução clara e correta, publicável como gabarito comentado com pequenos ajustes.%
\begin{itemize}%
\item%
{[}minor{]} Na (d), o candidato chama 32\% + 0,97$\cdot$51\% $\approx$ 81,5\% de cota máxima ('no máximo'). Esse valor é a estimativa direta, já que quem está no grupo 'All' deportaria qualquer imigrante irregular. O teto lógico é 84\%, porque só os 16\% do 'None' ficam de fora. Deveria dizer 'cerca de 81,5\%' e apresentar 84\% como limite superior absoluto.%
\item%
{[}minor{]} Na (a), o candidato só diz que 'All' = 32\%. Não explica que os 51\% (o único número acima da metade) são do grupo 'Some', que defende deportações condicionadas, o contrário de 'irrestrita'. Era essa confusão que a alternativa explorava, e valia nomeá{-}la.%
\item%
{[}minor{]} Na (b), a comparação usa só os critérios de crime (97\% e 52\%). Não menciona que 14\% e 5\% também ficam abaixo de 'chegada recente' (44\%) e de 'ter emprego' (15\%). Com isso o argumento comparativo de que esse grupo é 'mais tolerante' com vínculos familiares ficaria completo, mostrando que esses estão entre os menores índices do gráfico.%
\end{itemize}%
\subsection*{Questão 57: nota 9,0 (aprovada, rodada 1 de 1)}%
\textbf{Resposta da resolução às cegas do avaliador:} Alternativa A. Os dois artistas tratam a arte como transgressora ("revolutionary" para Moura, "dangerous" para Morrison). Só Morrison diz quem ataca os artistas: "dictators", "people in office" e "people who want to control and deceive". Moura fala apenas de um vago "these people".\par %
Conferi o enunciado na página 20 do PDF (impressa como 18), e a transcrição está fiel. O candidato marcou (A), a mesma resposta da minha referência, e a resposta está correta.\newline%
\newline%
Os dois depoimentos concordam que a arte é transgressora: Moura a chama de 'revolutionary' e Morrison, de 'dangerous'. Só Morrison diz quem ataca os artistas: 'dictators', 'people in office' e 'people who want to control and deceive'.\newline%
\newline%
O candidato tratou bem as principais armadilhas:\newline%
{-} **(B):** percebeu que a relação de medo está invertida. Moura diz que os críticos não temem o que ele fala, e sim o que os artistas fazem.\newline%
{-} **(C):** identificou o falso cognato: 'deceive' significa enganar, não decepcionar. Também mostrou que quem engana são os opressores.\newline%
{-} **(D):** notou que a segunda parte é verdadeira (só Moura cita filmes), mas que a primeira é falsa, porque nenhum dos textos trata de uma 'luta contra a punição'.\newline%
{-} **(A):** explicou a síntese 'natureza transgressora' e não confundiu as fontes.\newline%
\newline%
As falhas são pequenas:\newline%
{-} Fala em 'poderosos' para os 'they' de Moura, que o texto não identifica.\newline%
{-} Explica de forma curta por que 'these people' não conta como 'tipo de pessoa'.\newline%
{-} Não comenta a ressalva 'though it might bother them'.\newline%
{-} Não avalia a primeira parte de (B) e de (C).\newline%
\newline%
É uma resolução clara, correta e enxuta, que daria um gabarito comentado publicável com pequenos ajustes.%
\begin{itemize}%
\item%
{[}minor{]} Na seção 'Ponto em comum', o candidato escreve que 'os poderosos têm medo dos artistas'. Para Moura, porém, quem tem medo são os 'these people' e 'they', e o próprio candidato diz depois que Moura não os identifica. Chamar esse grupo de 'poderosos' tira dos textos uma informação que eles não dão e entra em tensão com a análise da alternativa a). Ficaria mais correto escrever algo como 'quem se opõe aos artistas os teme ou se sente ameaçado por eles'.%
\item%
{[}minor{]} Na alternativa a), a justificativa de que 'these people' não é um 'tipo de pessoa' está curta. O candidato diz apenas 'sem dizer quem são'. Teria sido melhor dizer que Moura só identifica esse grupo pelo que eles dizem ('stay in your lane'), sem categoria social ou política. Também faltou dizer, na própria análise de a), que só Morrison relata ataques concretos ('murdered, imprisoned'). Ele cita isso no resumo inicial, mas não usa o dado ao avaliar a alternativa.%
\item%
{[}minor{]} Na alternativa b), o candidato não menciona a ressalva 'though it might bother them', que é o que pode confundir quem lê rápido. Também não avalia a primeira parte de b) e de c) (arte e política; arte e história). O erro dessas alternativas está mesmo na segunda parte, mas um gabarito comentado completo diria que a primeira parte se sustenta.%
\item%
{[}minor{]} O candidato traduz 'people in office' de dois jeitos: 'ocupantes de cargos de poder' no resumo e 'ocupantes de cargos públicos' na alternativa a). É só uma questão de padronizar o texto.%
\end{itemize}%
\subsection*{Questão 58: nota 9,0 (aprovada, rodada 1 de 1)}%
\textbf{Resposta da resolução às cegas do avaliador:} Alternativa c): a ilustração dos frutos vermelhos em forma de gota, cobertos de cerdas flexíveis, que é o urucum (Bixa orellana).\par %
Conferi o original (página 20 do PDF, página 18 impressa). O enunciado e as quatro ilustrações conferem com a transcrição: a) é uma vagem de cacau vinho com sulcos; b) são frutos redondos alaranjados, tipo tucumã; c) são duas cápsulas vermelhas em forma de gota, cobertas de cerdas (urucum); d) são cachos vermelhos abertos com sementes pretas em arilo branco (guaraná).\newline%
\newline%
A resposta final do candidato, (C), é a mesma da minha referência e está correta. A resolução traduz bem cada pista do texto: dye, teardrop{-}shaped, malleable bristles, fuzzy, insect bites. Também identifica o urucum (Bixa orellana, do tupi uru{-}ku) e estabelece um critério claro: a figura tem de ser vermelha, em forma de gota e coberta de cerdas. Com isso, analisa as quatro alternativas.\newline%
\newline%
As armadilhas principais foram tratadas: a cor e a origem amazônica não decidem, e o vocabulário foi entendido. Os problemas são pequenos. A casca do cacau em a) foi descrita como \textbackslash{}"lisa\textbackslash{}", quando a ilustração mostra uma casca sulcada e não lisa. E a distratora d) poderia ter sido mais bem refutada, porque o texto também fala em sementes e o uso das sementes do guaraná é outro.\newline%
\newline%
É uma resolução clara e completa, que pode ser publicada com uma pequena correção na descrição da alternativa a).%
\begin{itemize}%
\item%
{[}minor{]} Na alternativa a), o candidato diz que o fruto tem "casca lisa com sulcos longitudinais". A ilustração da página 20 do PDF mostra uma vagem de cacau com sulcos e casca texturizada, ou seja, não lisa. Isso não muda a resposta, mas a descrição está errada e deixa passar uma armadilha: o fruto tem casca "non{-}smooth" e cor avermelhada. O certo seria dizer que a casca é sulcada e não lisa, e que o fruto só é descartado por ser grande e alongado, sem forma de gota e sem cerdas.%
\item%
{[}minor{]} Na alternativa d), principal distratora, a resolução aponta só a falta de cerdas. Faltou dizer que a casca é lisa, que o fruto não tem forma de gota e que as sementes à mostra, embora lembrem o "seeds" do texto, não servem para pintura corporal (no guaraná, são usadas como estimulante). O candidato apresenta o critério (vermelho + gota + cerdas) mas não desmonta essa atração da alternativa.%
\item%
{[}minor{]} Na alternativa b), o candidato também afirma "casca lisa" sem que isso seja essencial. Os critérios que realmente decidem (forma redonda, cor laranja, sem cerdas) estão corretos.%
\end{itemize}%
\subsection*{Questão 59: nota 9,0 (aprovada, rodada 1 de 1)}%
\textbf{Resposta da resolução às cegas do avaliador:} B: Aleksic explica a função de algumas gírias na demarcação de território e identidade cultural ("signaled in{-}group status", "tool of power to build community and shared identity").\par %
Conferi o enunciado nas páginas 20--21 do PDF (páginas 18--19 do caderno), e a transcrição está fiel. O candidato marca (B), a mesma resposta da minha referência, e a resposta está correta: Aleksic explica que as gírias "signaled in{-}group status" e servem como "tool of power to build community and shared identity". Isso equivale a demarcar pertencimento (território simbólico) e identidade cultural. Todas as armadilhas foram tratadas. Em (a), a resolução aponta o verbo valorativo 'defende' e o tom crítico, sustentados pelas citações de apropriação e de perda de exclusividade. Em (c), refuta o 'direito de recusar' mostrando que a apropriação é inevitável no ciclo descrito. Em (d), identifica corretamente a inversão de agentes: quem tinha a exclusividade eram os grupos marginalizados, e os dominantes a destroem ao buscar a 'perceived coolness'. A ligação entre 'território' e pertencimento aparece, embora só na seção introdutória. As citações estão bem escolhidas e o texto está claro. Os problemas são pequenos: um asterisco sobrando em 'marginalizados*', o exagero de 'corresponde exatamente' no item (b) sem retomar ali a noção de território simbólico, e a falta de 'signals queer excitement' como reforço para identidade cultural. Com pequenos ajustes, é uma resolução publicável como gabarito comentado.%
\begin{itemize}%
\item%
{[}minor{]} Há um erro de digitação em \textbackslash{}emph\{marginalizados*\}: sobrou um asterisco. Basta removê{-}lo.%
\item%
{[}minor{]} Na lista, a justificativa de (b) é curta e diz que a alternativa 'corresponde exatamente' aos trechos citados. Isso é exagero, porque 'demarcação de território' é uma leitura interpretativa (território simbólico, a fronteira entre quem é e quem não é do grupo), e não uma paráfrase literal. A ligação aparece na seção de ideia central ('delimitar um espaço próprio do grupo (território)'), mas deveria ser repetida no item (b), dizendo que se trata de território simbólico ou de pertencimento, e não físico.%
\item%
{[}minor{]} Para reforçar o segundo elemento de (b), 'identidade cultural', a resolução poderia citar também 'it signals queer excitement'. Foi usado só 'shared identity'.%
\item%
{[}minor{]} Em (c), a resolução refuta o 'direito de recusar' com 'it's going to be taken too' e 'never{-}ending cycle', o que está correto. Não comenta, porém, o verbo 'ilustra': nenhum exemplo do texto mostra uma recusa. É um reforço opcional.%
\end{itemize}%
\subsection*{Questão 60: nota 9,5 (aprovada, rodada 1 de 1)}%
\textbf{Resposta da resolução às cegas do avaliador:} C: "A forma como assistentes de IA são desenvolvidos tem o potencial de transformar a intimidade em mercadoria." Isso corresponde a "AI companion products can commodify intimacy through emotionally manipulative design...".\par %
Conferi o enunciado na p. 21 do PDF (p. 19 do caderno): a transcrição está fiel ao texto e às quatro alternativas. O candidato e a minha referência chegam à mesma resposta, a (C), que está correta. Ela corresponde a "AI companion products can commodify intimacy through emotionally manipulative design". O candidato mapeia "design" para "a forma como são desenvolvidos", "can" para "tem o potencial" e "commodify" para "transformar em mercadoria".\newline%
\newline%
As quatro alternativas são analisadas com citação do texto como prova:\newline%
{-} (a) confunde a lente de análise (economia política e estudos culturais) com o assunto dos estereótipos, que são de raça e gênero.\newline%
{-} (b) contraria "structurally foreclose{[}s{]} the possibility of genuinely reciprocal relationships". O candidato também nota que "cruel companionship" é um conceito criado pelos autores, e não algo que os assistentes "reforçam".\newline%
{-} (d) tem a primeira parte correta, mas a conclusão é invertida: os assistentes recorrem a estereótipos ("draw on") e os reembalam ("repackage"), em vez de combatê{-}los.\newline%
\newline%
As armadilhas principais foram todas tratadas: foreclose, draw on/repackage, lente versus objeto, o conceito criado pelos autores, o modal can e a primeira parte verdadeira da (d). Os problemas são só de forma: uma citação com flexão alterada sem marcação, os conectivos "apesar de" e "portanto" não apontados explicitamente como o mecanismo das distratoras, e a ausência de comentário sobre a tradução "assistentes" versus "companions". A resolução tem qualidade de gabarito comentado.%
\begin{itemize}%
\item%
{[}minor{]} Na parte 'Ideia central', o candidato coloca entre aspas a forma "forecloses", como se fosse citação literal. O original traz "structurally foreclose" (sem {-}s, com sujeito plural). Na alternativa (b) ele usa "foreclose{[}s{]}" com colchete, o que é aceitável. Mesmo assim, uma citação direta deve reproduzir o texto exatamente ou marcar a alteração.%
\item%
{[}minor{]} Na (d), o candidato mostra que a conclusão contradiz o texto, mas não aponta o conectivo "portanto" como a origem da falsa relação de consequência. Também não comenta, na (b), que o "apesar de" inverte a lógica concessiva. Ficaria mais completo se dissesse de forma explícita que esses conectivos são o mecanismo das distratoras.%
\item%
{[}minor{]} O candidato não comenta que as alternativas traduzem "AI companions" como "assistentes de IA", enquanto ele próprio usa "companheiros de IA". Isso não afeta a resposta, mas uma observação de uma linha evitaria que o leitor achasse que essa diferença invalida a (c).%
\end{itemize}%
\subsection*{Questão 61: nota 9,3 (aprovada, rodada 1 de 1)}%
\textbf{Resposta da resolução às cegas do avaliador:} Alternativa D. O poema alterna verbos no passado e no particípio para os traumas coloniais (stripped, erased, buried, burned, bled, were buried) com verbos no presente simples para a resistência (breathes, births, dances, carry, wear, rise).\par %
O candidato e a minha referência marcam a mesma alternativa, D, e ela está correta. Conferi a transcrição com a página 21 do PDF (caderno p. 19) e ela é fiel.\newline%
\newline%
A resolução organiza bem o poema em dois blocos: passado e particípio para o trauma (stripped, erased, buried, burned, bled, were buried), presente simples para a resistência (breathes, births, dances, carry, am, wear, pulses, rise). Também identifica corretamente 'who's carried' como present perfect e aponta os marcadores de virada 'Yet' e 'And still'.\newline%
\newline%
As quatro alternativas são refutadas ou confirmadas com citação do texto:\newline%
{-} (a) troca a cruz pela coroa, e o candidato mostra isso pelo verso 'She's carried the cross \ldots{} but has not worn its crown'.\newline%
{-} (b) inverte a função dos tempos verbais.\newline%
{-} (c) tem a primeira parte verdadeira (há personificação), mas a segunda contradiz 'I carry her in my marrow' e 'I am her memory'. O candidato não aceita a alternativa por estar meio certa.\newline%
{-} (d) é confirmada pela alternância.\newline%
\newline%
A menção a 'Still I Rise', de Maya Angelou, é um acréscimo pertinente. As armadilhas principais (cruz $\times$ coroa, inversão em (b), (c) meio certa, present perfect, marcadores de contraste) estão todas tratadas. Faltam apenas detalhes: a ressalva sobre o particípio positivo 'crowned' e o comentário sobre 'infringido'. É uma resolução de qualidade de gabarito comentado.%
\begin{itemize}%
\item%
{[}minor{]} A resolução não trata a ressalva de que nem todo particípio do poema exprime trauma. 'A continent crowned by salt and sun' e 'we rise / crowned' são particípios de sentido positivo (o candidato cita 'crowned' mas só destaca 'rise'). Uma frase reconhecendo que (d) descreve o padrão predominante, e não uma regra absoluta, tornaria a defesa de (d) imune a essa objeção.%
\item%
{[}minor{]} 'has not worn' (present perfect negativo) aparece no verso de contraste, mas a resolução não comenta que essa forma, na prática, fecha o bloco do passado ('ainda não usou a coroa'), e não um trauma. A interpretação dada ('o passado é de dor; o presente, de reivindicação') está certa, só falta precisão gramatical.%
\item%
{[}minor{]} O 'infringido' da alternativa (a), desvio de norma do enunciado (o correto seria 'infligido'), não é mencionado. Isso não afeta o julgamento, porque (a) cai pelo conteúdo, e a resolução refuta corretamente pela oposição cruz/coroa.%
\end{itemize}%
\subsection*{Questão 62: nota 9,3 (aprovada, rodada 1 de 1)}%
\textbf{Resposta da resolução às cegas do avaliador:} Alternativa D: usar animais como propriedade ou produto fere dois interesses que todos os seres sencientes, humanos ou não, têm em comum: o interesse em viver ("not dying") e o interesse em não sofrer ("not suffering").\par %
Conferi as páginas 21 e 22 do PDF. A transcrição está fiel, e a resposta certa é (D). O candidato e minha referência chegaram a ela pelo mesmo caminho.\newline%
\newline%
Acertos do candidato:\newline%
{-} Resume a tese central e aponta os três trechos que sustentam a alternativa (d): os animais explorados são sencientes; todo ser senciente tem os interesses de não sofrer e de não morrer; o uso como propriedade atinge esses dois interesses.\newline%
{-} Em (a), refuta a alternativa pela senciência explícita dos animais explorados. Também observa que o direito de não ser propriedade decorre justamente da senciência.\newline%
{-} Em (b), aponta os dois erros. O primeiro é que 'criminalizar' extrapola o texto, que fala em 'make clear\ldots{} morally unjustifiable'. O segundo é que 'however humane' inverte a justificativa 'prática desumana'. Essa é a armadilha mais sutil, e ele a explica bem.\newline%
{-} Em (c), resolve corretamente a referência de 'the former' e 'the latter'.\newline%
{-} Em (d), relaciona 'humanos e animais em comum' a 'all sentient beings, human or nonhuman' e reconhece as paráfrases: 'produto' corresponde a 'property', 'fere' a 'implicates' e 'viver' a 'not dying'.\newline%
\newline%
Problemas, todos menores:\newline%
{-} Em (a), a citação omite 'emoções'.\newline%
{-} Em (c), não comenta a palavra 'carece'.\newline%
{-} Em (d), não comenta 'at least two'.\newline%
\newline%
A resolução é clara e bem fundamentada, e poderia ser publicada como gabarito comentado com esses pequenos retoques.%
\begin{itemize}%
\item%
{[}minor{]} Na alternativa (a), o candidato cita a ressalva como 'mesmo que alguns deles não sejam capazes de sentir', mas o original diz 'sentir emoções'. A citação ficou truncada e ele não explica por que 'sentir emoções' corresponde a ser senciente, que é o que o texto afirma sobre os animais explorados. O que deveria ser feito: citar a frase inteira e dizer em uma linha que a ressalva nega a senciência garantida por 'There is absolutely no doubt that the animals we routinely exploit are sentient'.%
\item%
{[}minor{]} Na alternativa (c), o candidato mostra a inversão entre 'the former' (julgar pessoas) e 'the latter' (julgar ações), mas não comenta a afirmação de que o movimento 'carece' desse princípio. O texto também não diz isso. Seria um segundo motivo, mais curto, para descartar a alternativa.%
\item%
{[}minor{]} Na alternativa (d), falta comentar 'at least two interests'. Era bom dizer que 'pelo menos dois' não invalida a alternativa, porque ela cita os dois interesses que o texto nomeia e não afirma que são os únicos. Fecharia a armadilha de quem acha que a alternativa restringe demais.%
\end{itemize}%
\subsection*{Questão 63: nota 9,5 (aprovada, rodada 1 de 1)}%
\textbf{Resposta da resolução às cegas do avaliador:} Alternativa A: d = c$\cdot$$\Delta$t = 3,0$\times$10$^{8}$ m/s $\times$ (1,3$\times$10$^{9}$ $\times$ 3,2$\times$10$^{7}$ s) = 1,248$\times$10$^2$$^{5}$ m $\approx$ 1,2$\times$10$^2$$^{5}$ m.\par %
A resposta final está certa: alternativa (a), 1,2$\times$10$^2$$^{5}$ m, igual à minha referência. Conferi o enunciado na p. 22 do PDF e a transcrição está correta. Refiz as contas em Python: 3,0e8 $\times$ 1,3e9 $\times$ 3,2e7 = 1,248e25. Os três distratores também conferem: (b) 3,9e17, sem converter anos em segundos; (c) 1,248e22, com bilhão = 10$^{6}$; (d) 3,9e14, com os dois erros juntos. O candidato declara as hipóteses: movimento uniforme com v $\approx$ c, bilhão = 10$^{9}$ e expansão do Universo desprezada. Converte o tempo para segundos (4,16$\times$10$^{1}$$^{6}$ s), verifica as unidades (m/s$\cdot$s = m), arredonda para 2 algarismos significativos e explica a origem de cada distrator. Assim, todas as armadilhas que listei foram tratadas: conversão para segundos, bilhão = 10$^{9}$, normalização da notação científica, hipótese do modelo, uso dos valores dados e explicação dos distratores. Há só duas ressalvas pequenas. A checagem por ano{-}luz repete a mesma conta em outra ordem e não é uma verificação independente. A separação entre mantissas e potências poderia estar escrita. A resolução serve como gabarito comentado.%
\begin{itemize}%
\item%
{[}minor{]} A 'verificação da ordem de grandeza' por ano{-}luz multiplica os mesmos três fatores (c, 3,2$\times$10$^{7}$ s e 1,3$\times$10$^{9}$), só que em outra ordem. Ela pega erros de normalização das potências, mas não é uma checagem independente, e dizer que 'confirma o resultado' é exagero. Uma conferência independente usaria um dado externo, por exemplo a ordem de grandeza conhecida da distância de GW150914, cerca de 410 Mpc $\approx$ 1,3$\times$10$^2$$^{5}$ m. Outra opção é apresentar o trecho apenas como releitura da conta.%
\item%
{[}minor{]} A conta aparece direto como 12,48$\times$10$^2$$^{4}$ m. Faltou separar mantissas (3,0$\cdot$1,3$\cdot$3,2 = 12,48) e potências (10$^{8}$$\cdot$10$^{9}$$\cdot$10$^{7}$ = 10$^2$$^{4}$), o que tornaria a normalização para 1,248$\times$10$^2$$^{5}$ explícita linha a linha. É um detalhe de apresentação; o resultado está certo.%
\end{itemize}%
\subsection*{Questão 64: nota 9,5 (aprovada, rodada 1 de 1)}%
\textbf{Resposta da resolução às cegas do avaliador:} Alternativa B. A aceleração da gravidade na superfície é g = GM/R$^2$. Ela não depende da massa m do corpo, é diretamente proporcional à massa M da Terra e inversamente proporcional ao quadrado do raio R da Terra.\par %
Conferi o enunciado na página 22 do PDF, e a transcrição está fiel. O candidato marcou B, que é a resposta correta e coincide com a minha referência.\newline%
\newline%
A dedução está correta. A Terra é tratada como esfericamente simétrica, com justificativa pelo teorema das cascas. A partir de F = GMm/R$^2$ e da 2ª Lei de Newton, ele obtém g = F/m = GM/R$^2$. O cancelamento de m é justificado explicitamente pela igualdade entre massa gravitacional e massa inercial.\newline%
\newline%
Também estão presentes as verificações:\newline%
{-} a análise dimensional, que dá N/kg = m/s$^2$;\newline%
{-} a ordem de grandeza, que recalculei e dá $\approx$ 9,81 m/s$^2$;\newline%
{-} a análise das quatro alternativas, que identifica a C como o distrator da força peso (proporcional a m) e a A e a D como erradas por inverterem a dependência em R$^2$.\newline%
\newline%
As principais armadilhas foram tratadas: confundir força com aceleração, o sentido da dependência em R, o princípio da equivalência e a simetria esférica. Ficaram só falhas menores:\newline%
{-} a frase de que o enunciado manda desprezar a rotação, que ele não diz;\newline%
{-} a falta de \textbackslash{}"com as demais grandezas fixas\textbackslash{}" nas proporcionalidades;\newline%
{-} nenhuma menção à aceleração relativa G(M+m)/R$^2$, que é desprezível.\newline%
\newline%
É uma resolução que serviria como gabarito comentado.%
\begin{itemize}%
\item%
{[}minor{]} A resolução diz que despreza a rotação da Terra "como faz o enunciado", mas o enunciado não manda desprezá{-}la. O que ele faz é pedir a aceleração "decorrente da ação da força gravitacional da Terra", e a rotação fica de fora por isso. Seria melhor dizer isso e acrescentar que, mesmo incluindo a rotação (g\_ef $\approx$ GM/R$^2$ - $\omega$$^2$R cos$^2$$\lambda$), g continua sem depender de m.%
\item%
{[}minor{]} As proporcionalidades são exemplificadas ("dobrando M, g dobra"; "dobrando R, g cai para 1/4"), mas falta dizer que g $\propto$ M vale com R fixo e que g $\propto$ 1/R$^2$ vale com M fixo. Esse ceteris paribus deveria estar escrito.%
\item%
{[}minor{]} A resolução não comenta que GM/R$^2$ é a aceleração do corpo em relação a um referencial inercial e que a aceleração relativa corpo--Terra seria G(M+m)/R$^2$. Essa diferença é desprezível porque m ? M e não muda a resposta, mas um gabarito comentado completo poderia mencioná{-}la em uma linha.%
\end{itemize}%
\subsection*{Questão 65: nota 9,5 (aprovada, rodada 1 de 1)}%
\textbf{Resposta da resolução às cegas do avaliador:} Alternativa d) 240 N. A tração no fio que sustenta o bloco é T = mg = 120 N. A polia móvel é puxada para baixo pelos dois ramos desse fio (2T = 240 N), e o fio da manopla, preso ao eixo dela e desviado só por polias fixas, transmite essa mesma força: F = 2mg = 240 N.\par %
As respostas finais coincidem: o candidato e a referência marcam (d) 240 N. Conferi a montagem na figura da p. 22. O fio da manopla passa pelas duas polias fixas e desce até o eixo da polia móvel. O segundo fio passa por cima da polia móvel, com um ramo preso à base e o outro sustentando o bloco. As setas T junto à polia apontam para baixo nos dois ramos. A resolução isola cada corpo na ordem certa. No bloco, T = mg = 120 N. Na polia móvel sem massa, T$_{1}$ = 2T = 240 N. No fio 1, que passa só por polias fixas ideais, F = T$_{1}$ = 240 N. As contas conferem: 12$\cdot$10 = 120 e 2$\cdot$120 = 240. As hipóteses estão explícitas: fios inextensíveis e sem massa, polias sem massa e sem atrito, equilíbrio na iminência. Cada igualdade de trações vem justificada por elas. Todas as armadilhas foram tratadas. A topologia foi lida corretamente, com o destaque de que a carga pende de um ramo e não do eixo. A falácia de 'a polia móvel divide por 2' (alternativa b) é refutada. O trecho horizontal entre as polias fixas não altera F. A massa da polia é declarada desprezível. A iminência é interpretada como equilíbrio, sem aceleração. As alternativas a, b e c são descartadas com o motivo de cada erro. Ressalvas menores: o candidato supõe uma normal 'do piso', embora a figura mostre o bloco suspenso, e falta a conferência cinemática por trabalho virtual. Com essas ressalvas pequenas, a resolução tem qualidade de gabarito comentado.%
\begin{itemize}%
\item%
{[}minor{]} No item 1, o candidato afirma que o bloco está apoiado no piso e sofre uma normal N. Na figura (p. 22), o bloco aparece suspenso, um pouco acima da base. O resultado T = mg não muda, porque 'iminência' significa equilíbrio com eventual normal nula, mas o certo seria escrever 'a eventual normal do apoio se anula' ou tratar o caso como simples limiar de equilíbrio, sem dar como certo um contato que a figura não mostra.%
\item%
{[}minor{]} O fator 2 só é justificado pelas forças. Não há a verificação cinemática ou por trabalho virtual: a polia sobe d, o bloco sobe 2d e a mão desce d, logo F$\cdot$d = mg$\cdot$2d. Ela não é obrigatória, mas reforçaria o argumento da 'desvantagem mecânica', que hoje é apenas qualitativo na seção de verificação.%
\end{itemize}%
\subsection*{Questão 66: nota 9,5 (aprovada, rodada 1 de 1)}%
\textbf{Resposta da resolução às cegas do avaliador:} A alternativa correta é a **C (4,0 N)**. O disco está deslizando, então o atrito é cinético: f\_at = $\mu$\_C$\cdot$m$\cdot$g = 0,02 $\cdot$ 20 $\cdot$ 10 = 4,0 N.\par %
Conferi a transcrição com o original (p. 22 do PDF, impressa como p. 20, Questão 66), e ela está fiel. Minha referência e o candidato chegam à mesma resposta, a alternativa C (4,0 N), e ela está correta: o disco está deslizando ('se desloca', 'durante seu deslocamento'), então o atrito é cinético, N = mg = 200 N e f\_at = $\mu$\_C$\cdot$N = 0,02$\cdot$200 = 4,0 N.\newline%
\newline%
Pontos fortes da resolução:\newline%
{-} Justifica explicitamente o uso do coeficiente cinético pelo movimento relativo entre disco e gelo.\newline%
{-} Explica que $\mu$\_E$\cdot$N é apenas o valor máximo do atrito estático, de modo que $\mu$\_E é um dado distrator.\newline%
{-} Obtém a normal pela 2ª lei de Newton na vertical, usando a pista horizontal.\newline%
{-} Explica a origem de cada alternativa errada: $\mu$\_C$\cdot$m e $\mu$\_E$\cdot$m esquecem g e saem com unidade errada; $\mu$\_E$\cdot$mg é o atrito estático máximo.\newline%
{-} Confere as unidades e a plausibilidade física: o atrito é 2\% do peso e a desaceleração é a = 0,2 m/s$^2$, ambos coerentes.\newline%
\newline%
Faltam dois detalhes menores, que não afetam a correção. A resolução não diz por que a varredura e a rotação são irrelevantes para o valor pedido, e não deixa explícito que as vassouras não exercem força sobre o disco.\newline%
\newline%
As contas estão certas: 0,02$\cdot$20$\cdot$10 = 4; 0,1$\cdot$200 = 20; 0,02$\cdot$20 = 0,4; 0,1$\cdot$20 = 2. A resolução está no nível de um gabarito comentado.%
\begin{itemize}%
\item%
{[}minor{]} A resolução não diz por que a varredura (e a rotação do disco) não muda o valor pedido. Faltou uma frase dizendo que o enunciado já dá o $\mu$\_C da situação considerada e que, no modelo de Coulomb, a rotação não altera o módulo. Isso fecharia a armadilha 'a varredura muda o atrito, então qual $\mu$ usar?'.%
\item%
{[}minor{]} A hipótese 'as únicas forças verticais são o peso e a normal' aparece, mas não é justificada contra a objeção natural do contexto. As vassouras atuam no gelo e não empurram o disco, e uma meia frase bastaria para dizer isso. O mesmo vale para a desaceleração a = $\mu$\_C$\cdot$g: ela supõe que o atrito é a única força horizontal, o que é correto, mas fica implícito.%
\end{itemize}%
\subsection*{Questão 67: nota 9,5 (aprovada, rodada 1 de 1)}%
\textbf{Resposta da resolução às cegas do avaliador:} T3 = 50 $^\circ$C, pois $\Delta$T = Q/(m$\cdot$c\_liq) = 80/(10$\cdot$0,4) = 20 $^\circ$C acima de 30 $^\circ$C. O gráfico certo sobe de 20 $^\circ$C para 30 $^\circ$C (fase sólida), fica num patamar em 30 $^\circ$C durante a fusão e depois sobe até 50 $^\circ$C em t = 230 s. Alternativa b).\par %
Conferi o enunciado e os quatro gráficos na página 23 do PDF, e a transcrição está fiel. O candidato marca a alternativa (B), a mesma da minha referência, e a resposta está correta. A conta $\Delta$T = 80/(10$\cdot$0,4) = 20 $^\circ$C está certa e as unidades fecham sem conversão (J ÷ (g$\cdot$J/(g$\cdot$$^\circ$C)) = $^\circ$C). Ele parte de T2 = 30 $^\circ$C e chega a T3 = 50 $^\circ$C, sem cair no erro de somar $\Delta$T a T1. Eliminou cada alternativa pelo motivo certo. A (d) cai porque mostra resfriamento. A (c) cai porque não tem patamar de fusão e porque chega a 50 $^\circ$C perto de 220 s, não em 230 s; conferi isso no gráfico original. A (a) cai porque terminar em 120 $^\circ$C exigiria Q = 360 J, e não 80 J. A (b) atende a tudo: subida, patamar em 30 $^\circ$C e subida até 50 $^\circ$C em t = 230 s. Todas as armadilhas foram tratadas. O patamar é apresentado como obrigatório e é o que derruba a (c). O Q de 80 J é aplicado só ao trecho líquido. A potência variável aparece explicitamente como motivo para não comparar as inclinações. O candidato não inventa calor específico do sólido nem calor latente, confere o instante final e menciona a quebra do eixo. A solução é bem organizada, tem conferência de unidades e uma conclusão clara. Só faltam duas pequenas explicitações de hipótese na redação. Poderia ser publicada como gabarito comentado.%
\begin{itemize}%
\item%
{[}minor{]} Ao falar da fusão, o candidato supõe sem dizer que a amostra é Ga puro e que a pressão é constante. É a hipótese padrão e o enunciado a autoriza implicitamente. Uma frase como "supondo Ga puro a pressão constante" deixaria a premissa explícita.%
\item%
{[}minor{]} A frase "recebe calor, portanto sua temperatura não pode diminuir" está certa no contexto, já que não há trabalho relevante e todo o calor recebido aumenta a energia interna. Ainda assim, poderia citar esse motivo: o calor vira calor sensível ou latente, e nos dois casos a temperatura não cai. É só uma questão de rigor na redação e não afeta o resultado.%
\end{itemize}%
\subsection*{Questão 68: nota 9,5 (aprovada, rodada 1 de 1)}%
\textbf{Resposta da resolução às cegas do avaliador:} i = P/U = 3,0/12 = 0,25 A, alternativa b).\par %
Conferi o enunciado na página 23 do PDF (questão 68) e a transcrição está correta. Os dados são P = 3,0 W e U = 12 V.\newline%
\newline%
A resposta do candidato é a alternativa b), i = 0,25 A, igual à minha referência e correta: i = P/U = 3,0/12 = 0,25 A.\newline%
\newline%
A resolução é muito boa:\newline%
{-} **Lei usada:** P = U$\cdot$i, deduzida da definição de potência (energia U$\cdot$$\Delta$q por intervalo $\Delta$t). O candidato observa corretamente que não é preciso supor que o roteador seja ôhmico.\newline%
{-} **Unidades:** confere que W/V = A.\newline%
{-} **Verificação:** refaz U$\cdot$i = 3,0 W e comenta que a ordem de grandeza é plausível.\newline%
{-} **Distratores:** analisa todos corretamente. A alternativa a) dá 0,36 W, cerca de 8 vezes menos que o correto. A c) é a razão invertida U/P e daria 48 W. A d) é o produto P$\cdot$U, que nem tem unidade de corrente. Todas essas contas estão certas.\newline%
\newline%
Das armadilhas que listei, foram tratadas a razão invertida, o produto, o erro de ordem de grandeza, a análise dimensional e a hipótese de corrente contínua. Ficou só implícita a de que os percentuais são apenas contexto: a resolução não os usa, mas não diz isso.\newline%
\newline%
Há dois detalhes menores: falta dizer que os percentuais são irrelevantes para a conta, e a hipótese de corrente contínua é assumida sem ser apresentada como suposição. A resolução está em nível de gabarito comentado.%
\begin{itemize}%
\item%
{[}minor{]} A resolução não diz que os percentuais do texto (8\% do consumo da casa e 1,5\% do consumo do país) são só contexto e não entram na conta. Ela também não os usa, então não há erro. Faltou uma frase dizendo que esses dados não interferem no cálculo.%
\item%
{[}minor{]} A resolução afirma que P = U$\cdot$i 'vale para qualquer dispositivo, ôhmico ou não'. Isso é verdade para corrente contínua e para valores instantâneos. Em corrente alternada, com valores eficazes, a potência média é U$\cdot$i$\cdot$cos$\varphi$. A ressalva já está implícita, porque o candidato declarou a alimentação como contínua. Mas essa hipótese de CC não está no enunciado, e o certo seria apresentá{-}la explicitamente como suposição, observando que ela não muda a resposta.%
\end{itemize}%
\subsection*{Questão 69: nota 9,5 (aprovada, rodada 1 de 1)}%
\textbf{Resposta da resolução às cegas do avaliador:} d) 30 mA. Com o alarme ativado, a chave fecha e os quatro resistores de 1,2 k$\Omega$ ficam em paralelo. Assim, R\_eq = 300 $\Omega$ e i = 9,0 V / 300 $\Omega$ = 30 mA.\par %
Conferi a figura no PDF original (p. 23). Ela mostra quatro ramos em paralelo, ligados aos mesmos dois nós, e a bateria de 9 V em série com o conjunto. A chave do ramo de R4 aparece aberta. A resposta do candidato, (d) 30 mA, coincide com a minha referência e está correta. Com o alarme ativado, a chave fecha, então R\_eq = 1200/4 = 300 $\Omega$ e i = 9,0/300 = 30 mA. Conferi as contas em python: 30; 7,5; 22,5; 1,875 mA. A resolução declara as hipóteses (bateria ideal, fios e chave sem resistência, resistores ôhmicos) e justifica a ligação em paralelo pelos nós comuns. Ela calcula R\_eq corretamente e confere o resultado pela lei dos nós (4 $\times$ 7,5 mA). Também verifica unidades e ordem de grandeza e explica a origem de cada distrator: série (a), um ramo só (b) e três ramos com a chave aberta (c). Todas as armadilhas que identifiquei foram tratadas. Restam apenas detalhes de redação: não diz explicitamente que a figura mostra a chave aberta e não mostra que 22,5 mA arredonda para 22 mA. É uma resolução publicável como gabarito comentado.%
\begin{itemize}%
\item%
{[}minor{]} A resolução não diz explicitamente que a figura mostra a chave ABERTA (estado normal) e que, por isso, a situação pedida é diferente da desenhada. Ela só diz que 'com o dispositivo intacto, a chave está aberta'. A conclusão está certa, mas para fechar a armadilha principal com clareza caberia uma frase como 'a figura mostra a chave aberta; com o alarme ativado ela se fecha'.%
\item%
{[}minor{]} No distrator (c), a resolução escreve 22 mA $\approx$ 9,0/400 A sem mostrar que 9,0/400 A = 22,5 mA, valor que o enunciado arredonda para 22 mA (dois algarismos significativos). Faltou explicitar o valor intermediário.%
\end{itemize}%
\subsection*{Questão 70: nota 9,5 (aprovada, rodada 1 de 1)}%
\textbf{Resposta da resolução às cegas do avaliador:} Alternativa D. Na ética do discurso, justiça e solidariedade estão estreitamente ligadas e têm a mesma raiz, que é a vulnerabilidade de seres que só se tornam indivíduos pela socialização. A justiça preserva os direitos dos indivíduos e a solidariedade preserva as relações intersubjetivas, isto é, o bem{-}estar da comunidade.\par %
A resposta é a alternativa D, a mesma da referência e correta. A transcrição confere com a página 24 do PDF. A resolução reconstrói a tese do trecho: os dois princípios têm a mesma raiz, que é a vulnerabilidade, e a individuação só ocorre pela socialização, de modo que defender o indivíduo obriga a defender a comunidade. Cada alternativa é justificada com o próprio texto. Na (c), o candidato identifica as duas inversões: a vulnerabilidade é ponto de partida e não algo a combater, e a individuação depende da socialização. Na (b), mostra que 'em sentido moderno' só caracteriza a justiça. Na (d), mostra que a divisão justiça/indivíduo $\times$ solidariedade/relações intersubjetivas é paráfrase fiel e que a 'relação estreita' vem da raiz comum. Não usa conhecimento externo sobre Habermas nem contesta dados do enunciado. Há três ressalvas menores: o 'em contrapartida' não é comentado na (a), a implicação mútua vai um pouco além da direção literal do texto, e a refutação da (b) poderia ser mais direta. É uma resolução de nível de gabarito comentado.%
\begin{itemize}%
\item%
{[}minor{]} Na alternativa (a), o candidato não trata o conectivo 'em contrapartida', o único ponto do texto em que a leitura de oposição poderia se apoiar. Faltou uma frase dizendo que ele só marca a diferença de foco (indivíduo $\times$ comunidade), e não um conflito sem solução.%
\item%
{[}minor{]} O candidato afirma que 'um não se sustenta sem o outro' e que 'defender um obriga a defender o outro'. O texto afirma explicitamente só uma direção: defender os direitos do indivíduo obriga a defender o bem{-}estar da comunidade. A volta é inferida de 'mesma raiz' e 'igualmente'. O ideal seria citar a direção literal e apresentar a recíproca como consequência da raiz comum.%
\item%
{[}minor{]} A refutação da alternativa (b) poderia ser mais direta: a conclusão de Habermas é o oposto da alternativa, porque a moral que defende o indivíduo é obrigada a defender a comunidade. Logo a liberdade subjetiva não impede a solidariedade. O argumento do candidato está correto, só um pouco menos incisivo.%
\end{itemize}%
\subsection*{Questão 71: nota 9,3 (aprovada, rodada 1 de 1)}%
\textbf{Resposta da resolução às cegas do avaliador:} C: a força desumaniza tanto quem a sofre (vira coisa ou cadáver) quanto quem a usa, que é "embriagado" por ela com a mesma dureza e nunca a possui de fato.\par %
Minha resolução de referência confirma a resposta do candidato: alternativa C. A justificativa é sólida e trata todas as armadilhas principais:\newline%
{-} (a) é refutada com 'julga possuí{-}la' e 'Ninguém a possui verdadeiramente', e o candidato acrescenta que o vencedor também é embriagado pela força.\newline%
{-} Em (c), a passagem 'com a mesma dureza {[}...{]} embriaga' é explicada como desumanização de quem exerce a força: perda do domínio de si e ilusão de posse.\newline%
{-} (b) é rejeitada porque Weil chama as palavras de Aquiles de 'razoáveis'.\newline%
{-} Em (d), o candidato mostra que bois e ovelhas são bens recuperáveis em oposição à vida irrecuperável, sem nenhuma hierarquia entre espécies.\newline%
\newline%
O texto está claro e bem organizado, e as citações são fiéis ao enunciado. Faltam só nuances menores:\newline%
{-} não aponta que a segunda metade de (b) é verdadeira, o que a torna atraente;\newline%
{-} não usa a introdução ('todos os espectros políticos') como reforço;\newline%
{-} não comenta a ambiguidade do 'Ela' no trecho adaptado.\newline%
\newline%
É uma resolução de qualidade de gabarito comentado, com pequenos pontos que poderiam ser enriquecidos.%
\begin{itemize}%
\item%
{[}minor{]} Na alternativa (b), o candidato não comenta que a segunda metade (a perda irreversível da vida humana) é verdadeira. Esse é justamente o chamariz da alternativa. Bastava uma frase dizendo que o erro está no núcleo da afirmação, a 'falta de humanidade', e não na consequência citada.%
\item%
{[}minor{]} Não usa a introdução do enunciado ('domina os homens de todos os espectros políticos durante uma guerra') para reforçar que a força atinge todos os lados. Isso fortaleceria tanto a escolha de (c) quanto a rejeição de (a). O candidato só cita a 'idolatria da força'.%
\item%
{[}minor{]} Na alternativa (a), faltou lembrar que na Ilíada a sorte das armas muda de lado, e por isso 'vencedor' e 'vencido' não são categorias fixas. O argumento central (ninguém possui a força, quem a usa só 'julga possuí{-}la') já basta. Fica como sugestão de enriquecimento.%
\item%
{[}minor{]} Ao dizer que a força 'que não mata' transforma o homem vivo em pedra, o candidato acerta em relação ao ensaio original. Mas o trecho adaptado diz 'A força que mata {[}...{]}. Ela transforma o homem que continua vivo em pedra'. Uma nota sobre essa ambiguidade do 'Ela' no texto adaptado deixaria a leitura mais precisa. Isso não afeta a resposta.%
\item%
{[}minor{]} A expressão 'vira seu instrumento' é uma leve extrapolação: o texto fala em 'embriaga' e 'ninguém a possui'. É uma paráfrase defensável e coerente com Weil, mas seria melhor ancorá{-}la mais diretamente no verbo 'embriagar' (perder a medida e o domínio de si), como o candidato faz no primeiro bloco.%
\end{itemize}%
\subsection*{Questão 72: nota 9,5 (aprovada, rodada 1 de 1)}%
\textbf{Resposta da resolução às cegas do avaliador:} Alternativa A. Segundo Dussel, o ego moderno não se forma só pelo cogito cartesiano. Ele também se define como "ego descobridor, conquistador, colonizador" do Outro, e essa Alteridade é "constitutiva da própria Modernidade".\par %
A transcrição confere com a página 24 do PDF. O candidato marcou A, a mesma alternativa da minha referência, e a resposta está certa.\newline%
\newline%
A resolução começa com uma síntese fiel do texto. Ela separa a gestação nas cidades medievais do nascimento em 1492. Mostra que o Outro é constitutivo da Modernidade e que o título fala em 'encobrimento', não em 'descobrimento'. Também deixa claro que o diálogo é uma tarefa proposta, e não algo que a Modernidade já tenha feito. Depois justifica cada alternativa com base no texto.\newline%
\newline%
{-} **(a):** a resolução explica que o cogito entra como conhecimento prévio, por ser o marco clássico do sujeito moderno. Mostra que a alternativa não nega o cogito ('não apenas... mas também') e a liga a 'ego descobridor, conquistador, colonizador' e à palavra 'constitutiva'. Isso responde à armadilha de descartar a alternativa por citar Descartes, que não aparece no trecho.\newline%
{-} **(b):** a resolução separa a filosofia do diálogo, que é tarefa futura, da filosofia moderna, que no texto é dominação. Também critica o termo 'descoberto'.\newline%
{-} **(c):** a resolução contrapõe 'escombros' a 'cidades livres, centros de enorme criatividade'. Aponta ainda que 'por definição' e 'progresso e retrocesso' não estão no texto.\newline%
{-} **(d):** a resolução recusa a ideia de mediação neutra para 'dirimir' a violência. Mostra que a filosofia da libertação toma o partido do oprimido.\newline%
\newline%
As falhas são todas menores:\newline%
{-} Em (c), não diz que 'conquista e violência' está certo e funciona como isca.\newline%
{-} Em (d), não aponta que a alternativa reduz o Outro aos 'colonizados'.\newline%
{-} Cita 'o africano' e o 'ego conquiro', que vêm da obra mais ampla de Dussel, sem avisar que isso está fora do trecho.\newline%
\newline%
A resolução tem qualidade de gabarito comentado.%
\begin{itemize}%
\item%
{[}minor{]} Alternativa (c): a resolução refuta bem 'dos escombros da Idade Média' com o trecho 'cidades europeias medievais, livres, centros de enorme criatividade'. Também aponta que 'por definição contraditória' e 'progresso e retrocesso' não estão no texto. Só não diz que 'conquista e violência' está de acordo com o texto e funciona como isca. Deveria registrar que essa parte é verdadeira e que o erro está no resto da frase. Isso deixaria claro por que a alternativa engana.%
\item%
{[}minor{]} Alternativa (d): a resolução desmonta bem a ideia de mediação neutra e o verbo 'dirimir'. Mas não aponta que a alternativa reduz 'o Outro' aos 'colonizados', quando o texto fala em 'oprimido, incomunicado, excluído, Outro'. É um argumento a mais que ficou de fora.%
\item%
{[}minor{]} Na 'Ideia central', a resolução diz que a Alteridade é 'o indígena, o africano, o colonizado'. O trecho não fala do africano. Isso vem da obra mais ampla de Dussel, e a resolução deveria avisar que é conhecimento externo. O mesmo vale para a menção ao 'ego conquiro', que está correta mas também vem de fora do trecho. Na conclusão, a expressão 'e de forma decisiva' vai um pouco além da alternativa, mas a palavra 'constitutiva' do texto sustenta essa leitura.%
\end{itemize}%
\end{document}